\documentclass{article}
\usepackage[utf8]{inputenc}
\usepackage[spanish]{babel}
\usepackage{geometry}
\usepackage{xcolor}
\usepackage{booktabs}
\usepackage{longtable}
\geometry{a4paper, margin=1in}

\begin{document}

\newpage
\section{Batería de Preguntas para Validación del Agente}

Este apartado contiene 50 preguntas diseñadas para evaluar los diferentes nodos del agente (rechazo amable, entrevistador, baremo, calendario, normativo, procedimental) en base a los documentos de la Universidad de Sevilla proporcionados. Las respuestas han sido generadas automáticamente por el agente real desplegado en la aplicación.

\subsection{Categoría 1: Rechazo Amable (Fuera de dominio) (1-8)}
\textit{Evalúa la capacidad del agente para detectar consultas fuera del dominio de la US y generar una respuesta cordial de rechazo.}
\begin{itemize}
  \item \textbf{P1.} ¿Cuál es la receta tradicional de la paella valenciana? \\
        \begin{itemize}
          \item \textbf{Tipo de consulta:} simple
          \item \textbf{Respuesta del Agente:}

                Lo siento, pero no puedo ayudarte con la receta tradicional de la paella valenciana. Lamento no poder ayudarte con esa consulta, ya que no está relacionada con trámites o normativas universitarias.

                Te recomendaría contactar con el departamento de cultura o el centro de idiomas de la universidad para obtener información sobre eventos culturales o actividades relacionadas con la gastronomía valenciana. También podrías buscar en línea o consultar con un experto en cocina para obtener la receta que estás buscando.
        \end{itemize}


  \item \textbf{P2.} ¿Me puedes recomendar una buena serie de Netflix para ver este fin de semana? \\
        \begin{itemize}
          \item \textbf{Tipo de consulta:} simple
          \item \textbf{Respuesta del Agente:}

                Lo siento, pero no puedo ayudarte con recomendaciones de series de Netflix. Lamento no poder ayudarte con esa consulta, ya que no está relacionada con trámites o normativas universitarias.

                Te recomendaría contactar con la biblioteca universitaria o con un servicio de entretenimiento del campus para obtener recomendaciones de contenido audiovisual. También podrías buscar en línea en sitios web de críticos de cine o de blogs de entretenimiento para encontrar recomendaciones de series y películas.
        \end{itemize}


  \item \textbf{P3.} ¿Cómo puedo configurar el Bluetooth de mi coche para conectar el móvil? \\
        \begin{itemize}
          \item \textbf{Tipo de consulta:} simple
          \item \textbf{Respuesta del Agente:}

                Lo siento, pero no puedo ayudarte con la configuración del Bluetooth de tu coche. Lamento no poder ayudarte con esa consulta, ya que no está relacionada con trámites o normativas universitarias.

                Te recomendaría contactar con el servicio técnico de tu coche o con el fabricante para obtener ayuda específica sobre cómo configurar el Bluetooth de tu vehículo. También podrías buscar en línea recursos y tutoriales que te puedan guiar a través del proceso.

                Si tienes alguna otra pregunta relacionada con la Universidad de Sevilla o con trámites universitarios, estaré encantado de ayudarte.
        \end{itemize}


  \item \textbf{P4.} ¿Qué equipo ganó la última Champions League? \\
        \begin{itemize}
          \item \textbf{Tipo de consulta:} simple
          \item \textbf{Respuesta del Agente:}

                Lo siento, pero no puedo ayudarte con esa consulta. Lamento no poder ayudarte con esa pregunta, ya que no está relacionada con trámites o normativas universitarias. La Champions League es un torneo de fútbol internacional y no es un tema que esté dentro de mi ámbito de asistencia.

                Te recomendaría contactar con un sitio web de deportes o un canal de televisión que cubra eventos deportivos para obtener la información que buscas. También podrías buscar en línea en sitios como UEFA o ESPN para encontrar la respuesta a tu pregunta.
        \end{itemize}


  \item \textbf{P5.} Necesito consejos para cuidar un gato recién adoptado, ¿me ayudas? \\
        \begin{itemize}
          \item \textbf{Tipo de consulta:} simple
          \item \textbf{Respuesta del Agente:}

                Lo siento, pero no puedo ayudarte con consejos para cuidar un gato recién adoptado. Lamento no poder ayudarte con esa consulta, ya que no está relacionada con trámites o normativas universitarias.

                Te recomendaría contactar con el departamento de bienestar animal de la ciudad o con un veterinario para obtener consejos y orientación sobre cómo cuidar a tu nuevo amigo. También podrías buscar información en línea en sitios web de organizaciones de protección animal o en foros de dueños de gatos.
        \end{itemize}


  \item \textbf{P6.} ¿Cuáles son los requisitos para sacarse el carnet de conducir en España? \\
        \begin{itemize}
          \item \textbf{Tipo de consulta:} simple
          \item \textbf{Respuesta del Agente:}

                ¡Hola! Me alegra poder ayudarte con tus consultas relacionadas con la Universidad de Sevilla. Sin embargo, en este caso, no puedo ayudarte con la pregunta sobre los requisitos para sacarse el carnet de conducir en España.

                Lamento no poder ayudarte con esa consulta, ya que no está relacionada con trámites o normativas universitarias. Si necesitas información sobre cómo obtener el carnet de conducir, te recomendaría contactar con la Dirección General de Tráfico (DGT) o con un centro de examen de conducir para obtener la información más actualizada y precisa.

                Si tienes alguna otra pregunta relacionada con la Universidad de Sevilla, como trámites de matrícula, horarios de clase, requisitos para cursar un grado, etc., estaré encantado de ayudarte. ¡No dudes en preguntar!
        \end{itemize}


  \item \textbf{P7.} Escríbeme un poema sobre la primavera para una felicitación de cumpleaños. \\
        \begin{itemize}
          \item \textbf{Tipo de consulta:} simple
          \item \textbf{Respuesta del Agente:}

                Lo siento, pero no puedo ayudarte con la solicitud de un poema sobre la primavera para una felicitación de cumpleaños. Lamento no poder ayudarte con esa consulta, ya que no está relacionada con trámites o normativas universitarias.

                Te recomendaría contactar con el departamento de cultura o comunicación de la universidad, ya que podrían tener recursos o sugerencias para crear contenido creativo como un poema para una felicitación de cumpleaños. También podrías buscar en línea o en redes sociales para encontrar ideas y ejemplos de poesías sobre la primavera que puedan inspirarte.
        \end{itemize}


  \item \textbf{P8.} ¿Qué diferencias hay entre un iPhone 15 y un Samsung Galaxy S24? \\
        \begin{itemize}
          \item \textbf{Tipo de consulta:} simple
          \item \textbf{Respuesta del Agente:}

                Lo siento, pero no puedo ayudarte con esa consulta. Lamento no poder ayudarte con esa pregunta, ya que no está relacionada con trámites o normativas universitarias. La Universidad de Sevilla no tiene información específica sobre los modelos de teléfonos móviles ni sus características.

                Te recomendaría contactar con un especialista en tecnología o con el servicio de atención al cliente de Apple o Samsung para obtener más información sobre las diferencias entre estos modelos. También podrías buscar en línea en sitios web de tecnología o foros de discusión para encontrar respuestas a tus preguntas.
        \end{itemize}


\end{itemize}

\subsection{Categoría 2: Entrevistador (9-16)}
\textit{Analiza el desempeño del agente cuando la consulta es ambigua y necesita pedir más datos al usuario.}
\begin{itemize}
  \item \textbf{P9.} Voy a ir a un congreso y quiero que me liquiden los gastos del viaje, ¿qué tengo que entregar documentalmente? \\
        \begin{itemize}
          \item \textbf{Tipo de consulta:} multi
          \item \textbf{Respuesta del Agente:}

                ⚠️ Requisitos Previos:
                Antes de iniciar el trámite de liquidación de gastos, asegúrate de tener todos los documentos necesarios y de haber cumplido con los requisitos previos, como la tramitación de la Comisión de Servicios o Autorización de Desplazamiento, y la comprobación de que la actividad a desarrollar es un gasto elegible en la convocatoria del proyecto.

                A continuación, te presento los pasos a seguir para la liquidación de gastos del viaje a un congreso:

                1. \textbf{Obtén el certificado de asistencia}: Debes obtener un certificado de asistencia emitido por la entidad organizadora del congreso. La mera inscripción no acredita haber asistido.
                2. \textbf{Recopila los justificantes de gasto y pago}: Debes recopilar todos los justificantes de gasto y pago, como facturas, tickets, billetes electrónicos, etc. Asegúrate de que las facturas reflejen tus datos fiscales y no los de la Universidad de Sevilla.
                3. \textbf{Rellena el formulario de liquidación de gastos}: Debes rellenar el \textbf{Formulario de liquidación de gastos de viaje} y firmarlo a través de \textbf{Portafirmas}, tanto por ti como por el Investigador Principal (IP).
                4. \textbf{Prepara la documentación adicional}: Debes preparar la siguiente documentación adicional:
                * \textbf{Memoria justificativa}: Una memoria justificativa que describa las actividades realizadas y su relación con el proyecto.
                * \textbf{Documentos acreditativos}: Documentos acreditativos de la actividad realizada, como el certificado de asistencia.
                * \textbf{Facturas y justificantes de pago}: Facturas y justificantes de pago de los gastos realizados.
                5. \textbf{Presenta la solicitud de abono}: Debes presentar la solicitud de abono de la liquidación de gastos a través del \textbf{Registro Electrónico General}, dirigido a la Subdirección General de Gestión Económica de la Investigación. Asegúrate de adjuntar todos los documentos necesarios, incluyendo el \textbf{Formulario de liquidación de gastos}, la \textbf{Memoria justificativa}, los \textbf{Documentos acreditativos}, las \textbf{Facturas y justificantes de pago}, y cualquier otro documento requerido.
                6. \textbf{Verifica la documentación}: Verifica que la documentación presentada sea completa y correcta, y que se ajuste a los requisitos establecidos en la \textbf{Guía completa de liquidación de viajes} y en la \textbf{Miniguía para viajes por razón del servicio}.

                Recuerda que es importante seguir las instrucciones y requisitos establecidos en la \textbf{Guía completa de liquidación de viajes} y en la \textbf{Miniguía para viajes por razón del servicio} para asegurarte de que el trámite se realice de manera correcta y eficiente.
          \item \textbf{Interacción Usuario:} Soy PDI (Personal Docente e Investigador).
          \item \textbf{Respuesta Final:}

                ⚠️ Requisitos Previos:
                Antes de iniciar el trámite de liquidación de gastos, asegúrate de tener todos los documentos necesarios y de haber cumplido con los requisitos previos, como la tramitación de la Comisión de Servicios o Autorización de Desplazamiento, y la comprobación de que la actividad a desarrollar es un gasto elegible en la convocatoria del proyecto.

                A continuación, te presento los pasos a seguir para la liquidación de gastos del viaje a un congreso:

                1. \textbf{Obtén el certificado de asistencia}: Debes obtener un certificado de asistencia emitido por la entidad organizadora del congreso. La mera inscripción no acredita haber asistido.
                2. \textbf{Recopila los justificantes de gasto y pago}: Debes recopilar todos los justificantes de gasto y pago, como facturas, tickets, billetes electrónicos, etc. Asegúrate de que las facturas reflejen tus datos fiscales y no los de la Universidad de Sevilla.
                3. \textbf{Rellena el formulario de liquidación de gastos}: Debes rellenar el \textbf{Formulario de liquidación de gastos de viaje} y firmarlo a través de \textbf{Portafirmas}, tanto por ti como por el Investigador Principal (IP).
                4. \textbf{Prepara la documentación adicional}: Debes preparar la siguiente documentación adicional:
                * \textbf{Memoria justificativa}: Una memoria justificativa que describa las actividades realizadas y su relación con el proyecto.
                * \textbf{Documentos acreditativos}: Documentos acreditativos de la actividad realizada, como el certificado de asistencia.
                * \textbf{Facturas y justificantes de pago}: Facturas y justificantes de pago de los gastos realizados.
                5. \textbf{Presenta la solicitud de abono}: Debes presentar la solicitud de abono de la liquidación de gastos a través del \textbf{Registro Electrónico General}, dirigido a la Subdirección General de Gestión Económica de la Investigación. Asegúrate de adjuntar todos los documentos necesarios, incluyendo el \textbf{Formulario de liquidación de gastos}, la \textbf{Memoria justificativa}, los \textbf{Documentos acreditativos}, las \textbf{Facturas y justificantes de pago}, y cualquier otro documento requerido.
                6. \textbf{Verifica la documentación}: Verifica que la documentación presentada sea completa y correcta, y que se ajuste a los requisitos establecidos en la \textbf{Guía completa de liquidación de viajes} y en la \textbf{Miniguía para viajes por razón del servicio}.

                Recuerda que es importante seguir las instrucciones y requisitos establecidos en la \textbf{Guía completa de liquidación de viajes} y en la \textbf{Miniguía para viajes por razón del servicio} para asegurarte de que el trámite se realice de manera correcta y eficiente.
          \item \textbf{Contexto utilizado:}
                \begin{scriptsize}
                  FUENTE: guia-completa-liquidacion-de-viajes-v2.pdf | PÁGINA: 19
                  \# Subdirección de Gestión Económica de la Investigación

                  \#\# 4.3 EVIDENCIA DE LA REALIZACIÓN DE LA ACTIVIDAD (MOTIVO DEL VIAJE)

                  Los Organismos financiadores exigen prueba documental de que la actividad o evento ha tenido lugar para acreditar que el gasto es elegible.

                  A continuación, se señalan los más comunes según la naturaleza de la actividad:

                  - \textbf{Asistencia a un congreso, seminario, jornada o similar}: debe presentar certificado de asistencia emitido por la entidad organizadora del evento. La mera inscripción no acredita haber asistido.
                  - \textbf{Participación como ponente en un congreso, seminario, jornada o similar}: copia del programa donde se pueda comprobar su nombre y datos del evento.
                  - \textbf{Asistencia a cursos de formación}: certificado de asistencia y programa donde se especifique el número de horas o créditos ECTS del curso.
                  - \textbf{Reuniones científicas o de coordinación/seguimiento del proyecto}: La misma memoria justificativa puede servir a estos efectos si describe la temática tratada en la reunión y su relación con el proyecto.
                  - \textbf{Trabajos de campo}: memoria del itinerario y actividad desarrollada firmada por el investigador responsable del proyecto.
                  - \textbf{Estancias}: Documento expedido por el responsable del centro donde haya tenido lugar la misma con indicación de las actividades desarrolladas y/o personas con las que ha colaborado durante la estancia.

                  \#\# 4.4 FACTURAS ABONADAS POR EL INTERESADO, ASÍ COMO JUSTIFICANTES DE LOS PAGOS EFECTUADOS

                  Si ha abonado personalmente facturas o tickets, debe presentar los correspondientes originales. En el punto 5 de esta guía se le indica cómo debe proceder según se trate de un original en soporte electrónico o en soporte papel.

                  Recuerde que las facturas deben reflejar los datos fiscales de la persona que realiza el viaje, en ningún caso los de la US, además de incluir la fecha y descripción de los conceptos facturados (nombre de la persona que viaja, itinerario del trayecto, número de noches, precio por noche, etc.). Dado que los conceptos genéricos son motivo de reintegro, asegúrese de que las facturas no contienen errores. De lo contrario, el personal de la Subdirección de Gestión Económica de la Investigación le dará traslado de la incidencia para la subsanación oportuna, con el consiguiente retraso en la tramitación.

                  En los tickets/facturas simplificadas vigile que figura la fecha y hora de emisión.

                  Por cada factura, ticket, billete electrónico, u otro documento justificativo de gasto debe presentar el correspondiente justificante del pago efectuado personalmente, ya que está obligado a acreditar ante el Organismo financiador los gastos soportados, ya sea mediante tarjeta, transferencia o en efectivo.

                  En caso de “Pagos en efectivo”, debe presentar el recibo del proveedor firmado por éste, con identificación de los datos del sujeto deudor que efectúa el pago, es decir, del beneficiario de la liquidación.

                  En el caso concreto de gastos de locomoción, deberá presentar los siguientes documentos según el medio de transporte público utilizado, junto con el correspondiente justificante del pago materializado (movimiento o cargo bancario):

                  \#\#\# TREN / AVE:

                  - Factura
                  - Billete electrónico

                  VERSIÓN 2 OCTUBRE 2024    PÁGINA 19 | 21

                  ---

                  FUENTE: Miniguía viajes personal propio US.pdf | PÁGINA: 4
                  \# Miniguía para viajes por razón del servicio

                  \#\# Alojamiento en congresos u otras actividades con sede oficial

                  El alojamiento en uno de los hoteles propuestos por la organización se podrá reembolsar, aunque supere los límites generales. Si el gasto se imputa a una subvención externa, solo se podrá cargar hasta el límite fijado por la convocatoria. La diferencia, si la hubiera, podrá imputarse a una orgánica financiada con fondos propios. Será necesario adjuntar a la liquidación la documentación que acredite que el alojamiento pertenece a la oferta oficial del evento.

                  Si se prevé algún supuesto de carácter extraordinario, se requiere solicitar justificadamente autorización extraordinaria previa de Gerencia mediante este formulario: https://forms.office.com/e/KiNDfV0YF2.

                  - La solicitud de autorización extraordinaria no implica en ningún caso la aceptación automática de las condiciones expresadas por el comisionado o comisionada.
                  - Debe presentarse siempre con carácter previo al viaje e incluir todos los datos requeridos: destino, fechas, motivación, orgánica, presupuesto comparado de agencias (si aplica) y demás extremos relevantes, así como documentación justificativa que acredite la situación extraordinaria o no contemplada por la Instrucción.
                  - La revisión se realizará, atendiendo a criterios de legalidad, proporcionalidad y principio de economía.
                  - Las solicitudes presentadas a posteriori serán susceptibles de denegación, conforme al artículo 9.2 de la Instrucción.

                  \#\# Durante el viaje: justificantes y certificado de asistencia

                  Guarda todos los justificantes de gasto y pago, que serán necesarios para su liquidación posterior. Solicita certificado de asistencia o similar en todos los casos.

                  ----

                  Gerencia
                  Universidad de Sevilla

                  ---

                  FUENTE: guia-breve-liquidacion.pdf | PÁGINA: 2
                  \# Subdirección de Gestión Económica de la Investigación

                  \#\# AL TERMINAR LA ACTIVIDAD

                  \#\#\# 1. LIQUIDACIÓN DE GASTOS

                  Rellene el Formulario (INCLUIR EL ENLACE DE LA WEB) de liquidación de gastos de viaje. En la Guía (INCLUIR EL ENLACE DE LA WEB) para la cumplimentación de solicitudes de abono encontrará la información detallada para la correcta cumplimentación de este formulario.

                  Una vez terminado, debe ser firmado a través de Portafirmas, tanto por el IP como por la persona que realiza la actividad.

                  Podrá renunciar parcialmente a los gastos de la liquidación tal y como se le especifica en detalle en el punto 3.4 de la guía completa.

                  \#\#\# 2. DOCUMENTACIÓN ADICIONAL A PRESENTAR

                  Además de la Comisión de Servicio/Autorización de desplazamiento y del formulario de liquidación de gastos será necesario presentar la siguiente documentación:

                  - Una memoria justificativa en la que describa las actividades realizadas indicando la relación de éstas con alguno/s de los objetivos del proyecto. También deberá especificar el vínculo con el proyecto de la persona que realiza el viaje, así como su situación laboral. Puede consultar todos los detalles necesarios en el punto 4.2 de la guía completa.

                  - Documentación acreditativa de:
                  - La actividad realizada (certificado de asistencia, documento firmado por responsable del centro donde haya tenido lugar la estancia, etc.). Ver detalles en punto 4.3 de la guía completa.
                  - Los gastos de alojamiento y/o locomoción efectuados por la persona que ha realizado la actividad (facturas, tickets, billetes electrónicos, etc.), y trazabilidad bancaria de los mismos. Ver detalles en punto 4.4 de la guía completa.
                  - Si procede, informe de conformidad firmado por el IP de las facturas de alojamiento y/o locomoción derivadas de reservas de crédito a través de agencias de viaje del acuerdo marco vigente.
                  - Si procede, documento de desglose de la renuncia parcial de los gastos de liquidación en pdf.

                  VERSIÓN 2, OCTUBRE 2024    PÁGINA 2 | 4

                  ---

                  FUENTE: guia-breve-liquidacion.pdf | PÁGINA: 1
                  \# Subdirección de Gestión Económica de la Investigación

                  \#\# GUÍA BREVE
                  \#\#\# Solicitudes de liquidación de gastos de viajes

                  \#\#\# Introducción

                  En este documento se presenta una versión reducida de la Guía para la cumplimentación de solicitudes de abono de liquidaciones de gastos de viaje en el ámbito de la investigación, con el fin de facilitar los trámites necesarios para la cumplimentación del formulario de liquidación de gastos y de especificar la documentación necesaria para la realización y justificación de la actividad, así como el procedimiento de presentación de la solicitud de abono.

                  \#\#\# ANTES DE REALIZAR LA ACTIVIDAD

                  - Deberá tramitar la Comisión de Servicios con su licencia por estudios a personal docente (PDI funcionario o laboral) o Autorización de Desplazamiento (Personal Investigador o Técnico Contratado). Puede consultar detalles en el punto 4.1 de la guía completa.

                  - Deberá comprobar que la actividad a desarrollar figura como un gasto elegible en la convocatoria del proyecto con el que pretende justificarla, que la persona que la va a llevar a cabo guarda el vínculo que exige el organismo financiador, y que dispone de saldo suficiente en la orgánica del proyecto.

                  - Si va a alojarse haciendo uso de compañías que ofrecen servicios de reservas online, asegúrese de que podrá obtener factura por el gasto total realizado.

                  - Si va a reservar el medio de transporte y/o el alojamiento a través de agencias de viaje del acuerdo marco vigente, solicite la reserva de crédito correspondiente.

                  \#\#\# DURANTE LA ACTIVIDAD

                  - Conserve tarjetas de embarque si el viaje se ha realizado en avión.
                  - Pida certificados o documentos que acrediten su participación en la actividad.
                  - Pida facturas o tickets del alojamiento y/o medios de locomoción utilizados donde figuren sus datos fiscales, no los de la US. En facturas/tickets simplificados asegúrese de que aparecen la fecha y hora de emisión.

                  ---

                  FUENTE: guia-breve-liquidacion.pdf | PÁGINA: 3
                  \# Subdirección de Gestión Económica de la Investigación

                  \#\# 3. SOLICITUD DEL ABONO DE LA LIQUIDACIÓN DE GASTOS

                  A continuación, deberá proceder a la solicitud de abono de la liquidación de los gastos realizados mediante instancia efectuada a través del Registro Electrónico General, dirigido a la Subdirección General de Gestión Económica de la Investigación.

                  \textbf{IMPORTANTE:} Cada entrada de Registro debe corresponder a una sola solicitud de abono, la cual solo puede incluir una sola liquidación de gastos y sus correspondientes documentos obligatorios. Estos documentos son:

                  - Comisión Servicio o Autorización Desplazamiento con todas las firmas requeridas.
                  - Formulario de liquidación de gastos con las firmas requeridas.
                  - Memoria justificativa de la actividad realizada.
                  - Documentos acreditativos de la actividad realizada.
                  - Facturas abonadas y la trazabilidad bancaria de las mismas.
                  - Si se va a realizar un pago a terceros deberá adjuntar también el informe de conformidad firmado por el IP.
                  - Si ha solicitado renuncia parcial a parte del gasto deberá adjuntar también el documento de desglose de la renuncia en pdf.

                  Dado que actualmente el Registro Electrónico General solo permite subir un máximo de cinco archivos, organice la documentación de tal modo que cada archivo corresponda a lo siguiente:

                  1. Comisión de servicios u autorización de desplazamiento
                  2. Formulario de liquidación de dietas
                  3. Memoria justificativa
                  4. Resto de documentos (agrupados en uno o dos archivos, sabiendo que cada archivo solo puede tener un tamaño máximo de 10 Mb, y el tamaño máximo de los cinco ficheros es de 15 Mb)

                  Si su documentación incluye originales en papel deberá enviarlos por correo interno a la Subdirección General de Gestión Económica, junto con una copia de la solicitud de registro presentada.

                  \#\# CONSULTAS

                  Además de la guía completa, puede consultar también el documento de preguntas frecuentes (INCLUIR EL ENLACE DE LA WEB) elaborado por el Vicerrectorado de Investigación o dirigirse al personal del Servicio de Fiscalización y Pagos:

                  - Proyectos Estatales y Fundaciones:
                  - Cristina Guerra Aguilar-Galindo, cristina-ag@us.es 954 48 13 86
                  - Responsable Proyectos Comunidad Autónoma, Europeos o Internacionales:
                  - Rosa María López Valiente, rosi@us.es, 954 48 12 85

                  \textbf{VERSIÓN 2, OCTUBRE 2024        PÁGINA 3 | 4}
                \end{scriptsize}
        \end{itemize}


  \item \textbf{P10.} Quiero que me reconozcan los créditos que hice en mis anteriores estudios, ¿cómo funciona exactamente el proceso? \\
        \begin{itemize}
          \item \textbf{Tipo de consulta:} multi
          \item \textbf{Respuesta del Agente:}

                ⚠️ Requisitos Previos:
                Antes de iniciar el proceso de reconocimiento de créditos, es importante que estés admitido y matriculado en los estudios para los que solicitas el reconocimiento de créditos. Además, debes tener en cuenta que el reconocimiento de créditos se basa en la adecuación entre las competencias y conocimientos adquiridos en las materias superadas y los previstos en el plan de estudios del título de destino.

                A continuación, te presento los pasos a seguir para solicitar el reconocimiento de créditos:

                1. \textbf{Reúne la documentación necesaria}: Debes aportar la documentación justificativa de los créditos obtenidos y su contenido académico, indicando los módulos, materias o asignaturas que consideras superados. La documentación necesaria incluye:
                * \textbf{Certificado académico} que acredite los estudios universitarios cursados y superados.
                * \textbf{Plan de estudios} del título universitario que has cursado anteriormente.
                * \textbf{Programas de las asignaturas} que has superado.
                2. \textbf{Presenta la solicitud de reconocimiento de créditos}: Debes presentar la solicitud de reconocimiento de créditos en la \textbf{Secretaría del Centro} responsable del título para el que solicitas el reconocimiento, dentro de los plazos indicados en el \textbf{calendario académico} de cada curso.
                3. \textbf{Rellena el }formulario de solicitud\textbf{}: Debes rellenar el \textbf{formulario de solicitud de reconocimiento de créditos}, que puedes encontrar en la \textbf{Secretaría Virtual} de la Universidad de Sevilla o en la \textbf{SEVIUS}, y adjuntar la documentación necesaria.
                4. \textbf{La Comisión de Reconocimiento de Créditos evalúa la solicitud}: La \textbf{Comisión de Reconocimiento de Créditos} del Centro evaluará tu solicitud y emitirá una propuesta de resolución. La Comisión puede solicitar informe a los \textbf{Departamentos Universitarios} responsables de las enseñanzas objeto de reconocimiento.
                5. \textbf{Resolución del Decano o Director del Centro}: El \textbf{Decano o Director del Centro} dictará resolución sobre tu solicitud de reconocimiento de créditos, previa propuesta de la \textbf{Comisión de Reconocimiento de Créditos}. La resolución debe dictarse y notificarse en un plazo máximo de tres meses desde la presentación de la solicitud.
                6. \textbf{Notificación de la resolución}: La resolución se notificará al interesado, y si es favorable, se reflejará en el \textbf{expediente académico} del estudiante.

                Recuerda que, en caso de desestimación, puedes interponer \textbf{recurso de alzada} ante el \textbf{Rector} de la Universidad de Sevilla.
          \item \textbf{Interacción Usuario:} Fueron en otra universidad española oficial y no acabé la carrera.
          \item \textbf{Respuesta Final:}

                ⚠️ Requisitos Previos:
                Antes de iniciar el proceso de reconocimiento de créditos, es importante que estés admitido y matriculado en los estudios para los que solicitas el reconocimiento de créditos en la Universidad de Sevilla. Además, debes tener en cuenta que el reconocimiento de créditos se basa en la adecuación entre las competencias y conocimientos adquiridos en las materias superadas y los previstos en el plan de estudios del título de destino.

                A continuación, te presento los pasos a seguir para solicitar el reconocimiento de créditos:

                1. \textbf{Reúne la documentación necesaria}: Debes aportar la documentación justificativa de los créditos obtenidos y su contenido académico, indicando los módulos, materias o asignaturas que consideras superados. La documentación necesaria incluye:
                * \textbf{Certificado académico} que acredite los estudios universitarios cursados y superados en la otra universidad española oficial.
                * \textbf{Plan de estudios} del título universitario que has cursado anteriormente.
                * \textbf{Programas de las asignaturas} que has superado.
                2. \textbf{Presenta la solicitud de reconocimiento de créditos}: Debes presentar la solicitud de reconocimiento de créditos en la \textbf{Secretaría del Centro} responsable del título para el que solicitas el reconocimiento, dentro de los plazos indicados en el \textbf{calendario académico} de cada curso.
                3. \textbf{Rellena el }formulario de solicitud\textbf{}: Debes rellenar el \textbf{formulario de solicitud de reconocimiento de créditos}, que puedes encontrar en la \textbf{Secretaría Virtual} de la Universidad de Sevilla o en la \textbf{SEVIUS}, y adjuntar la documentación necesaria.
                4. \textbf{La Comisión de Reconocimiento de Créditos evalúa la solicitud}: La \textbf{Comisión de Reconocimiento de Créditos} del Centro evaluará tu solicitud y emitirá una propuesta de resolución. La Comisión puede solicitar informe a los \textbf{Departamentos Universitarios} responsables de las enseñanzas objeto de reconocimiento.
                5. \textbf{Resolución del Decano o Director del Centro}: El \textbf{Decano o Director del Centro} dictará resolución sobre tu solicitud de reconocimiento de créditos, previa propuesta de la \textbf{Comisión de Reconocimiento de Créditos}. La resolución debe dictarse y notificarse en un plazo máximo de tres meses desde la presentación de la solicitud.
                6. \textbf{Notificación de la resolución}: La resolución se notificará al interesado, y si es favorable, se reflejará en el \textbf{expediente académico} del estudiante.

                Recuerda que, en caso de desestimación, puedes interponer \textbf{recurso de alzada} ante el \textbf{Rector} de la Universidad de Sevilla.
          \item \textbf{Contexto utilizado:}
                \begin{scriptsize}
                  FUENTE: Normativa Reconocimiento y Transferencia de creditos.pdf | PÁGINA: 7
                  \# Textos consolidados

                  \#\# CAPÍTULO III
                  \#\#\# Reconocimiento de créditos en enseñanzas oficiales de Máster

                  \#\#\#\# Artículo 11. A partir de otros títulos de Grado, Máster o Doctorado.
                  11.1. Las solicitudes de reconocimiento de créditos basadas en estudios superados en títulos universitarios de Grado, Máster -ya sean de Programas Oficiales de Postgrado regulados por el Real Decreto 56/2005 o de títulos de Master desarrollados al amparo del Real Decreto 1393/2007- o periodo de formación específico del Doctorado –Real Decreto 1393/2007 y, en su caso, los derivados del Real Decreto 99/2011- se resolverán teniendo en cuenta la adecuación entre las competencias y conocimientos adquiridos en las materias superadas y los previstos en el plan de estudios del título de destino, indicándose las materias o asignaturas que se considerarán superadas por el interesado y que, por lo tanto, no estarán obligados a cursar.

                  11.2. Para la resolución de estas solicitudes se tendrán en cuenta los siguientes criterios:
                  1. En el caso de solicitudes de reconocimiento de créditos entre títulos oficiales de Máster que habiliten para el ejercicio de profesiones reguladas, para los que el Gobierno haya establecido las condiciones a las que han de adecuarse los planes de estudios, se reconocerán los créditos de los módulos definidos en la correspondiente norma reguladora que hayan sido superados por el estudiante. En caso de no haberse superado íntegramente un determinado módulo, el reconocimiento se llevará a cabo por materias o asignaturas en función de las competencias y conocimientos asociados a las mismas.
                  2. En el caso de títulos de Máster en proceso de extinción por la implantación de nuevos planes de estudios, la adaptación de los estudiantes a estos últimos se basará en el reconocimiento de créditos previsto en la tabla de adaptación incluida en la correspondiente memoria de verificación del título de Máster en cuestión.
                  3. En los procesos de adaptación de estudiantes a los nuevos planes de los títulos de Máster deberá garantizarse que la situación académica de aquellos no resulte perjudicada. A tal efecto, las materias, asignaturas o créditos superados que no tengan equivalencia en las correspondientes al plan de estudios de destino se incorporarán en el expediente del estudiante como créditos genéricos de carácter optativo. Si, aún así resultaran excedentes, los créditos restantes se podrán incorporar al expediente como créditos transferidos, a petición del interesado y siempre que se trate de materias o asignaturas completas.
                  4. El número de créditos reconocidos con cargo a la optatividad no podrá superar el número de créditos optativos exigido por el plan de estudios del título de destino.
                  5. En ningún caso podrán ser objeto de reconocimiento los Trabajos Fin de Máster.

                  \#\#\#\# Artículo 12. A partir de títulos de la anterior ordenación universitaria.
                  12.1. Las solicitudes de reconocimiento de créditos basadas en estudios superados en títulos universitarios de la anterior ordenación universitaria, Arquitecto, Ingeniero, Licenciado o periodo de docencia del doctorado, se resolverán teniendo en cuenta la adecuación entre las competencias y conocimientos adquiridos en las materias superadas y los previstos en el plan de estudios del título de destino, indicándose las materias o asignaturas que se considerarán superadas por el interesado y que, por lo tanto, no estarán obligados a cursar.

                  12.2. En ningún caso podrán ser objeto de reconocimiento los Trabajos Fin de Máster.

                  Este documento es de carácter informativo y no tiene valor jurídico.

                  ---

                  FUENTE: Normativa Reconocimiento y Transferencia de creditos.pdf | PÁGINA: 12
                  \# Textos consolidados

                  \#\# Artículo 20. Órganos responsables.

                  \#\#\# 20.1.
                  En la forma que determine cada Centro se constituirá en cada uno de ellos una Comisión de Reconocimiento de Créditos en la que estarán representados los diferentes sectores de la comunidad universitaria, que estará presidida por el Decano o Director, o Vicedecano o Subdirector en quien delegue y de la que formarán parte el Secretario y el Responsable de Administración del Centro.

                  \#\#\# 20.2.
                  En el caso de la Escuela Internacional de Posgrado se constituirá una Comisión de Reconocimiento de Créditos específica para los títulos de Máster vinculados a la misma, con representación de los diferentes sectores de la comunidad universitaria y con la composición que determine la Dirección de la misma, de la que formará parte, en todo caso, la persona responsable de la administración de la Escuela Internacional de Posgrado.

                  \#\#\# 20.3.
                  Serán funciones de la Comisión de Reconocimiento de Créditos:
                  1. Analizar las solicitudes presentadas sobre reconocimiento de créditos a partir de estudios universitarios cursados, estudios superiores no universitarios o a partir de experiencia laboral o profesional acreditada por los interesados y realizar la propuesta de resolución correspondiente.
                  2. En los casos de estudios previos cursados, solicitar informe a los Departamentos Universitarios responsables de las enseñanzas objeto de reconocimiento sobre la adecuación entre las competencias y conocimientos adquiridos por el interesado y los previstos en el plan de estudios para el que se solicita el reconocimiento. Este informe deberá ser evacuado en el plazo máximo de 15 días y no tendrá carácter vinculante. De no emitirse en el plazo señalado se proseguirán las actuaciones de la Comisión.
                  3. Para la valoración de la experiencia laboral y profesional aportada por los interesados, la Comisión, tras el estudio de la documentación presentada, podrá acordar la realización de una evaluación de los conocimientos y capacidades de los solicitantes para determinar la adquisición de las competencias alegadas. Esta evaluación podrá consistir en entrevistas profesionales, pruebas de competencia, demostraciones prácticas en situaciones similares a las de los puestos desempeñados u otros medios similares y para su realización se podrá contar con la asistencia de especialistas de los Departamentos correspondientes.
                  4. En los supuestos en que puedan reconocerse automáticamente créditos obtenidos en otras titulaciones de Grado de la misma o distintas ramas de conocimiento, en titulaciones oficiales de Máster o en otros títulos de enseñanza superior, esta Comisión elaborará tablas de reconocimiento de créditos que serán públicas y que permitirán a los estudiantes conocer anticipadamente las asignaturas, materias o módulos que le serían automáticamente reconocidos ante una hipotética solicitud.
                  5. Emitir informes sobre los contenidos de los recursos administrativos que se interpongan ante el Rector contra las resoluciones de reconocimiento de créditos basadas en las solicitudes indicadas en el apartado 1 anterior.
                  6. Cualesquiera otras funciones que pudieran asignársele en las disposiciones de desarrollo de esta norma.

                  \#\#\# 20.4.
                  No será necesaria la intervención de la Comisión de Reconocimiento de Créditos y se aprobarán de oficio con carácter automático las solicitudes de reconocimiento de créditos que correspondan a alguno de los supuestos que conlleven el reconocimiento automático, así como las que se deriven del acuerdo de estudios firmado por el estudiante y el Centro con ocasión del disfrute de una plaza de movilidad de los programas “SICUE”, “Erasmus” o similares.

                  \#\#\# 20.5.
                  Corresponderá al Decano o Director del Centro correspondiente o a la persona responsable de la Escuela Internacional de Posgrado, en su caso, dictar resolución, previa propuesta de la Comisión de Reconocimiento de Créditos, salvo en el supuesto previsto en el apartado anterior. La resolución,

                  ---

                  FUENTE: Normativa Reconocimiento y Transferencia de creditos.pdf | PÁGINA: 13
                  \# Textos consolidados

                  que en caso desestimatorio debe ser motivada académicamente, deberá dictarse y notificarse en un plazo máximo de tres meses desde la presentación de la solicitud.

                  20.6. El vencimiento del plazo sin haberse notificado resolución expresa legitimará al interesado para entenderla desestimada por silencio administrativo.

                  20.7. Contra las resoluciones del Decano o Director del Centro se podrá interponer recurso de alzada ante el Rector de la Universidad de Sevilla.

                  20.8. A efecto de la tramitación del procedimiento se declaran inhábiles los periodos no lectivos previstos en el calendario académico de cada curso.

                  \#\# Artículo 21. Solicitudes de transferencia de créditos.

                  Los expedientes de transferencia de créditos se tramitarán a petición del interesado. A estos efectos, los estudiantes que se incorporen a un nuevo estudio, mediante escrito dirigido al Decano o Director del Centro y en los plazos que se establezcan en el calendario académico de cada curso, indicarán si han cursado anteriormente otros estudios universitarios oficiales sin haberlos finalizado, aportando, en caso de no tratarse de estudios de la Universidad de Sevilla, la documentación justificativa que proceda de entre la contemplada en el artículo 19.1.

                  \#\# Artículo 22. Efectos del reconocimiento y transferencia de créditos.

                  22.1. En el proceso de reconocimiento quedarán reflejados de forma explícita aquellos módulos, materias o asignaturas que no deberán ser cursadas por el estudiante. Se entenderá en este caso que dichos módulos, materias o asignaturas ya han sido superadas, no serán susceptibles de nueva evaluación y se reflejarán en el expediente del estudiante como, módulos, materias o asignaturas reconocidas, indicándose el origen del reconocimiento.

                  22.2. En todo caso, el reconocimiento de créditos se referirá, al menos, a unidades de matrícula completas, es decir, no se podrá realizar el reconocimiento parcial de una asignatura.

                  22.3. Cuando la resolución del procedimiento dé lugar al reconocimiento de créditos optativos, el número de créditos reconocidos se minorará del número de créditos optativos exigido por el correspondiente plan de estudios y se reflejará en el expediente del estudiante como créditos optativos reconocidos, indicándose el origen del reconocimiento. En todo caso, el número de créditos optativos reconocidos no podrá superar el número de créditos exigido por el plan de estudios en cuestión.

                  22.4. En los casos procedentes, tras el proceso de reconocimiento de créditos, se permitirá a los interesados la ampliación de su matrícula en los términos recogidos en las Normas de Matrícula de cada curso académico.

                  22.5. La calificación de las asignaturas o, en su caso, de los créditos superados como consecuencia de un proceso de reconocimiento será equivalente a la calificación de las materias o asignaturas que han dado origen a éste. En caso necesario, se realizará la media ponderada cuando varias materias o asignaturas conlleven el reconocimiento de una sola en la titulación de destino.

                  22.6. Las calificaciones se reflejarán en el expediente académico en los términos recogidos en el Real Decreto 1125/2003, de 5 de septiembre, por el que se establece el sistema europeo de créditos y el sistema de calificaciones en las titulaciones universitarias de carácter oficial y validez en todo el territorio nacional.

                  22.7. No obstante lo anterior, cuando en el expediente académico de origen sólo se haga referencia a las calificaciones cualitativas, se transformarán en calificaciones numéricas, teniendo en cuenta la siguiente tabla de equivalencias:

                  Este documento es de carácter informativo y no tiene valor jurídico.

                  ---

                  FUENTE: Normativa Reconocimiento y Transferencia de creditos.pdf | PÁGINA: 4
                  \# Textos consolidados

                  superado íntegramente un determinado módulo, el reconocimiento se llevará a cabo por materias o asignaturas en función de las competencias y conocimientos asociados a las mismas.

                  7) En el caso de títulos oficiales de Grado que habiliten para el ejercicio de profesiones reguladas en España, para los que el Gobierno haya establecido las condiciones a las que han de adecuarse los planes de estudios, se reconocerán los créditos de los módulos definidos en la correspondiente norma reguladora que hayan sido superados por el estudiante. En caso de no haberse superado íntegramente un determinado módulo, el reconocimiento se llevará a cabo por materias o asignaturas en función de las competencias y conocimientos asociados a las mismas.

                  \#\# Artículo 5. A partir de títulos de Máster Universitario.

                  5.1. Las solicitudes de reconocimiento de créditos basadas en estudios superados en títulos oficiales de Máster Universitario (tanto los regulados por el RD 56/2005, como por el RD 1393/2007) o periodo de formación específico del Doctorado se resolverán teniendo en cuenta la adecuación entre las competencias y conocimientos adquiridos en las materias superadas y los previstos en el plan de estudios del título de destino, indicándose las materias o asignaturas que se considerarán superadas por el interesado y que, por lo tanto, no estarán obligados a cursar.

                  5.2. En ningún caso podrán ser objeto de reconocimiento los Trabajos Fin de Grado.

                  \#\# Artículo 6. A partir de títulos de la anterior ordenación universitaria.

                  6.1. Las solicitudes de reconocimiento de créditos basadas en estudios superados en títulos universitarios de la anterior ordenación universitaria, Arquitecto Técnico, Diplomado, Ingeniero Técnico, Maestro, Licenciado, Arquitecto, Ingeniero o periodo de docencia del doctorado, se resolverán teniendo en cuenta la adecuación entre las competencias y conocimientos adquiridos en las materias superadas y los previstos en el plan de estudios del título de destino, indicándose las materias o asignaturas que se considerarán superadas por el interesado y que, por lo tanto, no estarán obligados a cursar.

                  6.2. Para la resolución de estas solicitudes se tendrán en cuenta los siguientes criterios:

                  1) Cuando las competencias y conocimientos no estén explicitados o no puedan deducirse del plan de estudios de origen del estudiante se tomarán como referencia el número de créditos y/o los contenidos de las materias o asignaturas cursadas.

                  2) En el caso de títulos en proceso de extinción por la implantación de los nuevos títulos de Grado, la adaptación de los estudiantes a estos últimos se basará en el reconocimiento de créditos previsto en la tabla de adaptación incluida en la correspondiente memoria de verificación del título de Grado en cuestión.

                  3) En los procesos de adaptación de estudiantes a los nuevos planes de los títulos de Grado deberá garantizarse que la situación académica de aquellos no resulte perjudicada. A tal efecto, las materias, asignaturas o créditos superados que no tengan equivalencia en las correspondientes al plan de estudios de Grado se incorporarán en el expediente del estudiante como créditos genéricos de carácter optativo. Si, aún así resultaran excedentes, los créditos restantes se podrán incorporar al expediente como créditos transferidos, a petición del interesado y siempre que se trate de materias o asignaturas completas.

                  4) El número de créditos reconocidos con cargo a la optatividad no podrá superar el número de créditos optativos exigido por el plan de estudios del título de destino.

                  5) En ningún caso podrán ser objeto de reconocimiento los Trabajos Fin de Grado.

                  Este documento es de carácter informativo y no tiene valor jurídico.

                  ---

                  FUENTE: Normativa Reconocimiento y Transferencia de creditos.pdf | PÁGINA: 10
                  \# Textos consolidados

                  15.7. De los acuerdos de estudios que se establezcan se enviará copia a los Servicios Centrales del Rectorado que corresponda.
                  15.8. Con carácter general lo dispuesto en estas normas será de aplicación a la movilidad para dobles titulaciones sin perjuicio de las previsiones contenidas en los convenios respectivos.
                  15.9. Resultarán igualmente de aplicación las normas que eventualmente se aprueben por los órganos nacionales o internacionales competentes para cada programa específico de movilidad.

                  \#\# CAPÍTULO V
                  \#\#\# Transferencia de créditos
                  \textbf{Artículo 16. Definición.}
                  La transferencia de créditos implica que, en los documentos académicos oficiales acreditativos de las enseñanzas seguidas por cada estudiante, se incluirán la totalidad de los créditos obtenidos en enseñanzas universitarias oficiales cursadas con anterioridad, en ésta u otra universidad, que no hayan conducido a la finalización de sus estudios con la consiguiente obtención de un título oficial.

                  \textbf{Artículo 17. Aplicación.}
                  Los créditos correspondientes a materias o asignaturas previamente superadas por el estudiante, en enseñanzas universitarias no concluidas y que no puedan ser objeto de reconocimiento, serán transferidos a su expediente en los estudios a los que ha accedido con la calificación de origen y se reflejarán en los documentos académicos oficiales acreditativos de los estudios seguidos por el mismo, así como en el Suplemento Europeo al Título.

                  \#\# CAPÍTULO VI
                  \#\#\# Tramitación
                  \textbf{Artículo 18. Solicitudes de reconocimiento de créditos.}
                  18.1. Los expedientes de reconocimiento de créditos se tramitarán a solicitud del interesado, quién deberá aportar la documentación justificativa de los créditos obtenidos y su contenido académico, indicando los módulos, materias o asignaturas que considere superados.
                  18.2. Será requisito imprescindible estar admitido y matriculado en los correspondientes estudios, salvo en los casos de cambios de estudios oficiales de Grado, según lo dispuesto en la Resolución Rectoral por la que se regula la admisión a los títulos de Grado de la Universidad de Sevilla de los estudiantes que han iniciado anteriormente otros estudios universitarios.
                  18.3. Las solicitudes de reconocimiento de créditos tendrán su origen en materias o asignaturas realmente cursadas y superadas, en ningún caso se referirán a materias o asignaturas previamente reconocidas, convalidadas o adaptadas.
                  18.4. Las solicitudes se presentarán en la Secretaría del Centro responsable del título para el que se solicita el reconocimiento dentro de los plazos indicados en el calendario académico de cada curso.

                  \textbf{Artículo 19. Documentación acreditativa.}
                  19.1. En el caso de estudios universitarios cursados, estudios superiores no universitarios u otros estudios no oficiales, se aportará la siguiente documentación:

                  Este documento es de carácter informativo y no tiene valor jurídico.
                  10
                \end{scriptsize}
        \end{itemize}


  \item \textbf{P11.} Tengo un par de artículos científicos clave, ¿cómo se valora esto puntualmente en el Baremo de Contratación? \\
        \begin{itemize}
          \item \textbf{Tipo de consulta:} multi
          \item \textbf{Respuesta del Agente:}

                \textbf{Requisitos mínimos excluyentes:}
                * No se establecen requisitos mínimos excluyentes específicos para la valoración de artículos científicos en el Baremo de Contratación.

                \textbf{Criterios de evaluación para artículos científicos:}

                \vspace{0.4cm}
                \begin{table}[h!]\centering\renewcommand{\arraystretch}{1.3}
                  \begin{tabular}{|p{0.45\textwidth}|p{0.45\textwidth}|}\hline
                    Criterio                                                   & Detalles                                                                                                                                                                         \\\hline
                    Publicación en revistas de prestigio                       & Se valoran como revistas de prestigio aquellas que acreditan ISSN y están localizadas en posiciones preferentes de Journal Citation Reports (JCR) y SCImago Citation Ranks (SCR) \\\hline
                    Citas validadas en bases de datos de reconocido prestigio  & 0,1 puntos por cada cita acreditada en bases de datos como SCOPUS, WOS o Google Académico                                                                                        \\\hline
                    Posición del artículo en los diferentes índices de impacto & La comisión utiliza la mejor posición para realizar la baremación, considerando listados como JCR, SJR o Anexo I                                                                 \\\hline
                    Año de publicación del artículo                            & La edición de la clasificación de referencia para la baremación es la del año de publicación del artículo, o en su defecto, la inmediata anterior                                \\\hline
                  \end{tabular}\end{table}
                \vspace{0.4cm}

                \textbf{Desglose de la puntuación:}
                * La puntuación exacta para cada artículo científico dependerá de factores como la posición del artículo en los índices de impacto, el número de citas validadas y la calidad de la revista en la que se publicó.
                * La comisión de contratación utiliza la aplicación Bareprof para comprobar automáticamente la posición del artículo en los diferentes índices de impacto considerados en el baremo.
                * El candidato debe introducir correctamente el ISSN de la revista y el año de publicación del artículo para que el proceso de baremación sea efectivo.

                \textbf{Nota importante:}
                * La comisión de contratación puede modificar el ISSN en caso de que se compruebe que el identificador introducido por el candidato es incorrecto.
                * La documentación acreditativa de los méritos alegados debe ser proporcionada por el candidato para que la comisión pueda validar los méritos y asignar la puntuación correspondiente.
          \item \textbf{Interacción Usuario:} Están publicados en revistas del Q1 de JCR.
          \item \textbf{Respuesta Final:}

                Los artículos científicos clave que tienes están publicados en revistas del Q1 de Journal Citation Reports (JCR). Según el contexto proporcionado, se valoran como revistas de prestigio aquellas que acreditan ISSN y están localizadas en posiciones preferentes de Journal Citation Reports (JCR) y SCImago Citation Ranks (SCR).

                En cuanto a la valoración puntual en el Baremo de Contratación, se consideran los siguientes criterios:

                * Publicación en revistas de prestigio
                * Citas validadas en bases de datos de reconocido prestigio (0,1 puntos por cada cita acreditada)
                * Posición del artículo en los diferentes índices de impacto
                * Año de publicación del artículo

                La comisión de contratación utiliza la aplicación Bareprof para comprobar automáticamente la posición del artículo en los diferentes índices de impacto considerados en el baremo. Para que este proceso sea efectivo, debes introducir correctamente tanto el ISSN de la revista como el año de publicación del artículo.

                No se establecen requisitos mínimos excluyentes específicos para la valoración de artículos científicos en el Baremo de Contratación. La puntuación exacta para cada artículo científico dependerá de factores como la posición del artículo en los índices de impacto, el número de citas validadas y la calidad de la revista en la que se publicó.

                Es importante mencionar que la documentación acreditativa de los méritos alegados debe ser proporcionada para que la comisión pueda validar los méritos y asignar la puntuación correspondiente. La comisión de contratación puede modificar el ISSN en caso de que se compruebe que el identificador introducido por el candidato es incorrecto.

                En cuanto a los plazos y fechas, no se especifican fechas concretas en el contexto proporcionado. Sin embargo, se menciona que la edición de la clasificación de referencia para la baremación es la del año de publicación del artículo, o en su defecto, la inmediata anterior. También se menciona que la Revisión 4 del documento de ayuda al candidato es de \textbf{marzo de 2025}, y la Revisión 3 es de \textbf{diciembre de 2024}.
          \item \textbf{Contexto utilizado:}
                \begin{scriptsize}
                  FUENTE: 27831\_instrucciones-baremacinin-r03.pdf | PÁGINA: 2
                  4.2. RENDIMIENTO ACADÉMICO (BLOQUE I DEL BAREMO) ..........................................................................................18
                  4.2.1. Expediente Académico (epígrafe 1.1 del baremo) ..........................................................................................18
                  4.2.2. Título de doctor (epígrafe 1.2 del baremo) ...................................................................................................21
                  4.2.3. Otros títulos académicos (epígrafe 1.3 del baremo) ......................................................................................22
                  4.2.4. Premios académicos y menciones de calidad (epígrafe 1.4 del baremo) ........................................................22
                  4.3. ACTIVIDAD DOCENTE (BLOQUE II DEL BAREMO) ................................................................................................23
                  4.3.1. Docencia en asignaturas regladas universitarias (epígrafe 2.1 del baremo) ..................................................23
                  4.3.2. Contratos como PDI universitario (epígrafe 2.2 del baremo) .......................................................................23
                  4.3.3. Dirección de tesis doctorales (epígrafe 2.3 del baremo) ..............................................................................24
                  4.3.4. Otros méritos de docencia universitaria (epígrafe 2.4 del baremo) .................................................................24
                  4.4. ACTIVIDAD INVESTIGADORA (BLOQUE III DEL BAREMO) ....................................................................................24
                  4.4.1. Contratos y becas en procesos competitivos (epígrafe 3.1 del baremo) ........................................................25
                  4.4.2. Libros (o monografías) científicas (epígrafe 3.2 del baremo) .......................................................................26
                  4.4.3. Capítulos de Libro (epígrafe 3.3 del baremo) ...............................................................................................27
                  4.4.4. Artículos científicos de Revista (epígrafe 3.4 del baremo) ..........................................................................27
                  4.4.5. Participación en proyectos de investigación (epígrafe 3.5 del baremo) ..........................................................27
                  4.4.6. Trabajos presentados en congresos (epígrafe 3.6 del baremo) ...................................................................28
                  4.4.7. Comités científicos, técnicos o asesores (epígrafe 3.7 del baremo) ...............................................................28
                  4.4.8. Conocimiento de idiomas relevantes (epígrafe 3.8 del baremo) ..................................................................28
                  4.5. TRANSFERENCIA DEL CONOCIMIENTO (BLOQUE IV DEL BAREMO) ....................................................................29
                  4.5.1. Contratos de transferencia (epígrafe 4.1 del baremo) ..................................................................................29
                  4.5.2. Obra civil o de edificación, o exposiciones (epígrafe 4.2 del baremo) ............................................................29
                  4.5.3. Patentes (epígrafe 4.3 del baremo) ...............................................................................................................29
                  4.5.4. Spin-off (epígrafe 4.4 del baremo) ................................................................................................................29

                  ---

                  FUENTE: BAREMO 2024.pdf | PÁGINA: 21
                  \textbf{Complemento (aditivo, y hasta un máximo de 5 puntos):}

                  \vspace{0.4cm}
                  \begin{table}[h!]\centering\renewcommand{\arraystretch}{1.3}
                    \begin{tabular}{|p{0.45\textwidth}|p{0.45\textwidth}|}\hline
                      Descripción                                                                                      & Puntos por cita acreditada \\\hline
                      Citas validadas en base de datos de reconocido prestigio (e.g., SCOPUS, WOS o Google Académico). & 0,1 por cita acreditada    \\\hline
                    \end{tabular}\end{table}
                  \vspace{0.4cm}

                  \textbf{Nota 1:} Las aportaciones presentadas deberán haber sido publicadas en los años para los que se solicita la evaluación. No se tendrán en cuenta en ningún caso publicaciones aceptadas pero no publicadas a excepción de aquellas publicaciones que acrediten oficialmente el identificador de publicación DOI, HANDLE, PMID u equivalente y aceptado por la comunidad científica y validado por la comisión de contratación. Se valoran como revistas de prestigio (tanto de suscripción como de acceso abierto e híbridas, y tanto en formato papel como digital) las revistas que acreditando ISSN estén localizadas en posiciones preferentes de Journal Citation Reports (JCR) y SCImago Citation Ranks (SCR), sin distinción entre ambos catálogos. La edición de la clasificación de referencia para la baremación es la del año de publicación del artículo, o en su defecto la inmediata anterior (en caso de no haberse aún publicado la nueva edición). Para artículos aceptados y no publicados, y sin perjuicio de lo señalado anteriormente en relación con los artículos aceptados pero no publicados, será la última edición publicada con fecha fin del plazo de la convocatoria de concurso.

                  ---

                  FUENTE: 27831\_instrucciones-baremacinin-r03.pdf | PÁGINA: 28
                  \# Bareprof

                  \#\# Instrucciones de baremación

                  \#\#\# 4.4.3 Capítulos de Libro (epígrafe 3.3 del baremo)

                  - No se consideran capítulos de libros las contribuciones a congresos editadas en libros de actas o proceedings.
                  - Además de los capítulos de libros propiamente dichos, el baremo incluye en este apartado fichas técnicas de catálogos en publicaciones, reseñas y aportaciones en diccionarios comunes o científicos, con una valoración del 10\% de la que corresponde al capítulo de libro. Si la aportación corresponde a uno de los tipos citados, pero el candidato ha indicado como tipo de producción “Capítulo de libro”, la Comisión no debe validar el mérito.
                  - El editor de CVN no incorpora ninguna opción para indicar como tipo de producción las fichas técnicas de catálogos en publicaciones. En la Guía del Ayuda al Candidato se indica que, en tales casos, debe seleccionar “Reseña” en el campo Tipo de Producción. También sería válido si seleccionaran “Diccionario científico” o “Diccionario común”, pues la valoración es la misma, pero no así “Capítulo de libro”; en este último caso, el mérito no debe validarse.
                  - Para la asignación de la afinidad al mérito, consulte las consideraciones realizadas en la introducción al apartado 4.4.

                  \#\#\# 4.4.4 Artículos científicos de Revista (epígrafe 3.4 del baremo)

                  - Los datos de citas de un artículo científico deben ser introducidos por el candidato. La aplicación Bareprof no comprueba tales datos de citas.
                  - El baremo no establece distinción entre las fuentes de citas, valorando el mayor número de citas validado por la comisión.
                  - La aplicación Bareprof comprueba automáticamente la posición del artículo en los distintos índices de impacto considerados en el baremo. Para que este proceso sea efectivo, el candidato ha debido introducir correctamente tanto el ISSN de la revista como el año de publicación del artículo. La comisión está facultada para modificar el ISSN en caso de que se compruebe que el identificador introducido por el candidato es incorrecto.
                  - En pantalla se mostrará la posición de la revista en los diferentes listados automatizados (JCR, SJR o Anexo I). Para cada índice, se utilizará el listado que corresponde al año de publicación del artículo o, en su defecto, al año inmediatamente anterior si el primero no se encuentra aún disponible a la finalización del periodo de admisión de solicitudes. Bareprof usará la mejor posición para realizar la baremación.
                  - Para la asignación de la afinidad, consulte las consideraciones realizadas en la introducción al apartado 4.4.

                  \#\#\# 4.4.5 Participación en proyectos de investigación (epígrafe 3.5 del baremo)

                  - La pertenencia a grupos de investigación no es un mérito considerado en el Baremo. Muchos candidatos incluyen estos méritos como proyectos de investigación, que no deben ser validados por la Comisión.
                  - Los méritos englobados en este epígrafe del baremo no admiten coeficiente de afinidad, por lo que la Comisión NO puede invalidarlos alegando falta de afinidad con la plaza. Si el mérito ha sido consignado en el apartado correcto y se encuentra acreditado documentalmente, la Comisión debe validarlo. Por este mismo motivo, no es preciso que el candidato justifique la afinidad del.

                  ----

                  Revisión 3 (diciembre, 2024) 27

                  ---

                  FUENTE: 27831\_instrucciones-baremacinin-r03.pdf | PÁGINA: 2
                  1. INTRODUCCIÓN .........................................................................................................................................................3
                  2. DOCUMENTOS DE REFERENCIA .................................................................................................................................3
                  3. DESCRIPCIÓN DEL PROCESO DE BAREMACIÓN ..........................................................................................................4
                  3.1. COMPOSICIÓN DE LAS COMISIONES DE CONTRATACIÓN ..........................................................................................4
                  3.2. CONSTITUCIÓN DE LA COMISIÓN DE CONTRATACIÓN..........................................................................................5
                  3.3. BAREMACIÓN .........................................................................................................................................................6
                  3.3.1. Documentación acreditativa de los méritos alegados........................................................................................7
                  3.3.2. Proceso de baremación ...................................................................................................................................9
                  3.3.2.1. Validación del mérito................................................................................................................................11
                  3.3.2.2. Determinación de la afinidad del mérito ................................................................................................11
                  3.3.2.3. Edición del mérito......................................................................................................................................12
                  3.3.3. Generación del acta .......................................................................................................................................14
                  3.3.4. Periodo de reclamaciones ................................................................................................................................15
                  3.3.5. En caso de duda...............................................................................................................................................15
                  4. CONSIDERACIONES SOBRE LA BAREMACIÓN ...........................................................................................................15
                  4.1. CONSIDERACIONES DE CARÁCTER GENERAL ........................................................................................................15
                  4.1.1. Documentación acreditativa de los méritos....................................................................................................15
                  4.1.2. Duplicidad de méritos......................................................................................................................................16
                  4.1.3. Reubicación de méritos ...................................................................................................................................17
                  4.1.4. Entidades públicas y entidades privadas .........................................................................................................17
                  4.1.5. Ramas de conocimiento y áreas de conocimiento ..........................................................................................17
                  4.1.6. Documentación acreditativa en idioma distinto al castellano........................................................................17
                  4.2. RENDIMIENTO ACADÉMICO (BLOQUE I DEL BAREMO) ..........................................................................................18
                  4.2.1. Expediente Académico (epígrafe 1.1 del baremo) ..........................................................................................18
                  4.2.2. Título de doctor (epígrafe 1.2 del baremo) ...................................................................................................21

                  ---

                  FUENTE: 27831\_instrucciones-baremacinin-r03.pdf | PÁGINA: 32
                  \# Instrucciones de baremación

                  Por tanto, la comisión debe cumplimentar numéricamente los campos Seminarios, Tesis, Libro y Artículo, atendiendo a lo indicado por el candidato en el campo Resultados relevantes. Por ejemplo, si el candidato señala en Resultados relevantes que ha escrito un artículo (e identifica cuál y con quién lo ha escrito), en el campo Artículo la comisión debe incluir el valor numérico 1.

                  \vspace{0.4cm}
                  \begin{table}[h!]\centering\renewcommand{\arraystretch}{1.3}
                    \begin{tabular}{|p{0.45\textwidth}|p{0.45\textwidth}|p{0.45\textwidth}|p{0.45\textwidth}|}\hline
                      Seminarios & Tesis & Libro & Artículo \\\hline
                      4          &       &       & 3        \\\hline
                    \end{tabular}\end{table}
                  \vspace{0.4cm}

                  Solo se valora el material firmado con algún miembro del equipo de investigación de la Universidad o del Centro de Investigación donde se realiza la estancia, o la dirección (o codirección) de tesis doctorales a doctorandos de la Universidad o Centro de Investigación donde se realiza la estancia y durante la realización de esta. La documentación aportada por el candidato deberá acreditarlo. Por otra parte, los méritos indicados por el candidato en el campo Resultados relevantes para la valoración de la estancia de investigación, pueden haber sido incluidos en otros apartados del CVN, sin que por ello se considere que existe duplicidad de méritos.

                  La comisión debe tener presente que los campos Seminarios, Tesis, Libro y Artículo solo admiten valores numéricos.

                  - El campo Resultados relevantes no es de cumplimentación obligatoria, por lo que es posible que aparezca en blanco en la baremación, ya sea por olvido del candidato o por la inexistencia de méritos computables. En el primer caso, si la documentación acreditativa aportada por el candidato pone de manifiesto la existencia de méritos aplicables en este apartado, la comisión podrá cumplimentar a partir de ellos los campos Seminarios, Tesis, Libro y Artículo, aplicando el mismo criterio a todos los candidatos.

                  - La instrucción de 17-10-19 sobre los requisitos para el reconocimiento de las estancias de investigación dentro del mérito preferente establece que la duración mínima de la estancia es de 1 mes. Bareprof no considerará las estancias de duración inferior a esta, sin que sea necesario que la comisión las invalide.

                  - No se considera que exista duplicidad de méritos cuando el candidato ha disfrutado de un contrato de investigación que conlleva la realización de la investigación, parcial o totalmente, en una Universidad o Centro de Investigación distinto de aquel que ha dado lugar al contrato (por ejemplo, convocatorias Marie Sklowdoska Curie o Margarita Salas). En tales casos, el candidato podrá hacer constar el contrato en el epígrafe 3.1 del baremo y la estancia de investigación como mérito preferente, siempre que se reúnan los requisitos exigidos en tales apartados.

                  - En aquellos casos en que la estancia no esté sufragada con fondos propios, la Comisión deberá comprobar que el programa indicado por el candidato está destinado a financiar estancias con finalidad investigadora y no de otro tipo (como sucede con el programa Erasmus+, destinado a promover estancias formativas o de docencia).

                  Revisión 3 (diciembre, 2024) 31
                \end{scriptsize}
        \end{itemize}


  \item \textbf{P12.} He leído que tengo derecho a una exención para no pagar los precios públicos de la matrícula, ¿me la pueden tramitar? \\
        \begin{itemize}
          \item \textbf{Tipo de consulta:} multi
          \item \textbf{Respuesta del Agente:}

                Según el artículo 28 de la normativa de matrícula de la Universidad de Sevilla, los estudiantes pueden tener derecho a una exención o bonificación en la matrícula si cumplen con ciertas condiciones establecidas en el decreto por el que se determinan los precios públicos de las universidades públicas de Andalucía.

                Para tramitar la exención, es necesario que el estudiante marque todas las posibles exenciones a las que pueda tener derecho en el momento de la formalización de la matrícula. Sin embargo, hay algunas exenciones que se aplican de forma automática, como la aplicación de matrícula de honor obtenida en estudios universitarios oficiales cursados en la Universidad de Sevilla y la bonificación del 99\% por créditos aprobados en primera matrícula.

                En cuanto a la documentación necesaria para acreditar la exención, el estudiante deberá aportar la documentación correspondiente en el Buzón de Documentación (DOMUS) habilitado en la Secretaría Virtual de la Universidad de Sevilla (SEVIUS) en el plazo de diez días hábiles a contar desde el día siguiente a la finalización del periodo en el que haya formalizado su matrícula.

                Es importante destacar que la Universidad de Sevilla puede consultar telemáticamente la condición de familia numerosa, la condición de persona con discapacidad o la condición de víctima de actos de terrorismo del estudiante que alegue dicha condición en su matrícula, salvo que se oponga a ello.

                En resumen, para tramitar la exención para no pagar los precios públicos de la matrícula, es necesario que el estudiante:

                1. Marque las posibles exenciones a las que pueda tener derecho en el momento de la formalización de la matrícula.
                2. Aporte la documentación correspondiente en el Buzón de Documentación (DOMUS) habilitado en la Secretaría Virtual de la Universidad de Sevilla (SEVIUS) en el plazo establecido.
                3. Permita la consulta telemática de la condición que alegue, salvo que se oponga a ello.

                Es recomendable que el estudiante consulte con la secretaría del centro o con la Universidad de Sevilla para obtener más información sobre el procedimiento y los requisitos necesarios para tramitar la exención.
          \item \textbf{Interacción Usuario:} Tengo reconocida una discapacidad del 35\%.
          \item \textbf{Respuesta Final:}

                ⚠️ Requisitos Previos:
                Antes de iniciar el trámite, asegúrate de tener reconocida tu discapacidad del 35\% y de tener acceso a la \textbf{Secretaría Virtual de la Universidad de Sevilla (SEVIUS)} y al \textbf{Buzón de Documentación (DOMUS)}.

                Para tramitar la exención para no pagar los precios públicos de la matrícula en la Universidad de Sevilla, sigue los siguientes pasos:

                1. \textbf{Marca las posibles exenciones}: En el momento de la formalización de la matrícula, marca la exención correspondiente a tu condición de discapacidad del 35\%. Esto te permitirá iniciar el proceso de tramitación de la exención.
                2. \textbf{Aporta la documentación correspondiente}: Debes aportar la documentación que acredite tu condición de discapacidad del 35\% en el \textbf{Buzón de Documentación (DOMUS)} habilitado en la \textbf{Secretaría Virtual de la Universidad de Sevilla (SEVIUS)}. La documentación debe ser aportada en el plazo de 10 días hábiles a contar desde el día siguiente a la finalización del periodo en el que hayas formalizado tu matrícula.
                3. \textbf{Permite la consulta telemática}: La Universidad de Sevilla puede consultar telemáticamente tu condición de discapacidad, salvo que te opongas a ello. Si no te opones, la Universidad de Sevilla verificará tu condición de discapacidad de forma automática.
                4. \textbf{Verificación de la documentación}: La secretaría del centro verificará la documentación aportada para comprobar que cumples con los requisitos para la exención. Si la documentación es correcta, se aplicará la exención correspondiente.
                5. \textbf{Notificación de la resolución}: La Universidad de Sevilla te notificará la resolución sobre la exención. Si la exención es aprobada, no tendrás que pagar los precios públicos de la matrícula.

                Recuerda que, si te opones a la consulta telemática, debes acreditar documentalmente tu condición de discapacidad y seguir el procedimiento establecido para ello. Puedes consultar más información en los enlaces proporcionados por la Universidad de Sevilla, como [https://osi.us.es](https://osi.us.es) y [https://osi.us.es/sites/osi/files/doc/c.i.gestionestudiantesoficialesModif1.12.5.20CT.pdf](https://osi.us.es/sites/osi/files/doc/c.i.gestionestudiantesoficialesModif1.12.5.20CT.pdf).
          \item \textbf{Contexto utilizado:}
                \begin{scriptsize}
                  FUENTE: R.R. NORMAS DE MATRICULA GRADO Y MASTER UNIVERSITARIO 2025-2026.pdf | PÁGINA: 23
                  \# UNIVERSIDAD D SEVII

                  \#\# C)EST

                  \#\#\# Artículo 28. Reducción de precios. Exenciones y bonificaciones

                  Los estudiantes no abonarán en parte o en su totalidad los precios públicos por servicios académicos y/o administrativos, en su caso, cuando en ellos concurra/n alguna/s circunstancia/s que genere/n un derecho a la exención o bonificación en la matrícula, según lo establecido en el decreto por el que se determinan los precios públicos de las universidades públicas de Andalucía.

                  El estudiante deberá marcar todas las posibles exenciones a las que pudiera tener derecho en el momento de la formalización de la matrícula, excepto la aplicación de matrícula de honor obtenida en estudios universitarios oficiales cursados en la Universidad de Sevilla y la bonificación del 99\% por créditos aprobados en primera matrícula, al ser automáticas. Si el estudiante obtiene el reconocimiento de persona con discapacidad con posterioridad al período de su matrícula, pero con efectos de vigencia que incluya dicho periodo, aunque no marcara dicha exención en el momento de formalizar su matrícula, se le aplicará la exención correspondiente.

                  Cuando el estudiante haya marcado en automatrícula la exención de ser beneficiario de la condición de familia numerosa, de tener la condición de persona con discapacidad o de ser beneficiario de derechos por ser víctima de actos de terrorismo, la Universidad de Sevilla lo consultará telemáticamente, salvo que se opusiera a ello. De no poderse verificar telemáticamente alguna de estas condiciones, el estudiante deberá aportar la documentación acreditativa en el Buzón de Documentación (DOMUS) habilitado en la Secretaría Virtual de la Universidad de Sevilla (SEVIUS).

                  La documentación acreditativa del resto de las exenciones alegadas, en su caso, será aportada igualmente en el Buzón de Documentación (DOMUS) habilitado en la Secretaría Virtual de la Universidad de Sevilla (SEVIUS) en el plazo de diez días hábiles a contar desde el día siguiente a la finalización del periodo en el que haya formalizado su matrícula, para su verificación por la secretaría del centro.

                  Las exenciones y bonificaciones son las siguientes:

                  1. Estudiantes becarios de la convocatoria general del Ministerio de Educación, Formación Profesional y Deportes (con cargo a los Presupuestos Generales del Estado). La beca de matrícula sólo cubrirá los créditos en primera matrícula de los estudiantes becarios. Abonarán íntegramente los precios de los créditos o asignaturas en segunda o sucesivas matrículas y los precios públicos por servicios administrativos y el seguro escolar. Los estudiantes que soliciten esta beca podrán formalizar su matrícula sin el previo pago de los precios públicos por servicios académicos de los que se matriculen por primera vez. No obstante, las secretarías de los centros podrán requerir cautelarmente el abono de dichos precios públicos a quienes previsiblemente no cumplan los requisitos establecidos en la convocatoria de estas becas.

                  2. Los estudiantes becados por el Gobierno Vasco no estarán obligados a abonar el importe de la contraprestación por servicios académicos que sea cubierto por la beca.

                  3. Aplicación de matrícula de honor, premios y medallas:
                  a) La matrícula de honor en segundo curso de Bachillerato o Ciclo Formativo de Formación Profesional de Grado Superior, la obtención de Premio Extraordinario de Bachillerato o de Formación Profesional de Grado Superior o los premios de las Olimpiadas de Matemáticas, Física o Química u otras olimpiadas acreditadas de ámbito nacional o internacional, otorgará el derecho a exención total de los precios por servicios académicos, por una sola vez, en el primer curso de enseñanzas de grado.

                  Código Seguro De Verificación: `eqOjhZcmL0p6VJlZiTH/Yw==`
                  Fecha: 27/06/2025
                  Firmado Por: MIGUEL ANGEL CASTRO ARROYO
                  Url De Verificación: [https://pfirma.us.es/verifirma/code/eqOjhZcmL0p6VJlZiTH\%2FYw\%3D\%3D](https://pfirma.us.es/verifirma/code/eqOjhZcmL0p6VJlZiTH\%2FYw\%3D\%3D)
                  Página: 19/50

                  ---

                  FUENTE: R.R. NORMAS DE MATRICULA GRADO Y MASTER UNIVERSITARIO 2025-2026.pdf | PÁGINA: 24
                  \# UNIVERSIDAD D SEVII

                  \#\# C)EST

                  Se acreditará mediante certificado donde conste la matrícula de honor o con documento acreditativo de la obtención del premio extraordinario o la Medalla en las Olimpiadas, según corresponda.

                  Se aplicará a los precios públicos por la prestación de servicios académicos y solo se abonarán los precios públicos por servicios administrativos y el seguro escolar.

                  \#\#\# b) Matrícula de honor en estudios oficiales universitarios:

                  Los estudiantes que hayan obtenido matrícula de honor en alguna/s asignatura/s en el último curso universitario realizado tendrán derecho a la exención correspondiente, por una sola vez, en la siguiente matrícula de la misma titulación.

                  Se aplicará también a las siguientes situaciones:

                  1. En el caso de adaptación a un plan de estudios, se aplicará el derecho a la exención correspondiente en el nuevo plan de estudios matriculado. Solo se tendrán en cuenta las matrículas de honor obtenidas en el último curso académico del plan de origen.

                  2. Acceso a estudios de máster universitario si la matrícula de honor se ha obtenido en los estudios que le han dado acceso a dicho máster.

                  3. En caso de traslado de un centro adscrito a uno propio para cursar la misma titulación.

                  4. Estudiantes procedentes de otra universidad. Si las matrículas de honor obtenidas corresponden a asignaturas cursadas en la misma titulación en otra universidad, serán computadas a estos efectos, en los términos previstos en este apartado.

                  En los supuestos previstos en los epígrafes anteriores, conforme a lo dispuesto en el decreto por el que se determinan los precios públicos de las universidades públicas de Andalucía, la bonificación en el importe de la matrícula será equivalente al número de créditos en que obtuvo dicha matrícula de honor, multiplicado por el precio establecido para el crédito en primera matrícula de los estudios en que se matricule, determinado según sus circunstancias personales y familiares.

                  No se podrá aplicar el derecho a esta exención en el caso de que la calificación de matrícula de honor sea consecuencia de convalidación de asignaturas o reconocimiento de créditos.

                  Se deberán abonar íntegramente los precios públicos por servicios administrativos y el seguro escolar.

                  En el supuesto de no acreditar esta exención, previamente requerida, se eliminará de su matrícula mediante resolución dictada y notificada al efecto.

                  \#\# 4. Beneficiarios del título de familia numerosa:

                  a) Si fuera de categoría general, abonarán el 50\% de todos los precios públicos. El seguro escolar, en su caso, lo abonarán íntegramente.

                  b) Si fuera de categoría especial, tendrán derecho a la exención total de todos los precios públicos. Solo abonarán el seguro escolar.

                  ----

                  Código Seguro De Verificación: `eqOjhZcmL0p6VJlZiTH/Yw==`
                  Fecha: `27/06/2025`
                  Firmado Por: `MIGUEL ANGEL CASTRO ARROYO`
                  Url De Verificación: [https://pfirma.us.es/verifirma/code/eqOjhZcmL0p6VJlZiTH\%2FYw\%3D\%3D](https://pfirma.us.es/verifirma/code/eqOjhZcmL0p6VJlZiTH\%2FYw\%3D\%3D)
                  Página: `20/50`

                  ---

                  FUENTE: R.R. NORMAS DE MATRICULA GRADO Y MASTER UNIVERSITARIO 2025-2026.pdf | PÁGINA: 25
                  \# UNIVERSIDAD DE SEVILLA

                  \#\# Exención por familia numerosa

                  El estudiante tendrá derecho a la aplicación de la exención por familia numerosa cuando esté en posesión del título de familia numerosa con vigencia en el periodo de matrícula de que se trate. También tendrá derecho a la exención si está en trámite su renovación o expedición en dicho periodo.

                  Los beneficios se extenderán durante todo el curso académico, pudiéndose aplicar a los servicios académicos y administrativos recogidos en el anexo del decreto por el que se determinan los precios públicos de las universidades públicas de Andalucía. Una vez acreditada la condición de familia numerosa, no se exigirá una nueva acreditación para la expedición, entre otros, de certificados y/o títulos durante el curso académico.

                  Aquellos estudiantes que en el periodo que les corresponda de matrícula no tuvieran la condición de familia numerosa, pero la obtienen con posterioridad, podrán aplicarla para nuevas actuaciones administrativas que impliquen el pago de precios públicos, como pueden ser, entre otros, certificaciones académicas o el abono de los derechos de expedición del título, previa acreditación. No tendrán derecho a aplicarla en la matrícula ordinaria al no tener dicha condición en ese periodo.

                  Conforme lo establecido en el artículo 28.2 de la ley 39/2015, de 1 de octubre, del Procedimiento Administrativo Común de las Administraciones Públicas, la Universidad de Sevilla consultará telemáticamente el título de familia numerosa del estudiante que alegue dicha condición en su matrícula, salvo que se opusiera a ello.

                  De no poderse verificar telemáticamente la condición de familia numerosa cuando realice su automatrícula, el estudiante deberá aportar en el Buzón de Documentación (DOMUS) habilitado en la Secretaría Virtual de la Universidad de Sevilla (SEVIUS) el título de familia numerosa vigente en el periodo ordinario de matrícula o en el que le haya sido asignado si es posterior a este o el resguardo de haber solicitado su expedición o renovación en dicho periodo. En este último supuesto el título expedido o renovado deberá aportarse antes del 31 de diciembre y tener efectos de vigencia en el periodo en el que haya realizado su matrícula. En caso de ser beneficiario del título de familia numerosa de categoría general, será necesario presentar el carné de familia numerosa junto con la resolución de reconocimiento de dicha condición en la que conste el número de miembros cuando en el carné no aparezca la relación completa de beneficiarios del título.

                  También podrá acreditarse mediante certificado, según modelo normalizado, emitido por la Consejería competente en el que se acredite que, de acuerdo con lo establecido en la disposición transitoria quinta de la ley 26/2015, el estudiante tenía la condición de familia numerosa en fecha 1 de enero de 2015.

                  Las secretarías de los centros podrán hacer uso de la comprobación telemática, con independencia de la aportación de la documentación por el estudiante.

                  En todo caso, el centro deberá realizar las comprobaciones telemáticas a aquellos estudiantes a los que no haya sido posible la comprobación en la automatrícula y no hayan aportado en el Buzón de Documentación (DOMUS) la documentación exigida, antes de requerir la documentación acreditativa de la condición de familia numerosa.

                  En el supuesto de no acreditar dicha condición, previamente requerida, se eliminará de su matrícula mediante resolución dictada y notificada al efecto.

                  En caso de oposición a la consulta telemática, el estudiante deberá acreditar documentalmente su condición de familia numerosa.

                  ----

                  \textbf{Código Seguro De Verificación:} eqOjhZcmL0p6VJlZiTH/Yw==
                  \textbf{Fecha:} 27/06/2025
                  \textbf{Firmado Por:} MIGUEL ANGEL CASTRO ARROYO
                  \textbf{Url De Verificación:} Verificación
                  \textbf{Página:} 21/50

                  ---

                  FUENTE: R.R. NORMAS DE MATRICULA GRADO Y MASTER UNIVERSITARIO 2025-2026.pdf | PÁGINA: 27
                  \# UNIVERSIDAD DE SEVILLA

                  \#\# C)EST

                  \#\#\# 6. Víctimas de actos de terrorismo
                  El estudiante podrá aplicar la exención por tener la condición de beneficiario de derechos por ser víctima de actos de terrorismo cuando tenga esta condición en el periodo de matrícula de que se trate.

                  Tendrán derecho a la exención total de precios públicos por la prestación de servicios académicos y administrativos en los estudios conducentes a la obtención de un título universitario oficial en primera y sucesivas matrículas. En cualquier caso, deberán abonar el seguro escolar, si procede.

                  Conforme lo establecido en el artículo 28.2 de la ley 39/2015, de 1 de octubre, del Procedimiento Administrativo Común de las Administraciones Públicas, la Universidad de Sevilla consultará telemáticamente la condición de beneficiario de derechos por ser víctima de actos de terrorismo del estudiante que la alegue en su matrícula, salvo que se opusiera a ello.

                  De no poderse verificar telemáticamente la condición de beneficiario de derechos por ser víctima de actos de terrorismo cuando realice su automatrícula, el estudiante deberá aportar en el Buzón de Documentación (DOMUS) habilitado en la Secretaría Virtual de la Universidad de Sevilla (SEVIUS) la documentación que acredite documentalmente dicha condición, mediante informe preceptivo del Ministerio del Interior, o bien mediante el reconocimiento por parte de la Administración General del Estado de pensión extraordinaria por actos de terrorismo, o por sentencia judicial firme, según se establece en el artículo 4 bis 3, de la Ley 32/1999, de 8 de octubre, de solidaridad con las víctimas de actos de terrorismo. En caso de no disponer de la documentación citada, se acompañarán los documentos que se consideren justificativos de reunir las condiciones señaladas. En este caso, el centro aceptará la matrícula condicional y remitirá el expediente al Vicerrectorado de Estudiantes para su estudio y resolución.

                  Las secretarías de los centros podrán hacer uso de la comprobación telemática, con independencia de la aportación de la documentación por el estudiante.

                  En todo caso, el centro deberá realizar las comprobaciones telemáticas a aquellos estudiantes a los que no haya sido posible la comprobación en la automatrícula y no hayan aportado en el Buzón de Documentación (DOMUS) la documentación exigida, antes de requerir la documentación acreditativa de la condición de víctima de actos de terrorismo.

                  En el supuesto de no acreditar dicha condición, previamente requerida, se eliminará de su matrícula mediante resolución dictada y notificada al efecto.

                  En caso de oposición a la consulta telemática, el estudiante deberá acreditar documentalmente su condición de víctima de actos de terrorismo.

                  El estudiante podrá ejercer el derecho de oposición a la consulta indicando las circunstancias que motivan dicha oposición, según los procedimientos establecidos para ello que podrá consultar en:

                  - https://osi.us.es
                  - https://osi.us.es/sites/osi/files/doc/c.i.gestionestudiantesoficialesModif1.12.5.20CT.pdf

                  En este supuesto, no se aplicará la exención de víctima de actos de terrorismo hasta tanto la secretaría del centro verifique la documentación aportada por el interesado en el acto de oposición a la consulta telemática.

                  \#\#\# 7. Obtención de Premio Extraordinario Fin de Grado o de Enseñanzas Artísticas Superiores
                  Otorgará la exención total de los precios públicos por servicios académicos, por una sola vez, en el primer curso de enseñanzas conducentes a la obtención de otro título universitario oficial.

                  ----

                  \textbf{Código Seguro De Verificación:} eqOjhZcmL0p6VJlZiTH/Yw==
                  \textbf{Fecha:} 27/06/2025
                  \textbf{Firmado Por:} MIGUEL ANGEL CASTRO ARROYO
                  \textbf{Url De Verificación:} [https://pfirma.us.es/verifirma/code/eqOjhZcmL0p6VJlZiTH\%2FYw\%3D\%3D](https://pfirma.us.es/verifirma/code/eqOjhZcmL0p6VJlZiTH\%2FYw\%3D\%3D)
                  \textbf{Página:} 23/50

                  ---

                  FUENTE: R.R. NORMAS DE MATRICULA GRADO Y MASTER UNIVERSITARIO 2025-2026.pdf | PÁGINA: 28
                  \# UNIVERSIDAD D SEVII

                  C)EST

                  En el caso de estudiantes procedentes de otras universidades o centros de enseñanzas artísticas superiores, la obtención del Premio Extraordinario se acreditará mediante la aportación del acta de concesión del premio y declaración responsable de no haber aplicado con anterioridad los beneficios del premio obtenido en otros estudios universitarios oficiales.

                  En el supuesto de no acreditar esta exención, previamente requerida, se eliminará de su matrícula mediante resolución dictada y notificada al efecto.

                  \#\# 8. Víctimas de violencia de género:
                  (Ley Orgánica 1/2004, de 28 de diciembre, de Medidas de Protección Integral contra la Violencia de Género y Ley 13/2007, de 26 de noviembre, de Medidas de Prevención y Protección Integral contra la Violencia de Género).

                  Estarán exentos de efectuar el pago de los precios públicos por servicios administrativos y académicos en primera y sucesivas matrículas. Deben abonar el seguro escolar, si procede.

                  La acreditación de dicha circunstancia podrá alegarse a lo largo del curso académico, presentando cualquiera de los documentos que se relacionan en el artículo 30 de la Ley 13/2007, de 26 de noviembre, de medidas de prevención y protección integral contra la violencia de género.

                  En el supuesto de no acreditar esta exención, previamente requerida, se eliminará de su matrícula mediante resolución dictada y notificada al efecto.

                  La exención será extensible a los hijos descendientes menores de 25 años y dependientes económicamente, aportando también el libro de familia.

                  \#\# 9. Personas beneficiarias del ingreso mínimo vital o de la renta mínima de inserción social en Andalucía:
                  Están exentas del pago de los precios públicos por servicios académicos universitarios para la realización de estudios conducentes a la obtención de títulos de carácter oficial las personas beneficiarias de la prestación del ingreso mínimo vital, de conformidad con lo previsto en la disposición adicional séptima de la Ley 19/2021, de 20 de diciembre, por la que se establece el ingreso mínimo vital. Además del reconocimiento de la exención anterior, también estarán exentos los beneficiarios de la renta mínima de inserción social en Andalucía.

                  Las personas que se acojan a cualquiera de las dos exenciones previstas en este apartado deberán acreditar dicha condición en el periodo de matrícula de que se trate.

                  En el supuesto de no acreditar esta exención, previamente requerida, se eliminará de su matrícula mediante resolución dictada y notificada al efecto.

                  \#\# 10. Bonificación por créditos superados en primera matrícula:
                  El alumnado de centros propios y adscritos de las universidades públicas andaluzas, en los términos regulados por el decreto por el que se determinan los precios públicos de las universidades públicas de Andalucía, tendrá derecho a una bonificación equivalente al 99\% del precio de los créditos superados en primera matrícula para sus estudios de grado y máster universitario.

                  \#\#\# 10.1. Requisitos para la bonificación:
                  a) En los estudios de grado: Haber estado matriculado en centros propios o adscritos de las universidades públicas andaluzas durante el curso 2024-2025, tomándose en consideración los créditos correspondientes a asignaturas superadas en primera matrícula en dicho curso.

                  b) En la primera matriculación de estudios de máster universitario: Haber estado matriculado en centros propios o adscritos de las universidades públicas andaluzas o en enseñanzas.

                  ----

                  Código Seguro De Verificación: `eqOjhZcmL0p6VJlZiTH/Yw==`
                  Fecha: `27/06/2025`
                  Firmado Por: `MIGUEL ANGEL CASTRO ARROYO`
                  Url De Verificación: [https://pfirma.us.es/verifirma/code/eqOjhZcmL0p6VJlZiTH\%2FYw\%3D\%3D](https://pfirma.us.es/verifirma/code/eqOjhZcmL0p6VJlZiTH\%2FYw\%3D\%3D)
                  Página: `24/50`
                \end{scriptsize}
        \end{itemize}


  \item \textbf{P13.} Llevo dando algunas clases universitarias durante varios cuatrimestres, ¿cómo se computa esto exactamente en el baremo? \\
        \begin{itemize}
          \item \textbf{Tipo de consulta:} multi
          \item \textbf{Respuesta del Agente:}

                \textbf{Requisitos Mínimos Excluyentes:}
                * Experiencia docente durante un mínimo de 4 años de dedicación a tiempo completo, o su equivalente a tiempo parcial, para figuras de contratados predoctorales y postdoctorales.
                * Un mínimo de 12 ECTS o 120 horas impartidas para figuras de contratados predoctorales y postdoctorales.
                * Un mínimo de 24 ECTS o 240 horas impartidas para el resto de figuras de profesorado.
                * Experiencia docente mínima de 10 años de dedicación a tiempo completo, o su equivalente a tiempo parcial, y un mínimo de 140 ECTS para Catedrático de Universidad (CU).

                \textbf{Criterios de Evaluación:}
                \vspace{0.4cm}
                \begin{table}[h!]\centering\renewcommand{\arraystretch}{1.3}
                  \begin{tabular}{|p{0.45\textwidth}|p{0.45\textwidth}|}\hline
                    Criterio                                                & Detalles                                                                                                                                                             \\\hline
                    Experiencia Docente                                     & Se valorará la experiencia docente en función de la dedicación a tiempo completo o parcial, y el número de horas o ECTS impartidos.                                  \\\hline
                    Dedicación Docente                                      & Se considerará dedicación a tiempo completo cuando se alcancen las 6 horas de clase semanal (180 horas de capacidad docente anual).                                  \\\hline
                    Pluralidad, Interdisciplinariedad y Complejidad Docente & Se valorará la complejidad de la docencia en función del número y diversidad de asignaturas impartidas, y la impartición de docencia en distintos idiomas o centros. \\\hline
                    Recursos Educativos                                     & Se valorarán los materiales docentes y recursos educativos disponibles en acceso abierto y sin restricciones.                                                        \\\hline
                    Actividades de Formación a lo Largo de la Vida          & Se valorarán las actividades de formación y actualización docente.                                                                                                   \\\hline
                    Calidad de la Actividad Docente e Innovación            & Se valorará la calidad de la docencia y la innovación en la enseñanza.                                                                                               \\\hline
                    Tutorización Docente                                    & Se valorará la tutorización reglada y no reglada.                                                                                                                    \\\hline
                  \end{tabular}\end{table}
                \vspace{0.4cm}

                \textbf{Puntuación:}
                La puntuación se calculará en función de los criterios de evaluación, con un máximo de 100 puntos. El baremo se divide en varios apartados, incluyendo:

                * Experiencia Docente (25-45 puntos)
                + Dedicación Docente (25-35 puntos)
                + Pluralidad, Interdisciplinariedad y Complejidad Docente (0-6 puntos)
                + Recursos Educativos (0-5 puntos)
                + Actividades de Formación a lo Largo de la Vida (0-5 puntos)
                * Calidad de la Actividad Docente e Innovación (5-40 puntos)
                + Calidad de la Actividad Docente (5-15 puntos)
                + Proyectos de Innovación Docente (0-20 puntos)
                + Formación para la Mejora Docente Recibida (0-5 puntos)
                + Formación para la Mejora Docente Impartida (0-5 puntos)
                * Tutorización Docente (0-15 puntos)
                + Tutorización Reglada (0-15 puntos)
                + Otros Tipos de Tutorización (0-15 puntos)
                * Otros Méritos (0-10 puntos)

                \textbf{Cálculo de Puntos:}
                El cálculo de puntos se realizará en función de la experiencia docente y la dedicación a tiempo completo o parcial. Por ejemplo, si se ha impartido docencia durante 4 años a tiempo completo, se considerará un total de 24 ECTS o 240 horas impartidas, lo que supondría un máximo de 35 puntos en el apartado de Dedicación Docente.
          \item \textbf{Interacción Usuario:} Fueron más de 1200 horas en asignaturas de Máster Universitario Oficial.
          \item \textbf{Respuesta Final:}

                \textbf{Requisitos Mínimos Excluyentes:}
                * Experiencia docente durante un mínimo de 4 años de dedicación a tiempo completo, o su equivalente a tiempo parcial, para figuras de contratados predoctorales y postdoctorales.
                * Un mínimo de 12 ECTS o 120 horas impartidas para figuras de contratados predoctorales y postdoctorales.
                * Un mínimo de 24 ECTS o 240 horas impartidas para el resto de figuras de profesorado.
                * Experiencia docente mínima de 10 años de dedicación a tiempo completo, o su equivalente a tiempo parcial, y un mínimo de 140 ECTS para Catedrático de Universidad (CU).

                \textbf{Criterios de Evaluación:}
                \vspace{0.4cm}
                \begin{table}[h!]\centering\renewcommand{\arraystretch}{1.3}
                  \begin{tabular}{|p{0.45\textwidth}|p{0.45\textwidth}|}\hline
                    Criterio                                                & Detalles                                                                                                                                                             \\\hline
                    Experiencia Docente                                     & Se valorará la experiencia docente en función de la dedicación a tiempo completo o parcial, y el número de horas o ECTS impartidos.                                  \\\hline
                    Dedicación Docente                                      & Se considerará dedicación a tiempo completo cuando se alcancen las 6 horas de clase semanal (180 horas de capacidad docente anual).                                  \\\hline
                    Pluralidad, Interdisciplinariedad y Complejidad Docente & Se valorará la complejidad de la docencia en función del número y diversidad de asignaturas impartidas, y la impartición de docencia en distintos idiomas o centros. \\\hline
                    Recursos Educativos                                     & Se valorarán los materiales docentes y recursos educativos disponibles en acceso abierto y sin restricciones.                                                        \\\hline
                    Actividades de Formación a lo Largo de la Vida          & Se valorarán las actividades de formación y actualización docente.                                                                                                   \\\hline
                    Calidad de la Actividad Docente e Innovación            & Se valorará la calidad de la docencia y la innovación en la enseñanza.                                                                                               \\\hline
                    Tutorización Docente                                    & Se valorará la tutorización reglada y no reglada.                                                                                                                    \\\hline
                  \end{tabular}\end{table}
                \vspace{0.4cm}

                \textbf{Puntuación:}
                La puntuación se calculará en función de los criterios de evaluación, con un máximo de 100 puntos. El baremo se divide en varios apartados, incluyendo:

                * Experiencia Docente (25-45 puntos)
                + Dedicación Docente (25-35 puntos)
                + Pluralidad, Interdisciplinariedad y Complejidad Docente (0-6 puntos)
                + Recursos Educativos (0-5 puntos)
                + Actividades de Formación a lo Largo de la Vida (0-5 puntos)
                * Calidad de la Actividad Docente e Innovación (5-40 puntos)
                + Calidad de la Actividad Docente (5-15 puntos)
                + Proyectos de Innovación Docente (0-20 puntos)
                + Formación para la Mejora Docente Recibida (0-5 puntos)
                + Formación para la Mejora Docente Impartida (0-5 puntos)
                * Tutorización Docente (0-15 puntos)
                + Tutorización Reglada (0-15 puntos)
                + Otros Tipos de Tutorización (0-15 puntos)
                * Otros Méritos (0-10 puntos)

                \textbf{Cálculo de Puntos:}
                El cálculo de puntos se realizará en función de la experiencia docente y la dedicación a tiempo completo o parcial. Por ejemplo, si se ha impartido docencia durante 4 años a tiempo completo, se considerará un total de 24 ECTS o 240 horas impartidas, lo que supondría un máximo de 35 puntos en el apartado de Dedicación Docente.

                \textbf{Ejemplo de Cálculo de Puntos para más de 1200 horas de Docencia en Máster Universitario Oficial:}

                * Se considerará un máximo de 1200 horas de docencia oficial.
                * La puntuación se calculará en función del número de horas impartidas, con un máximo de 50 puntos.
                * El complemento multiplicativo por docencia en Máster Universitario Oficial es de 1,50.
                * Si se han impartido más de 1200 horas, se aplicará un factor de escala para calcular la puntuación final.

                Por ejemplo, si se han impartido 1500 horas de docencia en Máster Universitario Oficial, la puntuación sería:

                * 1500 horas x 1,00 (coeficiente por hora de docencia) = 150 puntos
                * Factor de escala: 1200/1500 = 0,8
                * Puntuación final: 150 x 0,8 = 120 puntos

                Sin embargo, como el máximo de puntos para este apartado es de 50, la puntuación final sería de 50 puntos.

                Es importante tener en cuenta que la puntuación final dependerá de la evaluación de la comisión y del baremo establecido por la Universidad de Sevilla. Es recomendable consultar el baremo y las instrucciones de evaluación para obtener más información sobre el cálculo de puntos y los criterios de evaluación.
          \item \textbf{Contexto utilizado:}
                \begin{scriptsize}
                  FUENTE: CRITERIOS DE EVALUACIÓN Y REQUISITOS MÍNIMOS.pdf | PÁGINA: 10
                  \# CATEDRÁTICA O CATEDRÁTICO DE UNIVERSIDAD (CU)

                  \vspace{0.4cm}
                  \begin{table}[h!]\centering\renewcommand{\arraystretch}{1.3}
                    \begin{tabular}{|p{0.45\textwidth}|p{0.45\textwidth}|}\hline
                      Mínimo & Máximo \\\hline
                      50     & 100    \\\hline
                      puntos & puntos \\\hline
                    \end{tabular}\end{table}
                  \vspace{0.4cm}

                  \#\# 2. ACTIVIDAD DOCENTE

                  \vspace{0.4cm}
                  \begin{table}[h!]\centering\renewcommand{\arraystretch}{1.3}
                    \begin{tabular}{|p{0.45\textwidth}|p{0.45\textwidth}|p{0.45\textwidth}|}\hline
                      Actividad                                                           & Mínimo & Máximo \\\hline
                      2.1. Experiencia docente                                            & 25     & 45     \\\hline
                      2.1.1. Dedicación docente                                           & 25     & 35     \\\hline
                      2.1.2. Pluralidad, interdisciplinariedad y complejidad docente      & 0      & 6      \\\hline
                      2.1.3. Recursos educativos                                          & 0      & 5      \\\hline
                      2.1.4. Actividades de formación a lo largo de la vida               & 0      & 5      \\\hline
                      2.2. Calidad de la actividad docente e innovación                   & 5      & 40     \\\hline
                      2.2.1. Calidad de la actividad docente                              & 5      & 15     \\\hline
                      2.2.2. Proyectos de innovación docente                              & 0      & 20     \\\hline
                      2.2.3. Formación para la mejora docente recibida                    & 0      & 5      \\\hline
                      2.2.4. Formación para la mejora docente impartida                   & 0      & 15     \\\hline
                      2.3. Tutorización docente                                           & 5      & 15     \\\hline
                      2.3.1. Tutorización reglada                                         & 0      & 15     \\\hline
                      2.3.2. Otros tipos de tutorización (no incluida en encargo docente) & 0      & 15     \\\hline
                      2.4. Otros méritos*                                                 & 0      & 10     \\\hline
                    \end{tabular}\end{table}
                  \vspace{0.4cm}

                  * Los puntos obtenidos en este apartado, de carácter opcional, se suman a los de los apartados anteriores, pero en ningún caso podrá superarse el máximo (100 puntos).

                  \#\# 2.1. EXPERIENCIA DOCENTE

                  \#\#\# 2.1.1. Dedicación docente

                  En el caso de PTU deberá acreditarse experiencia docente durante un mínimo de 4 años de dedicación a tiempo completo, o su equivalente a tiempo parcial, y, además:

                  - En el caso de figuras de contratados predoctorales y postdoctorales, incluyendo el periodo de orientación postdoctoral, un mínimo de 12 ECTS o 120 horas impartidas. El disfrute de un contrato postdoctoral se considerará dedicación a tiempo completo. El disfrute de una ayuda predoctoral que tenga asignada capacidad docente tutelada y se cubra totalmente se considerará en todos los casos dedicación a tiempo parcial.

                  - Para el resto de las figuras de profesorado, un mínimo de 24 ECTS o 240 horas impartidas.

                  - En caso de que la experiencia docente se haya acumulado combinando distintas modalidades (esto es: predoctoral/postdoctoral y resto de figuras de profesorado), el mínimo de ECTS u horas se calculará proporcionalmente al periodo de tiempo de duración de cada una de las figuras contractuales disfrutadas.

                  En el caso de CU deberá acreditarse una experiencia docente mínima de 10 años de dedicación a tiempo completo, o su equivalente a tiempo parcial, y un mínimo de 140 ECTS.

                  DIRECCIÓN DE ANECA – 20/03/2024

                  ---

                  FUENTE: 27831\_instrucciones-baremacinin-r03.pdf | PÁGINA: 24
                  \# Bareprof
                  \#\# Instrucciones de baremación

                  \#\#\# 4.3 Actividad docente (bloque II del baremo)

                  \#\#\#\# 4.3.1 Docencia en asignaturas regladas universitarias (epígrafe 2.1 del baremo)

                  - En este apartado el candidato reflejará la docencia realmente impartida (no la nominal de su contrato), acreditada mediante el certificado de actividades docentes o documento equivalente.
                  - La docencia debe introducirse obligatoriamente en horas. Si la comisión observa que el candidato ha introducido el número de créditos, podrá modificarlo para reflejar el valor correcto.
                  - La docencia recogida en este apartado podrá corresponder tanto a la actividad del candidato como PDI (epígrafe 2.2 del baremo) como a su actividad docente desarrollada como Personal Investigador (epígrafe 3.1 del baremo).
                  - Las Prácticas Externas se validan como actividad docente del candidato siempre que este actúe como tutor académico; el tutor profesional ejerce la labor en calidad de responsable de la entidad colaboradora y no de PDI. Su cómputo se corresponde con lo reflejado en el certificado de actividades docentes. Por ejemplo, la Universidad de Sevilla reconoce al tutor académico 0,1 créditos por cada 6 créditos de que conste la asignatura en el Plan de Estudios.
                  - Se valorará también en este apartado la dirección o tutela de trabajos académicos que aparezca recogida y cuantificada (en número de horas o créditos) en el certificado de actividades docentes, con independencia de que la labor de tutorización sea también valorada en el apartado 2.4.1 del baremo (dirección o tutela de trabajos académicos en titulaciones universitarias oficiales). Esto no supone duplicidad en el mérito, pues se valora, por un lado, la actividad docente presencial reconocida y, por otro, la propia labor de tutorización.

                  \#\#\#\# 4.3.2 Contratos como PDI universitario (epígrafe 2.2 del baremo)

                  - En este apartado el candidato registrará su actividad como PDI universitario mediante contrato laboral o pertenencia a los cuerpos docentes universitarios. Los contratos laborales que se contemplan corresponden a las modalidades de contratación (ordinarias y extraordinarias) recogidas en los apartados 1, 2 y 3 del artículo 14 del I Convenio Colectivo del PDI Laboral de las Universidad Públicas de Andalucía (BOJA 92/2008, de 9 de mayo) o sus análogas en otros Convenios, así como las categorías análogas de PDI de Universidades privadas.
                  - Profesor Ayudante
                  - Profesor Ayudante Doctor
                  - Profesor Colaborador
                  - Profesor Contratado Doctor
                  - Profesor Asociado
                  - Profesor Asociado de concierto con las instituciones sanitarias
                  - Profesor Sustituto Interino

                  Las figuras englobadas dentro de los cuerpos docentes universitarios son las siguientes:
                  - Profesor Titular de Escuela Universitaria
                  - Catedrático de Escuela Universitaria
                  - Profesor Titular de Universidad
                  - Catedrático de Universidad
                  - Figuras análogas en otros ámbitos internacionales

                  Revisión 3 (diciembre, 2024) 23

                  ---

                  FUENTE: CRITERIOS DE EVALUACIÓN Y REQUISITOS MÍNIMOS.pdf | PÁGINA: 11
                  \# Evaluación de Méritos Docentes

                  \#\# 1. Contratos de Profesorado

                  - Los contratos de profesorado asociado y los de profesorado sustituto interino se considerarán dedicación a tiempo completo cuando alcancen las 6 horas de clase semanal (180 horas de capacidad docente anual).
                  - El resto de los contratos de profesorado asociado y profesorado sustituto interino se considerará dedicación a tiempo parcial.
                  - Si la duración del contrato no es anual o por el curso académico completo, se contabilizará la parte proporcional.
                  - Los contratos de profesorado tutor de la UNED (y de figuras similares sin responsabilidad docente plena) se considerarán dedicación a tiempo parcial.
                  - En el caso de los contratos de profesorado vinculado (Ciencias de la Salud), cada 1,5 años de actividad clínica como especialista con compromiso docente con una universidad (colaborador clínico docente o figura similar) se considerará 1 año de dedicación a tiempo completo.

                  \#\# 2. Evaluación de Méritos

                  - En la evaluación de los méritos de actividad docente realizada por las personas que se encuentren en alguna de las situaciones especiales, derivadas de permisos, excedencias o bajas, enumeradas en el apartado primero (Disposiciones generales) de la Resolución de la Directora de ANECA de 28 de febrero de 2024, se disminuirán las exigencias de méritos relativos a años de docencia y horas lectivas impartidas de manera proporcional a la duración de los períodos de permiso, excedencia o baja acreditados.
                  - En todos los casos debe tratarse de docencia universitaria oficial (grado, postgrado o doctorado), o su equivalente en un sistema universitario extranjero.

                  \#\# 2.1. Criterios de Evaluación

                  \#\#\# 2.1.1. Pluralidad, Interdisciplinariedad y Complejidad Docente

                  - Se valorará la complejidad de la docencia en función del número y diversidad de asignaturas impartidas, de la docencia en asignaturas de primer curso o en grupos numerosos (entendidos como aquellos que tienen una ratio alumnado/profesorado elevada con relación a la ratio media del área o la titulación), así como la impartición de docencia en distintos idiomas, o en centros o campus diferentes al de adscripción, entre otros factores.
                  - Así, en los casos que queden suficientemente justificados, se elevará la puntuación obtenida en el apartado 2.1.1.
                  - Además, en todos los casos y en aras de considerar la interdisciplinariedad, se valorará también la docencia que trascienda las áreas de conocimiento asignadas a la Comisión de acreditación que realice la evaluación.

                  \#\#\# 2.1.2. Recursos Educativos

                  - En el caso de PTU podrá aportarse un máximo de 2 contribuciones.
                  - En el caso de CU podrá aportarse un máximo de 5 contribuciones.
                  - Se valorarán preferentemente los materiales docentes y los recursos educativos que estén disponibles en acceso abierto y sin restricciones a través de repositorios institucionales.

                  ----

                  \textbf{DIRECCIÓN DE ANECA – 20/03/2024}

                  ---

                  FUENTE: 27831\_instrucciones-baremacinin-r03.pdf | PÁGINA: 2
                  4.2. RENDIMIENTO ACADÉMICO (BLOQUE I DEL BAREMO) ..........................................................................................18
                  4.2.1. Expediente Académico (epígrafe 1.1 del baremo) ..........................................................................................18
                  4.2.2. Título de doctor (epígrafe 1.2 del baremo) ...................................................................................................21
                  4.2.3. Otros títulos académicos (epígrafe 1.3 del baremo) ......................................................................................22
                  4.2.4. Premios académicos y menciones de calidad (epígrafe 1.4 del baremo) ........................................................22
                  4.3. ACTIVIDAD DOCENTE (BLOQUE II DEL BAREMO) ................................................................................................23
                  4.3.1. Docencia en asignaturas regladas universitarias (epígrafe 2.1 del baremo) ..................................................23
                  4.3.2. Contratos como PDI universitario (epígrafe 2.2 del baremo) .......................................................................23
                  4.3.3. Dirección de tesis doctorales (epígrafe 2.3 del baremo) ..............................................................................24
                  4.3.4. Otros méritos de docencia universitaria (epígrafe 2.4 del baremo) .................................................................24
                  4.4. ACTIVIDAD INVESTIGADORA (BLOQUE III DEL BAREMO) ....................................................................................24
                  4.4.1. Contratos y becas en procesos competitivos (epígrafe 3.1 del baremo) ........................................................25
                  4.4.2. Libros (o monografías) científicas (epígrafe 3.2 del baremo) .......................................................................26
                  4.4.3. Capítulos de Libro (epígrafe 3.3 del baremo) ...............................................................................................27
                  4.4.4. Artículos científicos de Revista (epígrafe 3.4 del baremo) ..........................................................................27
                  4.4.5. Participación en proyectos de investigación (epígrafe 3.5 del baremo) ..........................................................27
                  4.4.6. Trabajos presentados en congresos (epígrafe 3.6 del baremo) ...................................................................28
                  4.4.7. Comités científicos, técnicos o asesores (epígrafe 3.7 del baremo) ...............................................................28
                  4.4.8. Conocimiento de idiomas relevantes (epígrafe 3.8 del baremo) ..................................................................28
                  4.5. TRANSFERENCIA DEL CONOCIMIENTO (BLOQUE IV DEL BAREMO) ....................................................................29
                  4.5.1. Contratos de transferencia (epígrafe 4.1 del baremo) ..................................................................................29
                  4.5.2. Obra civil o de edificación, o exposiciones (epígrafe 4.2 del baremo) ............................................................29
                  4.5.3. Patentes (epígrafe 4.3 del baremo) ...............................................................................................................29
                  4.5.4. Spin-off (epígrafe 4.4 del baremo) ................................................................................................................29

                  ---

                  FUENTE: CRITERIOS DE EVALUACIÓN Y REQUISITOS MÍNIMOS.pdf | PÁGINA: 9
                  \# Actividad Docente

                  En todos los casos, se valorará el grado de internacionalización, de innovación, de apertura a la sociedad y de multidisciplinariedad o interdisciplinariedad de las contribuciones presentadas.

                  Las personas solicitantes que acrediten trayectorias de investigación excepcionales, por haber obtenido financiación del Consejo Europeo de Investigación (ERC) en sus programas Consolidator Grant, Advanced Grant o Synergy Grant (o en otras convocatorias competitivas individuales internacionales de prestigio y tasa de concesión comparables a los de estos programas) y hayan desarrollado su carrera principalmente en una universidad extranjera en la que el cómputo y los instrumentos de medición de la calidad de la actividad docente resulten difíciles de trasladar al sistema español, o bien en una institución no universitaria dedicada a la investigación, podrán obtener la acreditación justificando el encargo docente desempeñado, o bien actividades docentes de tutorización y dirección de trabajos de investigación, respectivamente.

                  En caso de adaptación del puesto de trabajo con reducción de las horas lectivas o modalización del cumplimiento de las obligaciones docentes debido a una situación transitoria de pérdida de capacidades o de protección de la maternidad si existe riesgo laboral o una situación reconocida de discapacidad, las exigencias de méritos docentes se modularán de acuerdo con la relevancia de la adaptación del puesto de trabajo acreditada.

                  \#\# PROFESORA O PROFESOR TITULAR UNIVERSIDAD (PTU)

                  \vspace{0.4cm}
                  \begin{table}[h!]\centering\renewcommand{\arraystretch}{1.3}
                    \begin{tabular}{|p{0.45\textwidth}|p{0.45\textwidth}|p{0.45\textwidth}|}\hline
                      Actividad Docente                                                   & Mínimo puntos & Máximo puntos \\\hline
                      2. ACTIVIDAD DOCENTE                                                & 50            & 100           \\\hline
                      2.1. Experiencia docente                                            & 25            & 45            \\\hline
                      2.1.1. Dedicación docente                                           & 25            & 35            \\\hline
                      2.1.2. Pluralidad, interdisciplinariedad y complejidad docente      & 0             & 6             \\\hline
                      2.1.3. Recursos educativos                                          & 0             & 5             \\\hline
                      2.1.4. Actividades de formación a lo largo de la vida               & 0             & 5             \\\hline
                      2.2. Calidad de la actividad docente e innovación                   & 5             & 40            \\\hline
                      2.2.1. Calidad de la actividad docente                              & 5             & 15            \\\hline
                      2.2.2. Proyectos de innovación docente                              & 0             & 20            \\\hline
                      2.2.3. Formación para la mejora docente recibida                    & 0             & 15            \\\hline
                      2.2.4. Formación para la mejora docente impartida                   & 0             & 5             \\\hline
                      2.3. Tutorización docente                                           & 0             & 15            \\\hline
                      2.3.1. Tutorización reglada                                         & 0             & 15            \\\hline
                      2.3.2. Otros tipos de tutorización (no incluida en encargo docente) & 0             & 15            \\\hline
                      2.4. Otros méritos*                                                 & 0             & 10            \\\hline
                    \end{tabular}\end{table}
                  \vspace{0.4cm}

                  ----

                  DIRECCIÓN DE ANECA – 20/03/2024
                \end{scriptsize}
        \end{itemize}


  \item \textbf{P14.} Tengo una duda con mis exámenes, ¿cuántas convocatorias tengo en total para aprobar una asignatura matriculada? \\
        \begin{itemize}
          \item \textbf{Tipo de consulta:} multi
          \item \textbf{Respuesta del Agente:}

                Según la normativa de la Universidad de Sevilla, los estudiantes matriculados dispondrán de \textbf{seis convocatorias} para superar cada asignatura. No se computarán las convocatorias en las que el estudiante figure en acta con la mención de “No Presentado”.

                Además, si se agotan las \textbf{seis convocatorias} y el estudiante ha superado el primer año de permanencia, con más de \textbf{30 créditos} pendientes para finalizar sus estudios, dispondrá de una \textbf{primera convocatoria de gracia} para superar la asignatura, que se concederá de oficio sin necesidad de solicitarla.

                Es importante tener en cuenta que, si se agota la \textbf{primera convocatoria de gracia}, el estudiante puede solicitar al Rector una \textbf{segunda convocatoria de gracia}, que se resolverá favorablemente si el estudiante ha superado al menos el \textbf{50\%} de la carga lectiva global de la titulación. También se puede solicitar una \textbf{tercera convocatoria de gracia} si el estudiante ha superado el \textbf{75\%} de la carga lectiva global de la titulación.

                En resumen, el número total de convocatorias para aprobar una asignatura puede variar dependiendo de la situación del estudiante, pero inicialmente se dispone de \textbf{seis convocatorias}, con la posibilidad de obtener convocatorias de gracia adicionales bajo ciertas condiciones.
          \item \textbf{Interacción Usuario:} Ya superé los créditos del primer año, pero me quedan más de 30 créditos para terminar el grado.
          \item \textbf{Respuesta Final:}

                Según la normativa de la Universidad de Sevilla, establecida en la Resolución Rectoral de Normas de Matrícula de Grado y Máster Universitario 2025-2026, página 18, punto 2, los estudiantes matriculados dispondrán de \textbf{seis convocatorias} para superar cada asignatura.

                Además, según el punto 4 del mismo documento, si se agotan las seis convocatorias y el estudiante ha superado el primer año de permanencia, con más de \textbf{30 créditos} pendientes para finalizar sus estudios, dispondrá de una \textbf{primera convocatoria de gracia} para superar la asignatura, que se concederá de oficio sin necesidad de solicitarla.

                Por lo tanto, en tu caso, ya que has superado los créditos del primer año pero te quedan más de 30 créditos para terminar el grado, dispondrás de \textbf{siete convocatorias} en total para aprobar una asignatura matriculada: las \textbf{seis convocatorias} ordinarias y la \textbf{primera convocatoria de gracia}.

                Es importante tener en cuenta que, si se agota la primera convocatoria de gracia, podrás solicitar al Rector una \textbf{segunda convocatoria de gracia}, que se resolverá favorablemente si has superado al menos el \textbf{50\%} de la carga lectiva global de la titulación. También se puede solicitar una \textbf{tercera convocatoria de gracia} si has superado el \textbf{75\%} de la carga lectiva global de la titulación.
        \end{itemize}


  \item \textbf{P15.} Me quiero dar de baja y anular mi matrícula entera este año, ¿entonces me devolverán el dinero o todo lo que he pagado? \\
        \begin{itemize}
          \item \textbf{Tipo de consulta:} multi
          \item \textbf{Respuesta del Agente:}

                ⚠️ Requisitos Previos:
                Antes de iniciar el trámite de anulación de matrícula, es importante tener en cuenta los siguientes requisitos previos:
                - Debes haber realizado el pago de la matrícula correspondiente.
                - Debes cumplir con los plazos establecidos para la anulación de matrícula, que varían según la fecha de inicio de la docencia y el tipo de titulación.

                Para anular tu matrícula entera este año y solicitar la devolución del dinero pagado, sigue los siguientes pasos:

                1. \textbf{Verifica el plazo de anulación}: Comprueba si estás dentro del plazo establecido para la anulación de matrícula. Los plazos varían según la fecha de inicio de la docencia y el tipo de titulación. Puedes consultar los plazos en el artículo 11 de las Normas de Matrícula Grado y Máster Universitario 2025-2026.
                2. \textbf{Prepara la documentación necesaria}: Asegúrate de tener toda la documentación necesaria para solicitar la anulación de matrícula, incluyendo el \textbf{formulario de anulación de matrícula} y cualquier otro documento requerido por la Universidad de Sevilla.
                3. \textbf{Solicita la anulación de matrícula}: Dirígete al Decano/Director del centro responsable del título oficial de los estudios donde formalizaste tu matrícula y solicita la anulación de matrícula. Debes hacer constar la causa por la que solicitas la anulación.
                4. \textbf{Espera la resolución}: La Universidad de Sevilla resolverá tu solicitud de anulación de matrícula en primera instancia. Si se concede la anulación, se te devolverán los importes abonados, en su caso.
                5. \textbf{Verifica la devolución del dinero}: Si se concede la anulación de matrícula, la Universidad de Sevilla te devolverá el importe abonado en concepto de precios públicos por servicios académicos, excepto el importe de los precios públicos por servicios administrativos y el seguro escolar.

                Recuerda que la anulación de matrícula puede tener consecuencias académicas y económicas, por lo que es importante leer detenidamente las Normas de Matrícula Grado y Máster Universitario 2025-2026 y consultar con la Universidad de Sevilla si tienes alguna duda.
          \item \textbf{Interacción Usuario:} La solicito antes de finales de octubre del curso académico.
          \item \textbf{Respuesta Final:}

                ⚠️ Requisitos Previos:
                Antes de iniciar el trámite de anulación de matrícula, es importante tener en cuenta los siguientes requisitos previos:
                - Debes haber realizado el pago de la matrícula correspondiente.
                - Debes cumplir con los plazos establecidos para la anulación de matrícula, que varían según la fecha de inicio de la docencia y el tipo de titulación.

                Para anular tu matrícula entera este año y solicitar la devolución del dinero pagado antes de finales de octubre del curso académico, sigue los siguientes pasos:

                1. \textbf{Verifica el plazo de anulación}: Comprueba si estás dentro del plazo establecido para la anulación de matrícula. Según el artículo 11 de las Normas de Matrícula Grado y Máster Universitario 2025-2026, los plazos son:
                * Titulaciones de grado: hasta el 16 de septiembre de 2025.
                * Titulaciones de máster universitario que inicien la docencia el 6 de octubre: hasta el 15 de octubre de 2025.
                * Titulaciones de máster universitario que inicien la docencia el 22 de septiembre: hasta el 30 de septiembre de 2025.
                * Cuando la matrícula se haya formalizado con posterioridad al inicio del periodo docente del primer cuatrimestre: 7 días hábiles a contar desde el día de la finalización del periodo en el que haya realizado su matrícula.
                2. \textbf{Prepara la documentación necesaria}: Asegúrate de tener toda la documentación necesaria para solicitar la anulación de matrícula, incluyendo el \textbf{formulario de anulación de matrícula} y cualquier otro documento requerido por la Universidad de Sevilla.
                3. \textbf{Solicita la anulación de matrícula}: Dirígete al Decano/Director del centro responsable del título oficial de los estudios donde formalizaste tu matrícula y solicita la anulación de matrícula. Debes hacer constar la causa por la que solicitas la anulación.
                4. \textbf{Espera la resolución}: La Universidad de Sevilla resolverá tu solicitud de anulación de matrícula en primera instancia. Si se concede la anulación, se te devolverán los importes abonados, en su caso.
                5. \textbf{Verifica la devolución del dinero}: Si se concede la anulación de matrícula, la Universidad de Sevilla te devolverá el importe abonado en concepto de precios públicos por servicios académicos, excepto el importe de los precios públicos por servicios administrativos y el seguro escolar.

                Recuerda que la anulación de matrícula puede tener consecuencias académicas y económicas, por lo que es importante leer detenidamente las Normas de Matrícula Grado y Máster Universitario 2025-2026 y consultar con la Universidad de Sevilla si tienes alguna duda.
          \item \textbf{Contexto utilizado:}
                \begin{scriptsize}
                  FUENTE: R.R. NORMAS DE MATRICULA GRADO Y MASTER UNIVERSITARIO 2025-2026.pdf | PÁGINA: 11
                  \# UNIVERSIDAD D SEVII

                  C)EST

                  5. No se concederán anulaciones parciales de matrícula a solicitud del estudiante, excepto por coincidencia no prevista de horarios de clase y exámenes y las situaciones contempladas en estas normas. No obstante, por causas académicas sobrevenidas que imposibiliten la prestación del servicio, los centros podrán realizar de oficio anulaciones parciales de matrícula, ofreciendo a los estudiantes la posibilidad de matricular otras asignaturas o bien la devolución de los importes abonados en los créditos anulados, en su caso.

                  6. La matrícula se anulará de oficio cuando se trate de error, omisión, falseamiento de los datos consignados en esta o cuando se modifique la situación administrativa o académica del estudiante, o las matrículas condicionadas que no deban producir efectos.

                  Asimismo se anularán de oficio las matrículas cuando el estudiante no aporte la documentación requerida o por falta de pago.

                  7. Los centros anularán de oficio las matrículas efectuadas por estudiantes que se matriculen en otra titulación por haber obtenido plaza en los procedimientos de admisión para estudiantes de nuevo ingreso en esta u otra universidad pública, siempre que no se acredite una autorización para la simultaneidad de estudios.

                  \#\# Artículo 11. Supuestos de anulación de matrícula y sus efectos

                  1. Se podrá solicitar la anulación de matrícula ordinaria al inicio del curso académico en los siguientes plazos:

                  - Titulaciones de grado: hasta el 16 de septiembre de 2025.
                  - Titulaciones de máster universitario que inicien la docencia el 6 de octubre: hasta el 15 de octubre de 2025.
                  - Titulaciones de máster universitario que inicien la docencia el 22 de septiembre: hasta el 30 de septiembre de 2025.
                  - Cuando la matrícula se haya formalizado con posterioridad al inicio del periodo docente del primer cuatrimestre: 7 días hábiles a contar desde el día de la finalización del periodo en el que haya realizado su matrícula.

                  La concesión de la anulación supondrá:

                  a) La devolución de los importes abonados, en su caso.

                  b) La matrícula no se computará a efectos de convocatorias de examen ni de permanencia del estudiante, ni se contabilizará para futuras matrículas a efectos de aplicación de recargos.

                  c) La pérdida de la plaza adjudicada si el estudiante accedió por el procedimiento de admisión en el curso 2025-2026.

                  2. La anulación de la matrícula por haber obtenido plaza en esta u otra universidad pública (en el procedimiento o sistema de admisión que tenga establecido) en otra titulación de grado, máster universitario o estudios de doctorado donde hubiese sido admitido el estudiante con posterioridad a su matriculación en la Universidad de Sevilla, no tendrá ningún efecto académico ni económico si se solicita con anterioridad al 31 de octubre de 2025, en el caso de titulaciones de grado, y al 31 de diciembre de 2025 en las titulaciones de máster universitario.

                  Código Seguro De Verificación: `eqOjhZcmL0p6VJlZiTH/Yw==`
                  Fecha: 27/06/2025
                  Firmado Por: MIGUEL ANGEL CASTRO ARROYO
                  Url De Verificación: [https://pfirma.us.es/verifirma/code/eqOjhZcmL0p6VJlZiTH\%2FYw\%3D\%3D](https://pfirma.us.es/verifirma/code/eqOjhZcmL0p6VJlZiTH\%2FYw\%3D\%3D)
                  Página: 7/50

                  ---

                  FUENTE: R.R. NORMAS DE MATRICULA GRADO Y MASTER UNIVERSITARIO 2025-2026.pdf | PÁGINA: 12
                  \# UNIVERSIDAD DE SEVILLA

                  \#\# C)EST

                  En este supuesto procederá:

                  1. La devolución de los importes abonados, en su caso. Si la nueva plaza obtenida lo es de un centro propio de la Universidad de Sevilla, el importe abonado en la primera titulación podrá ser compensado en la nueva.

                  2. La pérdida de la plaza adjudicada en el procedimiento de admisión en el curso 2025-2026 en la titulación que abandona.

                  3. La anulación de matrícula por traslado de expediente a otros estudios de la Universidad de Sevilla no tendrá ningún efecto académico ni económico. Procederá la devolución de los importes abonados, en su caso. Si la nueva plaza obtenida lo es de un centro propio de la Universidad de Sevilla, el importe abonado en la primera titulación podrá ser compensado en la nueva.

                  4. La anulación de matrícula solicitada hasta el 30 de noviembre de 2025, como consecuencia de traslado de expediente a otra universidad o por haber obtenido plaza en un ciclo formativo de grado superior, dará lugar a la devolución de los importes abonados, en su caso.

                  5. El estudiante que en el presente curso académico esté al corriente de todos los pagos de su matrícula, incluso los no vencidos en caso de fraccionamiento, podrá solicitar la anulación de su matrícula hasta el 31 de marzo, con las siguientes consecuencias:

                  a) No procederá la devolución de los importes abonados en concepto de precios públicos por servicios académicos y administrativos y el seguro escolar.

                  b) Serán de aplicación los recargos económicos en la siguiente matrícula que se realice en los mismos estudios.

                  c) No se computará la matrícula realizada a efectos académicos. En consecuencia, a efectos de beca, solo se tendrá en cuenta el rendimiento académico obtenido en el curso anterior del que solicita la anulación de matrícula.

                  d) Si el estudiante accedió por el procedimiento de admisión en el curso 2025-2026, perderá la plaza adjudicada, sin perjuicio de poder volver a solicitar la admisión en cursos posteriores en la misma o en distinta titulación.

                  e) En caso de que el estudiante haya formalizado la matrícula con la exención de becario de carácter general, debe solicitar la renuncia a la solicitud de beca y abonar el importe de la matrícula como requisito para la anulación por esta causa.

                  6. El plazo para solicitar la anulación de la matrícula formalizada en el periodo general de ampliación será del 2 al 6 de febrero de 2026. La concesión de la anulación de la ampliación de matrícula supondrá la devolución del importe de los precios públicos abonados por servicios académicos, en su caso, y la matrícula no se contabilizará a efectos de aplicación de recargos, en cursos posteriores.

                  7. El estudiante afectado por la imposibilidad de obtener el visado de entrada en España o por otra causa no imputable al estudiante, podrá solicitar la anulación total de su matrícula, con las siguientes consecuencias:

                  a) Procederá la devolución de los importes abonados, en su caso.

                  b) El estudiante perderá la plaza adjudicada por el procedimiento de admisión en el curso 2025-2026.

                  ----

                  \textbf{Código Seguro De Verificación}: eqOjhZcmL0p6VJlZiTH/Yw==
                  \textbf{Fecha}: 27/06/2025
                  \textbf{Firmado Por}: MIGUEL ANGEL CASTRO ARROYO
                  \textbf{Url De Verificación}: Verificación
                  \textbf{Página}: 8/50

                  ---

                  FUENTE: R.R. NORMAS DE MATRICULA GRADO Y MASTER UNIVERSITARIO 2025-2026.pdf | PÁGINA: 10
                  \# UNIVERSIDAD DE SEVILLA

                  C)EST

                  5. En el acto de ampliación de matrícula serán de aplicación las reducciones de precios públicos previstas en el artículo 28 de estas normas siempre que se haya acreditado el derecho a la deducción o reducción de precios en el periodo ordinario, o bien, podrán aplicarse nuevas deducciones en el periodo de ampliación acreditadas conforme establece el citado artículo.

                  6. La falta de pago de la ampliación de matrícula determinará la anulación de la ampliación.

                  7. Con carácter general, la anulación de la matrícula ordinaria determinará también la anulación de la ampliación de matrícula, en su caso.

                  \#\# Artículo 9. Bloqueo de expediente

                  Como actuación previa a la anulación de matrícula, el centro bloqueará el expediente de aquellos estudiantes que no hayan aportado la documentación establecida en estas normas, o que se encuentren en situación de impago, una vez notificado el requerimiento correspondiente o cuando se produzca el rechazo automático del mismo por el transcurso del plazo establecido.

                  Este bloqueo tendrá las siguientes consecuencias sobre el expediente de la titulación a la que corresponda:

                  - No podrán emitirse certificaciones académicas.
                  - El estudiante no podrá solicitar su título universitario oficial.
                  - El estudiante no tendrá acceso a la plataforma de Enseñanza Virtual de la Universidad de Sevilla.

                  El estudiante dispondrá de 15 días hábiles siguientes a la notificación del requerimiento, o en su caso, desde que se produzca el rechazo automático de la misma por el transcurso del plazo establecido, para aportar la documentación requerida o abonar las cantidades pendientes, según proceda, en cuyo caso el centro procederá al desbloqueo de su expediente.

                  En el supuesto de que el estudiante no aporte la documentación requerida en el plazo indicado o no abone las cantidades pendientes de pago, el centro procederá a la anulación de su matrícula en el curso actual, con los efectos que se establecen en el artículo 11 de estas normas.

                  \#\# Artículo 10. Anulación de matrícula

                  1. La anulación de matrícula supone el cese de los efectos académicos de la matrícula realizada con la consiguiente pérdida de los derechos a evaluación, manteniéndose la obligación de satisfacer los precios públicos salvo que proceda su devolución en los términos previstos reglamentariamente.

                  2. La anulación de la matrícula en un curso conllevará la anulación de los créditos o asignaturas reconocidos en ese curso sin derecho a devolución de los importes satisfechos, lo que no impedirá que el estudiante pueda volver a solicitarlo en cursos posteriores conforme a los procedimientos correspondientes.

                  3. Las solicitudes de anulación de matrícula se dirigirán al Decano/Director del centro responsable del título oficial de los estudios donde formalice su matrícula por cualquiera de los medios establecidos en las disposiciones vigentes, haciendo constar la causa por la que se solicita, y serán resueltas en primera instancia por el centro.

                  4. No se concederán anulaciones de matrícula a los estudiantes con alguna calificación incorporada en actas ya cerradas definitivamente, en alguna de las asignaturas matriculadas en el curso académico para el que solicita la anulación.

                  Código Seguro De Verificación: eqOjhZcmL0p6VJlZiTH/Yw==
                  Fecha: 27/06/2025
                  Firmado Por: MIGUEL ANGEL CASTRO ARROYO
                  Url De Verificación: [https://pfirma.us.es/verifirma/code/eqOjhZcmL0p6VJlZiTH\%2FYw\%3D\%3D](https://pfirma.us.es/verifirma/code/eqOjhZcmL0p6VJlZiTH\%2FYw\%3D\%3D)
                  Página: 6/50

                  ---

                  FUENTE: R.R. NORMAS DE MATRICULA GRADO Y MASTER UNIVERSITARIO 2025-2026.pdf | PÁGINA: 13
                  \# UNIVERSIDAD D SEVII

                  C)EST

                  c) No se aplicarán recargos en siguientes cursos respecto a las asignaturas anuladas.

                  d) En el supuesto de obtención del visado con posterioridad al inicio del curso académico, podrá solicitar, previa aportación de la documentación acreditativa de la fecha de obtención de este, la anulación parcial de aquellas asignaturas cuya docencia se haya impartido total o parcialmente con anterioridad a la obtención del visado, salvo que hubiera alguna calificación incorporada en acta.

                  8. Excepcionalmente, el estudiante podrá solicitar la anulación total o parcial de matrícula de asignaturas no calificadas en acta por enfermedad grave sobrevenida, documentalmente acreditada, que le impida continuar los estudios.

                  Por enfermedad grave se entenderá cualquier dolencia o lesión física o psíquica que incapacite temporalmente para la ocupación o actividad habitual del estudiante durante un periodo continuado mínimo de tres meses en el curso académico para el que solicita la anulación y que requiera tratamiento o intervención clínica de cirugía mayor. El solicitante deberá aportar informe clínico expedido por el sistema público de salud español. Del mismo modo, en caso de pertenencia obligatoria a una mutualidad, el informe será expedido por la misma.

                  Las solicitudes se presentarán por cualquiera de los medios establecidos en las disposiciones vigentes dirigidas al Vicerrectorado de Estudiantes. La entrega por el estudiante de la documentación médica se podrá realizar en sobre cerrado e irá acompañada de un consentimiento expreso para el tratamiento de los datos incluidos en esta documentación, cuyo fin exclusivo será la gestión por la Universidad de Sevilla de su solicitud de anulación de matrícula por enfermedad grave sobrevenida. Por el Vicerrectorado de Estudiantes se remitirá dicha documentación a la Comisión Técnica de Expertos para la Valoración de los Supuestos Determinantes de la Anulación de la Matrícula por Enfermedad Grave Sobrevenida, comisión de carácter médico, para la emisión del correspondiente informe. La referida comisión podrá recabar cuantos otros informes estimen oportunos.

                  En el supuesto de resolverse favorablemente la solicitud de anulación de matrícula, los efectos serán:

                  a) La devolución del importe abonado en concepto de precios públicos por servicios académicos de las asignaturas anuladas, a excepción del importe de los precios públicos por servicios administrativos y el seguro escolar.

                  b) No se computará la matrícula a efectos de convocatorias de examen ni de permanencia del estudiante, ni se contabilizará para futuras matrículas a efectos de aplicación de recargos.

                  c) El estudiante no perderá la plaza adjudicada en el procedimiento de admisión en el curso 2025-2026.

                  9. En los supuestos de anulación de oficio por falta de algún requisito no subsanable y no imputable al interesado, procederá, en su caso, la devolución de los importes abonados, en su caso. Si el estudiante accedió por el procedimiento de admisión en el curso 2025-2026, perderá la plaza adjudicada.

                  10. Cuando la anulación de matrícula sea consecuencia de sanción disciplinaria:

                  a) No procederá la devolución de los importes abonados en concepto de precios públicos por servicios académicos y administrativos y el seguro escolar, ni se eximirá de la obligación de satisfacer las cantidades pendientes de pago.

                  b) Serán de aplicación los recargos económicos en la siguiente matrícula que realice en los mismos estudios.

                  Código Seguro De Verificación: eqOjhZcmL0p6VJlZiTH/Yw==
                  Fecha: 27/06/2025
                  Firmado Por: MIGUEL ANGEL CASTRO ARROYO
                  Url De Verificación: [https://pfirma.us.es/verifirma/code/eqOjhZcmL0p6VJlZiTH\%2FYw\%3D\%3D](https://pfirma.us.es/verifirma/code/eqOjhZcmL0p6VJlZiTH\%2FYw\%3D\%3D)
                  Página: 9/50

                  ---

                  FUENTE: R.R. NORMAS DE MATRICULA GRADO Y MASTER UNIVERSITARIO 2025-2026.pdf | PÁGINA: 14
                  \# UNIVERSIDAD D SEVII

                  \#\#\# C)EST

                  c) No se computará la matrícula realizada a efectos académicos.

                  d) En caso de que la sanción suponga la anulación total de la matrícula y el estudiante hubiera accedido por el procedimiento de admisión en el curso 2025-2026, perderá la plaza adjudicada.

                  11. En los supuestos de anulación por falta de pago o por falta de aportación de la documentación requerida, los efectos serán:

                  a) No procederá la devolución de los importes abonados en concepto de precios públicos por servicios académicos y administrativos y el seguro escolar, ni se eximirá de la obligación de satisfacer las cantidades pendientes de pago.

                  b) El abono de las cantidades pendientes de pago no podrá ser objeto de aplazamiento, fraccionamiento ni domiciliación bancaria.

                  c) Serán de aplicación los recargos económicos en la siguiente matrícula que realicen en los mismos estudios.

                  d) No se computará la matrícula realizada a efectos académicos y supondrá la pérdida de las calificaciones incorporadas en acta.

                  e) Si el estudiante accedió por el procedimiento de admisión en el curso 2025-2026, perderá la plaza adjudicada.

                  f) La pérdida de la condición de becario de la convocatoria general y, en consecuencia, la supresión de dicha exención en su matrícula.

                  \#\#\# Sección Tercera: Simultaneidad de estudios. Estudiantes visitantes

                  A) Simultaneidad de estudios

                  \#\#\#\# Artículo 12. De la matrícula para simultanear estudios

                  En materia de simultaneidad de estudios se estará a lo dispuesto en la regulación contenida en los Acuerdos de la Comisión de Distrito Único Universitario de Andalucía y en las demás normas de aplicación.

                  \#\#\#\# Artículo 13. Requisitos

                  1. Los estudiantes que deseen simultanear estudios desde su inicio deberán obtener plaza en ambas titulaciones en el procedimiento de admisión, para lo que deberán indicarlo expresamente en el mismo. Al matricularse, el estudiante deberá comunicar su situación de simultaneidad expresamente en los centros correspondientes. En el expediente académico que se abra en cada titulación se hará constar tal circunstancia.

                  2. Solo en los casos de simultaneidad entre estudios de grado ya iniciados en la Universidad de Sevilla con otros en los que haya obtenido plaza en otra universidad española, el estudiante deberá abonar los derechos de traslado de expediente a efectos de simultaneidad en el centro de origen y entregar dicho resguardo en el centro de destino. El incumplimiento de este requisito dará lugar a la anulación de la matrícula. Los centros implicados deberán dejar constancia del traslado por simultaneidad en el expediente del estudiante.

                  Código Seguro De Verificación: `eqOjhZcmL0p6VJlZiTH/Yw==`
                  Fecha: 27/06/2025
                  Firmado Por: MIGUEL ANGEL CASTRO ARROYO
                  Url De Verificación: [https://pfirma.us.es/verifirma/code/eqOjhZcmL0p6VJlZiTH\%2FYw\%3D\%3D](https://pfirma.us.es/verifirma/code/eqOjhZcmL0p6VJlZiTH\%2FYw\%3D\%3D)
                  Página: 10/50
                \end{scriptsize}
        \end{itemize}


  \item \textbf{P16.} ¿Estoy completamente obligado a realizar la evaluación continua de una asignatura con mis compañeros de proyecto? \\
        \begin{itemize}
          \item \textbf{Tipo de consulta:} multi
          \item \textbf{Respuesta del Agente:}

                Según la normativa reguladora de la evaluación y calificación de las asignaturas de la Universidad de Sevilla, la evaluación continua es un componente importante de la evaluación de las competencias, conocimientos y capacidades adquiridas por los estudiantes.

                En el artículo 6.1 de la normativa, se establece que los sistemas de evaluación pueden basarse en actividades de evaluación continua, entre otras opciones. Además, en el artículo 11.1, se define la evaluación continua como la participación en las clases lectivas, la realización de prácticas, trabajos y otras pruebas que se realicen en presencia de un profesor.

                Sin embargo, no hay una disposición explícita que establezca que los estudiantes estén obligados a realizar la evaluación continua con sus compañeros de proyecto. En el artículo 9.1, se establece que los estudiantes tienen derecho a optar entre las distintas posibilidades de evaluación que se contemplen en los proyectos docentes, lo que sugiere que pueden haber diferentes opciones de evaluación continua.

                En el artículo 4.3, se establece que los estudiantes tienen el derecho y el deber de participar en las diversas actividades académicas por las que hayan de ser evaluados. Sin embargo, no se especifica que deban realizar la evaluación continua con sus compañeros de proyecto.

                En resumen, aunque la evaluación continua es un componente importante de la evaluación, no hay una disposición explícita que establezca que los estudiantes estén obligados a realizar la evaluación continua con sus compañeros de proyecto. Es posible que los proyectos docentes contemplen diferentes opciones de evaluación continua, y los estudiantes pueden optar por la que mejor se adapte a sus necesidades y preferencias.

                Es recomendable que consultes con el profesor o coordinador de la asignatura para obtener más información sobre las opciones de evaluación continua disponibles y cómo se aplicarán en tu caso específico.
          \item \textbf{Interacción Usuario:} Sí, el caso es que mi actividad laboral actual me impide asistir a clase regularmente.
          \item \textbf{Respuesta Final:}

                Lamento no poder ayudarte con tu consulta, ya que no está relacionada con trámites o normativas universitarias. Sin embargo, te recomendaría contactar con el departamento de orientación académica o con el profesor de la asignatura para discutir posibles soluciones para tu situación. Es posible que puedan ofrecerte alternativas para seguir participando en la asignatura de manera remota o proporcionarte información sobre las políticas de asistencia y participación en clase.
          \item \textbf{Contexto utilizado:}
                \begin{scriptsize}
                  FUENTE: NORMATIVA REGUL.pdf | PÁGINA: 2
                  \# Evaluación Continua y Exámenes

                  4. Las actividades de evaluación continua y los exámenes de cada estudiante deberán ser evaluados y calificados en su integridad.

                  5. Cuando, con motivo del ejercicio por los estudiantes del derecho al paro académico u otra causa de fuerza mayor, no hubiere sido posible cubrir todos los contenidos de una asignatura en el periodo lectivo, el alcance de la evaluación será fijado mediante acuerdo entre los estudiantes de los grupos afectados, representados por sus delegados, y los profesores de la asignatura en dichos grupos. En caso de no haber acuerdo, resolverá la Comisión de Docencia del Departamento, oídos los delegados y los profesores afectados.

                  \#\# Artículo 5. Estudiantes con necesidades académicas especiales

                  La evaluación de las competencias, conocimientos y capacidades adquiridas por los estudiantes de la Universidad de Sevilla que, en virtud de lo dispuesto en el Reglamento general de estudiantes, tengan necesidades académicas especiales, se llevará a cabo de acuerdo con lo que se establezca, a este respecto, en la evaluación a que aluden los artículos 45 del Reglamento general de actividades docentes y 26.2.b) del Reglamento general de estudiantes.

                  \# CAPÍTULO II. SISTEMAS DE EVALUACIÓN

                  \#\# SECCIÓN 1ª. SISTEMAS DE EVALUACIÓN

                  \#\#\# Artículo 6. Sistemas de evaluación

                  1. Los diversos sistemas de evaluación de las competencias, conocimientos y capacidades adquiridas por los estudiantes incluidos en el programa de la asignatura podrán basarse en algunos de los siguientes elementos:
                  a) Actividades de evaluación continua.
                  b) Exámenes, parciales o finales.

                  2. Los sistemas podrán contemplar una relación de requisitos específicos como, por ejemplo, la realización de exámenes u otro tipo de pruebas, la asistencia a un número mínimo de horas de clases prácticas, la realización obligatoria de trabajos, proyectos o prácticas de laboratorio y la participación en seminarios.

                  3. La asistencia a las clases teóricas podrá puntuar de manera positiva en la ponderación de la calificación final, aunque no podrá exigirse como requisito ineludible para superar la asignatura. Las faltas de asistencia a la clases teóricas no podrán puntuar negativamente en la ponderación de la calificación final.

                  \#\#\# Artículo 7. Inclusión del sistema de evaluación en los proyectos docentes

                  1. En los proyectos docentes de cada curso académico se incluirá el sistema concreto de evaluación y calificación seleccionado de entre los sistemas propuestos en el programa de la asignatura. En dicho sistema deben constar los criterios de calificación de todas las actividades de evaluación continua y exámenes parciales y finales que se contemplen, así como su ponderación en la calificación final según la convocatoria de que se trate.

                  ---

                  FUENTE: NORMATIVA REGUL.pdf | PÁGINA: 4
                  \# SECCIÓN 2ª. ACTIVIDADES DE EVALUACIÓN CONTINUA

                  \#\# Artículo 11. Actividades de evaluación continua

                  1. Las actividades de evaluación continua mencionadas en el apartado 1 del artículo 6 comprenden:
                  a) La participación en las clases lectivas, tanto teóricas como prácticas, incluida la asistencia y defensa de ponencias y trabajos en seminarios.
                  b) La realización de prácticas informáticas, clínicas, jurídicas, de laboratorio, de campo, en aulas multidisciplinares de arquitectura o ingeniería, etc.
                  c) Los trabajos presentados en relación con el contenido de la asignatura.
                  d) Otras pruebas que se realicen como, por ejemplo, pequeñas pruebas de control periódico de conocimientos.
                  e) Cualquier otra actividad de evaluación que se lleve a cabo en presencia de un profesor ante un grupo de impartición de la asignatura en un aula, sala de seminario, laboratorio, taller, etc.

                  2. Las actividades de evaluación continua presenciales se realizarán siempre dentro del horario asignado a las clases lectivas de la asignatura fijado en su plan de organización docente.

                  3. A efectos de lo dispuesto en los artículos 28.1.g) del EUS y 34 del Reglamento General de Estudiantes, por disposición del Decano o Director del Centro, o previa solicitud de la Delegación de Alumnos, la comisión de seguimiento del plan de estudios, o el órgano que, en su caso, acuerde el Centro, se reunirá al comienzo de cada cuatrimestre con el profesorado y los delegados de los estudiantes de un mismo curso para coordinar el calendario de las actividades de evaluación continua de las diferentes asignaturas y materias de dicho cuatrimestre.

                  4. Los estudiantes deberán ser informados de los resultados de las actividades de evaluación continua en el plazo máximo de 20 días lectivos desde la realización de la actividad.

                  \#\# Artículo 12. Conservación de las calificaciones de actividades de evaluación continua

                  Los proyectos docentes podrán contemplar la posibilidad de conservar, en el caso de los estudiantes que no hubieran superado la asignatura en el curso académico anterior, las calificaciones obtenidas en aquellas actividades de evaluación continua cuyo contenido sea muy similar, sin perjuicio del derecho que tienen estos estudiantes de volver a realizar dichas actividades en su totalidad.

                  ---

                  FUENTE: NORMATIVA REGUL.pdf | PÁGINA: 1
                  \# NORMATIVA REGULADORA DE LA EVALUACIÓN Y CALIFICACIÓN DE LAS ASIGNATURAS

                  Aprobado por el Acuerdo 6.1 del Consejo de Gobierno en la sesión de 29 de septiembre de 2009 y modificado por el Acuerdo 1.1 del Consejo de Gobierno en la sesión de 18 de marzo de 2010.

                  \#\# CAPÍTULO I. PRINCIPIOS GENERALES

                  \#\#\# Artículo 1. Objeto

                  Serán objeto de evaluación individual las competencias, conocimientos y capacidades adquiridas por los estudiantes en relación con los objetivos y los contenidos fijados en los programas de la asignaturas.

                  \#\#\# Artículo 2. Ámbito de aplicación

                  La presente normativa se dicta en desarrollo de lo previsto en el Estatuto de la Universidad de Sevilla (en adelante EUS) y en el Reglamento general de actividades docentes, y será de aplicación a las asignaturas contenidas en los planes de estudios de las enseñanzas oficiales de Grado y de Máster y del periodo de docencia del Doctorado, así como, hasta su extinción, en los de las titulaciones oficiales reguladas por normativas anteriores a la Ley Orgánica 4/2007, de 12 de abril, por la que se modifica la Ley Orgánica 6/2001, de 21 de diciembre, de Universidades (en adelante LO 4/2007 y LO 6/2001, respectivamente).

                  \#\#\# Artículo 3. Responsables de la evaluación

                  La responsabilidad de la evaluación de una asignatura corresponde, en los términos contemplados en los proyectos docentes, a los profesores de la asignatura que tengan reconocida plena capacidad docente, sin perjuicio del derecho de los estudiantes a ser evaluados por el tribunal específico de evaluación al que se refiere el artículo 10.

                  \#\#\# Artículo 4. Principios generales

                  1. Los estudiantes tienen derecho a la corrección objetiva de las pruebas, exámenes u otros medios de evaluación de las competencias, conocimientos y capacidades adquiridas, a conocer sus calificaciones detalladamente en términos literales y numéricos dentro de los plazos fijados, así como a la revisión e impugnación de aquéllas mediante los mecanismos de garantía que se desarrollan en esta normativa.

                  2. Los profesores tienen el deber de evaluar de manera objetiva el nivel de las competencias, conocimientos y capacidades adquiridas por los estudiantes, ateniéndose al sistema incluido en el proyecto docente.

                  3. Los estudiantes tienen el derecho y el deber de participar en las diversas actividades académicas por las que hayan de ser evaluados.

                  ---

                  FUENTE: NORMATIVA REGUL.pdf | PÁGINA: 3
                  \# Artículo 7. Evaluación

                  1. El profesorado responsable de la evaluación se atendrá al sistema incluido en el proyecto docente. Salvo circunstancias excepcionales, los proyectos docentes no podrán ser modificados; en cualquier caso, las modificaciones que afecten al sistema de evaluación requerirán el visto bueno de los Delegados de los grupos afectados.

                  2. Con independencia de que existan uno o varios proyectos docentes, los criterios de calificación habrán de ser uniformes para los distintos grupos de una misma asignatura.

                  \# Artículo 8. Aprobado por curso

                  1. Los sistemas de evaluación contemplarán la posibilidad de aprobar por curso una asignatura, sea cuatrimestral o anual, de manera previa al examen final, caso de que lo hubiere.

                  2. Si el sistema de evaluación de una asignatura anual se basa, primordialmente, en exámenes, deberán realizarse dos exámenes parciales. La superación de éstos y de los demás elementos de evaluación, en los términos contemplados en el sistema de evaluación y calificación de la asignatura, conllevará el aprobado por curso.

                  3. En el caso del apartado anterior, si no se ha aprobado la asignatura de manera previa al examen final pero se ha superado uno de los exámenes parciales, el estudiante podrá no examinarse de los contenidos correspondientes en el examen final de la primera convocatoria.

                  \# Artículo 9. Derecho a optar entre los sistemas de evaluación

                  1. En cada asignatura los estudiantes tendrán derecho a optar entre las distintas posibilidades de evaluación que se contemplen en los proyectos docentes; en cualquier caso, la calificación máxima que se pueda obtener no se podrá ver afectada por el procedimiento de evaluación utilizado.

                  2. En aquellas asignaturas donde exista libre elección de grupo por parte de los estudiantes, el derecho a optar entre los sistemas de evaluación se ejercerá mediante la elección del grupo.

                  3. En otras circunstancias, el ejercicio de este derecho será solicitado mediante escrito dirigido al Decano o Director del Centro, que remitirá una copia al Director del Departamento, en el plazo máximo de un mes desde el comienzo de las clases de la asignatura.

                  \# Artículo 10. Tribunales específicos de evaluación

                  1. Para cada asignatura, el Consejo de Departamento elegirá, en la misma sesión en que se apruebe el proyecto de plan de asignación de profesorado, un tribunal específico de evaluación, así como un tribunal suplente.

                  2. El tribunal estará formado por tres profesores con plena capacidad docente del área de conocimiento, o área afín, a la que está adscrita la asignatura.

                  3. Presidirá el tribunal el miembro de mayor categoría y antigüedad y actuará de secretario el miembro de menor categoría y antigüedad.

                  4. Los estudiantes tienen derecho a que la evaluación la realice el tribunal específico de evaluación de la asignatura. El ejercicio de este derecho será solicitado mediante escrito, debidamente.

                  ---

                  FUENTE: NORMATIVA REGUL.pdf | PÁGINA: 6
                  \# Artículo 17. Cambio de la fecha de evaluación en casos excepcionales

                  1. Previa solicitud tendrán derecho a que se les facilite la realización de actividades de evaluación continua y exámenes en fechas distintas de las previstas los estudiantes que se encuentren en alguna de las siguientes situaciones excepcionales:
                  a) Estudiantes que sean representantes, que podrán ejercer este derecho en los términos recogidos en el artículo 5.1.e) del Reglamento General de Estudiantes.
                  b) Estudiantes matriculados en asignaturas distintas cuyos exámenes finales coincidan en la misma fecha. En este caso, el derecho se refiere a uno de dichos exámenes, cuya fecha podrá ser cambiada por una sola vez.
                  c) Estudiantes que estén en situación de ingreso hospitalario en la fecha de la evaluación o tengan un familiar en primer grado de consanguinidad en dicha situación.
                  d) Fallecimiento de un familiar en primer grado de consanguinidad en la fecha del examen o en los cuatro días anteriores.
                  e) Estudiantes que sean deportistas de alto nivel o deportistas de alto rendimiento, que podrán ejercer este derecho en los términos recogidos en el artículo 32 del Reglamento General de Estudiantes. En este caso el derecho enunciado se refiere únicamente a los exámenes que coincidan con actividades deportivas regladas.
                  f) Estudiantes que necesiten compaginar los estudios con la actividad laboral en el caso, que deberá ser justificado documentalmente, de que el convenio colectivo de aplicación no recoja los términos del ejercicio del derecho a los permisos necesarios para concurrir a exámenes establecido en el artículo 23.1.a) de la Ley 8/1980, de 10 de marzo, del Estatuto de los Trabajadores.

                  2. En las situaciones recogidas en las letras b) y e) del apartado 1, la solicitud será dirigida, a través de la Secretaría del Centro, al coordinador de la asignatura con una antelación mínima de dos meses respecto de la fecha prevista para la realización del examen. En el caso de la letra b), si las asignaturas son de distinto curso, la solicitud se referirá a la asignatura situada en el curso superior. En las demás situaciones, la solicitud será dirigida al coordinador de la asignatura con la mayor brevedad y por el medio más adecuado que las circunstancias permitan.

                  3. El coordinador de la asignatura, oídos el estudiante y el profesorado responsable de la actividad de evaluación, establecerá y comunicará a los interesados la nueva fecha que, en todo caso, deberá ser posterior, al menos en tres días naturales, al día en que se produzca la comunicación. Salvo imposibilidad material, en el caso de que haya varias solicitudes de cambio de fecha de un mismo examen, la nueva fecha será la misma para todos los solicitantes.

                  \# Artículo 18. Vigilancia de los exámenes

                  1. El profesorado de una asignatura tiene el deber de estar presente y participar en la vigilancia de los exámenes de dicha asignatura. Salvo imposibilidad material, cada dependencia en la que se celebre un examen deberá contar con al menos un profesor de la asignatura.

                  2. El profesorado del Departamento tiene el deber de colaborar en la vigilancia de los exámenes de las asignaturas adscritas al Departamento.

                  3. La organización y coordinación de los turnos de vigilancia de los exámenes corresponderá al Director del Departamento, previa consulta con los coordinadores de las asignaturas, quien procurará un reparto equitativo acorde con la dedicación docente.
                \end{scriptsize}
        \end{itemize}


\end{itemize}

\subsection{Categoría 3: Baremo (17-26)}
\textit{Pruebas relativas a la extracción y aplicación de criterios de evaluación y baremación de méritos en Bareprof.}
\begin{itemize}
  \item \textbf{P17.} Para la plaza de Profesor Ayudante Doctor, ¿cuál es el máximo de puntos que puedo conseguir en el apartado de actividad investigadora según el baremo de 2024? \\
        \begin{itemize}
          \item \textbf{Tipo de consulta:} simple
          \item \textbf{Respuesta del Agente:}

                \#\#\# Requisitos mínimos excluyentes:
                No se mencionan requisitos mínimos excluyentes específicos para el apartado de actividad investigadora en el baremo proporcionado.

                \#\#\# Criterios de evaluación para el apartado de actividad investigadora:
                La actividad investigadora se evalúa según varios criterios, que se detallan a continuación:

                \vspace{0.4cm}
                \begin{table}[h!]\centering\renewcommand{\arraystretch}{1.3}
                  \begin{tabular}{|p{0.45\textwidth}|p{0.45\textwidth}|}\hline
                    Criterio                                                          & Detalles           \\\hline
                    Contratos (o becas) en procesos competitivos                      & Máximo: 10 puntos  \\\hline
                    Participación en proyectos de investigación                       & Máximo: 10 puntos  \\\hline
                    Trabajos presentados en congresos científicos o tecnológicos      & Máximo: 10 puntos  \\\hline
                    Publicaciones científicas (artículos, capítulos de libro, libros) & Máximo: 100 puntos \\\hline
                  \end{tabular}\end{table}
                \vspace{0.4cm}

                \#\#\# Desglose de los apartados puntuables:
                1. \textbf{Contratos (o becas) en procesos competitivos}:
                - \textbf{Máximo:} 10 puntos
                - \textbf{Tipos de contratos y puntuaciones}:
                \vspace{0.4cm}
                \begin{table}[h!]\centering\renewcommand{\arraystretch}{1.3}
                  \begin{tabular}{|p{0.45\textwidth}|p{0.45\textwidth}|}\hline
                    Tipo de ayuda                                                                                                                                                                                                       & Puntos por tipo de contrato/año     \\\hline
                    FPU, FPI, Planes propios de Universidad (predoctoral)                                                                                                                                                               & 3,0 puntos por tipo de contrato/año \\\hline
                    Otras convocatorias públicas competitivas (predoctoral)                                                                                                                                                             & 2,0 puntos por tipo de contrato/año \\\hline
                    MEC, Juan de la Cierva, Fullbright, Humboldt, Marie Curie, Contratos de acceso a Sistema Español de Ciencia y Tecnología, o similares de fundaciones e instituciones de reconocido prestigio en I+D+i (posdoctoral) & 4,0 puntos por tipo de contrato/año \\\hline
                    Otras convocatorias públicas competitivas (posdoctoral)                                                                                                                                                             & 3,0 puntos por tipo de contrato/año \\\hline
                    Ramón y Cajal, Miguel Servet, Marie Curie, o similares de fundaciones e instituciones de reconocido prestigio en I+D+i (contratos sénior de reincorporación)                                                        & 6,0 puntos por tipo de contrato/año \\\hline
                    Otras convocatorias públicas competitivas (contratos sénior de reincorporación)                                                                                                                                     & 4,0 puntos por tipo de contrato/año \\\hline
                  \end{tabular}\end{table}
                \vspace{0.4cm}

                2. \textbf{Participación en proyectos de investigación}:
                - \textbf{Máximo:} 10 puntos
                - \textbf{Tipos de proyectos y puntuaciones}:
                \vspace{0.4cm}
                \begin{table}[h!]\centering\renewcommand{\arraystretch}{1.3}
                  \begin{tabular}{|p{0.45\textwidth}|p{0.45\textwidth}|}\hline
                    Tipo de Proyecto        & Puntos por Proyecto     \\\hline
                    Proyecto competitivo    & 5,0 puntos por proyecto \\\hline
                    Proyecto no competitivo & 2,0 puntos por proyecto \\\hline
                  \end{tabular}\end{table}
                \vspace{0.4cm}
                - \textbf{Complementos}:
                - \textbf{Ámbito geográfico}: Internacional (3,0 x puntos por proyecto), Internacional no UE (2,5 x puntos por proyecto), Nacional (1,5 x puntos por proyecto), Autonómico (1,0 x puntos por proyecto), Otros (0,5 x puntos por proyecto)
                - \textbf{Duración}: Menor a 2 años (1,0 x puntos por proyecto), 2 años o más (2,0 x puntos por proyecto)
                - \textbf{Grado de contribución}: Coordinador (3,00 x puntos por proyecto), Miembro del equipo (2,00 x puntos por proyecto), Técnico (1,00 x puntos por proyecto), Titulado Universitario en formación (0,75 x puntos por proyecto), Otros (0,50 x puntos por proyecto)

                3. \textbf{Trabajos presentados en congresos científicos o tecnológicos}:
                - \textbf{Máximo:} 10 puntos
                - \textbf{Tipos de participación y puntuaciones}:
                \vspace{0.4cm}
                \begin{table}[h!]\centering\renewcommand{\arraystretch}{1.3}
                  \begin{tabular}{|p{0.45\textwidth}|p{0.45\textwidth}|}\hline
                    Tipo de Participación                      & Puntos por Participación      \\\hline
                    Participativo-Ponencia plenaria            & 2,00 puntos por participación \\\hline
                    Participativo-Ponencia invitada            & 1,50 puntos por participación \\\hline
                    Participativo-Ponencia oral (comunicación) & 1,00 puntos por participación \\\hline
                    Participativo-Póster                       & 0,75 puntos por participación \\\hline
                  \end{tabular}\end{table}
                \vspace{0.4cm}
                - \textbf{Complemento}:
                - \textbf{Ámbito geográfico}: Internacional (2,00 x puntos por participación), Nacional (1,00 x puntos por participación), Regional (0,50 x puntos por participación), Local (0,25 x puntos por participación)

                4. \textbf{Publicaciones científicas}:
                - \textbf{Máximo:} 100 puntos
                - \textbf{Tipos de publicaciones y puntuaciones}:
                \vspace{0.4cm}
                \begin{table}[h!]\centering\renewcommand{\arraystretch}{1.3}
                  \begin{tabular}{|p{0.45\textwidth}|p{0.45\textwidth}|}\hline
                    Tipo de Publicación             & Puntos por Publicación             \\\hline
                    Artículos en revistas indexadas & Puntuación según baremo específico \\\hline
                    Capítulos de libro              & Puntuación según baremo específico \\\hline
                    Libros                          & Puntuación según baremo específico \\\hline
                  \end{tabular}\end{table}
                \vspace{0.4cm}

                \#\#\# Ponderación del apartado de actividad investigadora:
                - La actividad investigadora tiene una ponderación del \textbf{35\%} en la evaluación global para la plaza de Profesor Ayudante Doctor, según la Tabla 1 de ponderaciones.

                \#\#\# Máximo de puntos en el apartado de actividad investigadora:
                - El máximo de puntos que se puede conseguir en este apartado es \textbf{100 puntos}, considerando todos los criterios de evaluación mencionados.

                En resumen, el máximo de puntos que puedes conseguir en el apartado de actividad investigadora para la plaza de Profesor Ayudante Doctor es \textbf{100 puntos}, distribuidos según los criterios de evaluación y las ponderaciones establecidas en el baremo de 2024.
          \item \textbf{Contexto utilizado:}
                \begin{scriptsize}
                  FUENTE: BAREMO 2024.pdf | PÁGINA: 19
                  \# Núm. 12/2024, de 11 de octubre

                  \#\# I. DISPOSICIONES GENERALES

                  \#\#\# I.5. Secretario General

                  \#\#\#\# 2.4.6. Participación (exclusivamente ponente) en congresos, jornadas o seminarios orientados a la formación teórico-práctica para la actividad docente:
                  - \textbf{Máximo:} 100
                  - \textbf{Ponderación:} 0,10
                  - (máximo 10 aportaciones)

                  \#\#\#\#\# 2.4.6.1. Tipo de participación:
                  \vspace{0.4cm}
                  \begin{table}[h!]\centering\renewcommand{\arraystretch}{1.3}
                    \begin{tabular}{|p{0.45\textwidth}|p{0.45\textwidth}|}\hline
                      Tipo de participación                             & Puntos por participación    \\\hline
                      Participativo - Ponencia plenaria                 & 20 puntos por participación \\\hline
                      Participativo - Ponencia invitada                 & 15 puntos por participación \\\hline
                      Participativo - Ponencia oral (comunicación oral) & 10 puntos por participación \\\hline
                      Participativo - Póster                            & 5 puntos por participación  \\\hline
                    \end{tabular}\end{table}
                  \vspace{0.4cm}

                  \textbf{Complemento 1 (multiplicativo):}
                  - \textbf{Ámbito geográfico del congreso:}
                  - Internacional: 1,50 x participación
                  - Nacional: 1,00 x participación
                  - Regional o Local: 0,50 x participación

                  \textbf{Complemento 2 (multiplicando el Complemento 1):}
                  - Factor de corrección de 1,5 para participación en congresos orientados a la formación docente correspondientes a los últimos 5 años naturales (incluido el año natural de fecha de convocatoria de la plaza).

                  \textbf{Nota 1:} El Departamento deberá especificar en el ANEXO I el número de autores a partir del cual se penaliza la puntuación del aspirante. El factor de corrección por número de autores será el siguiente: (a) de 1 a 3 autores no penaliza = 100 \% de la cuantificación; (b) más de 3 autores = (100 \%) / (nº autores / i], siendo i ≥ 3, y siendo i el número de autores a partir del cual (no incluyendo i) se penaliza. Asimismo, el Departamento deberá especificar si el orden de los autores es relevante, y el coeficiente corrector aplicable que en ningún caso será inferior a 0,8, y únicamente aplicable a las posiciones intermedias en caso de al menos 3 autores.

                  \#\# III. ACTIVIDAD INVESTIGADORA

                  - \textbf{Máximo:} 100 puntos

                  \#\#\# 3.1. Contratos (o becas, si procede) en procesos competitivos convocados por entidades públicas o privadas:
                  - \textbf{Máximo:} 10 puntos

                  \#\#\#\# 3.1.1. Finalidad Predoctoral (nombre de la ayuda):
                  \vspace{0.4cm}
                  \begin{table}[h!]\centering\renewcommand{\arraystretch}{1.3}
                    \begin{tabular}{|p{0.45\textwidth}|p{0.45\textwidth}|}\hline
                      Tipo de ayuda                               & Puntos por tipo de contrato/año     \\\hline
                      - FPU, FPI, Planes propios de Universidad   & 3,0 puntos por tipo de contrato/año \\\hline
                      - Otras convocatorias públicas competitivas & 2,0 puntos por tipo de contrato/año \\\hline
                    \end{tabular}\end{table}
                  \vspace{0.4cm}

                  \#\#\#\# 3.1.2. Finalidad Posdoctoral (nombre de la ayuda):
                  \vspace{0.4cm}
                  \begin{table}[h!]\centering\renewcommand{\arraystretch}{1.3}
                    \begin{tabular}{|p{0.45\textwidth}|p{0.45\textwidth}|}\hline
                      Tipo de ayuda                                                                                                                                                                                           & Puntos por tipo de contrato/año     \\\hline
                      - MEC, Juan de la Cierva, Fullbright, Humboldt, Marie Curie, Contratos de acceso a Sistema Español de Ciencia y Tecnología, o similares de fundaciones e instituciones de reconocido prestigio en I+D+i & 4,0 puntos por tipo de contrato/año \\\hline
                      - Otras convocatorias públicas competitivas                                                                                                                                                             & 3,0 puntos por tipo de contrato/año \\\hline
                    \end{tabular}\end{table}
                  \vspace{0.4cm}

                  \#\#\#\# 3.1.3. Contratos Sénior de reincorporación (nombre de la ayuda):
                  \vspace{0.4cm}
                  \begin{table}[h!]\centering\renewcommand{\arraystretch}{1.3}
                    \begin{tabular}{|p{0.45\textwidth}|p{0.45\textwidth}|}\hline
                      Tipo de ayuda                                                                                                            & Puntos por tipo de contrato/año     \\\hline
                      - Ramón y Cajal, Miguel Servet, Marie Curie, o similares de fundaciones e instituciones de reconocido prestigio en I+D+i & 6,0 puntos por tipo de contrato/año \\\hline
                      - Otras convocatorias públicas competitivas                                                                              & 4,0 puntos por tipo de contrato/año \\\hline
                    \end{tabular}\end{table}
                  \vspace{0.4cm}

                  ---

                  FUENTE: BAREMO 2024.pdf | PÁGINA: 11
                  \# Núm. 12/2024, de 11 de octubre

                  \#\# I. DISPOSICIONES GENERALES

                  \#\#\# I.5. Secretario General

                  \#\# BAREMO DE CONTRATACIÓN DE PROFESORADO AYUDANTE DOCTOR, PROFESORADO ASOCIADO Y PROFESORADO INTERINO

                  \#\#\# CONSIDERACIONES GENERALES
                  El baremo es aplicable a las figuras contractuales de Profesorado Ayudante Doctor, Profesorado Asociado y Profesorado Interino.

                  Apartado 1. Porcentaje de cada bloque temático según figura contractual (en adelante plazas).
                  En el Baremo se distinguen 6 bloques, sus ponderaciones y sus puntuaciones:
                  1. RENDIMIENTO ACADÉMICO
                  2. ACTIVIDAD DOCENTE
                  3. ACTIVIDAD INVESTIGADORA
                  4. TRANSFERENCIA DEL CONOCIMIENTO
                  5. ACTIVIDAD PROFESIONAL EXTERNA A LA UNIVERSIDAD
                  6. MÉRITO PREFERENTE

                  La calificación del candidato se valora según puntuación acumulada atendiendo a las ponderaciones de la Tabla 1, y dependiendo del tipo de contrato. En caso de empate a puntos entre aspirantes la preferencia entre ellos se resolverá en función de la puntuación alcanzada en el Bloque Temático VI. Mérito Preferente. En caso de no resolverse el empate se calculará la media aritmética de la puntuación (sobre 100) de Actividad Docente (Bloque II) y Actividad Investigadora (Bloque III) en las plazas de Ayudante Doctor; en las plazas de Profesor Asociado la media aritmética se calculará con las puntuaciones (sobre 100) de la Actividad Docente (Bloque II) y Actividad Profesional Externa a la Universidad (Bloque V), y en el caso del Profesorado Interino la media aritmética se calculará con las puntuaciones (sobre 100) del Rendimiento Académico (Bloque I) y Actividad Docente (Bloque II).

                  \#\#\# TABLA 1. Ponderación de cada bloque temático según tipo de contrato.

                  \vspace{0.4cm}
                  \begin{table}[h!]\centering\renewcommand{\arraystretch}{1.3}
                    \begin{tabular}{|p{0.45\textwidth}|p{0.45\textwidth}|p{0.45\textwidth}|p{0.45\textwidth}|p{0.45\textwidth}|}\hline
                      BLOQUE TEMÁTICO                                   & TIPO DE CONTRATO/PLAZA (*) & PSI \% & ASO \% & AYD \% \\\hline
                      I. RENDIMIENTO ACADÉMICO                          &                            & 30     & 5      & 10     \\\hline
                      II. ACTIVIDAD DOCENTE                             &                            & 25     & 10     & 30     \\\hline
                      III. ACTIVIDAD INVESTIGADORA                      &                            & 25     & 10     & 35     \\\hline
                      IV. TRANSFERENCIA DEL CONOCIMIENTO                &                            & 5      & 15     & 5      \\\hline
                      V. ACTIVIDAD PROFESIONAL EXTERNA A LA UNIVERSIDAD &                            & 10     & 55     & 5      \\\hline
                      VI. MÉRITO PREFERENTE                             &                            & 5      & 5      & 15     \\\hline
                      TOTAL                                             &                            & 100    & 100    & 100    \\\hline
                    \end{tabular}\end{table}
                  \vspace{0.4cm}

                  (*) PSI: Profesor Interino; ASO: Profesor Asociado; AYD: Profesor Ayudante Doctor.

                  ---

                  FUENTE: BAREMO 2024.pdf | PÁGINA: 4
                  \# Núm. 12/2024, de 11 de octubre

                  \#\# I. DISPOSICIONES GENERALES

                  \#\#\# I.5. Secretario General

                  3. La regulación concreta para la valoración del mérito preferente será la que se establezca en el Baremo o en las normas que, en su caso, se dicten en desarrollo de esta normativa.

                  Por su parte, en el punto VI. Mérito Preferente del Baremo aprobado junto con la normativa de contratación figura:

                  “\textbf{Profesorado Ayudante Doctor:} El aspirante logra 15 puntos atendiendo al siguiente reparto: 10 puntos por su habilitación, acreditación o ser funcionario en el cuerpo docente de Titular de Universidad o Catedrático de Universidad, o hasta 10 puntos (acumulables al punto 1 y hasta 15 puntos) por sus estancias de investigación en universidades u otros centros de investigación externos a la Universidad, hasta 18 meses de estancia, y conforme a las instrucciones dictadas por la Universidad de Sevilla”.

                  Visto lo anterior, procede una actualización tanto del artículo 10.1 de la Normativa de contratación de Profesorado Ayudante Doctor, de Profesorado Asociado y de Profesorado Sustituto como del párrafo 3 del apartado VI del baremo a aplicar.

                  \#\# TÍTULO I

                  \#\#\# Disposiciones generales

                  \textbf{Artículo 1. Objeto de régimen jurídico aplicable.}

                  1. El objeto de la presente normativa es la regulación en la Universidad de Sevilla de los procedimientos de convocatoria y provisión mediante concurso público de las siguientes plazas de personal docente e investigador: Profesorado Ayudante Doctor, Profesorado Asociado y Profesorado Interino.
                  2. Los citados procedimientos se desarrollarán en el marco de las previsiones normativas establecidas en la Ley Orgánica 6/2001, de 21 de diciembre, el Texto Refundido de la Ley Andaluza de Universidades aprobado por Decreto Legislativo 1/2013, el Estatuto de la Universidad de Sevilla, el I Convenio Colectivo del Personal Docente e Investigador en régimen laboral de las Universidades Públicas andaluzas, la presente normativa y, en lo no previsto en estas normas, por el régimen de profesorado que sea de aplicación y por las disposiciones que regulan con carácter general el ingreso en la Administración Pública.
                  3. En materia de procedimiento, y en todo lo no recogido específicamente en esta normativa, se estará a lo dispuesto en la Ley 39/2015, de 1 de octubre, del Procedimiento Administrativo Común de las Administraciones Públicas y en la Ley 40/2015, de 1 de octubre, de Régimen Jurídico del Sector Público.

                  \textbf{Artículo 2. Principios.}

                  1. La selección del personal a que se refiere la presente normativa se efectuará mediante concurso público con respeto a los principios constitucionales de igualdad, mérito, capacidad y publicidad.

                  \#\# TÍTULO II

                  \#\#\# Tramitación previa y actuación de las comisiones de selección

                  \textbf{Artículo 3. Composición, sorteo y nombramiento de las Comisiones de Contratación.}

                  1. La composición y vigencia de las Comisiones de Contratación de las plazas de Profesor Ayudante Doctor y Profesor Asociado serán las establecidas en el artículo 90 del Estatuto de la Universidad de Sevilla.

                  ---

                  FUENTE: BAREMO 2024.pdf | PÁGINA: 23
                  \# Núm. 12/2024, de 11 de octubre

                  \#\# I. DISPOSICIONES GENERALES

                  \#\#\# I.5. Secretario General

                  \#\#\#\# 3.5. Participación en proyectos de investigación
                  financiados ejecutados o en curso, y resueltos por instituciones de investigación acreditadas en nuestra comunidad científica:
                  \textbf{Máximo:} 10 puntos
                  (máximo 5 proyectos con una duración mínima de 1 año cada proyecto)

                  \#\#\#\#\# 3.5.1. Tipo de proyecto:
                  \textbf{No admite coeficiente de afinidad.}

                  \vspace{0.4cm}
                  \begin{table}[h!]\centering\renewcommand{\arraystretch}{1.3}
                    \begin{tabular}{|p{0.45\textwidth}|p{0.45\textwidth}|}\hline
                      Tipo de Proyecto        & Puntos por Proyecto     \\\hline
                      Proyecto competitivo    & 5,0 puntos por proyecto \\\hline
                      Proyecto no competitivo & 2,0 puntos por proyecto \\\hline
                    \end{tabular}\end{table}
                  \vspace{0.4cm}

                  \textbf{Complemento 1 (multiplicativo) – Ámbito geográfico:}

                  \vspace{0.4cm}
                  \begin{table}[h!]\centering\renewcommand{\arraystretch}{1.3}
                    \begin{tabular}{|p{0.45\textwidth}|p{0.45\textwidth}|}\hline
                      Tipo de Proyecto               & Puntos por Proyecto       \\\hline
                      Proyecto internacional (UE)    & 3,0 x puntos por proyecto \\\hline
                      Proyecto Internacional (no UE) & 2,5 x puntos por proyecto \\\hline
                      Proyecto nacional              & 1,5 x puntos por proyecto \\\hline
                      Proyecto autonómico            & 1,0 x puntos por proyecto \\\hline
                      Otros                          & 0,5 x puntos por proyecto \\\hline
                    \end{tabular}\end{table}
                  \vspace{0.4cm}

                  \textbf{Complemento 2 (multiplicando al Complemento 1) – Duración:}

                  \vspace{0.4cm}
                  \begin{table}[h!]\centering\renewcommand{\arraystretch}{1.3}
                    \begin{tabular}{|p{0.45\textwidth}|p{0.45\textwidth}|}\hline
                      Duración                              & Puntos por Proyecto       \\\hline
                      De duración inferior a 2 años         & 1,0 x puntos por proyecto \\\hline
                      De duración igual o superior a 2 años & 2,0 x puntos por proyecto \\\hline
                    \end{tabular}\end{table}
                  \vspace{0.4cm}

                  \textbf{Complemento 3 (multiplicando al Complemento 1) – Grado de contribución:}

                  \vspace{0.4cm}
                  \begin{table}[h!]\centering\renewcommand{\arraystretch}{1.3}
                    \begin{tabular}{|p{0.45\textwidth}|p{0.45\textwidth}|}\hline
                      Grado de Contribución                                  & Puntos por Proyecto        \\\hline
                      Coordinador del proyecto (total, científico o gerente) & 3,00 x puntos por proyecto \\\hline
                      Miembro del equipo (investigador)                      & 2,00 x puntos por proyecto \\\hline
                      Técnico                                                & 1,00 x puntos por proyecto \\\hline
                      Titulado Universitario en formación                    & 0,75 x puntos por proyecto \\\hline
                      Otros                                                  & 0,50 x puntos por proyecto \\\hline
                    \end{tabular}\end{table}
                  \vspace{0.4cm}

                  \textbf{Nota 1:} Se precisa que el carácter Internacional, Nacional, Autonómico u Otro del proyecto científico o tecnológico no se basará en la mera precisión formal que se indique en el nombre del proyecto sino en el ámbito real de la convocatoria, quedando al juicio técnico de la comisión la distinción.

                  \#\#\#\# 3.6. Trabajos presentados en congresos científicos o tecnológicos:
                  \textbf{Máximo:} 10 puntos
                  (máximo 10 participaciones)

                  \#\#\#\#\# 3.6.1. Tipo de participación:

                  \vspace{0.4cm}
                  \begin{table}[h!]\centering\renewcommand{\arraystretch}{1.3}
                    \begin{tabular}{|p{0.45\textwidth}|p{0.45\textwidth}|}\hline
                      Tipo de Participación                      & Puntos por Participación      \\\hline
                      Participativo-Ponencia plenaria            & 2,00 puntos por participación \\\hline
                      Participativo-Ponencia invitada            & 1,50 puntos por participación \\\hline
                      Participativo-Ponencia oral (comunicación) & 1,00 puntos por participación \\\hline
                      Participativo-Póster                       & 0,75 puntos por participación \\\hline
                    \end{tabular}\end{table}
                  \vspace{0.4cm}

                  \textbf{Complemento 1 (multiplicativo): Ámbito geográfico del congreso:}

                  \vspace{0.4cm}
                  \begin{table}[h!]\centering\renewcommand{\arraystretch}{1.3}
                    \begin{tabular}{|p{0.45\textwidth}|p{0.45\textwidth}|}\hline
                      Ámbito Geográfico & Puntos por Participación        \\\hline
                      Internacional     & 2,00 x puntos por participación \\\hline
                      Nacional          & 1,00 x puntos por participación \\\hline
                      Regional          & 0,50 x puntos por participación \\\hline
                      Local             & 0,25 x puntos por participación \\\hline
                    \end{tabular}\end{table}
                  \vspace{0.4cm}

                  \textbf{Complemento 2 (multiplicando al complemento 1): Periodicidad:}

                  ---

                  FUENTE: BAREMO 2024.pdf | PÁGINA: 18
                  \# Núm. 12/2024, de 11 de octubre

                  \#\# I. DISPOSICIONES GENERALES

                  \#\#\# I.5. Secretario General

                  \#\#\#\# 2.4.4. En calidad de autor o coautor de publicaciones docentes o de carácter pedagógico, libros (libros de texto y manuales), artículos, u otros equivalentes de difusión del conocimiento en el dominio docente:
                  - \textbf{Máximo:} 100 (máximo 5 aportaciones en libros, 5 cursos online de al menos 10 horas de docencia no oficial no computados en apartado 2.4.2, 10 aportaciones en capítulos de libro y 10 artículos de revista no computados en el apartado III. Actividad Investigadora)
                  - \textbf{Ponderación:} 0,20

                  \#\#\#\#\# 2.4.4.1. Tipo de soporte:
                  \vspace{0.4cm}
                  \begin{table}[h!]\centering\renewcommand{\arraystretch}{1.3}
                    \begin{tabular}{|p{0.45\textwidth}|p{0.45\textwidth}|}\hline
                      Tipo de soporte                                                                                         & Puntos por unidad       \\\hline
                      Libros o manuales                                                                                       & 40 puntos por libro     \\\hline
                      Capítulos de libro                                                                                      & 10 puntos por capítulo  \\\hline
                      Artículos de revista                                                                                    & 20 puntos por artículo  \\\hline
                      Juegos didácticos, software, etc.                                                                       & 20 puntos por material  \\\hline
                      Actividades relevantes de difusión de la docencia en medios de comunicación audiovisual o equivalentes. & 20 puntos por actividad \\\hline
                    \end{tabular}\end{table}
                  \vspace{0.4cm}

                  \textbf{Complemento (multiplicativo):}
                  - Factor de corrección de 1,5 para materiales docentes publicados (o aceptados) en los últimos 5 años naturales (incluido el año natural de fecha de convocatoria de la plaza).
                  - \textbf{1,5 x material docente}

                  \textbf{Nota 1:} Se valorarán únicamente los libros, capítulos de libro y artículos cuyas editoriales o revistas se correspondan con los admitidos en los apartados 3.2, 3.3 y 3.4, respectivamente, de este Baremo, con la aclaración de que en este caso las publicaciones deben tener carácter docente y pedagógico. No se valora ningún mérito que no acredite ISBN o depósito legal comprobable mediante las bases de datos oficiales para tales fines. Los tipos de materiales (i.e., cursos online, juegos didácticos, software, difusión, etc.) que no puedan acreditar identificador por su naturaleza deberán acreditar algún otro identificador de publicación o aceptación oficialmente validado, o indicios verificables por la comisión de contratación.

                  \textbf{Nota 2:} El Departamento deberá especificar en el ANEXO I el número de autores a partir del cual se penaliza la puntuación del aspirante. El factor de corrección por número de autores será el siguiente:
                  - (a) de 1 a 3 autores no penaliza = 100 \% de la cuantificación;
                  - (b) más de 3 autores = (100 \%) / (nº autores / i], siendo i ≥ 3, y siendo i el número de autores a partir del cual (no incluyendo i) se penaliza. Asimismo, el Departamento deberá especificar si el orden de los autores es relevante, y el coeficiente corrector aplicable que en ningún caso será inferior a 0,8, y únicamente aplicable a las posiciones intermedias en caso de al menos 3 autores.

                  \#\#\#\# 2.4.5. Curso de formación e innovación docente:
                  - \textbf{Máximo:} 100 (máximo 100 horas de formación e innovación docente)
                  - \textbf{Ponderación:} 0,10

                  \#\#\#\#\# 2.4.5.1. Entidad organizadora:
                  \vspace{0.4cm}
                  \begin{table}[h!]\centering\renewcommand{\arraystretch}{1.3}
                    \begin{tabular}{|p{0.45\textwidth}|p{0.45\textwidth}|}\hline
                      Tipo de curso                                                                & Puntos por 10 h. docencia no oficial    \\\hline
                      Cursos organizados por Universidades o Centros de investigación de prestigio & 10 puntos por 10 h. docencia no oficial \\\hline
                      Cursos organizados por otras entidades                                       & 5 puntos por 10 h. docencia no oficial  \\\hline
                    \end{tabular}\end{table}
                  \vspace{0.4cm}

                  \textbf{Complemento (multiplicativo):}
                  - Factor de corrección de 1,5 para cursos de formación e innovación docente correspondientes a los últimos 5 cursos académicos (incluido el curso académico de fecha de convocatoria de la plaza).
                  - \textbf{1,5 x 10 h. docencia no oficial}
                \end{scriptsize}
        \end{itemize}


  \item \textbf{P18.} En el baremo de 2024, ¿con qué coeficiente de afinidad se valoran automáticamente los artículos publicados en revistas indexadas en JCR Q1 y Q2? \\
        \begin{itemize}
          \item \textbf{Tipo de consulta:} simple
          \item \textbf{Respuesta del Agente:}

                No existen requisitos mínimos excluyentes mencionados en el contexto proporcionado.

                A continuación, se detallan los criterios de evaluación y puntuación para los artículos publicados en revistas indexadas en JCR Q1 y Q2:

                \#\#\# Criterios de evaluación para artículos científicos de revista

                \vspace{0.4cm}
                \begin{table}[h!]\centering\renewcommand{\arraystretch}{1.3}
                  \begin{tabular}{|p{0.45\textwidth}|p{0.45\textwidth}|}\hline
                    Criterio                   & Detalles                                                                                                                                                                                 \\\hline
                    Posición de la publicación & Se valoran según su posición en los índices de impacto (JCR, SJR, ANEXO I)                                                                                                               \\\hline
                    Número de autores          & Se aplica un factor de corrección por número de autores, que puede penalizar la puntuación                                                                                               \\\hline
                    Afinidad                   & Se aplica un coeficiente de afinidad según la relación del artículo con el perfil investigador de la plaza convocada o las líneas estratégicas de investigación del área de conocimiento \\\hline
                  \end{tabular}\end{table}
                \vspace{0.4cm}

                \#\#\# Puntuación para artículos científicos de revista

                \vspace{0.4cm}
                \begin{table}[h!]\centering\renewcommand{\arraystretch}{1.3}
                  \begin{tabular}{|p{0.45\textwidth}|p{0.45\textwidth}|}\hline
                    Posición de la publicación                       & Puntos por artículo      \\\hline
                    Decil 1 del JCR                                  & 12,0 puntos por artículo \\\hline
                    Q1 del JCR o SJR                                 & 10,0 puntos por artículo \\\hline
                    Q2 del JCR o SJR, o ANEXO I (nivel I)            & 8,0 puntos por artículo  \\\hline
                    Q3 del JCR o SJR, o ANEXO I (nivel II)           & 6,0 puntos por artículo  \\\hline
                    Q4 del JCR o SJR, o ANEXO I (nivel III)          & 4,0 puntos por artículo  \\\hline
                    Artículos no incluidos ni en JCR, SJR ni ANEXO I & 0,5 puntos por artículo  \\\hline
                  \end{tabular}\end{table}
                \vspace{0.4cm}

                \#\#\# Coeficiente de afinidad

                \vspace{0.4cm}
                \begin{table}[h!]\centering\renewcommand{\arraystretch}{1.3}
                  \begin{tabular}{|p{0.45\textwidth}|p{0.45\textwidth}|p{0.45\textwidth}|}\hline
                    Tipo de Afinidad         & Coeficiente & Descripción                                                                                                            \\\hline
                    Afinidad muy alta        & 1,0         & Méritos correspondientes al perfil investigador de la plaza convocada                                                  \\\hline
                    Afinidad alta            & 0,75        & Méritos relativos a las líneas de investigación del área de conocimiento a que corresponde la plaza                    \\\hline
                    Afinidad media           & 0,50        & Méritos relativos a las áreas de conocimiento afines al departamento y área de conocimiento a que corresponde la plaza \\\hline
                    Afinidad baja o muy baja & 0,25 o 0,0  & Méritos relativos a alguna de las líneas de investigación de otras áreas de conocimiento                               \\\hline
                  \end{tabular}\end{table}
                \vspace{0.4cm}

                En cuanto a la pregunta del usuario, los artículos publicados en revistas indexadas en JCR Q1 y Q2 se valoran con un coeficiente de afinidad que depende de la relación del artículo con el perfil investigador de la plaza convocada o las líneas estratégicas de investigación del área de conocimiento. Sin embargo, no se especifica un coeficiente de afinidad automático para estos artículos. El coeficiente de afinidad se determina según la tipología de afinidad establecida en el baremo, que se aplica según la relación del artículo con el perfil investigador de la plaza convocada o las líneas estratégicas de investigación del área de conocimiento.

                Por lo tanto, no se puede determinar un coeficiente de afinidad automático para los artículos publicados en revistas indexadas en JCR Q1 y Q2 sin conocer la relación del artículo con el perfil investigador de la plaza convocada o las líneas estratégicas de investigación del área de conocimiento.
          \item \textbf{Contexto utilizado:}
                \begin{scriptsize}
                  FUENTE: BAREMO 2024.pdf | PÁGINA: 21
                  \# Núm. 12/2024, de 11 de octubre

                  \#\# I. DISPOSICIONES GENERALES

                  \#\#\# I.5. Secretario General

                  - Afinidad media: se aplicará el punto 2) a cada área de conocimiento afín (sin distinción de universidad) según listado oficial de la Universidad de Sevilla de comisiones juzgadoras a cuerpos docentes. No deben incorporarse las líneas al ANEXO I, siendo motivadas durante el proceso de baremación. El aspirante deberá motivar entre sus indicios de calidad las líneas estratégicas de investigación a que queda adscrito el resultado alcanzado con su mérito. Debe incluir su motivación en el campo (de su CVN Editor de FECYT): Resultados relevantes del apartado Actividad científica y tecnológica - Publicaciones, documentos científicos y técnicos.

                  \textbf{Nota 4:} El Departamento deberá especificar en el ANEXO I el número de autores a partir del cual se penaliza la puntuación del aspirante. El factor de corrección por número de autores será el siguiente: (a) de 1 a 3 autores no penaliza = 100 \% de la cuantificación; (b) más de 3 autores = (100 \%) / (n.º autores / i], siendo i ≥ 3, y siendo i el número de autores a partir del cual (no incluyendo i) se penaliza. Asimismo, el Departamento deberá especificar si el orden de los autores es relevante, y el coeficiente corrector aplicable que en ningún caso será inferior a 0,8, y únicamente aplicable a las posiciones intermedias en caso de al menos 3 autores.

                  \#\#\# 3.3. Capítulos de Libro: Máximo: 10 puntos
                  (máximo 10 capítulos)

                  \#\#\#\# 3.3.1. Posición de la publicación:

                  \vspace{0.4cm}
                  \begin{table}[h!]\centering\renewcommand{\arraystretch}{1.3}
                    \begin{tabular}{|p{0.45\textwidth}|p{0.45\textwidth}|}\hline
                      Posición de la publicación                     & Puntos por capítulo de libro     \\\hline
                      Q1 en el SPI                                   & 4,0 puntos por capítulo de libro \\\hline
                      Q2 en el SPI, o ANEXO I (nivel I, N1)          & 3,0 puntos por capítulo de libro \\\hline
                      Q3 en el SPI, o ANEXO I (nivel II, N2)         & 2,0 puntos por capítulo de libro \\\hline
                      Q4 en el SPI, o ANEXO I (nivel III, N3)        & 1,0 punto por capítulo de libro  \\\hline
                      Otros libros no incluidos ni en SPI ni ANEXO I & 0,5 puntos por capítulo de libro \\\hline
                    \end{tabular}\end{table}
                  \vspace{0.4cm}

                  \textbf{Complemento 1 (multiplicativo):}

                  \vspace{0.4cm}
                  \begin{table}[h!]\centering\renewcommand{\arraystretch}{1.3}
                    \begin{tabular}{|p{0.45\textwidth}|p{0.45\textwidth}|}\hline
                      Descripción                                                                                                   & Puntos por capítulo de libro       \\\hline
                      Fichas técnicas de catálogos en publicaciones, reseñas, o aportaciones en diccionarios comunes o científicos. & 0,1 x puntos por capítulo de libro \\\hline
                    \end{tabular}\end{table}
                  \vspace{0.4cm}

                  \textbf{Nota 1:} Se aplican las Notas del apartado 3.2. Libros.
                  \textbf{Nota 2:} En particular, no se considerarán capítulos de libros las contribuciones a congresos editadas en libros de actas o proceedings, o aquellas aportaciones cuyo contenido sea esencialmente tesis doctorales (salvo, y como excepción a juicio de la comisión de contratación, las publicadas en editoriales de prestigio según Nota 1), las recopilaciones legislativas y jurisprudenciales así como sus combinaciones, trabajos fin de estudios, o trabajos similares.

                  \#\#\# 3.4. Artículos científicos de Revista con proceso anónimo de revisión por pares: Máximo: 40 puntos
                  (máximo 10 artículos)

                  \#\#\#\# 3.4.1. Posición de la publicación:

                  \vspace{0.4cm}
                  \begin{table}[h!]\centering\renewcommand{\arraystretch}{1.3}
                    \begin{tabular}{|p{0.45\textwidth}|p{0.45\textwidth}|}\hline
                      Posición de la publicación                       & Puntos por artículo      \\\hline
                      Decil 1 del JCR                                  & 12,0 puntos por artículo \\\hline
                      Q1 del JCR o SJR                                 & 10,0 puntos por artículo \\\hline
                      Q2 del JCR o SJR, o ANEXO I (nivel I)            & 8,0 puntos por artículo  \\\hline
                      Q3 del JCR o SJR, o ANEXO I (nivel II)           & 6,0 puntos por artículo  \\\hline
                      Q4 del JCR o SJR, o ANEXO I (nivel III)          & 4,0 puntos por artículo  \\\hline
                      Artículos no incluidos ni en JCR, SJR ni ANEXO I & 0,5 puntos por artículo  \\\hline
                    \end{tabular}\end{table}
                  \vspace{0.4cm}

                  \textbf{Complemento (aditivo, y hasta un máximo de 5 puntos):}

                  ---

                  FUENTE: 27831\_instrucciones-baremacinin-r03.pdf | PÁGINA: 28
                  \# Bareprof

                  \#\# Instrucciones de baremación

                  \#\#\# 4.4.3 Capítulos de Libro (epígrafe 3.3 del baremo)

                  - No se consideran capítulos de libros las contribuciones a congresos editadas en libros de actas o proceedings.
                  - Además de los capítulos de libros propiamente dichos, el baremo incluye en este apartado fichas técnicas de catálogos en publicaciones, reseñas y aportaciones en diccionarios comunes o científicos, con una valoración del 10\% de la que corresponde al capítulo de libro. Si la aportación corresponde a uno de los tipos citados, pero el candidato ha indicado como tipo de producción “Capítulo de libro”, la Comisión no debe validar el mérito.
                  - El editor de CVN no incorpora ninguna opción para indicar como tipo de producción las fichas técnicas de catálogos en publicaciones. En la Guía del Ayuda al Candidato se indica que, en tales casos, debe seleccionar “Reseña” en el campo Tipo de Producción. También sería válido si seleccionaran “Diccionario científico” o “Diccionario común”, pues la valoración es la misma, pero no así “Capítulo de libro”; en este último caso, el mérito no debe validarse.
                  - Para la asignación de la afinidad al mérito, consulte las consideraciones realizadas en la introducción al apartado 4.4.

                  \#\#\# 4.4.4 Artículos científicos de Revista (epígrafe 3.4 del baremo)

                  - Los datos de citas de un artículo científico deben ser introducidos por el candidato. La aplicación Bareprof no comprueba tales datos de citas.
                  - El baremo no establece distinción entre las fuentes de citas, valorando el mayor número de citas validado por la comisión.
                  - La aplicación Bareprof comprueba automáticamente la posición del artículo en los distintos índices de impacto considerados en el baremo. Para que este proceso sea efectivo, el candidato ha debido introducir correctamente tanto el ISSN de la revista como el año de publicación del artículo. La comisión está facultada para modificar el ISSN en caso de que se compruebe que el identificador introducido por el candidato es incorrecto.
                  - En pantalla se mostrará la posición de la revista en los diferentes listados automatizados (JCR, SJR o Anexo I). Para cada índice, se utilizará el listado que corresponde al año de publicación del artículo o, en su defecto, al año inmediatamente anterior si el primero no se encuentra aún disponible a la finalización del periodo de admisión de solicitudes. Bareprof usará la mejor posición para realizar la baremación.
                  - Para la asignación de la afinidad, consulte las consideraciones realizadas en la introducción al apartado 4.4.

                  \#\#\# 4.4.5 Participación en proyectos de investigación (epígrafe 3.5 del baremo)

                  - La pertenencia a grupos de investigación no es un mérito considerado en el Baremo. Muchos candidatos incluyen estos méritos como proyectos de investigación, que no deben ser validados por la Comisión.
                  - Los méritos englobados en este epígrafe del baremo no admiten coeficiente de afinidad, por lo que la Comisión NO puede invalidarlos alegando falta de afinidad con la plaza. Si el mérito ha sido consignado en el apartado correcto y se encuentra acreditado documentalmente, la Comisión debe validarlo. Por este mismo motivo, no es preciso que el candidato justifique la afinidad del.

                  ----

                  Revisión 3 (diciembre, 2024) 27

                  ---

                  FUENTE: BAREMO 2024.pdf | PÁGINA: 22
                  \# Núm. 12/2024, de 11 de octubre

                  \#\# I. DISPOSICIONES GENERALES

                  \#\#\# I.5. Secretario General

                  Asimismo, y sin perjuicio de lo establecido en el primer párrafo de la Nota 1, el Departamento puede elaborar un listado complementario incluido en su ANEXO I de un máximo de 10 revistas (identificadas por su ISSN) que no estando contenidas en los listados JCR y SJR de aplicación en el año de la publicación del mérito (o aceptación si procede) necesariamente cuenten con sello de calidad FECYT, o estén indexadas en bases de datos citadas expresa y literalmente (independientemente de a qué rama o campo de conocimiento se haga referencia) en el último documento publicado por la Agencia Nacional de Evaluación de la Calidad y Acreditación (ANECA): PRINCIPIOS Y ORIENTACIONES PARA LA APLICACIÓN DE LOS CRITERIOS DE EVALUACIÓN (se aplica en todos los casos el documento de Profesor Ayudante Doctor), y también el último documento publicado por la Comisión Nacional Evaluadora de la Actividad Investigadora (CNEAI) para la evaluación de la actividad investigadora de los profesores universitarios: LOS CRITERIOS ESPECÍFICOS APROBADOS PARA CADA UNO DE LOS CAMPOS DE EVALUACIÓN.

                  Finalmente se indica expresamente que en caso de que la comisión de contratación valore como ‘revista seudocientífica o depredadora’ el mérito aportado por el aspirante, previa motivación de la comisión de contratación, deberá puntuarse con coeficiente de afinidad igual a 0.

                  \textbf{Nota 2:} En relación con el coeficiente de afinidad aplicable se considera la siguiente tipología:
                  - \textbf{Afinidad muy alta (1,0):} méritos correspondientes al perfil investigador de la plaza convocada, o en su defecto a las líneas estratégicas de investigación del área de conocimiento del departamento de la Universidad de Sevilla solicitante de la plaza, pudiéndose completar en el Anexo I.
                  - \textbf{Afinidad alta (0,75):} méritos relativos a las líneas de investigación del área de conocimiento a que corresponde la plaza –sin consideración de universidad de adscripción del área de conocimiento, pudiéndose completar en el Anexo I.
                  - \textbf{Afinidad media (0,50):} méritos relativos a las áreas de conocimiento (sin distinción de universidad) afines al departamento y área de conocimiento a que corresponde la plaza según listado oficial de la Universidad de Sevilla de comisiones juzgadoras a cuerpos docentes.
                  - \textbf{Afinidad baja o muy baja (0,25 o 0,0):} méritos relativos a alguna de las líneas de investigación de otras áreas de conocimiento. El coeficiente de afinidad aplicable deberá ser 0,25 (baja) o 0,0 (muy baja), motivándose por la comisión de contratación. Los artículos publicados en revista seudocientífica o depredadora son merecedores en todos los casos de un coeficiente de afinidad 0,0. La comisión de contratación deberá motivar de modo preciso su consideración de seudocientífica o depredadora.

                  ---

                  FUENTE: BAREMO 2024.pdf | PÁGINA: 24
                  \# Núm. 12/2024, de 11 de octubre

                  \#\# I. DISPOSICIONES GENERALES

                  \#\#\# I.5. Secretario General

                  Nota 1: Se valoran los congresos que respeten el último documento publicado por la Agencia Nacional de Evaluación de la Calidad y Acreditación (ANECA): PRINCIPIOS Y ORIENTACIONES PARA LA APLICACIÓN DE LOS CRITERIOS DE EVALUACIÓN (se aplica en todos los casos el documento de Profesor Ayudante Doctor). Asimismo, se precisa que el carácter Internacional, Nacional, Regional o Local del Congreso científico o tecnológico no se basará en la mera precisión formal que se indique en el nombre del congreso sino en su ámbito real de participación, quedando al juicio técnico de la comisión la distinción.

                  Nota 2: En relación con el coeficiente de afinidad aplicable se considera la siguiente tipología:

                  \vspace{0.4cm}
                  \begin{table}[h!]\centering\renewcommand{\arraystretch}{1.3}
                    \begin{tabular}{|p{0.45\textwidth}|p{0.45\textwidth}|p{0.45\textwidth}|}\hline
                      Tipo de Afinidad         & Coeficiente & Descripción                                                                                                                                                                                                                                                                                                                                                                                                                                                   \\\hline
                      Afinidad muy alta        & 1,0         & Méritos correspondientes al perfil investigador de la plaza convocada, o en su defecto a las líneas estratégicas de investigación del área de conocimiento del departamento de la Universidad de Sevilla solicitante de la plaza, pudiéndose completar en el Anexo I.                                                                                                                                                                                         \\\hline
                      Afinidad alta            & 0,75        & Méritos relativos a las líneas de investigación del área de conocimiento a que corresponde la plaza –sin consideración de universidad de adscripción del área de conocimiento, pudiéndose completar en el Anexo I.                                                                                                                                                                                                                                            \\\hline
                      Afinidad media           & 0,50        & Méritos relativos a las áreas de conocimiento (sin distinción de universidad) afines al departamento y área de conocimiento a que corresponde la plaza según listado oficial de la Universidad de Sevilla de comisiones juzgadoras a cuerpos docentes.                                                                                                                                                                                                        \\\hline
                      Afinidad baja o muy baja & 0,25 o 0,0  & Méritos relativos a alguna de las líneas de investigación de otras áreas de conocimiento. El coeficiente de afinidad aplicable deberá ser 0,25 (baja) o 0,0 (muy baja), motivándose por la comisión de contratación. Los congresos seudocientíficos o depredadores son merecedores en todos los casos de un coeficiente de afinidad igual a 0,0. La comisión de contratación deberá motivar de modo preciso su consideración de seudocientífico o depredador. \\\hline
                    \end{tabular}\end{table}
                  \vspace{0.4cm}

                  Las líneas de investigación se describen del modo siguiente:

                  ---

                  FUENTE: BAREMO 2024.pdf | PÁGINA: 18
                  \# Núm. 12/2024, de 11 de octubre

                  \#\# I. DISPOSICIONES GENERALES

                  \#\#\# I.5. Secretario General

                  \#\#\#\# 2.4.4. En calidad de autor o coautor de publicaciones docentes o de carácter pedagógico, libros (libros de texto y manuales), artículos, u otros equivalentes de difusión del conocimiento en el dominio docente:
                  - \textbf{Máximo:} 100 (máximo 5 aportaciones en libros, 5 cursos online de al menos 10 horas de docencia no oficial no computados en apartado 2.4.2, 10 aportaciones en capítulos de libro y 10 artículos de revista no computados en el apartado III. Actividad Investigadora)
                  - \textbf{Ponderación:} 0,20

                  \#\#\#\#\# 2.4.4.1. Tipo de soporte:
                  \vspace{0.4cm}
                  \begin{table}[h!]\centering\renewcommand{\arraystretch}{1.3}
                    \begin{tabular}{|p{0.45\textwidth}|p{0.45\textwidth}|}\hline
                      Tipo de soporte                                                                                         & Puntos por unidad       \\\hline
                      Libros o manuales                                                                                       & 40 puntos por libro     \\\hline
                      Capítulos de libro                                                                                      & 10 puntos por capítulo  \\\hline
                      Artículos de revista                                                                                    & 20 puntos por artículo  \\\hline
                      Juegos didácticos, software, etc.                                                                       & 20 puntos por material  \\\hline
                      Actividades relevantes de difusión de la docencia en medios de comunicación audiovisual o equivalentes. & 20 puntos por actividad \\\hline
                    \end{tabular}\end{table}
                  \vspace{0.4cm}

                  \textbf{Complemento (multiplicativo):}
                  - Factor de corrección de 1,5 para materiales docentes publicados (o aceptados) en los últimos 5 años naturales (incluido el año natural de fecha de convocatoria de la plaza).
                  - \textbf{1,5 x material docente}

                  \textbf{Nota 1:} Se valorarán únicamente los libros, capítulos de libro y artículos cuyas editoriales o revistas se correspondan con los admitidos en los apartados 3.2, 3.3 y 3.4, respectivamente, de este Baremo, con la aclaración de que en este caso las publicaciones deben tener carácter docente y pedagógico. No se valora ningún mérito que no acredite ISBN o depósito legal comprobable mediante las bases de datos oficiales para tales fines. Los tipos de materiales (i.e., cursos online, juegos didácticos, software, difusión, etc.) que no puedan acreditar identificador por su naturaleza deberán acreditar algún otro identificador de publicación o aceptación oficialmente validado, o indicios verificables por la comisión de contratación.

                  \textbf{Nota 2:} El Departamento deberá especificar en el ANEXO I el número de autores a partir del cual se penaliza la puntuación del aspirante. El factor de corrección por número de autores será el siguiente:
                  - (a) de 1 a 3 autores no penaliza = 100 \% de la cuantificación;
                  - (b) más de 3 autores = (100 \%) / (nº autores / i], siendo i ≥ 3, y siendo i el número de autores a partir del cual (no incluyendo i) se penaliza. Asimismo, el Departamento deberá especificar si el orden de los autores es relevante, y el coeficiente corrector aplicable que en ningún caso será inferior a 0,8, y únicamente aplicable a las posiciones intermedias en caso de al menos 3 autores.

                  \#\#\#\# 2.4.5. Curso de formación e innovación docente:
                  - \textbf{Máximo:} 100 (máximo 100 horas de formación e innovación docente)
                  - \textbf{Ponderación:} 0,10

                  \#\#\#\#\# 2.4.5.1. Entidad organizadora:
                  \vspace{0.4cm}
                  \begin{table}[h!]\centering\renewcommand{\arraystretch}{1.3}
                    \begin{tabular}{|p{0.45\textwidth}|p{0.45\textwidth}|}\hline
                      Tipo de curso                                                                & Puntos por 10 h. docencia no oficial    \\\hline
                      Cursos organizados por Universidades o Centros de investigación de prestigio & 10 puntos por 10 h. docencia no oficial \\\hline
                      Cursos organizados por otras entidades                                       & 5 puntos por 10 h. docencia no oficial  \\\hline
                    \end{tabular}\end{table}
                  \vspace{0.4cm}

                  \textbf{Complemento (multiplicativo):}
                  - Factor de corrección de 1,5 para cursos de formación e innovación docente correspondientes a los últimos 5 cursos académicos (incluido el curso académico de fecha de convocatoria de la plaza).
                  - \textbf{1,5 x 10 h. docencia no oficial}
                \end{scriptsize}
        \end{itemize}


  \item \textbf{P19.} Dentro de la experiencia docente universitaria en el baremo, ¿cuántos puntos se otorgan por cada hora impartida para el puesto de Ayudante Doctor en el límite general? \\
        \begin{itemize}
          \item \textbf{Tipo de consulta:} simple
          \item \textbf{Respuesta del Agente:}

                \#\#\# Requisitos Mínimos Excluyentes

                * No se mencionan requisitos mínimos excluyentes específicos para la evaluación de la experiencia docente universitaria en el baremo proporcionado.

                \#\#\# Criterios de Evaluación para Experiencia Docente Universitaria

                \vspace{0.4cm}
                \begin{table}[h!]\centering\renewcommand{\arraystretch}{1.3}
                  \begin{tabular}{|p{0.45\textwidth}|p{0.45\textwidth}|}\hline
                    Criterio                                                       & Detalles         \\\hline
                    Experiencia Docente                                            & Máximo 50 puntos \\\hline
                    2.1.1. Dedicación Docente                                      & 25-35 puntos     \\\hline
                    2.1.2. Pluralidad, Interdisciplinariedad y Complejidad Docente & 0-6 puntos       \\\hline
                    2.1.3. Recursos Educativos                                     & 0-5 puntos       \\\hline
                    2.1.4. Actividades de Formación a lo Largo de la Vida          & 0-5 puntos       \\\hline
                  \end{tabular}\end{table}
                \vspace{0.4cm}

                \#\#\# Puntos por Hora Impartida para Ayudante Doctor

                Según el baremo, para el puesto de Ayudante Doctor, se otorgan \textbf{1 punto por cada 10 horas de docencia}. Esto se calcula considerando que el máximo de puntos para la experiencia docente es 50 y se consideran hasta 1.200 horas de docencia oficial. Por lo tanto, el cálculo es el siguiente:

                - \textbf{1 punto / 10 horas} = \textbf{0,1 puntos por hora}.

                Sin embargo, es importante tener en cuenta que este cálculo se basa en la interpretación de la información proporcionada y podría requerir ajustes según la normativa específica o contextos adicionales no mencionados en el texto proporcionado.

                \#\#\# Complementos a la Puntuación

                Además, existen complementos que pueden multiplicar la puntuación obtenida por la experiencia docente, como:

                - \textbf{Complemento 1 (multiplicativo):} Docencia en Máster Universitario (oficial) puede multiplicar la puntuación por 1,50.
                - \textbf{Complemento 2 (multiplicando al complemento 1):} Docencia en inglés u otras lenguas diferentes a las que son objeto de estudio puede multiplicar la puntuación por 1,50.

                Estos complementos pueden aumentar significativamente la puntuación final por la experiencia docente, dependiendo de las características específicas de la docencia impartida.
          \item \textbf{Contexto utilizado:}
                \begin{scriptsize}
                  FUENTE: BAREMO 2024.pdf | PÁGINA: 15
                  \# Núm. 12/2024, de 11 de octubre

                  \#\# I. DISPOSICIONES GENERALES

                  \#\#\# I.5. Secretario General

                  1.3. Otros títulos académicos:                                                            \textbf{Máximo: 10 puntos}
                  - 1.3.1. Titulación Universitaria de Grado o equivalente (como Licenciado/a, Diplomado/a, Arquitecto/a, Arquitecto/a técnico/a, Ingeniero/a o Ingeniero/a técnico/a), y Titulación Universitaria de Máster (Oficial):
                  - APROBADO: 5,0-6,9 nota media                            \textbf{3 puntos}
                  - NOTABLE: 7,0-8,9 nota media                             \textbf{6 puntos}
                  - SOBRESALIENTE: 9,0-10,0 nota media                      \textbf{10 puntos}
                  - 1.3.2. Título de doctor (tesis doctoral)                  \textbf{10 puntos}

                  \textbf{Nota 1:} En relación con el coeficiente de afinidad y la homologación de títulos se aplicarán las Notas del apartado 1.1.

                  1.4. Premios académicos y menciones de calidad:                                           \textbf{Máximo: 5 puntos}
                  - 1.4.1. Premio (Nacional) Fin de Carrera de Educación Universitaria: Primer premio           \textbf{5 puntos por premio}           No admite coeficiente de afinidad.
                  - 1.4.2. Premio extraordinario de grado o equivalente        \textbf{4 puntos por premio}
                  - 1.4.3. Premio Extraordinario de Doctorado                  \textbf{4 puntos por premio}
                  - 1.4.4. Doctorado internacional y europeo                   \textbf{2 puntos por indicio}
                  - 1.4.5. Mención de calidad (o ‘hacia la excelencia’)        \textbf{1 punto por indicio}
                  - 1.4.6. Tesis en régimen de cotutela                        \textbf{1 punto por indicio}

                  \#\# II. ACTIVIDAD DOCENTE                                                                     \textbf{Máximo: 100 puntos}

                  2.1. Docencia en asignaturas regladas universitarias:                                     \textbf{Máximo: 50 puntos}
                  (máximo 1.200 horas de docencia oficial)
                  - 2.1.1. Titulación Universitaria de Grado o equivalente (como Licenciado/a, Diplomado/a, Arquitecto/a, Arquitecto/a técnico/a, Ingeniero/a o Ingeniero/a técnico/a), y Titulación Universitaria de Máster (Oficial):
                  \textbf{1,00 x 10 horas de docencia}
                  En el caso de que el aspirante acredite docencia superior a 1.200 horas, su puntuación se escalará proporcionalmente a 1.200 horas.

                  \textbf{Complemento 1 (multiplicativo):}                             \textbf{1,50 x puntuación obtenida}
                  Docencia en Máster Universitario (oficial)                  en 2.1.1 en asignaturas de Máster Universitario (oficial)

                  \textbf{Complemento 2 (multiplicando al complemento 1):}
                  Docencia en inglés u otras lenguas diferentes a las que son objeto de estudio                                      \textbf{1,50 x puntuación obtenida}
                  en 2.1.1 en asignaturas en inglés u otras lenguas diferentes a las que son objeto de estudio

                  \textbf{Nota 1:} Los Trabajos Fin de Carrera (Grado o equivalente, y Máster) no serán evaluados en este apartado. Se remite al aspirante al apartado 2.4.1.
                  \textbf{Nota 2:} Se considera la siguiente tipología de asignaturas:
                  - Afinidad muy alta (1,0): asignaturas cuyo contenido corresponda al perfil de la plaza convocada, o en su defecto las asignaturas impartidas por el área de conocimiento del departamento solicitante de la plaza.
                  - Afinidad alta  (0,75): asignaturas impartidas por el área de conocimiento solicitante sin consideración de universidad de impartición, e incluyendo con ello las asignaturas no necesariamente impartidas por el área de conocimiento del departamento de la Universidad de Sevilla convocante de la plaza.
                  - Afinidad media  (0,50): asignaturas impartidas por áreas de conocimiento afines al departamento y área de conocimiento a que corresponde la plaza según listado oficial de la Universidad de Sevilla de comisiones juzgadoras a cuerpos docentes, y sin consideración de universidad de impartición.
                  - Afinidad baja o muy baja      (0,25 o 0,0): asignaturas no pertenecientes a ninguno de los tipos anteriores. El coeficiente de afinidad aplicable deberá ser 0,0 o 0,25, motivándose por la comisión de contratación.

                  ---

                  FUENTE: CRITERIOS DE EVALUACIÓN Y REQUISITOS MÍNIMOS.pdf | PÁGINA: 10
                  \# CATEDRÁTICA O CATEDRÁTICO DE UNIVERSIDAD (CU)

                  \vspace{0.4cm}
                  \begin{table}[h!]\centering\renewcommand{\arraystretch}{1.3}
                    \begin{tabular}{|p{0.45\textwidth}|p{0.45\textwidth}|}\hline
                      Mínimo & Máximo \\\hline
                      50     & 100    \\\hline
                      puntos & puntos \\\hline
                    \end{tabular}\end{table}
                  \vspace{0.4cm}

                  \#\# 2. ACTIVIDAD DOCENTE

                  \vspace{0.4cm}
                  \begin{table}[h!]\centering\renewcommand{\arraystretch}{1.3}
                    \begin{tabular}{|p{0.45\textwidth}|p{0.45\textwidth}|p{0.45\textwidth}|}\hline
                      Actividad                                                           & Mínimo & Máximo \\\hline
                      2.1. Experiencia docente                                            & 25     & 45     \\\hline
                      2.1.1. Dedicación docente                                           & 25     & 35     \\\hline
                      2.1.2. Pluralidad, interdisciplinariedad y complejidad docente      & 0      & 6      \\\hline
                      2.1.3. Recursos educativos                                          & 0      & 5      \\\hline
                      2.1.4. Actividades de formación a lo largo de la vida               & 0      & 5      \\\hline
                      2.2. Calidad de la actividad docente e innovación                   & 5      & 40     \\\hline
                      2.2.1. Calidad de la actividad docente                              & 5      & 15     \\\hline
                      2.2.2. Proyectos de innovación docente                              & 0      & 20     \\\hline
                      2.2.3. Formación para la mejora docente recibida                    & 0      & 5      \\\hline
                      2.2.4. Formación para la mejora docente impartida                   & 0      & 15     \\\hline
                      2.3. Tutorización docente                                           & 5      & 15     \\\hline
                      2.3.1. Tutorización reglada                                         & 0      & 15     \\\hline
                      2.3.2. Otros tipos de tutorización (no incluida en encargo docente) & 0      & 15     \\\hline
                      2.4. Otros méritos*                                                 & 0      & 10     \\\hline
                    \end{tabular}\end{table}
                  \vspace{0.4cm}

                  * Los puntos obtenidos en este apartado, de carácter opcional, se suman a los de los apartados anteriores, pero en ningún caso podrá superarse el máximo (100 puntos).

                  \#\# 2.1. EXPERIENCIA DOCENTE

                  \#\#\# 2.1.1. Dedicación docente

                  En el caso de PTU deberá acreditarse experiencia docente durante un mínimo de 4 años de dedicación a tiempo completo, o su equivalente a tiempo parcial, y, además:

                  - En el caso de figuras de contratados predoctorales y postdoctorales, incluyendo el periodo de orientación postdoctoral, un mínimo de 12 ECTS o 120 horas impartidas. El disfrute de un contrato postdoctoral se considerará dedicación a tiempo completo. El disfrute de una ayuda predoctoral que tenga asignada capacidad docente tutelada y se cubra totalmente se considerará en todos los casos dedicación a tiempo parcial.

                  - Para el resto de las figuras de profesorado, un mínimo de 24 ECTS o 240 horas impartidas.

                  - En caso de que la experiencia docente se haya acumulado combinando distintas modalidades (esto es: predoctoral/postdoctoral y resto de figuras de profesorado), el mínimo de ECTS u horas se calculará proporcionalmente al periodo de tiempo de duración de cada una de las figuras contractuales disfrutadas.

                  En el caso de CU deberá acreditarse una experiencia docente mínima de 10 años de dedicación a tiempo completo, o su equivalente a tiempo parcial, y un mínimo de 140 ECTS.

                  DIRECCIÓN DE ANECA – 20/03/2024

                  ---

                  FUENTE: BAREMO 2024.pdf | PÁGINA: 4
                  \# Núm. 12/2024, de 11 de octubre

                  \#\# I. DISPOSICIONES GENERALES

                  \#\#\# I.5. Secretario General

                  3. La regulación concreta para la valoración del mérito preferente será la que se establezca en el Baremo o en las normas que, en su caso, se dicten en desarrollo de esta normativa.

                  Por su parte, en el punto VI. Mérito Preferente del Baremo aprobado junto con la normativa de contratación figura:

                  “\textbf{Profesorado Ayudante Doctor:} El aspirante logra 15 puntos atendiendo al siguiente reparto: 10 puntos por su habilitación, acreditación o ser funcionario en el cuerpo docente de Titular de Universidad o Catedrático de Universidad, o hasta 10 puntos (acumulables al punto 1 y hasta 15 puntos) por sus estancias de investigación en universidades u otros centros de investigación externos a la Universidad, hasta 18 meses de estancia, y conforme a las instrucciones dictadas por la Universidad de Sevilla”.

                  Visto lo anterior, procede una actualización tanto del artículo 10.1 de la Normativa de contratación de Profesorado Ayudante Doctor, de Profesorado Asociado y de Profesorado Sustituto como del párrafo 3 del apartado VI del baremo a aplicar.

                  \#\# TÍTULO I

                  \#\#\# Disposiciones generales

                  \textbf{Artículo 1. Objeto de régimen jurídico aplicable.}

                  1. El objeto de la presente normativa es la regulación en la Universidad de Sevilla de los procedimientos de convocatoria y provisión mediante concurso público de las siguientes plazas de personal docente e investigador: Profesorado Ayudante Doctor, Profesorado Asociado y Profesorado Interino.
                  2. Los citados procedimientos se desarrollarán en el marco de las previsiones normativas establecidas en la Ley Orgánica 6/2001, de 21 de diciembre, el Texto Refundido de la Ley Andaluza de Universidades aprobado por Decreto Legislativo 1/2013, el Estatuto de la Universidad de Sevilla, el I Convenio Colectivo del Personal Docente e Investigador en régimen laboral de las Universidades Públicas andaluzas, la presente normativa y, en lo no previsto en estas normas, por el régimen de profesorado que sea de aplicación y por las disposiciones que regulan con carácter general el ingreso en la Administración Pública.
                  3. En materia de procedimiento, y en todo lo no recogido específicamente en esta normativa, se estará a lo dispuesto en la Ley 39/2015, de 1 de octubre, del Procedimiento Administrativo Común de las Administraciones Públicas y en la Ley 40/2015, de 1 de octubre, de Régimen Jurídico del Sector Público.

                  \textbf{Artículo 2. Principios.}

                  1. La selección del personal a que se refiere la presente normativa se efectuará mediante concurso público con respeto a los principios constitucionales de igualdad, mérito, capacidad y publicidad.

                  \#\# TÍTULO II

                  \#\#\# Tramitación previa y actuación de las comisiones de selección

                  \textbf{Artículo 3. Composición, sorteo y nombramiento de las Comisiones de Contratación.}

                  1. La composición y vigencia de las Comisiones de Contratación de las plazas de Profesor Ayudante Doctor y Profesor Asociado serán las establecidas en el artículo 90 del Estatuto de la Universidad de Sevilla.

                  ---

                  FUENTE: BAREMO 2024.pdf | PÁGINA: 3
                  \# Núm. 12/2024, de 11 de octubre

                  \#\# I. DISPOSICIONES GENERALES

                  \#\#\# I.5. Secretario General

                  El artículo 12 de la referida normativa establecía: “El umbral mínimo se fija en 40 puntos para profesores ayudantes doctores, 30 puntos para profesores asociados y 20 puntos para profesores interinos. Ningún candidato cuya puntuación sea inferior a la establecida por dicho umbral podrá ser propuesto para la adjudicación de la plaza”. Idéntica gradación venía reflejada en el apartado 2 del baremo, que establecía el umbral mínimo para poder ser contratado en 20 puntos para profesorado interino, 30 para profesorado asociado y 40 para profesorado ayudante doctor.

                  Con posterioridad, por Acuerdo 7.5 de Consejo de Gobierno de 22 de abril de 2020, se modificó el referido artículo 12 de la normativa en el sentido siguiente: “El umbral mínimo se fija en 40 puntos para profesores ayudantes doctores, 30 puntos para profesores asociados y 20 puntos para profesores interinos, o en caso de acuerdo del Consejo de Departamento tras la publicación de las actas resultantes del proceso de baremación y con validez hasta la resolución del siguiente concurso público de méritos para la contratación de profesorado interino en la Universidad de Sevilla, 10 puntos. Ningún candidato cuya puntuación sea inferior a la establecida por dicho umbral podrá ser propuesto para la adjudicación de la plaza”. Asimismo, el apartado 2 del baremo quedó modificado de la siguiente forma: “El umbral mínimo queda establecido en el art. 12 de esta normativa”.

                  Dado el tiempo transcurrido y las incidencias acaecidas en las listas priorizadas de candidatos resultantes de los diferentes procesos selectivos, procede efectuar una modificación de la normativa relativa al umbral mínimo necesario para la contratación de profesorado sustituto interino.

                  \#\# MOTIVACIÓN DE LA MODIFICACIÓN

                  APROBADA POR ACUERDO 6.6/CG 26-9-24

                  El artículo 78 de la Ley Orgánica 2/2023, de 22 de marzo, del Sistema Universitario regula la figura del Profesor Ayudante Doctor dentro de las diferentes categorías de personal docente e investigador laboral. Dicho artículo dispone que las Universidades podrán contratar bajo esta modalidad a las personas que ostenten el título de Doctor sin necesidad de acreditación.

                  La Disposición transitoria tercera del referido cuerpo legal establece en su apartado 1 que la acreditación vigente de Profesor Ayudante Doctor o de la figura equivalente en la normativa autonómica se considerará como un mérito preferente durante los cuatro años posteriores a la aprobación de esta ley orgánica, a efectos del acceso a la figura de Profesor Ayudante Doctor.

                  El Consejo de Gobierno de la Universidad de Sevilla, por Acuerdo 6.7 en sesión de 27 de febrero de 2019 (posteriormente modificado por Acuerdo 7.5/CG 22-4-20 y Acuerdo 10.6/CG 1-3-24) procedió a la aprobación de la Normativa sobre procedimiento de contratación de profesorado ayudante doctor, profesorado asociado y profesorado interino, junto con el baremo establecido para dichas figuras contractuales.

                  En el artículo 10 de dicha normativa se dispone lo siguiente:

                  1. Según lo dispuesto en el artículo 48.3 de la Ley 6/2001, Orgánica de Universidades, se considerará mérito preferente estar acreditado para participar en los concursos de acceso a cuerpos docentes universitarios.
                  2. El artículo 50 del mismo texto legal dispone que para la contratación de profesores ayudantes doctores se considerará mérito preferente la estancia del candidato en universidades o centros de investigación de reconocido prestigio, españoles o extranjeros, distintos de la universidad que lleve a cabo la contratación.

                  ---

                  FUENTE: BAREMO 2024.pdf | PÁGINA: 11
                  \# Núm. 12/2024, de 11 de octubre

                  \#\# I. DISPOSICIONES GENERALES

                  \#\#\# I.5. Secretario General

                  \#\# BAREMO DE CONTRATACIÓN DE PROFESORADO AYUDANTE DOCTOR, PROFESORADO ASOCIADO Y PROFESORADO INTERINO

                  \#\#\# CONSIDERACIONES GENERALES
                  El baremo es aplicable a las figuras contractuales de Profesorado Ayudante Doctor, Profesorado Asociado y Profesorado Interino.

                  Apartado 1. Porcentaje de cada bloque temático según figura contractual (en adelante plazas).
                  En el Baremo se distinguen 6 bloques, sus ponderaciones y sus puntuaciones:
                  1. RENDIMIENTO ACADÉMICO
                  2. ACTIVIDAD DOCENTE
                  3. ACTIVIDAD INVESTIGADORA
                  4. TRANSFERENCIA DEL CONOCIMIENTO
                  5. ACTIVIDAD PROFESIONAL EXTERNA A LA UNIVERSIDAD
                  6. MÉRITO PREFERENTE

                  La calificación del candidato se valora según puntuación acumulada atendiendo a las ponderaciones de la Tabla 1, y dependiendo del tipo de contrato. En caso de empate a puntos entre aspirantes la preferencia entre ellos se resolverá en función de la puntuación alcanzada en el Bloque Temático VI. Mérito Preferente. En caso de no resolverse el empate se calculará la media aritmética de la puntuación (sobre 100) de Actividad Docente (Bloque II) y Actividad Investigadora (Bloque III) en las plazas de Ayudante Doctor; en las plazas de Profesor Asociado la media aritmética se calculará con las puntuaciones (sobre 100) de la Actividad Docente (Bloque II) y Actividad Profesional Externa a la Universidad (Bloque V), y en el caso del Profesorado Interino la media aritmética se calculará con las puntuaciones (sobre 100) del Rendimiento Académico (Bloque I) y Actividad Docente (Bloque II).

                  \#\#\# TABLA 1. Ponderación de cada bloque temático según tipo de contrato.

                  \vspace{0.4cm}
                  \begin{table}[h!]\centering\renewcommand{\arraystretch}{1.3}
                    \begin{tabular}{|p{0.45\textwidth}|p{0.45\textwidth}|p{0.45\textwidth}|p{0.45\textwidth}|p{0.45\textwidth}|}\hline
                      BLOQUE TEMÁTICO                                   & TIPO DE CONTRATO/PLAZA (*) & PSI \% & ASO \% & AYD \% \\\hline
                      I. RENDIMIENTO ACADÉMICO                          &                            & 30     & 5      & 10     \\\hline
                      II. ACTIVIDAD DOCENTE                             &                            & 25     & 10     & 30     \\\hline
                      III. ACTIVIDAD INVESTIGADORA                      &                            & 25     & 10     & 35     \\\hline
                      IV. TRANSFERENCIA DEL CONOCIMIENTO                &                            & 5      & 15     & 5      \\\hline
                      V. ACTIVIDAD PROFESIONAL EXTERNA A LA UNIVERSIDAD &                            & 10     & 55     & 5      \\\hline
                      VI. MÉRITO PREFERENTE                             &                            & 5      & 5      & 15     \\\hline
                      TOTAL                                             &                            & 100    & 100    & 100    \\\hline
                    \end{tabular}\end{table}
                  \vspace{0.4cm}

                  (*) PSI: Profesor Interino; ASO: Profesor Asociado; AYD: Profesor Ayudante Doctor.
                \end{scriptsize}
        \end{itemize}


  \item \textbf{P20.} Si estoy aplicando a Profesor Asociado según las instrucciones del baremo, ¿cómo se calculan exactamente los puntos correspondientes a la docencia oficial de más de 1200 horas? \\
        \begin{itemize}
          \item \textbf{Tipo de consulta:} simple
          \item \textbf{Respuesta del Agente:}

                No existen requisitos mínimos excluyentes mencionados en el contexto proporcionado para la evaluación de la docencia oficial de más de 1200 horas.

                A continuación, se presentan los criterios de evaluación para la docencia en asignaturas regladas universitarias (epígrafe 2.1 del baremo):

                \vspace{0.4cm}
                \begin{table}[h!]\centering\renewcommand{\arraystretch}{1.3}
                  \begin{tabular}{|p{0.45\textwidth}|p{0.45\textwidth}|}\hline
                    Criterio                                        & Detalles                                                                                                      \\\hline
                    Docencia en asignaturas regladas universitarias & Máximo 50 puntos, máximo 1.200 horas de docencia oficial                                                      \\\hline
                    Puntuación por hora de docencia                 & 1,00 x 10 horas de docencia                                                                                   \\\hline
                    Complemento 1 (multiplicativo)                  & 1,50 x puntuación obtenida para docencia en Máster Universitario (oficial)                                    \\\hline
                    Complemento 2 (multiplicando al complemento 1)  & 1,50 x puntuación obtenida para docencia en inglés u otras lenguas diferentes a las que son objeto de estudio \\\hline
                    Afinidad de la asignatura                       & Afinidad muy alta (1,0), afinidad alta (0,75), afinidad media (0,50), afinidad baja o muy baja (0,25 o 0,0)   \\\hline
                  \end{tabular}\end{table}
                \vspace{0.4cm}

                Para calcular los puntos correspondientes a la docencia oficial de más de 1200 horas, se deben seguir los siguientes pasos:

                1. \textbf{Selección de horas}: Seleccionar en primer lugar las horas de máxima afinidad. Si con eso se alcanzan las 1200 horas, no incluir ninguna más. Si no se llega a 1200 horas, seguir incluyendo horas en orden descendente de puntuación por afinidad, hasta completar las 1200 horas.
                2. \textbf{Cálculo de la puntuación}: Calcular la puntuación total teniendo en cuenta la afinidad de las horas y los complementos multiplicativos (docencia en Máster Universitario y docencia en inglés u otras lenguas).
                3. \textbf{Escalado individual}: Si el candidato supera las 1200 horas, se calculará su puntuación total y se aplicará un factor de escala para ajustar la puntuación a la escala de 0 a 50 puntos.

                Por ejemplo, si un candidato ha impartido 1500 horas de docencia en asignaturas regladas universitarias, con afinidad muy alta, y 300 horas de docencia en Máster Universitario, la puntuación se calcularía de la siguiente manera:

                * Puntuación por hora de docencia: 1,00 x 10 horas de docencia = 10 puntos por hora
                * Puntuación total por 1200 horas de docencia: 1200 horas x 10 puntos/hora = 12000 puntos
                * Puntuación total por 300 horas de docencia en Máster Universitario: 300 horas x 10 puntos/hora x 1,50 (complemento 1) = 4500 puntos
                * Puntuación total: 12000 puntos + 4500 puntos = 16500 puntos
                * Factor de escala: 1200 horas / 1500 horas = 0,8
                * Puntuación final: 16500 puntos x 0,8 = 13200 puntos

                Sin embargo, como la puntuación máxima para este apartado es de 50 puntos, la puntuación final se ajustaría a 50 puntos.

                Es importante tener en cuenta que la comisión de contratación puede aplicar la regla de saturación si algún candidato supera la puntuación máxima establecida para un subapartado. En ese caso, se asignaría la máxima puntuación al candidato con la puntuación más alta y se disminuiría proporcionalmente la puntuación de los demás candidatos.
          \item \textbf{Contexto utilizado:}
                \begin{scriptsize}
                  FUENTE: 27830\_guia-ayuda-candidato-r04-3-.pdf | PÁGINA: 19
                  \# Guía de ayuda al candidato

                  \#\# Actividad profesional externa a la Universidad

                  En caso de duplicidad, es dispositivo de la comisión decidir si procede su consideración favorable (validación) y en qué único lugar se baremará, en aplicación de la discrecionalidad técnica de que goza. El criterio que haya adoptado la comisión para resolver la duplicidad deberá ser aplicado a todos los candidatos.

                  \#\# 5.4 Puntuaciones escaladas

                  Al cumplimentar el CVN para su uso en Bareprof, el candidato debe tener presente que el baremo de la Universidad de Sevilla contempla dos tipos de escalado en las puntuaciones de los distintos subapartados que componen el baremo. Uno de estos escalados corresponde a la denominada regla de saturación, que se aplica con carácter general. Dicha regla opera de la siguiente forma: cuando algún candidato supera la puntuación máxima establecida para un subapartado, se asigna la máxima puntuación al que la tuviera más alta y se disminuye proporcionalmente la de los demás candidatos.

                  Por ejemplo, supongamos que la puntuación máxima de un subapartado es 40 puntos y que hay un candidato que obtiene 50 puntos, siendo ésta la puntuación más alta que presentan los candidatos; en tal caso, todas las puntuaciones se multiplicarán por el factor de escala 0.8, resultante de la operación 40/50. De este modo, el candidato con máxima puntuación pasará a tener 50×0.8 = 40 puntos y un candidato que tuviera 40 puntos, pasaría a tener 40×0.8 = 32 puntos.

                  El otro tipo de escalado contemplado en el baremo es un escalado individual que se aplica en algunos subapartados específicos cuando se superan determinados valores máximos, como es el caso de la docencia en asignaturas regladas universitarias (epígrafe 2.1) o los contratos como PDI universitario (epígrafe 2.2). Para explicar su funcionamiento, consideraremos, por ejemplo, el caso del apartado 2.1, en el que se establece un límite de 1200 horas de docencia. Si un candidato supera las 1200 horas, se calculará su puntuación total teniendo en cuenta la afinidad de estas o su distinto valor si son horas de docencia en Máster y se realizará el escalado de esa puntuación en función del número de horas impartidas.

                  Veamos algunos ejemplos:

                  1. El candidato ha impartido 1500 horas en titulaciones de Grado, todas con afinidad muy alta. Su puntuación sería 150 puntos, que habrá que multiplicar por el factor de escala 0.8 (1200/1500=0.8). El resultado final sería 150*0.8 = 120, que es el mismo valor que obtendría si introdujera solo 1200 horas.
                  2. El candidato ha impartido 1500 horas en titulaciones de Grado, 1200 de ellas con afinidad muy alta y las otras 300 con afinidad media. La puntuación sería (1200/10) + (0.5\textit{300/10) = 135. El factor de escala sería el mismo de antes, por lo que la puntuación final sería 135}0.8 = 108. Sin embargo, si solo hubiera reflejado en el CVN las 1200 de máxima afinidad, tendría 120 puntos.

                  Por tanto, si cuenta con horas de distinta afinidad o distinto complemento multiplicativo, debe seleccionar en primer lugar las de máxima afinidad/complemento. Si con eso alcanza las 1200, no incluya ninguna más. Si no llega a 1200, siga incluyendo horas en orden descendente de puntuación por afinidad/complemento, hasta completar las 1200. Si incluye más, puede resultar penalizado, como se ha puesto de manifiesto en el segundo ejemplo. Posteriormente, si algún candidato en este subapartado superara los 50 puntos (valor máximo del mismo), se aplicaría la regla de saturación descrita anteriormente.

                  El escalado individual se aplicará, además, en el caso de los contratos como PDI universitario (epígrafe 2.2), la docencia universitaria no oficial (epígrafe 2.4.2), los cursos de formación e innovación docente (epígrafe 2.4.5) y las estancias de investigación (bloque VI).

                  Revisión 4 (marzo, 2025) 18

                  ---

                  FUENTE: BAREMO 2024.pdf | PÁGINA: 15
                  \# Núm. 12/2024, de 11 de octubre

                  \#\# I. DISPOSICIONES GENERALES

                  \#\#\# I.5. Secretario General

                  1.3. Otros títulos académicos:                                                            \textbf{Máximo: 10 puntos}
                  - 1.3.1. Titulación Universitaria de Grado o equivalente (como Licenciado/a, Diplomado/a, Arquitecto/a, Arquitecto/a técnico/a, Ingeniero/a o Ingeniero/a técnico/a), y Titulación Universitaria de Máster (Oficial):
                  - APROBADO: 5,0-6,9 nota media                            \textbf{3 puntos}
                  - NOTABLE: 7,0-8,9 nota media                             \textbf{6 puntos}
                  - SOBRESALIENTE: 9,0-10,0 nota media                      \textbf{10 puntos}
                  - 1.3.2. Título de doctor (tesis doctoral)                  \textbf{10 puntos}

                  \textbf{Nota 1:} En relación con el coeficiente de afinidad y la homologación de títulos se aplicarán las Notas del apartado 1.1.

                  1.4. Premios académicos y menciones de calidad:                                           \textbf{Máximo: 5 puntos}
                  - 1.4.1. Premio (Nacional) Fin de Carrera de Educación Universitaria: Primer premio           \textbf{5 puntos por premio}           No admite coeficiente de afinidad.
                  - 1.4.2. Premio extraordinario de grado o equivalente        \textbf{4 puntos por premio}
                  - 1.4.3. Premio Extraordinario de Doctorado                  \textbf{4 puntos por premio}
                  - 1.4.4. Doctorado internacional y europeo                   \textbf{2 puntos por indicio}
                  - 1.4.5. Mención de calidad (o ‘hacia la excelencia’)        \textbf{1 punto por indicio}
                  - 1.4.6. Tesis en régimen de cotutela                        \textbf{1 punto por indicio}

                  \#\# II. ACTIVIDAD DOCENTE                                                                     \textbf{Máximo: 100 puntos}

                  2.1. Docencia en asignaturas regladas universitarias:                                     \textbf{Máximo: 50 puntos}
                  (máximo 1.200 horas de docencia oficial)
                  - 2.1.1. Titulación Universitaria de Grado o equivalente (como Licenciado/a, Diplomado/a, Arquitecto/a, Arquitecto/a técnico/a, Ingeniero/a o Ingeniero/a técnico/a), y Titulación Universitaria de Máster (Oficial):
                  \textbf{1,00 x 10 horas de docencia}
                  En el caso de que el aspirante acredite docencia superior a 1.200 horas, su puntuación se escalará proporcionalmente a 1.200 horas.

                  \textbf{Complemento 1 (multiplicativo):}                             \textbf{1,50 x puntuación obtenida}
                  Docencia en Máster Universitario (oficial)                  en 2.1.1 en asignaturas de Máster Universitario (oficial)

                  \textbf{Complemento 2 (multiplicando al complemento 1):}
                  Docencia en inglés u otras lenguas diferentes a las que son objeto de estudio                                      \textbf{1,50 x puntuación obtenida}
                  en 2.1.1 en asignaturas en inglés u otras lenguas diferentes a las que son objeto de estudio

                  \textbf{Nota 1:} Los Trabajos Fin de Carrera (Grado o equivalente, y Máster) no serán evaluados en este apartado. Se remite al aspirante al apartado 2.4.1.
                  \textbf{Nota 2:} Se considera la siguiente tipología de asignaturas:
                  - Afinidad muy alta (1,0): asignaturas cuyo contenido corresponda al perfil de la plaza convocada, o en su defecto las asignaturas impartidas por el área de conocimiento del departamento solicitante de la plaza.
                  - Afinidad alta  (0,75): asignaturas impartidas por el área de conocimiento solicitante sin consideración de universidad de impartición, e incluyendo con ello las asignaturas no necesariamente impartidas por el área de conocimiento del departamento de la Universidad de Sevilla convocante de la plaza.
                  - Afinidad media  (0,50): asignaturas impartidas por áreas de conocimiento afines al departamento y área de conocimiento a que corresponde la plaza según listado oficial de la Universidad de Sevilla de comisiones juzgadoras a cuerpos docentes, y sin consideración de universidad de impartición.
                  - Afinidad baja o muy baja      (0,25 o 0,0): asignaturas no pertenecientes a ninguno de los tipos anteriores. El coeficiente de afinidad aplicable deberá ser 0,0 o 0,25, motivándose por la comisión de contratación.

                  ---

                  FUENTE: 27831\_instrucciones-baremacinin-r03.pdf | PÁGINA: 24
                  \# Bareprof
                  \#\# Instrucciones de baremación

                  \#\#\# 4.3 Actividad docente (bloque II del baremo)

                  \#\#\#\# 4.3.1 Docencia en asignaturas regladas universitarias (epígrafe 2.1 del baremo)

                  - En este apartado el candidato reflejará la docencia realmente impartida (no la nominal de su contrato), acreditada mediante el certificado de actividades docentes o documento equivalente.
                  - La docencia debe introducirse obligatoriamente en horas. Si la comisión observa que el candidato ha introducido el número de créditos, podrá modificarlo para reflejar el valor correcto.
                  - La docencia recogida en este apartado podrá corresponder tanto a la actividad del candidato como PDI (epígrafe 2.2 del baremo) como a su actividad docente desarrollada como Personal Investigador (epígrafe 3.1 del baremo).
                  - Las Prácticas Externas se validan como actividad docente del candidato siempre que este actúe como tutor académico; el tutor profesional ejerce la labor en calidad de responsable de la entidad colaboradora y no de PDI. Su cómputo se corresponde con lo reflejado en el certificado de actividades docentes. Por ejemplo, la Universidad de Sevilla reconoce al tutor académico 0,1 créditos por cada 6 créditos de que conste la asignatura en el Plan de Estudios.
                  - Se valorará también en este apartado la dirección o tutela de trabajos académicos que aparezca recogida y cuantificada (en número de horas o créditos) en el certificado de actividades docentes, con independencia de que la labor de tutorización sea también valorada en el apartado 2.4.1 del baremo (dirección o tutela de trabajos académicos en titulaciones universitarias oficiales). Esto no supone duplicidad en el mérito, pues se valora, por un lado, la actividad docente presencial reconocida y, por otro, la propia labor de tutorización.

                  \#\#\#\# 4.3.2 Contratos como PDI universitario (epígrafe 2.2 del baremo)

                  - En este apartado el candidato registrará su actividad como PDI universitario mediante contrato laboral o pertenencia a los cuerpos docentes universitarios. Los contratos laborales que se contemplan corresponden a las modalidades de contratación (ordinarias y extraordinarias) recogidas en los apartados 1, 2 y 3 del artículo 14 del I Convenio Colectivo del PDI Laboral de las Universidad Públicas de Andalucía (BOJA 92/2008, de 9 de mayo) o sus análogas en otros Convenios, así como las categorías análogas de PDI de Universidades privadas.
                  - Profesor Ayudante
                  - Profesor Ayudante Doctor
                  - Profesor Colaborador
                  - Profesor Contratado Doctor
                  - Profesor Asociado
                  - Profesor Asociado de concierto con las instituciones sanitarias
                  - Profesor Sustituto Interino

                  Las figuras englobadas dentro de los cuerpos docentes universitarios son las siguientes:
                  - Profesor Titular de Escuela Universitaria
                  - Catedrático de Escuela Universitaria
                  - Profesor Titular de Universidad
                  - Catedrático de Universidad
                  - Figuras análogas en otros ámbitos internacionales

                  Revisión 3 (diciembre, 2024) 23

                  ---

                  FUENTE: BAREMO 2024.pdf | PÁGINA: 11
                  \# Núm. 12/2024, de 11 de octubre

                  \#\# I. DISPOSICIONES GENERALES

                  \#\#\# I.5. Secretario General

                  \#\# BAREMO DE CONTRATACIÓN DE PROFESORADO AYUDANTE DOCTOR, PROFESORADO ASOCIADO Y PROFESORADO INTERINO

                  \#\#\# CONSIDERACIONES GENERALES
                  El baremo es aplicable a las figuras contractuales de Profesorado Ayudante Doctor, Profesorado Asociado y Profesorado Interino.

                  Apartado 1. Porcentaje de cada bloque temático según figura contractual (en adelante plazas).
                  En el Baremo se distinguen 6 bloques, sus ponderaciones y sus puntuaciones:
                  1. RENDIMIENTO ACADÉMICO
                  2. ACTIVIDAD DOCENTE
                  3. ACTIVIDAD INVESTIGADORA
                  4. TRANSFERENCIA DEL CONOCIMIENTO
                  5. ACTIVIDAD PROFESIONAL EXTERNA A LA UNIVERSIDAD
                  6. MÉRITO PREFERENTE

                  La calificación del candidato se valora según puntuación acumulada atendiendo a las ponderaciones de la Tabla 1, y dependiendo del tipo de contrato. En caso de empate a puntos entre aspirantes la preferencia entre ellos se resolverá en función de la puntuación alcanzada en el Bloque Temático VI. Mérito Preferente. En caso de no resolverse el empate se calculará la media aritmética de la puntuación (sobre 100) de Actividad Docente (Bloque II) y Actividad Investigadora (Bloque III) en las plazas de Ayudante Doctor; en las plazas de Profesor Asociado la media aritmética se calculará con las puntuaciones (sobre 100) de la Actividad Docente (Bloque II) y Actividad Profesional Externa a la Universidad (Bloque V), y en el caso del Profesorado Interino la media aritmética se calculará con las puntuaciones (sobre 100) del Rendimiento Académico (Bloque I) y Actividad Docente (Bloque II).

                  \#\#\# TABLA 1. Ponderación de cada bloque temático según tipo de contrato.

                  \vspace{0.4cm}
                  \begin{table}[h!]\centering\renewcommand{\arraystretch}{1.3}
                    \begin{tabular}{|p{0.45\textwidth}|p{0.45\textwidth}|p{0.45\textwidth}|p{0.45\textwidth}|p{0.45\textwidth}|}\hline
                      BLOQUE TEMÁTICO                                   & TIPO DE CONTRATO/PLAZA (*) & PSI \% & ASO \% & AYD \% \\\hline
                      I. RENDIMIENTO ACADÉMICO                          &                            & 30     & 5      & 10     \\\hline
                      II. ACTIVIDAD DOCENTE                             &                            & 25     & 10     & 30     \\\hline
                      III. ACTIVIDAD INVESTIGADORA                      &                            & 25     & 10     & 35     \\\hline
                      IV. TRANSFERENCIA DEL CONOCIMIENTO                &                            & 5      & 15     & 5      \\\hline
                      V. ACTIVIDAD PROFESIONAL EXTERNA A LA UNIVERSIDAD &                            & 10     & 55     & 5      \\\hline
                      VI. MÉRITO PREFERENTE                             &                            & 5      & 5      & 15     \\\hline
                      TOTAL                                             &                            & 100    & 100    & 100    \\\hline
                    \end{tabular}\end{table}
                  \vspace{0.4cm}

                  (*) PSI: Profesor Interino; ASO: Profesor Asociado; AYD: Profesor Ayudante Doctor.

                  ---

                  FUENTE: 27830\_guia-ayuda-candidato-r04-3-.pdf | PÁGINA: 28
                  \# Guía de ayuda al candidato

                  - En el campo Tipo de Programa, el candidato seleccionará el que corresponda a la titulación en que se ha impartido la docencia (Arquitectura, Arquitectura Técnica, Diplomatura, etc.). Es especialmente importante que este campo se cumplimente correctamente en el caso de la docencia en Máster Oficial, pues esta tiene asociada en el baremo un coeficiente multiplicador.
                  - Para introducir la docencia en titulaciones de Grado, se deberá seleccionar la opción Otros en el campo Tipo de Programa; al hacerlo, aparece un nuevo campo denominado Otros, en el que se deberá introducir el término Grado.
                  - En el campo Departamento se debe hacer constar el área de conocimiento, entre paréntesis tras el nombre del departamento. Por ejemplo:
                  \textbf{Departamento:} Filologías Integradas (Área de Conocimiento: Estudios de Asia Oriental)
                  - En los casos en que el certificado de actividades docentes no explicite la docencia impartida en la asignatura en número de horas, debe indicarse el equivalente en horas de los créditos impartidos. En ningún caso debe introducir el número de créditos; si así lo hiciera, este dato sería interpretado necesariamente como número de horas.
                  - Si la docencia se ha impartido en idioma distinto al oficial de su territorio de impartición (por ejemplo, docencia en francés en Universidad española), y con las salvaguardas hechas en el baremo, debe indicarlo expresamente en el campo Idioma de la asignatura. Para que la docencia en inglés u otro idioma se contabilice, el candidato ha de contar con una acreditación del idioma igual o superior al nivel B2 del Marco Común Europeo de Referencia para las Lenguas (o equivalente). Se exceptúa de la valoración de este mérito a aquellas plazas en que el uso del idioma acreditado sea consustancial con su docencia (por ejemplo, inglés en Filología Inglesa).
                  - El candidato puede incluir como mérito en este apartado las horas de actividad que consten en el certificado de actividades docentes asociadas a la dirección o tutela de trabajos académicos en titulaciones universitarias oficiales o a la tutela de prácticas externas.

                  \#\# Documentación acreditativa:

                  - El candidato debe incluir certificación oficial de la Universidad donde se expliciten los puestos ocupados (hoja de servicios oficial) y la docencia impartida oficialmente (certificado de actividades docentes).
                  - En los casos en que el certificado de actividades docentes no explicite el departamento y, en particular, el área de conocimiento de adscripción de las asignaturas impartidas, debe adjuntarse, o bien una hoja de servicios que permita validar a qué departamento y área pertenece el candidato en cada periodo docente, o cualquier otro documento que supla la necesidad descrita (por ejemplo, el programa de la asignatura publicado por la Universidad de Sevilla indica Departamento y Área de conocimiento responsable de la asignatura). El área de conocimiento desempeña un papel clave en la valoración del candidato y debe quedar debidamente acreditada su pertenencia.
                  - Si el candidato ha impartido la docencia de la asignatura en idioma distinto al oficial de su territorio de impartición, deberá incluir la acreditación del nivel del idioma (al menos B2 del Marco Común Europeo de Referencia para las Lenguas o equivalente) junto con un certificado del secretario del departamento responsable de la docencia (u órgano equivalente), de haber impartido la docencia en otro idioma (salvo que el certificado de actividades docentes indique el idioma en que se imparte la asignatura).

                  Revisión 4 (marzo, 2025) 27
                \end{scriptsize}
        \end{itemize}


  \item \textbf{P21.} Para Profesor Sustituto Interino, ¿es un requisito estrictamente excluyente sacar más de 0 puntos en la experiencia investigadora? \\
        \begin{itemize}
          \item \textbf{Tipo de consulta:} simple
          \item \textbf{Respuesta del Agente:}

                En relación con la pregunta planteada sobre si es un requisito estrictamente excluyente sacar más de 0 puntos en la experiencia investigadora para el Profesor Sustituto Interino, debemos analizar la normativa aplicable.

                Según el artículo 12 de la normativa sobre procedimiento de contratación de profesorado ayudante doctor, profesorado asociado y profesorado interino, el umbral mínimo se fija en 20 puntos para profesores interinos. Sin embargo, este umbral puede reducirse hasta un mínimo de 5 puntos en caso de acuerdo del Consejo de Departamento adoptado tras la publicación de las actas resultantes del proceso de baremación y con validez hasta la resolución del siguiente concurso público de méritos para la contratación de profesorado sustituto interino en la Universidad de Sevilla.

                En cuanto a la experiencia investigadora, el baremo establece que el mérito preferente para el Profesorado Interino se logra con 5 puntos, que pueden ser obtenidos por habilitación, acreditación o ser funcionario en el cuerpo docente de Titular de Universidad o Catedrático de Universidad, o por estancias de investigación en universidades u otros centros de investigación externos a la Universidad de Sevilla.

                No obstante, no se establece expresamente que obtener más de 0 puntos en la experiencia investigadora sea un requisito estrictamente excluyente para el Profesor Sustituto Interino. Lo que sí se establece es que el umbral mínimo de 20 puntos (o 5 puntos en caso de acuerdo del Consejo de Departamento) debe ser superado para poder ser propuesto para la adjudicación de la plaza.

                En consecuencia, aunque la experiencia investigadora es un mérito valorable en el proceso de selección, no es un requisito estrictamente excluyente que deba ser superado con una puntuación mayor que 0. Lo que realmente importa es alcanzar el umbral mínimo establecido para el Profesorado Interino, que puede variar según las circunstancias.

                Por lo tanto, la respuesta a la pregunta del usuario es que no es un requisito estrictamente excluyente sacar más de 0 puntos en la experiencia investigadora para el Profesor Sustituto Interino, siempre y cuando se supere el umbral mínimo establecido.
          \item \textbf{Contexto utilizado:}
                \begin{scriptsize}
                  FUENTE: BAREMO 2024.pdf | PÁGINA: 8
                  \# Núm. 12/2024, de 11 de octubre

                  \#\# I. DISPOSICIONES GENERALES

                  \#\#\# I.5. Secretario General

                  \#\# TÍTULO III
                  \#\#\# Desarrollo del concurso

                  \textbf{Artículo 11. Actuación de las Comisiones de Contratación.}
                  1. Una vez elevadas a definitivas las listas de aspirantes admitidos y excluidos al concurso, se comunicarán al Presidente de la Comisión de Contratación, debiéndose acusar recibo de las mismas. Asimismo, se le indicará la ubicación de los modelos oficiales en los que se documentará necesariamente el desarrollo del concurso en soporte informático incluida la documentación aportada por cada uno de los candidatos para su valoración.
                  2. La Comisión dispondrá de un plazo máximo de 25 días hábiles para resolver y elevar, en su caso, la correspondiente propuesta de adjudicación desde la comunicación referida por el Servicio de Programación Docente.

                  \textbf{Artículo 12. Umbral mínimo.}
                  El umbral mínimo se fija en 40 puntos para profesores ayudantes doctores, 30 puntos para profesores asociados y 20 puntos para profesorado sustituto interino.
                  El umbral para la contratación de profesorado sustituto interino podrá reducirse hasta un mínimo de 5 puntos, en caso de acuerdo del Consejo de Departamento adoptado tras la publicación de las actas resultantes del proceso de baremación y con validez hasta la resolución del siguiente concurso público de méritos para la contratación de profesorado sustituto interino en la Universidad de Sevilla. Ningún candidato cuya puntuación sea inferior a la establecida en dichos umbrales podrá ser propuesto para la adjudicación de la plaza.

                  \#\# TÍTULO IV
                  \#\#\# Resolución del concurso

                  \textbf{Artículo 13. Propuesta de adjudicación.}
                  1. Una vez asignada la puntuación a cada candidato, la Comisión de Contratación elaborará y firmará el Acta de Propuesta de Adjudicación, en la que propondrá:
                  - Al candidato que haya obtenido mayor puntuación, siempre que tenga una puntuación igual o superior al umbral mínimo, o bien
                  - La no provisión de la plaza si ningún candidato tiene una puntuación igual o superior al umbral mínimo.
                  2. Los Presidentes de las Comisiones deberán comunicar al Servicio de Programación Docente la correspondiente Acta de Propuesta de Adjudicación.
                  3. El Acta de Propuesta de Adjudicación será publicada en la sede en la que se establezca para las publicaciones oficiales por medios electrónicos de la Universidad de Sevilla. La publicación del Acta de Propuesta de Adjudicación servirá de notificación al adjudicatario de la plaza, de conformidad con lo dispuesto en el artículo 45 de la Ley 39/2015, de 1 de octubre, del Procedimiento Administrativo Común de las Administraciones Públicas, no efectuándose notificación personal alguna.
                  4. Contra estas propuestas, los interesados podrán interponer recurso de alzada ante el Sr. Rector de esta Universidad, en el plazo de un mes a partir del día siguiente a aquel en que tenga lugar la publicación.

                  ---

                  FUENTE: BAREMO 2024.pdf | PÁGINA: 3
                  \# Núm. 12/2024, de 11 de octubre

                  \#\# I. DISPOSICIONES GENERALES

                  \#\#\# I.5. Secretario General

                  El artículo 12 de la referida normativa establecía: “El umbral mínimo se fija en 40 puntos para profesores ayudantes doctores, 30 puntos para profesores asociados y 20 puntos para profesores interinos. Ningún candidato cuya puntuación sea inferior a la establecida por dicho umbral podrá ser propuesto para la adjudicación de la plaza”. Idéntica gradación venía reflejada en el apartado 2 del baremo, que establecía el umbral mínimo para poder ser contratado en 20 puntos para profesorado interino, 30 para profesorado asociado y 40 para profesorado ayudante doctor.

                  Con posterioridad, por Acuerdo 7.5 de Consejo de Gobierno de 22 de abril de 2020, se modificó el referido artículo 12 de la normativa en el sentido siguiente: “El umbral mínimo se fija en 40 puntos para profesores ayudantes doctores, 30 puntos para profesores asociados y 20 puntos para profesores interinos, o en caso de acuerdo del Consejo de Departamento tras la publicación de las actas resultantes del proceso de baremación y con validez hasta la resolución del siguiente concurso público de méritos para la contratación de profesorado interino en la Universidad de Sevilla, 10 puntos. Ningún candidato cuya puntuación sea inferior a la establecida por dicho umbral podrá ser propuesto para la adjudicación de la plaza”. Asimismo, el apartado 2 del baremo quedó modificado de la siguiente forma: “El umbral mínimo queda establecido en el art. 12 de esta normativa”.

                  Dado el tiempo transcurrido y las incidencias acaecidas en las listas priorizadas de candidatos resultantes de los diferentes procesos selectivos, procede efectuar una modificación de la normativa relativa al umbral mínimo necesario para la contratación de profesorado sustituto interino.

                  \#\# MOTIVACIÓN DE LA MODIFICACIÓN

                  APROBADA POR ACUERDO 6.6/CG 26-9-24

                  El artículo 78 de la Ley Orgánica 2/2023, de 22 de marzo, del Sistema Universitario regula la figura del Profesor Ayudante Doctor dentro de las diferentes categorías de personal docente e investigador laboral. Dicho artículo dispone que las Universidades podrán contratar bajo esta modalidad a las personas que ostenten el título de Doctor sin necesidad de acreditación.

                  La Disposición transitoria tercera del referido cuerpo legal establece en su apartado 1 que la acreditación vigente de Profesor Ayudante Doctor o de la figura equivalente en la normativa autonómica se considerará como un mérito preferente durante los cuatro años posteriores a la aprobación de esta ley orgánica, a efectos del acceso a la figura de Profesor Ayudante Doctor.

                  El Consejo de Gobierno de la Universidad de Sevilla, por Acuerdo 6.7 en sesión de 27 de febrero de 2019 (posteriormente modificado por Acuerdo 7.5/CG 22-4-20 y Acuerdo 10.6/CG 1-3-24) procedió a la aprobación de la Normativa sobre procedimiento de contratación de profesorado ayudante doctor, profesorado asociado y profesorado interino, junto con el baremo establecido para dichas figuras contractuales.

                  En el artículo 10 de dicha normativa se dispone lo siguiente:

                  1. Según lo dispuesto en el artículo 48.3 de la Ley 6/2001, Orgánica de Universidades, se considerará mérito preferente estar acreditado para participar en los concursos de acceso a cuerpos docentes universitarios.
                  2. El artículo 50 del mismo texto legal dispone que para la contratación de profesores ayudantes doctores se considerará mérito preferente la estancia del candidato en universidades o centros de investigación de reconocido prestigio, españoles o extranjeros, distintos de la universidad que lleve a cabo la contratación.

                  ---

                  FUENTE: 27831\_instrucciones-baremacinin-r03.pdf | PÁGINA: 5
                  \# Bareprof

                  \#\# Instrucciones de baremación

                  - Instrucción de 17-10-19 sobre los requisitos para el reconocimiento de las estancias de investigación dentro del mérito preferente en los concursos de profesorado ayudante doctor, asociado y sustituto interino (BOUS 14/2019, de 21 de noviembre).
                    [https://www.us.es/bous-numeros/numero-14-21-de-noviembre-de-2019](https://www.us.es/bous-numeros/numero-14-21-de-noviembre-de-2019)

                  - Resolución rectoral, de 24 marzo de 2021, por la que se emite una interpretación sobre la determinación del umbral mínimo aplicable en los procedimientos selectivos para plazas de Profesores Ayudantes Doctores, Profesores Asociados y Profesores Interinos establecido en el artículo 12 de la normativa sobre procedimiento de contratación de Profesorado Ayudante Doctor, Profesorado Asociado y Profesorado Interino, junto con el baremo establecido para dichas figuras contractuales de la Universidad de Sevilla (BOUS 1/2021, de 7 de abril).
                    [https://www.us.es/bous-numeros/numero-1-7-de-abril-de-2021](https://www.us.es/bous-numeros/numero-1-7-de-abril-de-2021)

                  - Resolución rectoral, de 20 de abril de 2022, por la que se determinan las áreas afines en el ámbito de los concursos de profesorado de la Universidad de Sevilla (BOUS 3/2022, de 19 de mayo).
                    [https://www.us.es/bous-numeros/numero-3-19-de-mayo-de-2022](https://www.us.es/bous-numeros/numero-3-19-de-mayo-de-2022)

                  Asimismo, los miembros de las comisiones de contratación deberán consultar los siguientes documentos disponibles en la plataforma Bareprof:
                  - Guía de ayuda al candidato.
                  - Anexo I correspondiente al departamento y área de conocimiento de la plaza a baremar.

                  \#\# 3 Descripción del proceso de baremación

                  \#\#\# 3.1 Composición de las comisiones de contratación

                  La composición de las comisiones de contratación correspondientes a concursos de Profesorado Ayudante Doctor se determina según lo dispuesto en el acuerdo 9.7/CG 23-7-74, por el que se aprueban las normas provisionales para adaptar a la Ley Orgánica 2/2023, de 22 de marzo, del Sistema Universitario, el modo de elección de las comisiones se selección de plazas de Profesorado Ayudante Doctor. Tales comisiones pueden consultarse en el tablón electrónico oficial de la Universidad de Sevilla ([https://www.us.es/tablon-virtual](https://www.us.es/tablon-virtual)).

                  En cuanto a los concursos de plazas de Profesorado Asociado y Profesorado Sustituto, la elección de las comisiones se realiza conforme a lo establecido en la Normativa sobre Procedimiento de Contratación de Profesorado Ayudante Doctor, Profesorado Asociado y Profesorado Interino (véase el apartado 2 de este documento). La composición de estas comisiones puede consultarse en la Secretaría Virtual de la Universidad de Sevilla, accediendo a la siguiente dirección:
                  [https://sevius4.us.es/?comisionescontratacion](https://sevius4.us.es/?comisionescontratacion)

                  Para realizar la consulta, debe seleccionarse el año del sorteo, el departamento, el área de conocimiento y la comisión (según el tipo de plaza a baremar).

                  Revisión 3 (diciembre, 2024)

                  ---

                  FUENTE: BAREMO 2024.pdf | PÁGINA: 27
                  \# Núm. 12/2024, de 11 de octubre

                  \#\# I. DISPOSICIONES GENERALES

                  \#\#\# I.5. Secretario General

                  4.4. Participar en la generación de Spin-off, con actividad económica contrastable durante al menos 2 años:
                  \textbf{Máximo: 10 puntos}
                  - 4.4.1. Participar en la generación de Spin-off: \textbf{5 puntos por participación}

                  \textbf{Complemento 1 (multiplicativo):}
                  \vspace{0.4cm}
                  \begin{table}[h!]\centering\renewcommand{\arraystretch}{1.3}
                    \begin{tabular}{|p{0.45\textwidth}|p{0.45\textwidth}|}\hline
                      Ámbito geográfico          & Puntos por mérito       \\\hline
                      Internacional (UE o no UE) & 2,0 x puntos por mérito \\\hline
                      Nacional                   & 1,5 x puntos por mérito \\\hline
                      Autonómico y Otros         & 1,0 x puntos por mérito \\\hline
                    \end{tabular}\end{table}
                  \vspace{0.4cm}

                  \textbf{Complemento 2 (multiplicando el complemento 1):}
                  \vspace{0.4cm}
                  \begin{table}[h!]\centering\renewcommand{\arraystretch}{1.3}
                    \begin{tabular}{|p{0.45\textwidth}|p{0.45\textwidth}|}\hline
                      Grado de contribución                                 & Puntos por mérito       \\\hline
                      Coordinador (o socio fundador)                        & 2,0 x puntos por mérito \\\hline
                      Investigador                                          & 1,5 x puntos por mérito \\\hline
                      Técnico, Titulado Universitario en formación, u Otros & 1,0 x puntos por mérito \\\hline
                    \end{tabular}\end{table}
                  \vspace{0.4cm}

                  \textbf{Nota 1:} El aspirante y la universidad (u organismo de pertenencia del investigador que se evalúa) deben formar parte de su capital social.

                  \#\# V. ACTIVIDAD PROFESIONAL EXTERNA A LA UNIVERSIDAD

                  \textbf{Máximo: 100 puntos}
                  5.1. Actividad profesional con trayectoria constatada y contrastada de calidad, y no acreditada en apartado IV.
                  (máximo 5 méritos de actividad profesional, y un mínimo de 6 meses continuados por actividad)
                  \textbf{Complemento (multiplicativo):}
                  - Factor de corrección de 1,5 para actividades profesionales correspondientes a los últimos 5 años naturales (incluido el año natural de fecha de convocatoria de la plaza).
                  \textbf{Nota 1:} Se contempla dentro de este apartado la docencia en Institutos de Enseñanza Secundaria.

                  \#\# VI. MÉRITO PREFERENTE

                  \textbf{Profesorado Interino:}
                  El aspirante logra 5 puntos atendiendo al siguiente reparto:
                  - 5 puntos por (1) su habilitación, acreditación o ser funcionario en el cuerpo docente de Titular de Universidad o Catedrático de Universidad, o (2) hasta 5 puntos (no acumulables al punto 1) por sus estancias de investigación en universidades u otros centros de investigación externos a la Universidad de Sevilla, hasta 18 meses de estancia, y conforme a las instrucciones dictadas por la Universidad de Sevilla.

                  \textbf{Profesorado Asociado:}
                  El aspirante logra 5 puntos atendiendo al siguiente reparto:
                  - 5 puntos por (1) su habilitación, acreditación o ser funcionario en el cuerpo docente de Titular de Universidad o Catedrático de Universidad, o (2) hasta 5 puntos (no acumulables al punto 1) por sus estancias de investigación en universidades u otros centros de investigación externos a la Universidad de Sevilla, hasta 18 meses de estancia, y conforme a las instrucciones dictadas por la Universidad de Sevilla.

                  \textbf{Profesorado Ayudante Doctor:}
                  El aspirante logra 15 puntos atendiendo al siguiente reparto:
                  - 10 puntos por su habilitación, acreditación o ser funcionario en el cuerpo docente de Titular de Universidad o Catedrático de Universidad, o por su acreditación a Profesor Ayudante Doctor o a Profesor Contratado Doctor/Profesor Permanente Laboral, o hasta 10 puntos (acumulables al punto 1 y hasta 15 puntos) por sus estancias de investigación en universidades u otros centros de investigación externos a la Universidad, hasta 18 meses de estancia, y conforme a las instrucciones dictadas por la Universidad de Sevilla.

                  En los tres casos anteriores, y en el supuesto de que el aspirante esté acreditado/a, habilitado/a o sea funcionario/a de los Cuerpos Docentes Universitarios en una rama de conocimiento distinta a la que pertenece la plaza convocada, los valores admitidos se reducirán en un 75 \%. En el ANEXO I, no obstante, se podrán señalar otras ramas de conocimiento merecedoras de la puntuación máxima. Esta apreciación solo es válida en los casos de que la acreditación no tenga carácter universal.

                  ---

                  FUENTE: BAREMO 2024.pdf | PÁGINA: 4
                  \# Núm. 12/2024, de 11 de octubre

                  \#\# I. DISPOSICIONES GENERALES

                  \#\#\# I.5. Secretario General

                  3. La regulación concreta para la valoración del mérito preferente será la que se establezca en el Baremo o en las normas que, en su caso, se dicten en desarrollo de esta normativa.

                  Por su parte, en el punto VI. Mérito Preferente del Baremo aprobado junto con la normativa de contratación figura:

                  “\textbf{Profesorado Ayudante Doctor:} El aspirante logra 15 puntos atendiendo al siguiente reparto: 10 puntos por su habilitación, acreditación o ser funcionario en el cuerpo docente de Titular de Universidad o Catedrático de Universidad, o hasta 10 puntos (acumulables al punto 1 y hasta 15 puntos) por sus estancias de investigación en universidades u otros centros de investigación externos a la Universidad, hasta 18 meses de estancia, y conforme a las instrucciones dictadas por la Universidad de Sevilla”.

                  Visto lo anterior, procede una actualización tanto del artículo 10.1 de la Normativa de contratación de Profesorado Ayudante Doctor, de Profesorado Asociado y de Profesorado Sustituto como del párrafo 3 del apartado VI del baremo a aplicar.

                  \#\# TÍTULO I

                  \#\#\# Disposiciones generales

                  \textbf{Artículo 1. Objeto de régimen jurídico aplicable.}

                  1. El objeto de la presente normativa es la regulación en la Universidad de Sevilla de los procedimientos de convocatoria y provisión mediante concurso público de las siguientes plazas de personal docente e investigador: Profesorado Ayudante Doctor, Profesorado Asociado y Profesorado Interino.
                  2. Los citados procedimientos se desarrollarán en el marco de las previsiones normativas establecidas en la Ley Orgánica 6/2001, de 21 de diciembre, el Texto Refundido de la Ley Andaluza de Universidades aprobado por Decreto Legislativo 1/2013, el Estatuto de la Universidad de Sevilla, el I Convenio Colectivo del Personal Docente e Investigador en régimen laboral de las Universidades Públicas andaluzas, la presente normativa y, en lo no previsto en estas normas, por el régimen de profesorado que sea de aplicación y por las disposiciones que regulan con carácter general el ingreso en la Administración Pública.
                  3. En materia de procedimiento, y en todo lo no recogido específicamente en esta normativa, se estará a lo dispuesto en la Ley 39/2015, de 1 de octubre, del Procedimiento Administrativo Común de las Administraciones Públicas y en la Ley 40/2015, de 1 de octubre, de Régimen Jurídico del Sector Público.

                  \textbf{Artículo 2. Principios.}

                  1. La selección del personal a que se refiere la presente normativa se efectuará mediante concurso público con respeto a los principios constitucionales de igualdad, mérito, capacidad y publicidad.

                  \#\# TÍTULO II

                  \#\#\# Tramitación previa y actuación de las comisiones de selección

                  \textbf{Artículo 3. Composición, sorteo y nombramiento de las Comisiones de Contratación.}

                  1. La composición y vigencia de las Comisiones de Contratación de las plazas de Profesor Ayudante Doctor y Profesor Asociado serán las establecidas en el artículo 90 del Estatuto de la Universidad de Sevilla.
                \end{scriptsize}
        \end{itemize}


  \item \textbf{P22.} ¿Cómo se valora la participación como Investigador Principal (IP) en proyectos de investigación competitivos para obtener la puntuación máxima en Bareprof? \\
        \begin{itemize}
          \item \textbf{Tipo de consulta:} simple
          \item \textbf{Respuesta del Agente:}

                No existen requisitos mínimos excluyentes para la participación como Investigador Principal (IP) en proyectos de investigación competitivos. Sin embargo, es importante tener en cuenta que el proyecto debe ser financiado, ejecutado o en curso, y resuelto por instituciones de investigación acreditadas en la comunidad científica.

                A continuación, se presentan los criterios de evaluación y puntuación para la participación como Investigador Principal (IP) en proyectos de investigación competitivos:

                \vspace{0.4cm}
                \begin{table}[h!]\centering\renewcommand{\arraystretch}{1.3}
                  \begin{tabular}{|p{0.45\textwidth}|p{0.45\textwidth}|}\hline
                    Criterio              & Detalles                                                                                                                                                                                                                                                     \\\hline
                    Tipo de proyecto      & Proyecto competitivo: 5,0 puntos por proyecto                                                                                                                                                                                                                \\\hline
                    Ámbito geográfico     & Internacional (UE): 3,0 x puntos por proyecto, Internacional (no UE): 2,5 x puntos por proyecto, Nacional: 1,5 x puntos por proyecto, Autonómico: 1,0 x puntos por proyecto, Otros: 0,5 x puntos por proyecto                                                \\\hline
                    Duración              & De duración inferior a 2 años: 1,0 x puntos por proyecto, De duración igual o superior a 2 años: 2,0 x puntos por proyecto                                                                                                                                   \\\hline
                    Grado de contribución & Coordinador del proyecto: 3,00 x puntos por proyecto, Miembro del equipo (investigador): 2,00 x puntos por proyecto, Técnico: 1,00 x puntos por proyecto, Titulado Universitario en formación: 0,75 x puntos por proyecto, Otros: 0,50 x puntos por proyecto \\\hline
                  \end{tabular}\end{table}
                \vspace{0.4cm}

                Para obtener la puntuación máxima en Bareprof, es importante tener en cuenta los siguientes aspectos:

                * La participación como Investigador Principal (IP) en proyectos de investigación competitivos se valora con un máximo de 10 puntos.
                * El número máximo de proyectos que se pueden valorar es de 5, con una duración mínima de 1 año cada proyecto.
                * La puntuación se calcula multiplicando el número de puntos por proyecto por los complementos correspondientes (ámbito geográfico, duración y grado de contribución).

                Es importante destacar que la comisión de evaluación tendrá en cuenta la documentación acreditativa presentada por el candidato y podrá solicitar información adicional para confirmar la veracidad de los méritos alegados.
          \item \textbf{Contexto utilizado:}
                \begin{scriptsize}
                  FUENTE: BAREMO 2024.pdf | PÁGINA: 23
                  \# Núm. 12/2024, de 11 de octubre

                  \#\# I. DISPOSICIONES GENERALES

                  \#\#\# I.5. Secretario General

                  \#\#\#\# 3.5. Participación en proyectos de investigación
                  financiados ejecutados o en curso, y resueltos por instituciones de investigación acreditadas en nuestra comunidad científica:
                  \textbf{Máximo:} 10 puntos
                  (máximo 5 proyectos con una duración mínima de 1 año cada proyecto)

                  \#\#\#\#\# 3.5.1. Tipo de proyecto:
                  \textbf{No admite coeficiente de afinidad.}

                  \vspace{0.4cm}
                  \begin{table}[h!]\centering\renewcommand{\arraystretch}{1.3}
                    \begin{tabular}{|p{0.45\textwidth}|p{0.45\textwidth}|}\hline
                      Tipo de Proyecto        & Puntos por Proyecto     \\\hline
                      Proyecto competitivo    & 5,0 puntos por proyecto \\\hline
                      Proyecto no competitivo & 2,0 puntos por proyecto \\\hline
                    \end{tabular}\end{table}
                  \vspace{0.4cm}

                  \textbf{Complemento 1 (multiplicativo) – Ámbito geográfico:}

                  \vspace{0.4cm}
                  \begin{table}[h!]\centering\renewcommand{\arraystretch}{1.3}
                    \begin{tabular}{|p{0.45\textwidth}|p{0.45\textwidth}|}\hline
                      Tipo de Proyecto               & Puntos por Proyecto       \\\hline
                      Proyecto internacional (UE)    & 3,0 x puntos por proyecto \\\hline
                      Proyecto Internacional (no UE) & 2,5 x puntos por proyecto \\\hline
                      Proyecto nacional              & 1,5 x puntos por proyecto \\\hline
                      Proyecto autonómico            & 1,0 x puntos por proyecto \\\hline
                      Otros                          & 0,5 x puntos por proyecto \\\hline
                    \end{tabular}\end{table}
                  \vspace{0.4cm}

                  \textbf{Complemento 2 (multiplicando al Complemento 1) – Duración:}

                  \vspace{0.4cm}
                  \begin{table}[h!]\centering\renewcommand{\arraystretch}{1.3}
                    \begin{tabular}{|p{0.45\textwidth}|p{0.45\textwidth}|}\hline
                      Duración                              & Puntos por Proyecto       \\\hline
                      De duración inferior a 2 años         & 1,0 x puntos por proyecto \\\hline
                      De duración igual o superior a 2 años & 2,0 x puntos por proyecto \\\hline
                    \end{tabular}\end{table}
                  \vspace{0.4cm}

                  \textbf{Complemento 3 (multiplicando al Complemento 1) – Grado de contribución:}

                  \vspace{0.4cm}
                  \begin{table}[h!]\centering\renewcommand{\arraystretch}{1.3}
                    \begin{tabular}{|p{0.45\textwidth}|p{0.45\textwidth}|}\hline
                      Grado de Contribución                                  & Puntos por Proyecto        \\\hline
                      Coordinador del proyecto (total, científico o gerente) & 3,00 x puntos por proyecto \\\hline
                      Miembro del equipo (investigador)                      & 2,00 x puntos por proyecto \\\hline
                      Técnico                                                & 1,00 x puntos por proyecto \\\hline
                      Titulado Universitario en formación                    & 0,75 x puntos por proyecto \\\hline
                      Otros                                                  & 0,50 x puntos por proyecto \\\hline
                    \end{tabular}\end{table}
                  \vspace{0.4cm}

                  \textbf{Nota 1:} Se precisa que el carácter Internacional, Nacional, Autonómico u Otro del proyecto científico o tecnológico no se basará en la mera precisión formal que se indique en el nombre del proyecto sino en el ámbito real de la convocatoria, quedando al juicio técnico de la comisión la distinción.

                  \#\#\#\# 3.6. Trabajos presentados en congresos científicos o tecnológicos:
                  \textbf{Máximo:} 10 puntos
                  (máximo 10 participaciones)

                  \#\#\#\#\# 3.6.1. Tipo de participación:

                  \vspace{0.4cm}
                  \begin{table}[h!]\centering\renewcommand{\arraystretch}{1.3}
                    \begin{tabular}{|p{0.45\textwidth}|p{0.45\textwidth}|}\hline
                      Tipo de Participación                      & Puntos por Participación      \\\hline
                      Participativo-Ponencia plenaria            & 2,00 puntos por participación \\\hline
                      Participativo-Ponencia invitada            & 1,50 puntos por participación \\\hline
                      Participativo-Ponencia oral (comunicación) & 1,00 puntos por participación \\\hline
                      Participativo-Póster                       & 0,75 puntos por participación \\\hline
                    \end{tabular}\end{table}
                  \vspace{0.4cm}

                  \textbf{Complemento 1 (multiplicativo): Ámbito geográfico del congreso:}

                  \vspace{0.4cm}
                  \begin{table}[h!]\centering\renewcommand{\arraystretch}{1.3}
                    \begin{tabular}{|p{0.45\textwidth}|p{0.45\textwidth}|}\hline
                      Ámbito Geográfico & Puntos por Participación        \\\hline
                      Internacional     & 2,00 x puntos por participación \\\hline
                      Nacional          & 1,00 x puntos por participación \\\hline
                      Regional          & 0,50 x puntos por participación \\\hline
                      Local             & 0,25 x puntos por participación \\\hline
                    \end{tabular}\end{table}
                  \vspace{0.4cm}

                  \textbf{Complemento 2 (multiplicando al complemento 1): Periodicidad:}

                  ---

                  FUENTE: 27830\_guia-ayuda-candidato-r04-3-.pdf | PÁGINA: 42
                  \# Guía de ayuda al candidato

                  - Para la cumplimentación del campo Resultados relevantes y, en particular, la consulta de líneas de investigación y áreas afines, véanse las notas correspondientes del apartado 6.5.2 (Libros), situadas en la página 36, que son igualmente aplicables a este caso.
                  - En caso de introducir más de un mérito en este apartado, el criterio de ordenación será la fecha de publicación, en orden descendente.

                  \#\# Documentación acreditativa:

                  - El candidato debe incluir en la documentación aportada copia de las páginas inicial y final del artículo y, en cualquier caso, aquella donde figure un resumen de este. Se deberá también incluir el índice del volumen o número de la revista donde ha sido publicado el artículo si dicha información no aparece reflejada en la primera página de este.
                  - Como prueba fehaciente del mérito no se considera la copia del manuscrito ni tampoco la carta de ningún editor informando de la aceptación del artículo.

                  La posición de la publicación en JCR/SJR se determina a partir del ISSN de la revista y del año de publicación del artículo. Es posible que un artículo se haya publicado en una revista que actualmente se encuentre en JCR/SJR, pero que no lo estuviera en el año de su publicación, por lo que Bareprof no determinará ninguna posición (salvo que estuviera contemplada en el Anexo I).

                  \#\# 6.5.5 Participación en proyectos de investigación

                  \#\#\# Participación en proyectos de investigación

                  3.5. Participación en proyectos de investigación financiados ejecutados o en curso, y resueltos por instituciones de investigación acreditadas en nuestra comunidad científica

                  - \textbf{Epígrafe baremo}: Participación en proyectos de investigación financiados ejecutados o en curso, y resueltos por instituciones de investigación acreditadas en nuestra comunidad científica
                  - \textbf{Pestaña del CVN}: Experiencia científica y tecnológica

                  \#\#\#\# Apartado

                  - Proyectos de I+D+i financiados en convocatorias competitivas de Administraciones o entidades públicas y privadas

                  \#\#\#\# Campos

                  - Nombre del proyecto
                  - Entidad financiadora
                  - Fecha de inicio
                  - Fecha de finalización
                  - Ámbito geográfico
                  - Grado de contribución
                  - Cód. según financiadora
                  - Resultados relevantes
                  - Aportación del solicitante

                  \#\#\# Indicaciones:

                  - Con independencia del carácter competitivo o no competitivo del proyecto, este se introducirá en el apartado indicado (Proyectos de I+D+i financiados en convocatorias competitivas de Administraciones o entidades públicas y privadas).

                  Revisión 4 (marzo, 2025) 41

                  ---

                  FUENTE: 27831\_instrucciones-baremacinin-r03.pdf | PÁGINA: 28
                  \# Bareprof

                  \#\# Instrucciones de baremación

                  \#\#\# 4.4.3 Capítulos de Libro (epígrafe 3.3 del baremo)

                  - No se consideran capítulos de libros las contribuciones a congresos editadas en libros de actas o proceedings.
                  - Además de los capítulos de libros propiamente dichos, el baremo incluye en este apartado fichas técnicas de catálogos en publicaciones, reseñas y aportaciones en diccionarios comunes o científicos, con una valoración del 10\% de la que corresponde al capítulo de libro. Si la aportación corresponde a uno de los tipos citados, pero el candidato ha indicado como tipo de producción “Capítulo de libro”, la Comisión no debe validar el mérito.
                  - El editor de CVN no incorpora ninguna opción para indicar como tipo de producción las fichas técnicas de catálogos en publicaciones. En la Guía del Ayuda al Candidato se indica que, en tales casos, debe seleccionar “Reseña” en el campo Tipo de Producción. También sería válido si seleccionaran “Diccionario científico” o “Diccionario común”, pues la valoración es la misma, pero no así “Capítulo de libro”; en este último caso, el mérito no debe validarse.
                  - Para la asignación de la afinidad al mérito, consulte las consideraciones realizadas en la introducción al apartado 4.4.

                  \#\#\# 4.4.4 Artículos científicos de Revista (epígrafe 3.4 del baremo)

                  - Los datos de citas de un artículo científico deben ser introducidos por el candidato. La aplicación Bareprof no comprueba tales datos de citas.
                  - El baremo no establece distinción entre las fuentes de citas, valorando el mayor número de citas validado por la comisión.
                  - La aplicación Bareprof comprueba automáticamente la posición del artículo en los distintos índices de impacto considerados en el baremo. Para que este proceso sea efectivo, el candidato ha debido introducir correctamente tanto el ISSN de la revista como el año de publicación del artículo. La comisión está facultada para modificar el ISSN en caso de que se compruebe que el identificador introducido por el candidato es incorrecto.
                  - En pantalla se mostrará la posición de la revista en los diferentes listados automatizados (JCR, SJR o Anexo I). Para cada índice, se utilizará el listado que corresponde al año de publicación del artículo o, en su defecto, al año inmediatamente anterior si el primero no se encuentra aún disponible a la finalización del periodo de admisión de solicitudes. Bareprof usará la mejor posición para realizar la baremación.
                  - Para la asignación de la afinidad, consulte las consideraciones realizadas en la introducción al apartado 4.4.

                  \#\#\# 4.4.5 Participación en proyectos de investigación (epígrafe 3.5 del baremo)

                  - La pertenencia a grupos de investigación no es un mérito considerado en el Baremo. Muchos candidatos incluyen estos méritos como proyectos de investigación, que no deben ser validados por la Comisión.
                  - Los méritos englobados en este epígrafe del baremo no admiten coeficiente de afinidad, por lo que la Comisión NO puede invalidarlos alegando falta de afinidad con la plaza. Si el mérito ha sido consignado en el apartado correcto y se encuentra acreditado documentalmente, la Comisión debe validarlo. Por este mismo motivo, no es preciso que el candidato justifique la afinidad del.

                  ----

                  Revisión 3 (diciembre, 2024) 27

                  ---

                  FUENTE: 27831\_instrucciones-baremacinin-r03.pdf | PÁGINA: 25
                  \# Instrucciones de baremación

                  - Los contratos de investigación corresponden al apartado 3.1 del baremo, Contratos de investigación (o becas, si procede) en procesos competitivos convocados por entidades públicas o privadas, por lo que el candidato no debe incluirlos en este apartado. En caso de que lo sean, la comisión no deberá validarlos.
                  - La comisión no debe modificar las fechas de finalización de los contratos en caso de que esta sea posterior a la fecha de finalización del periodo de admisión de solicitudes establecida en la convocatoria del concurso. La aplicación Bareprof computará los meses efectivamente disfrutados por el aspirante hasta dicha fecha fin de admisión de solicitudes. En caso de que el candidato no haya introducido correctamente algún dato, este deberá ser corregido por la comisión (vista la documentación aportada).
                  - Los méritos correspondientes a este epígrafe del baremo no admiten coeficiente de afinidad, por lo que la Comisión NO puede invalidarlos alegando falta de afinidad con la plaza. Si el mérito ha sido consignado en el apartado correcto y se encuentra acreditado documentalmente, la Comisión debe validarlo.

                  \#\# 4.3.3 Dirección de tesis doctorales (epígrafe 2.3 del baremo)

                  - La actividad de dirección solo se validará en caso de que la defensa pública de la tesis haya tenido lugar. No se considerarán las labores de dirección de tesis en desarrollo.
                  - Los méritos englobados en este epígrafe del baremo no admiten coeficiente de afinidad, por lo que la Comisión NO puede invalidarlos alegando falta de afinidad con la plaza. Si el mérito ha sido consignado en el apartado correcto y se encuentra acreditado documentalmente, la Comisión debe validarlo.

                  \#\# 4.3.4 Otros méritos de docencia universitaria (epígrafe 2.4 del baremo)

                  - La dirección de Trabajos Fin de Estudios (epígrafe 2.4.1 del baremo) no requiere de la defensa del trabajo para su validación, sino del mero inicio de la tutela.
                  - Los méritos correspondientes a los epígrafes 2.4.1, 2.4.3, 2.4.5 y 2.4.6 del baremo no admiten coeficiente de afinidad, por lo que la Comisión NO puede invalidarlos alegando falta de afinidad con la plaza. Si el mérito ha sido consignado en el apartado correcto y se encuentra acreditado documentalmente, la Comisión debe validarlo.

                  \#\# 4.4 Actividad investigadora (bloque III del baremo)

                  En aquellos méritos de este apartado que permiten coeficiente de afinidad (libros o monografías científicas, capítulos de libros, artículos científicos y trabajos presentados en congresos), la afinidad muy alta se otorga a méritos que correspondan al perfil investigador de la plaza convocada o, en su defecto, a las líneas estratégicas de investigación del área de conocimiento del departamento de la Universidad de Sevilla solicitante de la plaza, que, según el baremo de la Universidad de Sevilla, son las practicadas por los grupos de investigación, proyectos de investigación competitivos e institutos de investigación universitarios en los que participe el departamento y área de conocimiento convocante de la plaza. Tales líneas de investigación pueden consultarse en Bareprof y deberán ser tenidas en cuenta por la comisión para determinar la afinidad del mérito. Debe precisarse que el listado de líneas de investigación que recoge Bareprof no tiene carácter vinculante, por lo que, si un candidato acredita documentalmente que un mérito corresponde a una línea de investigación estratégica del departamento y área de conocimiento, la Comisión deberá.

                  Revisión 3 (diciembre, 2024) 24

                  ---

                  FUENTE: 27831\_instrucciones-baremacinin-r03.pdf | PÁGINA: 17
                  \# Bareprof

                  \#\# Instrucciones de baremación

                  No obstante, la comisión podrá, si así lo estima en el ejercicio de sus facultades, valerse de cualquier otro medio de prueba complementario y válido en derecho para decidir si procede validar o corregir el currículum del aspirante leído por Bareprof. Por ejemplo, puede recurrir a medios adicionales de prueba para confirmar si la convocatoria a que corresponde un contrato de investigación (personal investigador) es competitiva (podría corresponder en ese caso al apartado 3.1) o no competitiva (en su lugar, podría corresponder al bloque V si fuera actividad externa a la Institución universitaria). En este sentido, la normativa habilita a las comisiones a solicitar (excepcionalmente) al candidato documentación adicional a la ya debidamente entregada para confirmar la veracidad o adecuación de un mérito y siempre siguiendo los mismos criterios para todos los aspirantes. El requerimiento de dicha documentación aclaratoria, que deberá quedar reflejado en las actas de la comisión, se realizará mediante escrito oficial del presidente de la comisión remitido por correo postal certificado o mediante correo electrónico desde su cuenta de correo institucional @us.es. En ningún caso este procedimiento permitirá al candidato subsanar la documentación acreditativa que debió aportar al solicitar su participación en el concurso.

                  Finalmente, en caso de que la comisión no disponga de medio de prueba válido y existan dudas razonables de la existencia del mérito o del indicio de calidad que permite complementar su puntuación, se deberá actuar de la siguiente forma:
                  - Si no dispone de prueba fiable de la existencia del mérito, este no se debe validar.
                  - Si la existencia del indicio de calidad no está acreditada (pero sí el mérito):
                  a) Se procederá a validar el mérito.
                  b) Si el indicio de calidad no es obligatorio (es decir, no aparece recuadrado en rojo en Bareprof si no se incluye), se debe dejar en blanco.
                  c) Si el indicio de calidad es obligatorio, se debe elegir el complemento de menor puntuación según el baremo.

                  El cuadro Justificación debe rellenarse obligatoriamente tras cualquier corrección de un dato del aspirante por parte de la comisión. Su uso es opcional en cualquier otro caso.

                  \#\# 4.1.2 Duplicidad de méritos

                  Las Comisiones no baremarán el mismo mérito en dos apartados distintos del baremo (no admitiéndose, pues, la duplicidad de méritos) y podrán, asumida la buena fe del candidato, baremarlo exclusivamente en aquel apartado al que corresponde por su mayor idoneidad conceptual. No obstante, no se entenderá como duplicidad el caso de que un aspirante incluya como mérito su contrato como personal investigador en Contratos (o becas, si procede) en procesos competitivos convocados por entidades públicas o privadas (epígrafe 3.1 del baremo), al mismo tiempo que registra su docencia reglada impartida durante su contrato en Docencia en asignaturas regladas universitarias (epígrafe 2.1). Del mismo modo, el candidato no solo podrá alegar sus contratos como PDI (por ejemplo, sustituciones interinas), sino también la producción científica o docencia impartida durante la vigencia de tales contratos. Análogamente, no se considera que existe duplicidad de méritos en las tesis doctorales con respecto a la producción científica derivada de las mismas. Tampoco se considera que exista duplicidad de méritos cuando el candidato ha disfrutado de un contrato de investigación (epígrafe 3.1 del baremo) que conlleva la realización de la investigación, parcial o totalmente, en una Universidad o Centro de Investigación distinto de aquel que ha dado lugar al contrato (por ejemplo, convocatorias Marie Sklowdoska Curie o Margarita Salas) y, al mismo tiempo, aporta dicha estancia como mérito preferente.

                  Revisión 3 (diciembre, 2024) 16
                \end{scriptsize}
        \end{itemize}


  \item \textbf{P23.} En el apartado de transferencia del conocimiento del Baremo 2024, ¿cuántos puntos exactos se pueden conseguir como máximo mediante patentes en explotación? \\
        \begin{itemize}
          \item \textbf{Tipo de consulta:} simple
          \item \textbf{Respuesta del Agente:}

                No existen requisitos mínimos excluyentes para este apartado.

                A continuación, se detallan los criterios de evaluación para el apartado de patentes en explotación:

                \vspace{0.4cm}
                \begin{table}[h!]\centering\renewcommand{\arraystretch}{1.3}
                  \begin{tabular}{|p{0.45\textwidth}|p{0.45\textwidth}|}\hline
                    Criterio              & Detalles                                                                      \\\hline
                    Patentes de Invención & Ámbito geográfico: Internacional (4,0 puntos), Nacional (3,0 puntos)          \\\hline
                                          & Complemento 1: PCT (2,0 x puntos), No PCT (1,0 x puntos)                      \\\hline
                                          & Complemento 2: En explotación, en exclusiva (2,0 x puntos), No (0,5 x puntos) \\\hline
                    Modelos de Utilidad   & Ámbito geográfico: Internacional (4,0 puntos), Nacional (3,0 puntos)          \\\hline
                                          & Complemento 1: En explotación, en exclusiva (1,5 x puntos), No (0,5 x puntos) \\\hline
                    Diseños Industriales  & Ámbito geográfico: Internacional (3,0 puntos), Nacional (2,0 puntos)          \\\hline
                                          & Complemento 1: En explotación, en exclusiva (1,5 x puntos), No (0,5 x puntos) \\\hline
                  \end{tabular}\end{table}
                \vspace{0.4cm}

                Para calcular los puntos máximos que se pueden conseguir mediante patentes en explotación, debemos considerar los siguientes escenarios:

                1. \textbf{Patentes de Invención}:
                * Internacional con PCT y en explotación: 4,0 puntos x 2,0 (PCT) x 2,0 (en explotación) = 16,0 puntos
                * Nacional con PCT y en explotación: 3,0 puntos x 2,0 (PCT) x 2,0 (en explotación) = 12,0 puntos
                2. \textbf{Modelos de Utilidad}:
                * Internacional en explotación: 4,0 puntos x 1,5 (en explotación) = 6,0 puntos
                * Nacional en explotación: 3,0 puntos x 1,5 (en explotación) = 4,5 puntos
                3. \textbf{Diseños Industriales}:
                * Internacional en explotación: 3,0 puntos x 1,5 (en explotación) = 4,5 puntos
                * Nacional en explotación: 2,0 puntos x 1,5 (en explotación) = 3,0 puntos

                Considerando que se pueden presentar un máximo de 5 patentes, los puntos máximos que se pueden conseguir mediante patentes en explotación serían:

                * 5 patentes de invención internacionales con PCT y en explotación: 5 x 16,0 puntos = 80,0 puntos (superior al límite de 20 puntos para este apartado)
                * Combinación de patentes de invención, modelos de utilidad y diseños industriales en explotación: el máximo teórico sería 20 puntos, considerando la combinación de patentes que maximice la puntuación.

                En resumen, el máximo de puntos que se pueden conseguir mediante patentes en explotación es de \textbf{20 puntos}, siempre que se cumplan los requisitos y se maximice la puntuación con la combinación de patentes adecuada.
          \item \textbf{Contexto utilizado:}
                \begin{scriptsize}
                  FUENTE: BAREMO 2024.pdf | PÁGINA: 26
                  \textbf{Nota 1:} No deben registrarse en este apartado las empresas spin-off de I+D+i, que se registrarán en el apartado 4.4.
                  \textbf{Nota 2:} Únicamente se debe incluir en este apartado aquel mérito que demuestre su carácter de investigación capaz de generar conocimiento.

                  \#\#\#\# 4.3. Patentes:
                  \textbf{Máximo:} 20 puntos (máximo 5 patentes)

                  \vspace{0.4cm}
                  \begin{table}[h!]\centering\renewcommand{\arraystretch}{1.3}
                    \begin{tabular}{|p{0.45\textwidth}|p{0.45\textwidth}|}\hline
                      Concepto                                                 & Puntuación por mérito   \\\hline
                      4.3.1. Patentes de Invención:                            &                         \\\hline
                      Ámbito geográfico:                                       &                         \\\hline
                      - Internacional                                          & 4,0 puntos por mérito   \\\hline
                      - Nacional, y otros                                      & 3,0 puntos por mérito   \\\hline
                      \textbf{Complemento 1 (multiplicativo a 4.3.1):}         &                         \\\hline
                      - PCT                                                    & 2,0 x puntos por mérito \\\hline
                      - No PCT                                                 & 1,0 x puntos por mérito \\\hline
                      \textbf{Complemento 2 (multiplicando al complemento 1):} &                         \\\hline
                      En explotación, en exclusiva:                            &                         \\\hline
                      - Sí                                                     & 2,0 x puntos por mérito \\\hline
                      - No                                                     & 0,5 x puntos por mérito \\\hline
                      4.3.2. Modelos de Utilidad:                              &                         \\\hline
                      Ámbito geográfico:                                       &                         \\\hline
                      - Internacional                                          & 4,0 puntos por mérito   \\\hline
                      - Nacional, y otros                                      & 3,0 puntos por mérito   \\\hline
                      \textbf{Complemento 1 (multiplicativo a 4.3.2):}         &                         \\\hline
                      En explotación, en exclusiva:                            &                         \\\hline
                      - Sí                                                     & 1,5 x puntos por mérito \\\hline
                      - No                                                     & 0,5 x puntos por mérito \\\hline
                      4.3.3. Diseños Industriales:                             &                         \\\hline
                      Ámbito geográfico:                                       &                         \\\hline
                      - Internacional                                          & 3,0 puntos por mérito   \\\hline
                      - Nacional, y otros                                      & 2,0 puntos por mérito   \\\hline
                      \textbf{Complemento 1 (multiplicativo a 4.3.3):}         &                         \\\hline
                      En explotación, en exclusiva:                            &                         \\\hline
                      - Sí                                                     & 1,5 x puntos por mérito \\\hline
                      - No                                                     & 0,5 x puntos por mérito \\\hline
                    \end{tabular}\end{table}
                  \vspace{0.4cm}

                  ---

                  FUENTE: BAREMO 2024.pdf | PÁGINA: 26
                  \# Núm. 12/2024, de 11 de octubre

                  \#\# I. DISPOSICIONES GENERALES

                  \#\#\# I.5. Secretario General

                  \#\#\#\# 4.2. Obra civil o de edificación, o exposiciones:
                  \textbf{Máximo:} 35 (máximo 10 exposiciones) puntos

                  \vspace{0.4cm}
                  \begin{table}[h!]\centering\renewcommand{\arraystretch}{1.3}
                    \begin{tabular}{|p{0.45\textwidth}|p{0.45\textwidth}|}\hline
                      Concepto                                                 & Puntuación por mérito        \\\hline
                      4.2.1. Obras o exposiciones artísticas                   & 2,00 puntos por mérito       \\\hline
                      \textbf{Complemento 1 (multiplicativo):}                 &                              \\\hline
                      Autores:                                                 &                              \\\hline
                      - Individual (único autor)                               & 3,00 x puntuación por mérito \\\hline
                      - Colectiva (al menos dos autores)                       & 1,00 x puntuación por mérito \\\hline
                      \textbf{Complemento 2 (multiplicando el Complemento 1):} &                              \\\hline
                      Ámbito geográfico:                                       &                              \\\hline
                      - Internacional                                          & 2,00 x puntuación por mérito \\\hline
                      - Nacional, y otros                                      & 1,00 x puntuación por mérito \\\hline
                      \textbf{Complemento 3 (multiplicando el Complemento 1):} &                              \\\hline
                      Comisario de exposición                                  & 2,00 x puntuación por mérito \\\hline
                      \textbf{Complemento 4 (multiplicando el Complemento 1):} &                              \\\hline
                      Monográfica                                              & 1,50 x puntuación por mérito \\\hline
                      \textbf{Complemento 5 (multiplicando el Complemento 1):} &                              \\\hline
                      Catálogo                                                 & 1,50 x puntuación por mérito \\\hline
                    \end{tabular}\end{table}
                  \vspace{0.4cm}

                  \textbf{Nota 1:} No deben registrarse en este apartado las empresas spin-off de I+D+i, que se registrarán en el apartado 4.4.
                  \textbf{Nota 2:} Únicamente se debe incluir en este apartado aquel mérito que demuestre su carácter de investigación capaz de generar conocimiento.

                  \#\#\#\# 4.3. Patentes:
                  \textbf{Máximo:} 20 puntos (máximo 5 patentes)

                  ---

                  FUENTE: BAREMO 2024.pdf | PÁGINA: 25
                  \# Núm. 12/2024, de 11 de octubre

                  \#\# I. DISPOSICIONES GENERALES

                  \#\#\# I.5. Secretario General

                  \#\#\#\# 3.7. Comités científicos, técnicos o asesores:
                  \textbf{Máximo: 5 puntos}
                  (máximo 5 méritos por epígrafe)

                  \vspace{0.4cm}
                  \begin{table}[h!]\centering\renewcommand{\arraystretch}{1.3}
                    \begin{tabular}{|p{0.45\textwidth}|p{0.45\textwidth}|}\hline
                      Modalidad de la actividad                                                                                                                                                                                    & Puntuación            \\\hline
                      Participación en comités científicos o comités de revisión de artículos en revistas científicas o tecnológicas indexadas con factor de impacto JCR o SJR, o revistas contenidas en ANEXO I del apartado 3.4. & 1,0 punto por mérito  \\\hline
                      Organizador principal (presidente del comité) de Congresos de prestigio del apartado 3.6.                                                                                                                    & 3,0 puntos por mérito \\\hline
                      Pertenencia a Comités Científico y Organizador de Congresos de prestigio del apartado 3.6.                                                                                                                   & 1,0 punto por mérito  \\\hline
                    \end{tabular}\end{table}
                  \vspace{0.4cm}

                  \#\#\#\# 3.8. Conocimiento de idiomas relevantes para la investigación y la docencia:
                  \textbf{Máximo: 5 puntos}

                  \vspace{0.4cm}
                  \begin{table}[h!]\centering\renewcommand{\arraystretch}{1.3}
                    \begin{tabular}{|p{0.45\textwidth}|p{0.45\textwidth}|}\hline
                      Nivel                          & Puntuación         \\\hline
                      C2, o equivalente según idioma & 5 puntos por nivel \\\hline
                      C1, o equivalente según idioma & 3 puntos por nivel \\\hline
                      B2, o equivalente según idioma & 2 puntos por nivel \\\hline
                    \end{tabular}\end{table}
                  \vspace{0.4cm}

                  \textbf{Notas:}
                  1. El mérito alegado se valora incluso si la institución que lo otorga atribuye al certificado un periodo de validez temporal determinada y ya extinguido.
                  2. El aspirante puede alegar más de un idioma.
                  3. El Departamento deberá especificar en ANEXO I los idiomas a valorar. Se exceptúa de la valoración de este mérito a aquellas plazas en que el uso del idioma acreditado sea consustancial con su docencia (e.g., filologías).

                  \#\# IV. TRANSFERENCIA DEL CONOCIMIENTO
                  \textbf{Máximo: 100 puntos}

                  \#\#\#\# 4.1. Contratos de transferencia (artículo 83 de la LOU 6/2001) o convenios relacionados directa o indirectamente con los ámbitos materiales de las respectivas áreas de conocimiento con instituciones públicas o privadas en el ámbito cultural, social, técnico, económico y empresarial que contribuyen a la innovación del tejido económico y social, y no contemplados en el apartado V:\textit{}
                  \textbf{Máximo: 35 puntos}
                  (máximo 5 contratos de transferencia)

                  \vspace{0.4cm}
                  \begin{table}[h!]\centering\renewcommand{\arraystretch}{1.3}
                    \begin{tabular}{|p{0.45\textwidth}|p{0.45\textwidth}|}\hline
                      Grado de contribución               & Puntuación                  \\\hline
                      Coordinador                         & 20 puntos por coordinación  \\\hline
                      Investigador                        & 15 puntos por participación \\\hline
                      Técnico                             & 8 puntos por participación  \\\hline
                      Titulado Universitario en formación & 8 puntos por participación  \\\hline
                      Otros                               & 4 puntos por participación  \\\hline
                    \end{tabular}\end{table}
                  \vspace{0.4cm}

                  \textbf{Complemento 1 (multiplicativo):}
                  \vspace{0.4cm}
                  \begin{table}[h!]\centering\renewcommand{\arraystretch}{1.3}
                    \begin{tabular}{|p{0.45\textwidth}|p{0.45\textwidth}|}\hline
                      Ámbito geográfico          & Puntuación              \\\hline
                      Internacional (UE o no UE) & 3,0 x puntos por mérito \\\hline
                      Nacional                   & 2,0 x puntos por mérito \\\hline
                      Autonómico                 & 1,5 x puntos por mérito \\\hline
                      Otros                      & 1,0 x puntos por mérito \\\hline
                    \end{tabular}\end{table}
                  \vspace{0.4cm}

                  \textbf{Notas:}
                  1. Se contemplan en este apartado los cursos de formación (al amparo del art. 83 de la LOU o también convenios formalizados) en empresas o administraciones públicas (distintas de universidades u organismos de investigación).
                  2. Únicamente se debe incluir en este apartado aquel mérito que demuestre su carácter de investigación capaz de generar conocimiento.

                  ---

                  FUENTE: CRITERIOS DE EVALUACIÓN Y REQUISITOS MÍNIMOS.pdf | PÁGINA: 6
                  \# Resolución CNEAI-ANECA sobre Acreditación

                  \#\# 1.2.2. Transferencia e intercambio de conocimiento y actividad de carácter profesional

                  Para la acreditación a PTU, de las 8 contribuciones que pueden aportarse, un máximo de 3 puede corresponderse con resultados de transferencia e intercambio de conocimiento, que estén claramente diferenciados de los de investigación, docencia u otros méritos curriculares, y que atiendan a las características propias de cada ámbito del conocimiento.

                  Para la acreditación a CU, de las 20 contribuciones que pueden aportarse, un máximo de 5 puede corresponderse con estos resultados de transferencia e intercambio de conocimiento.

                  Se reconocerán automáticamente 10 puntos por transferencia e intercambio de conocimiento a quienes hayan obtenido un sexenio de transferencia CNEAI-ANECA. Podrán presentarse otras contribuciones para alcanzar el máximo en este apartado, siempre que no hayan sido ya presentadas para obtener el sexenio de transferencia.

                  En defecto de sexenio de transferencia, se podrán incluir méritos que respondan a alguna de las siguientes categorías:
                  - Actividades de carácter científico, técnico o artístico que generen valor económico o valor social a través de la transferencia e intercambio de conocimiento.
                  - Derechos de propiedad industrial e intelectual en explotación.
                  - Actividades con valor social, artístico o cultural.
                  - Participación en Cátedras institucionales y de empresa, o similares, que constituyan una fórmula de colaboración con instituciones públicas y privadas que, entre sus objetivos, lleven a cabo actividades de transferencia de conocimiento.
                  - Actividades en el ámbito del fomento de la cultura científica, tecnológica, de innovación y de la ciencia ciudadana.
                  - Asesoramiento científico a administraciones públicas.
                  - Asesoramiento legislativo, ponencia de leyes.
                  - Transferencia de conocimiento realizada al amparo de servicios desempeñados en entidades públicas o privadas diferentes de la gestión universitaria, por ejemplo, otras administraciones públicas, instituciones u organismos nacionales o internacionales.
                  - Derechos de propiedad industrial e intelectual derivados de la actividad investigadora.
                  - Participación en el equipo promotor de Empresas Basadas en el Conocimiento a partir de la actividad investigadora.
                  - Participación en la realización de ensayos clínicos, etc.

                  En todos los casos se valorará el tipo de participación (responsable o miembro de equipo), la duración, la financiación, la contratación de personal, el grado de innovación y de multidisciplinariedad o interdisciplinariedad, los resultados, productos o servicios resultantes, y el impacto social logrado.

                  \textbf{DIRECCIÓN DE ANECA – 20/03/2024}

                  ---

                  FUENTE: 27831\_instrucciones-baremacinin-r03.pdf | PÁGINA: 30
                  \# Instrucciones de baremación

                  - Los méritos englobados en este epígrafe del baremo no admiten coeficiente de afinidad, por lo que la Comisión NO puede invalidarlos alegando falta de afinidad con la plaza. Si el mérito ha sido consignado en el apartado correcto y se encuentra acreditado documentalmente, la Comisión debe validarlo.

                  \#\# 4.5 Transferencia del conocimiento (bloque IV del baremo)

                  \#\#\# 4.5.1 Contratos de transferencia (epígrafe 4.1 del baremo)

                  No se recogen indicaciones relativas a este apartado en la versión actual de este documento.

                  \#\#\# 4.5.2 Obra civil o de edificación, o exposiciones (epígrafe 4.2 del baremo)

                  No se recogen indicaciones relativas a este apartado en la versión actual de este documento.

                  \#\#\# 4.5.3 Patentes (epígrafe 4.3 del baremo)

                  Para que este mérito pueda validarse, debe acreditarse su concesión, no solo su solicitud.

                  \#\#\# 4.5.4 Spin-off (epígrafe 4.4 del baremo)

                  No se recogen indicaciones relativas a este apartado en la versión actual de este documento.

                  \#\# 4.6 Actividad profesional externa (bloque V del baremo)

                  - Este apartado se corresponde con Actividad profesional externa a la actividad académica dentro de la Universidad. Por tanto, se admiten los contratos en el ámbito universitario como Personal Técnico, de Gestión y de Administración y Servicios (aunque no como PDI).
                  - Se contempla dentro de este apartado la docencia en Institutos de Enseñanza Secundaria.

                  \#\# 4.7 Mérito preferente (bloque VI del baremo)

                  \#\#\# 4.7.1 Acreditación

                  - En los concursos correspondientes a plazas de Profesorado Asociado y Profesor Sustituto, las únicas acreditaciones que pueden validarse son las correspondientes a categorías de cuerpos docentes universitarios, es decir:
                  - Profesor Titular de Universidad
                  - Catedrático de Universidad
                  Las demás acreditaciones (Profesor Ayudante Doctor, Profesor Contratado Doctor) no deberán ser validadas.
                  - En lo que respecta a las plazas de Profesorado Ayudante Doctor, se admiten como mérito preferente, además de las acreditaciones anteriormente citadas, las correspondientes a Profesor Ayudante Doctor y a Profesor Contratado Doctor/Profesor Permanente Laboral. Esta distinción con respecto a las figuras citadas en el punto anterior tiene su origen en los cambios introducidos.

                  Revisión 3 (diciembre, 2024) 29
                \end{scriptsize}
        \end{itemize}


  \item \textbf{P24.} Según las instrucciones del baremo vigentes, ¿se puede valorar un mismo artículo simultáneamente en los apartados de investigación y docencia? \\
        \begin{itemize}
          \item \textbf{Tipo de consulta:} simple
          \item \textbf{Respuesta del Agente:}

                \textbf{Requisitos mínimos excluyentes:}
                * No se admitirá la duplicidad de méritos en dos apartados distintos del baremo.
                * Los méritos deben corresponder conceptualmente al apartado en el que se incluyan.

                \textbf{Criterios de evaluación:}

                \vspace{0.4cm}
                \begin{table}[h!]\centering\renewcommand{\arraystretch}{1.3}
                  \begin{tabular}{|p{0.45\textwidth}|p{0.45\textwidth}|}\hline
                    Criterio                & Detalles                                                                          \\\hline
                    Duplicidad de méritos   & No se admitirá la duplicidad de méritos en dos apartados distintos del baremo.    \\\hline
                    Idoneidad conceptual    & Los méritos deben corresponder conceptualmente al apartado en el que se incluyan. \\\hline
                    Actividad docente       & Se valorará la docencia realmente impartida, no la nominal del contrato.          \\\hline
                    Actividad investigadora & Se valorará la producción científica y la investigación realizada.                \\\hline
                  \end{tabular}\end{table}
                \vspace{0.4cm}

                \textbf{Puntuación:}

                * La puntuación se calculará según los criterios establecidos en el baremo vigente.
                * La comisión de contratación podrá solicitar documentación adicional para confirmar la veracidad o adecuación de un mérito.

                \textbf{Respuesta a la pregunta del usuario:}
                No, según las instrucciones del baremo vigentes, no se puede valorar un mismo artículo simultáneamente en los apartados de investigación y docencia, ya que se consideraría duplicidad de méritos. Sin embargo, se pueden valorar diferentes aspectos de un mismo trabajo, como la producción científica y la docencia impartida, siempre que se incluyan en apartados diferentes y no se considere duplicidad de méritos.

                \textbf{Ejemplos de valoración:}

                * Un artículo publicado en una revista científica se puede valorar en el apartado de investigación, pero no se puede valorar el mismo artículo en el apartado de docencia.
                * Un contrato de investigación se puede valorar en el apartado de investigación, y la docencia impartida durante ese contrato se puede valorar en el apartado de docencia, siempre que se cumplan los requisitos establecidos en el baremo.
          \item \textbf{Contexto utilizado:}
                \begin{scriptsize}
                  FUENTE: 27831\_instrucciones-baremacinin-r03.pdf | PÁGINA: 17
                  \# Bareprof

                  \#\# Instrucciones de baremación

                  No obstante, la comisión podrá, si así lo estima en el ejercicio de sus facultades, valerse de cualquier otro medio de prueba complementario y válido en derecho para decidir si procede validar o corregir el currículum del aspirante leído por Bareprof. Por ejemplo, puede recurrir a medios adicionales de prueba para confirmar si la convocatoria a que corresponde un contrato de investigación (personal investigador) es competitiva (podría corresponder en ese caso al apartado 3.1) o no competitiva (en su lugar, podría corresponder al bloque V si fuera actividad externa a la Institución universitaria). En este sentido, la normativa habilita a las comisiones a solicitar (excepcionalmente) al candidato documentación adicional a la ya debidamente entregada para confirmar la veracidad o adecuación de un mérito y siempre siguiendo los mismos criterios para todos los aspirantes. El requerimiento de dicha documentación aclaratoria, que deberá quedar reflejado en las actas de la comisión, se realizará mediante escrito oficial del presidente de la comisión remitido por correo postal certificado o mediante correo electrónico desde su cuenta de correo institucional @us.es. En ningún caso este procedimiento permitirá al candidato subsanar la documentación acreditativa que debió aportar al solicitar su participación en el concurso.

                  Finalmente, en caso de que la comisión no disponga de medio de prueba válido y existan dudas razonables de la existencia del mérito o del indicio de calidad que permite complementar su puntuación, se deberá actuar de la siguiente forma:
                  - Si no dispone de prueba fiable de la existencia del mérito, este no se debe validar.
                  - Si la existencia del indicio de calidad no está acreditada (pero sí el mérito):
                  a) Se procederá a validar el mérito.
                  b) Si el indicio de calidad no es obligatorio (es decir, no aparece recuadrado en rojo en Bareprof si no se incluye), se debe dejar en blanco.
                  c) Si el indicio de calidad es obligatorio, se debe elegir el complemento de menor puntuación según el baremo.

                  El cuadro Justificación debe rellenarse obligatoriamente tras cualquier corrección de un dato del aspirante por parte de la comisión. Su uso es opcional en cualquier otro caso.

                  \#\# 4.1.2 Duplicidad de méritos

                  Las Comisiones no baremarán el mismo mérito en dos apartados distintos del baremo (no admitiéndose, pues, la duplicidad de méritos) y podrán, asumida la buena fe del candidato, baremarlo exclusivamente en aquel apartado al que corresponde por su mayor idoneidad conceptual. No obstante, no se entenderá como duplicidad el caso de que un aspirante incluya como mérito su contrato como personal investigador en Contratos (o becas, si procede) en procesos competitivos convocados por entidades públicas o privadas (epígrafe 3.1 del baremo), al mismo tiempo que registra su docencia reglada impartida durante su contrato en Docencia en asignaturas regladas universitarias (epígrafe 2.1). Del mismo modo, el candidato no solo podrá alegar sus contratos como PDI (por ejemplo, sustituciones interinas), sino también la producción científica o docencia impartida durante la vigencia de tales contratos. Análogamente, no se considera que existe duplicidad de méritos en las tesis doctorales con respecto a la producción científica derivada de las mismas. Tampoco se considera que exista duplicidad de méritos cuando el candidato ha disfrutado de un contrato de investigación (epígrafe 3.1 del baremo) que conlleva la realización de la investigación, parcial o totalmente, en una Universidad o Centro de Investigación distinto de aquel que ha dado lugar al contrato (por ejemplo, convocatorias Marie Sklowdoska Curie o Margarita Salas) y, al mismo tiempo, aporta dicha estancia como mérito preferente.

                  Revisión 3 (diciembre, 2024) 16

                  ---

                  FUENTE: 27830\_guia-ayuda-candidato-r04-3-.pdf | PÁGINA: 18
                  \# Guía de ayuda al candidato

                  \#\# 5.2 Reubicación de méritos

                  Las comisiones de contratación no pueden reubicar los méritos que han sido incluidos en apartados a los que no corresponden conceptualmente, por lo que estos méritos no podrán ser validados. Por ejemplo, si el candidato ha incluido como mérito en el epígrafe 2.1 (Docencia en asignaturas regladas universitarias) su actividad docente en un Máster propio, ésta no será validada por la comisión, pues corresponde al epígrafe 2.4.2 (Docencia universitaria no oficial). Para evitar este tipo de inconvenientes, debe prestar mucha atención a la información que se proporciona para cada mérito en el apartado 6 de esta guía, donde se indica de forma precisa en qué apartado del CVN se debe consignar la información (así como los campos que deben cumplimentarse).

                  \vspace{0.4cm}
                  \begin{table}[h!]\centering\renewcommand{\arraystretch}{1.3}
                    \begin{tabular}{|p{0.45\textwidth}|p{0.45\textwidth}|}\hline
                      Contratos y becas en procesos competitivos &                                                                                                            \\\hline
                      Epigrafe baremo                            & 3.1. Contratos (o becas, si procede) en procesos competitivos convocados por entidades públicas o privadas \\\hline
                      Pestaña del CVN                            & Actividad científica y tecnológica                                                                         \\\hline
                      Apartado                                   & Ayudas y becas obtenidas                                                                                   \\\hline
                      Campos                                     & Nombre de la ayuda                                                                                         \\\hline
                                                                 & Entidad concesionaria                                                                                      \\\hline
                                                                 & Finalidad ["Predoctoral", "Posdoctoral"]                                                                   \\\hline
                                                                 & Fecha de concesión                                                                                         \\\hline
                                                                 & Fecha de finalización                                                                                      \\\hline
                    \end{tabular}\end{table}
                  \vspace{0.4cm}

                  \#\# 5.3 Duplicidad de méritos

                  Las Comisiones no baremarán el mismo mérito en dos apartados distintos del baremo (no admitiéndose, pues, la duplicidad de méritos) y podrán baremarlo exclusivamente en aquel apartado al que corresponde por su mayor idoneidad conceptual. No obstante, no se entenderá como duplicidad el caso de un aspirante que incluya como mérito su contrato como personal investigador en Contratos (o becas, si procede) en procesos competitivos convocados por entidades públicas o privadas (epígrafe 3.1 del baremo), al mismo tiempo que registra su docencia reglada impartida durante su contrato en Docencia en asignaturas regladas universitarias (epígrafe 2.1). Del mismo modo, el candidato no solo podrá alegar sus contratos como PDI (por ejemplo, como profesorado sustituto), sino también la producción científica o docencia impartida durante la vigencia de tales contratos. Análogamente, no se considera que existe duplicidad de méritos en las tesis doctorales con respecto a la producción científica derivada de las mismas. Tampoco se considera que exista duplicidad de méritos cuando el candidato ha disfrutado de un contrato de investigación (epígrafe 3.1 del baremo) que conlleva la realización de la investigación, parcial o totalmente, en una Universidad o Centro de Investigación distinto de aquel que ha dado lugar al contrato (por ejemplo, convocatorias Marie Sklowdoska Curie o Margarita Salas) y, al mismo tiempo, aporta dicha estancia como mérito preferente. Por el contrario, sí se entenderá como duplicidad que la ponencia expuesta en Trabajos presentados en congresos científicos o tecnológicos (apartado 3.6), sea también registrada como Capítulos de Libro (epígrafe 3.3), en tanto que se valora esencialmente el contenido del mérito y no solo el formato de publicación. También se entenderá como duplicidad el que un aspirante incluya su contrato de Investigación con una entidad no universitaria ajustado a Contratos (o becas, si procede) en procesos competitivos.

                  Revisión 4 (marzo, 2025) 17

                  ---

                  FUENTE: 27831\_instrucciones-baremacinin-r03.pdf | PÁGINA: 24
                  \# Bareprof
                  \#\# Instrucciones de baremación

                  \#\#\# 4.3 Actividad docente (bloque II del baremo)

                  \#\#\#\# 4.3.1 Docencia en asignaturas regladas universitarias (epígrafe 2.1 del baremo)

                  - En este apartado el candidato reflejará la docencia realmente impartida (no la nominal de su contrato), acreditada mediante el certificado de actividades docentes o documento equivalente.
                  - La docencia debe introducirse obligatoriamente en horas. Si la comisión observa que el candidato ha introducido el número de créditos, podrá modificarlo para reflejar el valor correcto.
                  - La docencia recogida en este apartado podrá corresponder tanto a la actividad del candidato como PDI (epígrafe 2.2 del baremo) como a su actividad docente desarrollada como Personal Investigador (epígrafe 3.1 del baremo).
                  - Las Prácticas Externas se validan como actividad docente del candidato siempre que este actúe como tutor académico; el tutor profesional ejerce la labor en calidad de responsable de la entidad colaboradora y no de PDI. Su cómputo se corresponde con lo reflejado en el certificado de actividades docentes. Por ejemplo, la Universidad de Sevilla reconoce al tutor académico 0,1 créditos por cada 6 créditos de que conste la asignatura en el Plan de Estudios.
                  - Se valorará también en este apartado la dirección o tutela de trabajos académicos que aparezca recogida y cuantificada (en número de horas o créditos) en el certificado de actividades docentes, con independencia de que la labor de tutorización sea también valorada en el apartado 2.4.1 del baremo (dirección o tutela de trabajos académicos en titulaciones universitarias oficiales). Esto no supone duplicidad en el mérito, pues se valora, por un lado, la actividad docente presencial reconocida y, por otro, la propia labor de tutorización.

                  \#\#\#\# 4.3.2 Contratos como PDI universitario (epígrafe 2.2 del baremo)

                  - En este apartado el candidato registrará su actividad como PDI universitario mediante contrato laboral o pertenencia a los cuerpos docentes universitarios. Los contratos laborales que se contemplan corresponden a las modalidades de contratación (ordinarias y extraordinarias) recogidas en los apartados 1, 2 y 3 del artículo 14 del I Convenio Colectivo del PDI Laboral de las Universidad Públicas de Andalucía (BOJA 92/2008, de 9 de mayo) o sus análogas en otros Convenios, así como las categorías análogas de PDI de Universidades privadas.
                  - Profesor Ayudante
                  - Profesor Ayudante Doctor
                  - Profesor Colaborador
                  - Profesor Contratado Doctor
                  - Profesor Asociado
                  - Profesor Asociado de concierto con las instituciones sanitarias
                  - Profesor Sustituto Interino

                  Las figuras englobadas dentro de los cuerpos docentes universitarios son las siguientes:
                  - Profesor Titular de Escuela Universitaria
                  - Catedrático de Escuela Universitaria
                  - Profesor Titular de Universidad
                  - Catedrático de Universidad
                  - Figuras análogas en otros ámbitos internacionales

                  Revisión 3 (diciembre, 2024) 23

                  ---

                  FUENTE: 27831\_instrucciones-baremacinin-r03.pdf | PÁGINA: 18
                  \# Instrucciones de baremación

                  Por el contrario, sí se entenderá como duplicidad que la ponencia expuesta en Trabajos presentados en congresos científicos o tecnológicos (apartado 3.6), sea también registrada como Capítulos de Libro (epígrafe 3.3), en tanto que se valora esencialmente el contenido del mérito y no solo el formato de publicación. También se entenderá como duplicidad que un aspirante incluya su contrato de Investigación con una entidad no universitaria ajustado a Contratos (o becas, si procede) en procesos competitivos convocados por entidades públicas o privadas (epígrafe 3.1.), y asimismo lo incluya en el bloque V. Actividad profesional externa a la Universidad. En caso de duplicidad de méritos, es dispositivo de la comisión decidir si procede su consideración favorable (validación) y en qué único lugar se baremará, en aplicación de la discrecionalidad técnica de que goza. El criterio que haya adoptado la comisión para resolver la duplicidad deberá ser aplicado a todos los candidatos. En estos casos en que se produzca la invalidación de un mérito, se deberá rellenar el campo Justificación para mejor conocimiento del aspirante de los motivos que llevan a la comisión a adoptar su decisión.

                  \#\# 4.1.3 Reubicación de méritos

                  La Comisión no validará aquellos méritos que se aleguen en apartados a los que no corresponden según idoneidad conceptual. Dichos méritos no pueden ser reubicados por la comisión. Por ejemplo, si el candidato ha incluido como mérito en el apartado 2.1 (Docencia en asignaturas regladas universitarias) su actividad docente en un Máster propio, ésta no deberá ser validada, pues corresponde al apartado 2.4.2 (Docencia universitaria no oficial).

                  \#\# 4.1.4 Entidades públicas y entidades privadas

                  El baremo no establece distinción entre organismos públicos y privados allí donde sea aplicable, salvo las menciones expresas señaladas en los documentos vigentes. Por tanto, el carácter público o privado de la entidad vinculada al mérito no puede influir en el nivel de afinidad asignado por la comisión.

                  \#\# 4.1.5 Ramas de conocimiento y áreas de conocimiento

                  Para la aplicación de las disposiciones del baremo y las precisiones que en él se hacen, las comisiones no deben considerar intercambiables los conceptos de rama de conocimiento (Artes y Humanidades, Ciencias, Ciencias de la Salud, Ciencias Sociales y Jurídicas e Ingeniería y Arquitectura) y área de conocimiento, en que las primeras se subdividen.
                    [https://investigacion.us.es/secretariado-investigacion/plan-propio/areas-ramas-conocimiento](https://investigacion.us.es/secretariado-investigacion/plan-propio/areas-ramas-conocimiento)

                  \#\# 4.1.6 Documentación acreditativa en idioma distinto al castellano

                  La certificación oficial para la acreditación de las estancias de investigación dentro del Mérito preferente debe aportarse en castellano o en el idioma en que esta haya sido expedida, junto con la traducción al castellano emitida por traductor jurado oficial, según dispone la Instrucción dada al efecto (consulte el apartado 2 de esta guía). Para el resto de la documentación acreditativa, la comisión deberá decidir para qué idiomas no será precisa la aportación de traducción jurada al castellano de los documentos acreditativos de los méritos aportados, criterio que deberá ser el mismo para todos los candidatos. Puede tratarse de idiomas de uso común, como el inglés, o bien de otros que sean consustanciales con la docencia de la plaza y que,

                  Revisión 3 (diciembre, 2024) 17

                  ---

                  FUENTE: BAREMO 2024.pdf | PÁGINA: 18
                  \# Núm. 12/2024, de 11 de octubre

                  \#\# I. DISPOSICIONES GENERALES

                  \#\#\# I.5. Secretario General

                  \#\#\#\# 2.4.4. En calidad de autor o coautor de publicaciones docentes o de carácter pedagógico, libros (libros de texto y manuales), artículos, u otros equivalentes de difusión del conocimiento en el dominio docente:
                  - \textbf{Máximo:} 100 (máximo 5 aportaciones en libros, 5 cursos online de al menos 10 horas de docencia no oficial no computados en apartado 2.4.2, 10 aportaciones en capítulos de libro y 10 artículos de revista no computados en el apartado III. Actividad Investigadora)
                  - \textbf{Ponderación:} 0,20

                  \#\#\#\#\# 2.4.4.1. Tipo de soporte:
                  \vspace{0.4cm}
                  \begin{table}[h!]\centering\renewcommand{\arraystretch}{1.3}
                    \begin{tabular}{|p{0.45\textwidth}|p{0.45\textwidth}|}\hline
                      Tipo de soporte                                                                                         & Puntos por unidad       \\\hline
                      Libros o manuales                                                                                       & 40 puntos por libro     \\\hline
                      Capítulos de libro                                                                                      & 10 puntos por capítulo  \\\hline
                      Artículos de revista                                                                                    & 20 puntos por artículo  \\\hline
                      Juegos didácticos, software, etc.                                                                       & 20 puntos por material  \\\hline
                      Actividades relevantes de difusión de la docencia en medios de comunicación audiovisual o equivalentes. & 20 puntos por actividad \\\hline
                    \end{tabular}\end{table}
                  \vspace{0.4cm}

                  \textbf{Complemento (multiplicativo):}
                  - Factor de corrección de 1,5 para materiales docentes publicados (o aceptados) en los últimos 5 años naturales (incluido el año natural de fecha de convocatoria de la plaza).
                  - \textbf{1,5 x material docente}

                  \textbf{Nota 1:} Se valorarán únicamente los libros, capítulos de libro y artículos cuyas editoriales o revistas se correspondan con los admitidos en los apartados 3.2, 3.3 y 3.4, respectivamente, de este Baremo, con la aclaración de que en este caso las publicaciones deben tener carácter docente y pedagógico. No se valora ningún mérito que no acredite ISBN o depósito legal comprobable mediante las bases de datos oficiales para tales fines. Los tipos de materiales (i.e., cursos online, juegos didácticos, software, difusión, etc.) que no puedan acreditar identificador por su naturaleza deberán acreditar algún otro identificador de publicación o aceptación oficialmente validado, o indicios verificables por la comisión de contratación.

                  \textbf{Nota 2:} El Departamento deberá especificar en el ANEXO I el número de autores a partir del cual se penaliza la puntuación del aspirante. El factor de corrección por número de autores será el siguiente:
                  - (a) de 1 a 3 autores no penaliza = 100 \% de la cuantificación;
                  - (b) más de 3 autores = (100 \%) / (nº autores / i], siendo i ≥ 3, y siendo i el número de autores a partir del cual (no incluyendo i) se penaliza. Asimismo, el Departamento deberá especificar si el orden de los autores es relevante, y el coeficiente corrector aplicable que en ningún caso será inferior a 0,8, y únicamente aplicable a las posiciones intermedias en caso de al menos 3 autores.

                  \#\#\#\# 2.4.5. Curso de formación e innovación docente:
                  - \textbf{Máximo:} 100 (máximo 100 horas de formación e innovación docente)
                  - \textbf{Ponderación:} 0,10

                  \#\#\#\#\# 2.4.5.1. Entidad organizadora:
                  \vspace{0.4cm}
                  \begin{table}[h!]\centering\renewcommand{\arraystretch}{1.3}
                    \begin{tabular}{|p{0.45\textwidth}|p{0.45\textwidth}|}\hline
                      Tipo de curso                                                                & Puntos por 10 h. docencia no oficial    \\\hline
                      Cursos organizados por Universidades o Centros de investigación de prestigio & 10 puntos por 10 h. docencia no oficial \\\hline
                      Cursos organizados por otras entidades                                       & 5 puntos por 10 h. docencia no oficial  \\\hline
                    \end{tabular}\end{table}
                  \vspace{0.4cm}

                  \textbf{Complemento (multiplicativo):}
                  - Factor de corrección de 1,5 para cursos de formación e innovación docente correspondientes a los últimos 5 cursos académicos (incluido el curso académico de fecha de convocatoria de la plaza).
                  - \textbf{1,5 x 10 h. docencia no oficial}
                \end{scriptsize}
        \end{itemize}


  \item \textbf{P25.} Para la figura de Profesor Asociado, ¿qué porcentaje de peso relativo tiene exactamente la experiencia profesional o laboral en la nota final del proceso? \\
        \begin{itemize}
          \item \textbf{Tipo de consulta:} simple
          \item \textbf{Respuesta del Agente:}

                \#\#\# Requisitos Mínimos Excluyentes:
                * Umbral mínimo para profesores asociados: 30 puntos.

                \#\#\# Criterios de Evaluación para Profesor Asociado:

                \vspace{0.4cm}
                \begin{table}[h!]\centering\renewcommand{\arraystretch}{1.3}
                  \begin{tabular}{|p{0.45\textwidth}|p{0.45\textwidth}|}\hline
                    Criterio                                       & Detalles       \\\hline
                    Rendimiento Académico                          & 5\% del total  \\\hline
                    Actividad Docente                              & 10\% del total \\\hline
                    Actividad Investigadora                        & 10\% del total \\\hline
                    Transferencia del Conocimiento                 & 15\% del total \\\hline
                    Actividad Profesional Externa a la Universidad & 55\% del total \\\hline
                    Mérito Preferente                              & 5\% del total  \\\hline
                  \end{tabular}\end{table}
                \vspace{0.4cm}

                \#\#\# Desglose de la Actividad Profesional Externa a la Universidad (Bloque V):

                La experiencia profesional o laboral se evalúa dentro del Bloque V, que tiene un peso relativo del 55\% en la nota final del proceso para la figura de Profesor Asociado. Este bloque se refiere a la actividad profesional externa a la universidad, lo que incluye experiencia laboral en el sector privado, instituciones públicas, o cualquier otra actividad profesional relevante fuera del ámbito académico universitario.

                \#\#\# Porcentaje de Peso Relativo de la Experiencia Profesional o Laboral:

                El porcentaje de peso relativo de la experiencia profesional o laboral en la nota final del proceso para la figura de Profesor Asociado es del 55\%. Esto significa que más de la mitad de la puntuación total se basa en la evaluación de la experiencia profesional o laboral del candidato, lo que destaca la importancia de esta experiencia en el proceso de selección para esta figura contractual.
          \item \textbf{Contexto utilizado:}
                \begin{scriptsize}
                  FUENTE: BAREMO 2024.pdf | PÁGINA: 11
                  \# Núm. 12/2024, de 11 de octubre

                  \#\# I. DISPOSICIONES GENERALES

                  \#\#\# I.5. Secretario General

                  \#\# BAREMO DE CONTRATACIÓN DE PROFESORADO AYUDANTE DOCTOR, PROFESORADO ASOCIADO Y PROFESORADO INTERINO

                  \#\#\# CONSIDERACIONES GENERALES
                  El baremo es aplicable a las figuras contractuales de Profesorado Ayudante Doctor, Profesorado Asociado y Profesorado Interino.

                  Apartado 1. Porcentaje de cada bloque temático según figura contractual (en adelante plazas).
                  En el Baremo se distinguen 6 bloques, sus ponderaciones y sus puntuaciones:
                  1. RENDIMIENTO ACADÉMICO
                  2. ACTIVIDAD DOCENTE
                  3. ACTIVIDAD INVESTIGADORA
                  4. TRANSFERENCIA DEL CONOCIMIENTO
                  5. ACTIVIDAD PROFESIONAL EXTERNA A LA UNIVERSIDAD
                  6. MÉRITO PREFERENTE

                  La calificación del candidato se valora según puntuación acumulada atendiendo a las ponderaciones de la Tabla 1, y dependiendo del tipo de contrato. En caso de empate a puntos entre aspirantes la preferencia entre ellos se resolverá en función de la puntuación alcanzada en el Bloque Temático VI. Mérito Preferente. En caso de no resolverse el empate se calculará la media aritmética de la puntuación (sobre 100) de Actividad Docente (Bloque II) y Actividad Investigadora (Bloque III) en las plazas de Ayudante Doctor; en las plazas de Profesor Asociado la media aritmética se calculará con las puntuaciones (sobre 100) de la Actividad Docente (Bloque II) y Actividad Profesional Externa a la Universidad (Bloque V), y en el caso del Profesorado Interino la media aritmética se calculará con las puntuaciones (sobre 100) del Rendimiento Académico (Bloque I) y Actividad Docente (Bloque II).

                  \#\#\# TABLA 1. Ponderación de cada bloque temático según tipo de contrato.

                  \vspace{0.4cm}
                  \begin{table}[h!]\centering\renewcommand{\arraystretch}{1.3}
                    \begin{tabular}{|p{0.45\textwidth}|p{0.45\textwidth}|p{0.45\textwidth}|p{0.45\textwidth}|p{0.45\textwidth}|}\hline
                      BLOQUE TEMÁTICO                                   & TIPO DE CONTRATO/PLAZA (*) & PSI \% & ASO \% & AYD \% \\\hline
                      I. RENDIMIENTO ACADÉMICO                          &                            & 30     & 5      & 10     \\\hline
                      II. ACTIVIDAD DOCENTE                             &                            & 25     & 10     & 30     \\\hline
                      III. ACTIVIDAD INVESTIGADORA                      &                            & 25     & 10     & 35     \\\hline
                      IV. TRANSFERENCIA DEL CONOCIMIENTO                &                            & 5      & 15     & 5      \\\hline
                      V. ACTIVIDAD PROFESIONAL EXTERNA A LA UNIVERSIDAD &                            & 10     & 55     & 5      \\\hline
                      VI. MÉRITO PREFERENTE                             &                            & 5      & 5      & 15     \\\hline
                      TOTAL                                             &                            & 100    & 100    & 100    \\\hline
                    \end{tabular}\end{table}
                  \vspace{0.4cm}

                  (*) PSI: Profesor Interino; ASO: Profesor Asociado; AYD: Profesor Ayudante Doctor.

                  ---

                  FUENTE: CRITERIOS DE EVALUACIÓN Y REQUISITOS MÍNIMOS.pdf | PÁGINA: 9
                  \# Actividad Docente

                  En todos los casos, se valorará el grado de internacionalización, de innovación, de apertura a la sociedad y de multidisciplinariedad o interdisciplinariedad de las contribuciones presentadas.

                  Las personas solicitantes que acrediten trayectorias de investigación excepcionales, por haber obtenido financiación del Consejo Europeo de Investigación (ERC) en sus programas Consolidator Grant, Advanced Grant o Synergy Grant (o en otras convocatorias competitivas individuales internacionales de prestigio y tasa de concesión comparables a los de estos programas) y hayan desarrollado su carrera principalmente en una universidad extranjera en la que el cómputo y los instrumentos de medición de la calidad de la actividad docente resulten difíciles de trasladar al sistema español, o bien en una institución no universitaria dedicada a la investigación, podrán obtener la acreditación justificando el encargo docente desempeñado, o bien actividades docentes de tutorización y dirección de trabajos de investigación, respectivamente.

                  En caso de adaptación del puesto de trabajo con reducción de las horas lectivas o modalización del cumplimiento de las obligaciones docentes debido a una situación transitoria de pérdida de capacidades o de protección de la maternidad si existe riesgo laboral o una situación reconocida de discapacidad, las exigencias de méritos docentes se modularán de acuerdo con la relevancia de la adaptación del puesto de trabajo acreditada.

                  \#\# PROFESORA O PROFESOR TITULAR UNIVERSIDAD (PTU)

                  \vspace{0.4cm}
                  \begin{table}[h!]\centering\renewcommand{\arraystretch}{1.3}
                    \begin{tabular}{|p{0.45\textwidth}|p{0.45\textwidth}|p{0.45\textwidth}|}\hline
                      Actividad Docente                                                   & Mínimo puntos & Máximo puntos \\\hline
                      2. ACTIVIDAD DOCENTE                                                & 50            & 100           \\\hline
                      2.1. Experiencia docente                                            & 25            & 45            \\\hline
                      2.1.1. Dedicación docente                                           & 25            & 35            \\\hline
                      2.1.2. Pluralidad, interdisciplinariedad y complejidad docente      & 0             & 6             \\\hline
                      2.1.3. Recursos educativos                                          & 0             & 5             \\\hline
                      2.1.4. Actividades de formación a lo largo de la vida               & 0             & 5             \\\hline
                      2.2. Calidad de la actividad docente e innovación                   & 5             & 40            \\\hline
                      2.2.1. Calidad de la actividad docente                              & 5             & 15            \\\hline
                      2.2.2. Proyectos de innovación docente                              & 0             & 20            \\\hline
                      2.2.3. Formación para la mejora docente recibida                    & 0             & 15            \\\hline
                      2.2.4. Formación para la mejora docente impartida                   & 0             & 5             \\\hline
                      2.3. Tutorización docente                                           & 0             & 15            \\\hline
                      2.3.1. Tutorización reglada                                         & 0             & 15            \\\hline
                      2.3.2. Otros tipos de tutorización (no incluida en encargo docente) & 0             & 15            \\\hline
                      2.4. Otros méritos*                                                 & 0             & 10            \\\hline
                    \end{tabular}\end{table}
                  \vspace{0.4cm}

                  ----

                  DIRECCIÓN DE ANECA – 20/03/2024

                  ---

                  FUENTE: BAREMO 2024.pdf | PÁGINA: 3
                  \# Núm. 12/2024, de 11 de octubre

                  \#\# I. DISPOSICIONES GENERALES

                  \#\#\# I.5. Secretario General

                  El artículo 12 de la referida normativa establecía: “El umbral mínimo se fija en 40 puntos para profesores ayudantes doctores, 30 puntos para profesores asociados y 20 puntos para profesores interinos. Ningún candidato cuya puntuación sea inferior a la establecida por dicho umbral podrá ser propuesto para la adjudicación de la plaza”. Idéntica gradación venía reflejada en el apartado 2 del baremo, que establecía el umbral mínimo para poder ser contratado en 20 puntos para profesorado interino, 30 para profesorado asociado y 40 para profesorado ayudante doctor.

                  Con posterioridad, por Acuerdo 7.5 de Consejo de Gobierno de 22 de abril de 2020, se modificó el referido artículo 12 de la normativa en el sentido siguiente: “El umbral mínimo se fija en 40 puntos para profesores ayudantes doctores, 30 puntos para profesores asociados y 20 puntos para profesores interinos, o en caso de acuerdo del Consejo de Departamento tras la publicación de las actas resultantes del proceso de baremación y con validez hasta la resolución del siguiente concurso público de méritos para la contratación de profesorado interino en la Universidad de Sevilla, 10 puntos. Ningún candidato cuya puntuación sea inferior a la establecida por dicho umbral podrá ser propuesto para la adjudicación de la plaza”. Asimismo, el apartado 2 del baremo quedó modificado de la siguiente forma: “El umbral mínimo queda establecido en el art. 12 de esta normativa”.

                  Dado el tiempo transcurrido y las incidencias acaecidas en las listas priorizadas de candidatos resultantes de los diferentes procesos selectivos, procede efectuar una modificación de la normativa relativa al umbral mínimo necesario para la contratación de profesorado sustituto interino.

                  \#\# MOTIVACIÓN DE LA MODIFICACIÓN

                  APROBADA POR ACUERDO 6.6/CG 26-9-24

                  El artículo 78 de la Ley Orgánica 2/2023, de 22 de marzo, del Sistema Universitario regula la figura del Profesor Ayudante Doctor dentro de las diferentes categorías de personal docente e investigador laboral. Dicho artículo dispone que las Universidades podrán contratar bajo esta modalidad a las personas que ostenten el título de Doctor sin necesidad de acreditación.

                  La Disposición transitoria tercera del referido cuerpo legal establece en su apartado 1 que la acreditación vigente de Profesor Ayudante Doctor o de la figura equivalente en la normativa autonómica se considerará como un mérito preferente durante los cuatro años posteriores a la aprobación de esta ley orgánica, a efectos del acceso a la figura de Profesor Ayudante Doctor.

                  El Consejo de Gobierno de la Universidad de Sevilla, por Acuerdo 6.7 en sesión de 27 de febrero de 2019 (posteriormente modificado por Acuerdo 7.5/CG 22-4-20 y Acuerdo 10.6/CG 1-3-24) procedió a la aprobación de la Normativa sobre procedimiento de contratación de profesorado ayudante doctor, profesorado asociado y profesorado interino, junto con el baremo establecido para dichas figuras contractuales.

                  En el artículo 10 de dicha normativa se dispone lo siguiente:

                  1. Según lo dispuesto en el artículo 48.3 de la Ley 6/2001, Orgánica de Universidades, se considerará mérito preferente estar acreditado para participar en los concursos de acceso a cuerpos docentes universitarios.
                  2. El artículo 50 del mismo texto legal dispone que para la contratación de profesores ayudantes doctores se considerará mérito preferente la estancia del candidato en universidades o centros de investigación de reconocido prestigio, españoles o extranjeros, distintos de la universidad que lleve a cabo la contratación.

                  ---

                  FUENTE: CRITERIOS DE EVALUACIÓN Y REQUISITOS MÍNIMOS.pdf | PÁGINA: 10
                  \# CATEDRÁTICA O CATEDRÁTICO DE UNIVERSIDAD (CU)

                  \vspace{0.4cm}
                  \begin{table}[h!]\centering\renewcommand{\arraystretch}{1.3}
                    \begin{tabular}{|p{0.45\textwidth}|p{0.45\textwidth}|}\hline
                      Mínimo & Máximo \\\hline
                      50     & 100    \\\hline
                      puntos & puntos \\\hline
                    \end{tabular}\end{table}
                  \vspace{0.4cm}

                  \#\# 2. ACTIVIDAD DOCENTE

                  \vspace{0.4cm}
                  \begin{table}[h!]\centering\renewcommand{\arraystretch}{1.3}
                    \begin{tabular}{|p{0.45\textwidth}|p{0.45\textwidth}|p{0.45\textwidth}|}\hline
                      Actividad                                                           & Mínimo & Máximo \\\hline
                      2.1. Experiencia docente                                            & 25     & 45     \\\hline
                      2.1.1. Dedicación docente                                           & 25     & 35     \\\hline
                      2.1.2. Pluralidad, interdisciplinariedad y complejidad docente      & 0      & 6      \\\hline
                      2.1.3. Recursos educativos                                          & 0      & 5      \\\hline
                      2.1.4. Actividades de formación a lo largo de la vida               & 0      & 5      \\\hline
                      2.2. Calidad de la actividad docente e innovación                   & 5      & 40     \\\hline
                      2.2.1. Calidad de la actividad docente                              & 5      & 15     \\\hline
                      2.2.2. Proyectos de innovación docente                              & 0      & 20     \\\hline
                      2.2.3. Formación para la mejora docente recibida                    & 0      & 5      \\\hline
                      2.2.4. Formación para la mejora docente impartida                   & 0      & 15     \\\hline
                      2.3. Tutorización docente                                           & 5      & 15     \\\hline
                      2.3.1. Tutorización reglada                                         & 0      & 15     \\\hline
                      2.3.2. Otros tipos de tutorización (no incluida en encargo docente) & 0      & 15     \\\hline
                      2.4. Otros méritos*                                                 & 0      & 10     \\\hline
                    \end{tabular}\end{table}
                  \vspace{0.4cm}

                  * Los puntos obtenidos en este apartado, de carácter opcional, se suman a los de los apartados anteriores, pero en ningún caso podrá superarse el máximo (100 puntos).

                  \#\# 2.1. EXPERIENCIA DOCENTE

                  \#\#\# 2.1.1. Dedicación docente

                  En el caso de PTU deberá acreditarse experiencia docente durante un mínimo de 4 años de dedicación a tiempo completo, o su equivalente a tiempo parcial, y, además:

                  - En el caso de figuras de contratados predoctorales y postdoctorales, incluyendo el periodo de orientación postdoctoral, un mínimo de 12 ECTS o 120 horas impartidas. El disfrute de un contrato postdoctoral se considerará dedicación a tiempo completo. El disfrute de una ayuda predoctoral que tenga asignada capacidad docente tutelada y se cubra totalmente se considerará en todos los casos dedicación a tiempo parcial.

                  - Para el resto de las figuras de profesorado, un mínimo de 24 ECTS o 240 horas impartidas.

                  - En caso de que la experiencia docente se haya acumulado combinando distintas modalidades (esto es: predoctoral/postdoctoral y resto de figuras de profesorado), el mínimo de ECTS u horas se calculará proporcionalmente al periodo de tiempo de duración de cada una de las figuras contractuales disfrutadas.

                  En el caso de CU deberá acreditarse una experiencia docente mínima de 10 años de dedicación a tiempo completo, o su equivalente a tiempo parcial, y un mínimo de 140 ECTS.

                  DIRECCIÓN DE ANECA – 20/03/2024

                  ---

                  FUENTE: 27831\_instrucciones-baremacinin-r03.pdf | PÁGINA: 30
                  \# Instrucciones de baremación

                  - Los méritos englobados en este epígrafe del baremo no admiten coeficiente de afinidad, por lo que la Comisión NO puede invalidarlos alegando falta de afinidad con la plaza. Si el mérito ha sido consignado en el apartado correcto y se encuentra acreditado documentalmente, la Comisión debe validarlo.

                  \#\# 4.5 Transferencia del conocimiento (bloque IV del baremo)

                  \#\#\# 4.5.1 Contratos de transferencia (epígrafe 4.1 del baremo)

                  No se recogen indicaciones relativas a este apartado en la versión actual de este documento.

                  \#\#\# 4.5.2 Obra civil o de edificación, o exposiciones (epígrafe 4.2 del baremo)

                  No se recogen indicaciones relativas a este apartado en la versión actual de este documento.

                  \#\#\# 4.5.3 Patentes (epígrafe 4.3 del baremo)

                  Para que este mérito pueda validarse, debe acreditarse su concesión, no solo su solicitud.

                  \#\#\# 4.5.4 Spin-off (epígrafe 4.4 del baremo)

                  No se recogen indicaciones relativas a este apartado en la versión actual de este documento.

                  \#\# 4.6 Actividad profesional externa (bloque V del baremo)

                  - Este apartado se corresponde con Actividad profesional externa a la actividad académica dentro de la Universidad. Por tanto, se admiten los contratos en el ámbito universitario como Personal Técnico, de Gestión y de Administración y Servicios (aunque no como PDI).
                  - Se contempla dentro de este apartado la docencia en Institutos de Enseñanza Secundaria.

                  \#\# 4.7 Mérito preferente (bloque VI del baremo)

                  \#\#\# 4.7.1 Acreditación

                  - En los concursos correspondientes a plazas de Profesorado Asociado y Profesor Sustituto, las únicas acreditaciones que pueden validarse son las correspondientes a categorías de cuerpos docentes universitarios, es decir:
                  - Profesor Titular de Universidad
                  - Catedrático de Universidad
                  Las demás acreditaciones (Profesor Ayudante Doctor, Profesor Contratado Doctor) no deberán ser validadas.
                  - En lo que respecta a las plazas de Profesorado Ayudante Doctor, se admiten como mérito preferente, además de las acreditaciones anteriormente citadas, las correspondientes a Profesor Ayudante Doctor y a Profesor Contratado Doctor/Profesor Permanente Laboral. Esta distinción con respecto a las figuras citadas en el punto anterior tiene su origen en los cambios introducidos.

                  Revisión 3 (diciembre, 2024) 29
                \end{scriptsize}
        \end{itemize}


  \item \textbf{P26.} En caso de que en la baremación dos solicitantes tengan exactamente la misma puntuación global para ser Ayudante Doctor, ¿cuál es el primer criterio de desempate de la Universidad? \\
        \begin{itemize}
          \item \textbf{Tipo de consulta:} simple
          \item \textbf{Respuesta del Agente:}

                \#\#\# Requisitos Mínimos Excluyentes

                *   No se mencionan requisitos mínimos excluyentes específicos en el texto proporcionado para el proceso de evaluación y baremación de méritos para la contratación de Ayudante Doctor en la Universidad de Sevilla.

                \#\#\# Criterios de Evaluación

                \vspace{0.4cm}
                \begin{table}[h!]\centering\renewcommand{\arraystretch}{1.3}
                  \begin{tabular}{|p{0.45\textwidth}|p{0.45\textwidth}|}\hline
                    Criterio                                       & Detalles                    \\\hline
                    Rendimiento Académico                          & 10\% de la puntuación total \\\hline
                    Actividad Docente                              & 30\% de la puntuación total \\\hline
                    Actividad Investigadora                        & 35\% de la puntuación total \\\hline
                    Transferencia del Conocimiento                 & 5\% de la puntuación total  \\\hline
                    Actividad Profesional Externa a la Universidad & 5\% de la puntuación total  \\\hline
                    Mérito Preferente                              & 15\% de la puntuación total \\\hline
                  \end{tabular}\end{table}
                \vspace{0.4cm}

                \#\#\# Desempate en Caso de Empate

                En caso de empate a puntos entre aspirantes, la preferencia entre ellos se resolverá en función de la puntuación alcanzada en el Bloque Temático VI, Mérito Preferente. Si el empate persiste, se calculará la media aritmética de la puntuación (sobre 100) de Actividad Docente (Bloque II) y Actividad Investigadora (Bloque III) para las plazas de Ayudante Doctor.

                \#\#\# Cálculo de Puntos

                La calificación del candidato se valora según puntuación acumulada atendiendo a las ponderaciones de la Tabla 1, y dependiendo del tipo de contrato. La puntuación final y última del candidato se redondeará al primer decimal inmediato superior en caso que el segundo decimal fuese igual o superior a 5 y, quedando en el correspondiente en caso que fuera igual o inferior a cuatro.

                \#\#\# Umbral Mínimo

                El umbral mínimo para poder ser contratado como Ayudante Doctor es de 40 puntos, según el artículo 12 de la normativa. Ningún candidato cuya puntuación sea inferior a la establecida por dicho umbral podrá ser propuesto para la adjudicación de la plaza.

                \#\#\# Mérito Preferente

                El mérito preferente para Ayudante Doctor incluye:

                *   Habilitación, acreditación o ser funcionario en el cuerpo docente de Titular de Universidad o Catedrático de Universidad (hasta 10 puntos).
                *   Estancias de investigación en universidades u otros centros de investigación externos a la Universidad (hasta 10 puntos, acumulables al punto anterior y hasta 15 puntos en total).

                \#\#\# Coeficiente de Afinidad

                La comisión de contratación puede aplicar un coeficiente de afinidad para reducir la puntuación de méritos que no correspondan al perfil de la plaza en concurso. El coeficiente de afinidad puede ser: 0,0 (muy bajo), 0,25 (bajo), 0,50 (medio), 0,75 (alto), 1,00 (muy alto).
          \item \textbf{Contexto utilizado:}
                \begin{scriptsize}
                  FUENTE: BAREMO 2024.pdf | PÁGINA: 11
                  \# Núm. 12/2024, de 11 de octubre

                  \#\# I. DISPOSICIONES GENERALES

                  \#\#\# I.5. Secretario General

                  \#\# BAREMO DE CONTRATACIÓN DE PROFESORADO AYUDANTE DOCTOR, PROFESORADO ASOCIADO Y PROFESORADO INTERINO

                  \#\#\# CONSIDERACIONES GENERALES
                  El baremo es aplicable a las figuras contractuales de Profesorado Ayudante Doctor, Profesorado Asociado y Profesorado Interino.

                  Apartado 1. Porcentaje de cada bloque temático según figura contractual (en adelante plazas).
                  En el Baremo se distinguen 6 bloques, sus ponderaciones y sus puntuaciones:
                  1. RENDIMIENTO ACADÉMICO
                  2. ACTIVIDAD DOCENTE
                  3. ACTIVIDAD INVESTIGADORA
                  4. TRANSFERENCIA DEL CONOCIMIENTO
                  5. ACTIVIDAD PROFESIONAL EXTERNA A LA UNIVERSIDAD
                  6. MÉRITO PREFERENTE

                  La calificación del candidato se valora según puntuación acumulada atendiendo a las ponderaciones de la Tabla 1, y dependiendo del tipo de contrato. En caso de empate a puntos entre aspirantes la preferencia entre ellos se resolverá en función de la puntuación alcanzada en el Bloque Temático VI. Mérito Preferente. En caso de no resolverse el empate se calculará la media aritmética de la puntuación (sobre 100) de Actividad Docente (Bloque II) y Actividad Investigadora (Bloque III) en las plazas de Ayudante Doctor; en las plazas de Profesor Asociado la media aritmética se calculará con las puntuaciones (sobre 100) de la Actividad Docente (Bloque II) y Actividad Profesional Externa a la Universidad (Bloque V), y en el caso del Profesorado Interino la media aritmética se calculará con las puntuaciones (sobre 100) del Rendimiento Académico (Bloque I) y Actividad Docente (Bloque II).

                  \#\#\# TABLA 1. Ponderación de cada bloque temático según tipo de contrato.

                  \vspace{0.4cm}
                  \begin{table}[h!]\centering\renewcommand{\arraystretch}{1.3}
                    \begin{tabular}{|p{0.45\textwidth}|p{0.45\textwidth}|p{0.45\textwidth}|p{0.45\textwidth}|p{0.45\textwidth}|}\hline
                      BLOQUE TEMÁTICO                                   & TIPO DE CONTRATO/PLAZA (*) & PSI \% & ASO \% & AYD \% \\\hline
                      I. RENDIMIENTO ACADÉMICO                          &                            & 30     & 5      & 10     \\\hline
                      II. ACTIVIDAD DOCENTE                             &                            & 25     & 10     & 30     \\\hline
                      III. ACTIVIDAD INVESTIGADORA                      &                            & 25     & 10     & 35     \\\hline
                      IV. TRANSFERENCIA DEL CONOCIMIENTO                &                            & 5      & 15     & 5      \\\hline
                      V. ACTIVIDAD PROFESIONAL EXTERNA A LA UNIVERSIDAD &                            & 10     & 55     & 5      \\\hline
                      VI. MÉRITO PREFERENTE                             &                            & 5      & 5      & 15     \\\hline
                      TOTAL                                             &                            & 100    & 100    & 100    \\\hline
                    \end{tabular}\end{table}
                  \vspace{0.4cm}

                  (*) PSI: Profesor Interino; ASO: Profesor Asociado; AYD: Profesor Ayudante Doctor.

                  ---

                  FUENTE: BAREMO 2024.pdf | PÁGINA: 4
                  \# Núm. 12/2024, de 11 de octubre

                  \#\# I. DISPOSICIONES GENERALES

                  \#\#\# I.5. Secretario General

                  3. La regulación concreta para la valoración del mérito preferente será la que se establezca en el Baremo o en las normas que, en su caso, se dicten en desarrollo de esta normativa.

                  Por su parte, en el punto VI. Mérito Preferente del Baremo aprobado junto con la normativa de contratación figura:

                  “\textbf{Profesorado Ayudante Doctor:} El aspirante logra 15 puntos atendiendo al siguiente reparto: 10 puntos por su habilitación, acreditación o ser funcionario en el cuerpo docente de Titular de Universidad o Catedrático de Universidad, o hasta 10 puntos (acumulables al punto 1 y hasta 15 puntos) por sus estancias de investigación en universidades u otros centros de investigación externos a la Universidad, hasta 18 meses de estancia, y conforme a las instrucciones dictadas por la Universidad de Sevilla”.

                  Visto lo anterior, procede una actualización tanto del artículo 10.1 de la Normativa de contratación de Profesorado Ayudante Doctor, de Profesorado Asociado y de Profesorado Sustituto como del párrafo 3 del apartado VI del baremo a aplicar.

                  \#\# TÍTULO I

                  \#\#\# Disposiciones generales

                  \textbf{Artículo 1. Objeto de régimen jurídico aplicable.}

                  1. El objeto de la presente normativa es la regulación en la Universidad de Sevilla de los procedimientos de convocatoria y provisión mediante concurso público de las siguientes plazas de personal docente e investigador: Profesorado Ayudante Doctor, Profesorado Asociado y Profesorado Interino.
                  2. Los citados procedimientos se desarrollarán en el marco de las previsiones normativas establecidas en la Ley Orgánica 6/2001, de 21 de diciembre, el Texto Refundido de la Ley Andaluza de Universidades aprobado por Decreto Legislativo 1/2013, el Estatuto de la Universidad de Sevilla, el I Convenio Colectivo del Personal Docente e Investigador en régimen laboral de las Universidades Públicas andaluzas, la presente normativa y, en lo no previsto en estas normas, por el régimen de profesorado que sea de aplicación y por las disposiciones que regulan con carácter general el ingreso en la Administración Pública.
                  3. En materia de procedimiento, y en todo lo no recogido específicamente en esta normativa, se estará a lo dispuesto en la Ley 39/2015, de 1 de octubre, del Procedimiento Administrativo Común de las Administraciones Públicas y en la Ley 40/2015, de 1 de octubre, de Régimen Jurídico del Sector Público.

                  \textbf{Artículo 2. Principios.}

                  1. La selección del personal a que se refiere la presente normativa se efectuará mediante concurso público con respeto a los principios constitucionales de igualdad, mérito, capacidad y publicidad.

                  \#\# TÍTULO II

                  \#\#\# Tramitación previa y actuación de las comisiones de selección

                  \textbf{Artículo 3. Composición, sorteo y nombramiento de las Comisiones de Contratación.}

                  1. La composición y vigencia de las Comisiones de Contratación de las plazas de Profesor Ayudante Doctor y Profesor Asociado serán las establecidas en el artículo 90 del Estatuto de la Universidad de Sevilla.

                  ---

                  FUENTE: 27831\_instrucciones-baremacinin-r03.pdf | PÁGINA: 5
                  \# Bareprof

                  \#\# Instrucciones de baremación

                  - Instrucción de 17-10-19 sobre los requisitos para el reconocimiento de las estancias de investigación dentro del mérito preferente en los concursos de profesorado ayudante doctor, asociado y sustituto interino (BOUS 14/2019, de 21 de noviembre).
                    [https://www.us.es/bous-numeros/numero-14-21-de-noviembre-de-2019](https://www.us.es/bous-numeros/numero-14-21-de-noviembre-de-2019)

                  - Resolución rectoral, de 24 marzo de 2021, por la que se emite una interpretación sobre la determinación del umbral mínimo aplicable en los procedimientos selectivos para plazas de Profesores Ayudantes Doctores, Profesores Asociados y Profesores Interinos establecido en el artículo 12 de la normativa sobre procedimiento de contratación de Profesorado Ayudante Doctor, Profesorado Asociado y Profesorado Interino, junto con el baremo establecido para dichas figuras contractuales de la Universidad de Sevilla (BOUS 1/2021, de 7 de abril).
                    [https://www.us.es/bous-numeros/numero-1-7-de-abril-de-2021](https://www.us.es/bous-numeros/numero-1-7-de-abril-de-2021)

                  - Resolución rectoral, de 20 de abril de 2022, por la que se determinan las áreas afines en el ámbito de los concursos de profesorado de la Universidad de Sevilla (BOUS 3/2022, de 19 de mayo).
                    [https://www.us.es/bous-numeros/numero-3-19-de-mayo-de-2022](https://www.us.es/bous-numeros/numero-3-19-de-mayo-de-2022)

                  Asimismo, los miembros de las comisiones de contratación deberán consultar los siguientes documentos disponibles en la plataforma Bareprof:
                  - Guía de ayuda al candidato.
                  - Anexo I correspondiente al departamento y área de conocimiento de la plaza a baremar.

                  \#\# 3 Descripción del proceso de baremación

                  \#\#\# 3.1 Composición de las comisiones de contratación

                  La composición de las comisiones de contratación correspondientes a concursos de Profesorado Ayudante Doctor se determina según lo dispuesto en el acuerdo 9.7/CG 23-7-74, por el que se aprueban las normas provisionales para adaptar a la Ley Orgánica 2/2023, de 22 de marzo, del Sistema Universitario, el modo de elección de las comisiones se selección de plazas de Profesorado Ayudante Doctor. Tales comisiones pueden consultarse en el tablón electrónico oficial de la Universidad de Sevilla ([https://www.us.es/tablon-virtual](https://www.us.es/tablon-virtual)).

                  En cuanto a los concursos de plazas de Profesorado Asociado y Profesorado Sustituto, la elección de las comisiones se realiza conforme a lo establecido en la Normativa sobre Procedimiento de Contratación de Profesorado Ayudante Doctor, Profesorado Asociado y Profesorado Interino (véase el apartado 2 de este documento). La composición de estas comisiones puede consultarse en la Secretaría Virtual de la Universidad de Sevilla, accediendo a la siguiente dirección:
                  [https://sevius4.us.es/?comisionescontratacion](https://sevius4.us.es/?comisionescontratacion)

                  Para realizar la consulta, debe seleccionarse el año del sorteo, el departamento, el área de conocimiento y la comisión (según el tipo de plaza a baremar).

                  Revisión 3 (diciembre, 2024)

                  ---

                  FUENTE: BAREMO 2024.pdf | PÁGINA: 12
                  \# Núm. 12/2024, de 11 de octubre

                  \#\# I. DISPOSICIONES GENERALES

                  \#\#\# I.5. Secretario General

                  \#\#\#\# Apartado 2. Puntuaciones de los méritos y apartados.

                  Los candidatos introducirán sus méritos (CV) a través del editor de CVN de la Fundación Española para la Ciencia y la Tecnología (FECYT), atendiendo a las instrucciones publicadas por el Vicerrectorado con competencias en profesorado. Con antelación a la convocatoria de plazas o convocatoria de profesorado interino se publicará en el medio electrónico oficial de la Universidad de Sevilla el documento de instrucciones para cumplimentar el CVN Editor de FECYT. En el caso de que FECYT modifique la estructura de su CVN se dictarán asimismo las instrucciones precisas por el Vicerrectorado con competencias en profesorado para la adaptación que corresponda.

                  Las Comisiones podrán excluir de la baremación un mérito que se alegue en dos apartados distintos del baremo, no valorando aquellos que se aleguen en apartados a que no correspondan. Aplicado el Baremo si alguno o algunos de los concursantes a una misma plaza obtuvieran puntuación en exceso sobre los máximos establecidos en alguno de los subapartados, se asignará la puntuación máxima admitida a quien la tuviese más alta, escalándose proporcionalmente la valoración de los demás concursantes. No se escalarán los valores obtenidos por los concursantes en los diferentes subapartados donde expresamente se indique que no procede, o se hagan otras precisiones.

                  La puntuación final y última del candidato se redondeará al primer decimal inmediato superior en caso que el segundo decimal fuese igual o superior a 5 y, quedando en el correspondiente en caso que fuera igual o inferior a cuatro. El umbral mínimo queda establecido en el artículo 12 de esta normativa.

                  En caso de que la Universidad de Sevilla considere necesario emitir interpretaciones al baremo destinadas a su obligado cumplimiento por los actores implicados en el proceso, se harán públicas bajo el tipo Instrucciones del Baremo.

                  \#\#\#\# Apartado 3. Coeficiente de afinidad.

                  Cuando la comisión de contratación considere que los méritos y su contenido alegados por el candidato no corresponden en todo o en parte al perfil de la plaza en concurso contenida en la resolución de convocatoria, la valoración de tales méritos se verá reducida en función de su menor afinidad; la reducción deberá ser debidamente motivada en las actas a elaborar por la comisión de contratación, y se aplicará en aquellos subapartados donde proceda según el baremo aplicable.

                  A tal efecto, la comisión de contratación fijará el coeficiente de afinidad en uno de los valores siguientes: 0,0 (muy bajo), 0,25 (bajo), 0,50 (medio), 0,75 (alto), 1,00 (muy alto). El coeficiente fijado por la comisión de contratación se multiplicará por la puntuación base dada al mérito, estableciéndose, en primer lugar, el coeficiente de afinidad y, en segundo lugar, la valoración del mérito según indicios de calidad. Para su aplicación deberán atenderse necesariamente las notas publicadas en cada apartado del baremo que lo especifique.

                  Asimismo, se establece que la comisión de contratación podrá solicitar al aspirante en cualquier momento del proceso de baremación documentación acreditativa y explicativa de sus méritos, estando el aspirante obligado a su remisión.

                  \#\#\#\# Apartado 4. Mérito preferente.

                  El candidato que alegue estar en posesión del mérito preferente al que se refiere el artículo 48.3 de la Ley Orgánica de Universidades (estar acreditado a los cuerpos docentes universitarios, i.e., Titular de Universidad o Catedrático de Universidad) o al que se refiere el Art. 50 de la Ley Orgánica de

                  ---

                  FUENTE: BAREMO 2024.pdf | PÁGINA: 3
                  \# Núm. 12/2024, de 11 de octubre

                  \#\# I. DISPOSICIONES GENERALES

                  \#\#\# I.5. Secretario General

                  El artículo 12 de la referida normativa establecía: “El umbral mínimo se fija en 40 puntos para profesores ayudantes doctores, 30 puntos para profesores asociados y 20 puntos para profesores interinos. Ningún candidato cuya puntuación sea inferior a la establecida por dicho umbral podrá ser propuesto para la adjudicación de la plaza”. Idéntica gradación venía reflejada en el apartado 2 del baremo, que establecía el umbral mínimo para poder ser contratado en 20 puntos para profesorado interino, 30 para profesorado asociado y 40 para profesorado ayudante doctor.

                  Con posterioridad, por Acuerdo 7.5 de Consejo de Gobierno de 22 de abril de 2020, se modificó el referido artículo 12 de la normativa en el sentido siguiente: “El umbral mínimo se fija en 40 puntos para profesores ayudantes doctores, 30 puntos para profesores asociados y 20 puntos para profesores interinos, o en caso de acuerdo del Consejo de Departamento tras la publicación de las actas resultantes del proceso de baremación y con validez hasta la resolución del siguiente concurso público de méritos para la contratación de profesorado interino en la Universidad de Sevilla, 10 puntos. Ningún candidato cuya puntuación sea inferior a la establecida por dicho umbral podrá ser propuesto para la adjudicación de la plaza”. Asimismo, el apartado 2 del baremo quedó modificado de la siguiente forma: “El umbral mínimo queda establecido en el art. 12 de esta normativa”.

                  Dado el tiempo transcurrido y las incidencias acaecidas en las listas priorizadas de candidatos resultantes de los diferentes procesos selectivos, procede efectuar una modificación de la normativa relativa al umbral mínimo necesario para la contratación de profesorado sustituto interino.

                  \#\# MOTIVACIÓN DE LA MODIFICACIÓN

                  APROBADA POR ACUERDO 6.6/CG 26-9-24

                  El artículo 78 de la Ley Orgánica 2/2023, de 22 de marzo, del Sistema Universitario regula la figura del Profesor Ayudante Doctor dentro de las diferentes categorías de personal docente e investigador laboral. Dicho artículo dispone que las Universidades podrán contratar bajo esta modalidad a las personas que ostenten el título de Doctor sin necesidad de acreditación.

                  La Disposición transitoria tercera del referido cuerpo legal establece en su apartado 1 que la acreditación vigente de Profesor Ayudante Doctor o de la figura equivalente en la normativa autonómica se considerará como un mérito preferente durante los cuatro años posteriores a la aprobación de esta ley orgánica, a efectos del acceso a la figura de Profesor Ayudante Doctor.

                  El Consejo de Gobierno de la Universidad de Sevilla, por Acuerdo 6.7 en sesión de 27 de febrero de 2019 (posteriormente modificado por Acuerdo 7.5/CG 22-4-20 y Acuerdo 10.6/CG 1-3-24) procedió a la aprobación de la Normativa sobre procedimiento de contratación de profesorado ayudante doctor, profesorado asociado y profesorado interino, junto con el baremo establecido para dichas figuras contractuales.

                  En el artículo 10 de dicha normativa se dispone lo siguiente:

                  1. Según lo dispuesto en el artículo 48.3 de la Ley 6/2001, Orgánica de Universidades, se considerará mérito preferente estar acreditado para participar en los concursos de acceso a cuerpos docentes universitarios.
                  2. El artículo 50 del mismo texto legal dispone que para la contratación de profesores ayudantes doctores se considerará mérito preferente la estancia del candidato en universidades o centros de investigación de reconocido prestigio, españoles o extranjeros, distintos de la universidad que lleve a cabo la contratación.
                \end{scriptsize}
        \end{itemize}


\end{itemize}

\subsection{Categoría 4: Calendario (27-34)}
\textit{Evalúa la precisión del agente al consultar plazos, fechas límite, festivos y periodos lectivos.}
\begin{itemize}
  \item \textbf{P27.} ¿Qué día empiezan oficialmente las clases del curso 2025-2026 para todos aquellos estudiantes de Grado y Máster Universitario? \\
        \begin{itemize}
          \item \textbf{Tipo de consulta:} simple
          \item \textbf{Respuesta del Agente:}

                El inicio de las clases para los estudiantes de Grado y Máster Universitario en el curso 2025-2026 es el siguiente:

                - Para los estudiantes de Grado, las clases comienzan oficialmente el \textbf{lunes 8 de septiembre de 2025}.
                - Para los estudiantes de Máster Universitario, las clases comienzan oficialmente el \textbf{lunes 6 de octubre de 2025}.

                Es importante destacar que, aunque la mayoría de los estudiantes de Máster Universitario comienzan el \textbf{6 de octubre de 2025}, existen excepciones para aquellas titulaciones que hayan cubierto sus plazas a \textbf{22 de septiembre de 2025}, las cuales podrán adaptar el inicio del periodo docente a su situación específica, siempre respetando la fecha límite de entrega de actas y requiriendo comunicación previa al Vicerrectorado de Estudiantes.
          \item \textbf{Contexto utilizado:}
                \begin{scriptsize}
                  FUENTE: calendario-academico-2025-2026-cg.pdf | PÁGINA: 3
                  \# CALENDARIO ACADÉMICO
                  \#\# CURSO 2025/2026

                  \#\#\# 4. SOLICITUDES DE RECONOCIMIENTO/TRANSFERENCIA DE CRÉDITOS EN LOS ESTUDIOS DE GRADO Y MÁSTER UNIVERSITARIO

                  La presentación de solicitudes de reconocimiento de créditos ECTS para asignaturas regladas, se realizará desde la fecha de la formalización de automatrícula/matrícula por el estudiante hasta el 22 de septiembre de 2025. Aquellos estudiantes que realicen su matrícula con posterioridad al 5 de septiembre dispondrán de un plazo de 10 días hábiles, contados a partir del día siguiente a la formalización de su matrícula, para la presentación de la solicitud.

                  En los casos de cambio de Universidad y/o estudios se realizará según lo dispuesto en la Resolución Rectoral por la que se regula la admisión a los títulos de grado de la Universidad de Sevilla de los estudiantes que han iniciado anteriormente estudios universitarios.

                  Las solicitudes de reconocimiento de créditos por participación en actividades universitarias culturales, deportivas, de representación estudiantil, solidarias y de cooperación, así como las solicitudes de reconocimiento basadas en experiencia profesional, el reconocimiento por la acreditación de otros niveles de idioma, podrán presentarse en cualquier momento del curso académico. Los referidos reconocimientos se podrán imputar al curso anterior, 2024/2025, siempre que la solicitud se presente antes del 8 de septiembre de 2025.

                  Las solicitudes de reconocimiento de prácticas extracurriculares por la asignatura “Prácticas Externas” podrán presentarse en cualquier momento del curso académico. No obstante, se podrán imputar al curso anterior, 2024/2025, siempre que la solicitud se presente hasta el 22 de septiembre de 2025, inclusive.

                  El plazo para solicitar transferencia de créditos coincidirá con el general para solicitar reconocimiento de créditos.

                  \#\#\# 5. CALENDARIO DE ACTIVIDADES DOCENTES Y EXÁMENES

                  \#\#\#\# 5.1. INICIO DE ACTIVIDADES DOCENTES

                  - El periodo docente de los estudiantes de grado dará comienzo el lunes 8 de septiembre de 2025.
                  - El periodo docente de los estudiantes de máster universitario dará comienzo el lunes 6 de octubre de 2025.

                  Aquellas titulaciones de máster universitario que hayan cubierto sus plazas a 22 (L) de septiembre de 2025 podrán adaptar el inicio del periodo docente a su situación específica.

                  Por otro lado, cuando no se hayan cubierto las plazas en la Fase 2, se podrá adaptar el inicio del periodo docente a su situación específica.

                  En ambos casos se habrá de respetar siempre la fecha límite de entrega de actas y se requerirá comunicación previa al Vicerrectorado de Estudiantes.

                  \#\#\#\# 5.2. CONSIDERACIONES GENERALES

                  La actividad docente tendrá una duración de 40 semanas, 30 de clases y 10 de exámenes. Las 30 semanas de clases se dividirán en 15 para cada cuatrimestre.

                  Acuerdo 4.2 CG/20-3-25 por el que se aprueba el Calendario Académico para el curso 2025-2026  Página 3 de 13

                  ---

                  FUENTE: calendario-academico-2025-2026-cg.pdf | PÁGINA: 4
                  \# CALENDARIO ACADÉMICO
                  \#\# CURSO 2025/2026

                  En los títulos de grado, con carácter general, se establecen tres semanas para la realización de exámenes de la primera convocatoria de las asignaturas cuatrimestrales impartidas en el primer cuatrimestre. No obstante, por necesidades organizativas, los Centros podrán establecer solamente dos semanas, previa comunicación del Vicerrectorado de Estudiantes.

                  Para el caso de los títulos de máster universitario, el período correspondiente a los exámenes de la primera convocatoria de las asignaturas del primer cuatrimestre será de una semana y no coincidirá con los períodos de realización de exámenes de los títulos de grado.

                  Por necesidades organizativas los Centros podrán ampliar los periodos de exámenes siempre que se respeten los plazos límites de entrega de actas, previa comunicación al Vicerrectorado de Estudiantes.

                  En los títulos de máster universitario, para la presentación de los trabajos fin de máster el periodo de evaluación de la convocatoria prevista en el mes de julio del curso 2024/2025, quedará ampliado hasta el 20 (J) de noviembre de 2025 y la entrega de actas tendrá como fecha límite el 28 (V) de noviembre de 2025.

                  La convocatoria de julio del trabajo fin de máster correspondiente al curso 2025/2026 se extenderá hasta el fin del plazo que se establezca en la 3ª convocatoria del curso siguiente.

                  \#\#\# 5.3. PERIODOS DE CLASES

                  \vspace{0.4cm}
                  \begin{table}[h!]\centering\renewcommand{\arraystretch}{1.3}
                    \begin{tabular}{|p{0.45\textwidth}|p{0.45\textwidth}|p{0.45\textwidth}|}\hline
                      PERÍODO DE CLASES & PRIMER CUATRIMESTRE                                            & SEGUNDO CUATRIMESTRE                                                 \\\hline
                      TÍTULOS DE GRADO  & Del 8 (L) de septiembre de 2025 al 19 (V) de diciembre de 2025 & Caso general: Del 26 (L) de enero 2026 al 22 (V) de mayo de 2026     \\\hline
                                        &                                                                & Caso excepcional: Del 19 (L) de enero 2026 al 15 (V) de mayo de 2026 \\\hline
                      TÍTULOS DE MÁSTER & Del 6 (L) de octubre de 2025 al 3 (M) de febrero de 2026       & Del 12 (J) de febrero de 2026 al 10 (X) de junio de 2026             \\\hline
                    \end{tabular}\end{table}
                  \vspace{0.4cm}

                  - Caso general: 3 semanas de realización de exámenes de las asignaturas cuatrimestrales impartidas en el primer cuatrimestre.
                  - Caso excepcional: 2 semanas de realización de exámenes de las asignaturas cuatrimestrales impartidas en el primer cuatrimestre.

                  Acuerdo 4.2 CG/20-3-25 por el que se aprueba el Calendario Académico para el curso 2025-2026  Página 4 de 13

                  ---

                  FUENTE: calendario-academico-2025-2026-cg.pdf | PÁGINA: 1
                  \# CALENDARIO ACADÉMICO
                  \#\# CURSO 2025/2026

                  \#\#\# GRADO, MÁSTER Y DOCTORADO

                  El presente Calendario se establece en aplicación de lo dispuesto en el artículo 31 del Reglamento General de Actividades Docentes, aprobado por Acuerdo del Claustro de 5/2/2009 y resulta de aplicación en el curso 2025/2026 a los estudios oficiales de grado y de máster universitario.

                  Asimismo, se incorpora el Calendario correspondiente a los estudios oficiales de doctorado (anexo II) y a los estudios impartidos en el Instituto de Idiomas (anexo III).

                  El Rector podrá dictar cuantas Resoluciones sean necesarias para el cumplimiento, desarrollo e interpretación del presente Acuerdo, singularmente para el caso de titulaciones conjuntas, másteres universitarios y programas de doctorado que así lo precisen.

                  \#\#\# ESTUDIOS DE GRADO Y MÁSTER

                  \#\#\#\# 1. CALENDARIO DE MATRÍCULA

                  \#\#\#\#\# 1.1. MATRÍCULA ORDINARIA PARA ESTUDIANTES QUE CONTINÚAN EN LOS MISMOS ESTUDIOS

                  Con independencia del procedimiento de automatrícula/matrícula seguido, el plazo estará comprendido entre el 9 (X) y el 31 (J) de julio de 2025 y entre el 1 (L) y el 5 (V) de septiembre de 2025.

                  El acceso escalonado a la automatrícula tendrá lugar del 9 (X) al 25 (V) de julio de 2025.

                  \#\#\#\#\# 1.2. MATRÍCULA ORDINARIA PARA ESTUDIANTES DE NUEVO INGRESO

                  Los plazos de matrícula en los títulos de grado y máster universitario, para estudiantes procedentes de preinscripción, vendrán determinados por el Calendario aprobado por la Comisión del Distrito Único Universitario Andaluz (anexo I).

                  Los estudiantes admitidos en la Fase de Extranjeros del procedimiento de ingreso en estudios de grado y en la Fase 1 del procedimiento de admisión en estudios de máster universitario, realizarán su prematrícula en los plazos determinados por el calendario aprobado por la Comisión del Distrito Único Universitario Andaluz y formalizarán su matrícula/automatrícula en los siguientes plazos:

                  - Del 3 (J) al 9 (X) de julio de 2025 en estudios de grado.
                  - Del 8 (M) al 10 (J) de julio de 2025 en estudios de máster universitario.

                  \#\#\#\#\# 1.3. AMPLIACIÓN DE MATRÍCULA

                  - Entre el 12 (L) y el 30 (V) de enero de 2026.
                  - El plazo general de ampliación se prorrogará para los trabajos fin de estudios hasta el 18 (X) de febrero de 2026 para las titulaciones de grado y hasta el 6 (V) de marzo de 2026 para las titulaciones de máster universitario salvo para aquellos Centros que no lo autoricen. Estos últimos deberán comunicarlo al Vicerrectorado de Estudiantes antes del 28 (V) de noviembre de 2025.

                  Acuerdo 4.2 CG/20-3-25 por el que se aprueba el Calendario Académico para el curso 2025-2026  Página 1 de 13

                  ---

                  FUENTE: calendario-academico-2025-2026-cg.pdf | PÁGINA: 2
                  \# CALENDARIO ACADÉMICO
                  \#\# CURSO 2025/2026

                  \#\#\# 1.4. SIMULTANEIDAD DE ESTUDIOS

                  El plazo general de matrícula en los estudios de grado y máster universitario coincidirá con el determinado por la Comisión de Distrito Único Universitario de Andalucía para los estudiantes de nuevo ingreso con plaza adjudicada al final del proceso de preinscripción.

                  La matrícula de aquellos estudiantes que accedan en régimen de simultaneidad, de acuerdo con lo previsto en la Resolución Rectoral de 16/06/2010, se llevará a cabo tras su admisión por el Centro.

                  \#\#\# 1.5. SOLICITUDES PARA CURSAR REQUISITOS FORMATIVOS COMPLEMENTARIOS PARA HOMOLOGACIÓN DE TÍTULOS EXTRANJEROS

                  El periodo de presentación de solicitudes para cursar requisitos formativos complementarios exigidos para la homologación de títulos extranjeros de educación superior a títulos españoles estará comprendido entre el 9 (X) y el 31 (J) de julio de 2025 y entre el 1 (L) y el 5 (V) de septiembre de 2025.

                  Igualmente, se podrán presentar solicitudes en el periodo de ampliación de matrícula, entre 12 (L) y el 30 (V) de enero de 2026.

                  \#\#\# 1.6. INSCRIPCIÓN Y MATRÍCULA DE ESTUDIANTES ENTRANTES DE MOVILIDAD INTERNACIONAL

                  Fechas de selección de asignaturas, prematrícula y matrícula anticipada:
                  - Primer semestre o curso completo: 15 (J) de mayo al 31 (J) de julio de 2025.
                  - Segundo semestre: 3 (L) de noviembre al 19 (V) de diciembre de 2025.

                  Fecha de asignación de grupos (matriculación efectiva):
                  - Primer semestre o curso completo: 1 (L) al 19 (V) de septiembre de 2025.
                  - Segundo semestre: 12 (L) de enero al 13 (V) de febrero de 2026.

                  \#\# 2. SOLICITUDES DE CAMBIO DE UNIVERSIDAD, CENTRO Y/O ESTUDIOS. CONVALIDACIÓN PARCIAL DE ESTUDIOS EXTRANJEROS

                  El período de presentación de solicitudes para cambios de Universidad, Centro y/o estudios estará comprendido entre el 1 (J) de mayo y el 31 (J) de julio de 2025. A efectos de cómputos de plazos en la resolución de estas solicitudes, se declara inhábil el mes de agosto.

                  En el caso de traslado para continuar estudios en las titulaciones conjuntas del Campus de Excelencia Internacional de Andalucía Tech el plazo de presentación de solicitudes en el Centro de destino será del 2 (L) al 30 (L) de junio de 2025.

                  \#\# 3. SOLICITUDES DE REANUDACIÓN DE ESTUDIOS ESTUDIANTES QUE TRASLADARON EXPEDIENTE

                  Los estudiantes que iniciaron estudios de grado o máster universitario y hubieran trasladado su expediente, podrán solicitar la reanudación de sus estudios en el centro entre el 1 y el 15 de septiembre de 2025, siempre que tengan un mínimo de 6 créditos superados/reconocidos y no hayan agotado las normas de permanencia en los estudios que pretenden reanudar. De ser estimada su solicitud, dispondrán de un plazo de diez días hábiles para formalizar su matrícula y gestionar el traslado de expediente, contados a partir de la fecha de notificación de la resolución de admisión.

                  Acuerdo 4.2 CG/20-3-25 por el que se aprueba el Calendario Académico para el curso 2025-2026  Página 2 de 13

                  ---

                  FUENTE: R.R. NORMAS DE MATRICULA GRADO Y MASTER UNIVERSITARIO 2025-2026.pdf | PÁGINA: 51
                  \# UNIVERSIDAD DE SEVILLA

                  \#\# B) AUTOMATRÍCULA ESTUDIANTES QUE CONTINÚAN SUS ESTUDIOS EN LA MISMA TITULACIÓN

                  \#\#\# 1. Plazos

                  Todos los estudiantes de la Universidad de Sevilla que continúan estudios de grado o de máster universitario, con las excepciones que se indican en la página 34, realizarán su matrícula para el curso 2025-2026 por el procedimiento de automatrícula.

                  \#\#\#\# 1.1. Estudiantes de grado

                  Para ordenar el acceso a la automatrícula y garantizar el correcto desarrollo del proceso, en el curso 2025-2026, se establece un plazo personalizado de matrícula para cada estudiante. Este plazo permite al estudiante matricularse desde el día y hora que le sean asignados hasta el 5 de septiembre (excluido agosto).

                  Cada estudiante podrá consultar en la Secretaría Virtual (https://sevius4.us.es/?exp\_punt\_auto) la fecha y hora a partir de las cuales puede acceder a la automatrícula (nunca antes).

                  Para más información sobre el sistema de acceso escalonado, debe consultarse en https://cat.us.es/acceso-escalonado.

                  \#\#\#\# PLAZO PERSONALIZADO DE ACCESO AUTOMATRÍCULA ESTUDIANTES DE GRADO

                  \vspace{0.4cm}
                  \begin{table}[h!]\centering\renewcommand{\arraystretch}{1.3}
                    \begin{tabular}{|p{0.45\textwidth}|p{0.45\textwidth}|p{0.45\textwidth}|p{0.45\textwidth}|p{0.45\textwidth}|p{0.45\textwidth}|p{0.45\textwidth}|p{0.45\textwidth}|}\hline
                      DESDE  & HASTA  & FECHA    & HORA  & DESDE  & HASTA  & FECHA    & HORA  \\\hline
                      14,401 & 30,000 & 9/07/25  & 09:00 & 10,080 & 10,155 & 17/07/25 & 15:00 \\\hline
                      13,536 & 14,400 & 9/07/25  & 12:00 & 9,997  & 10,079 & 18/07/25 & 09:00 \\\hline
                      12,987 & 13,535 & 9/07/25  & 15:00 & 9,915  & 9,996  & 18/07/25 & 12:00 \\\hline
                      12,584 & 12,986 & 10/07/25 & 09:00 & 9,826  & 9,914  & 18/07/25 & 15:00 \\\hline
                      12,263 & 12,583 & 10/07/25 & 12:00 & 9,734  & 9,825  & 21/07/25 & 09:00 \\\hline
                      12,009 & 12,262 & 10/07/25 & 15:00 & 9,636  & 9,733  & 21/07/25 & 12:00 \\\hline
                      11,781 & 12,008 & 11/07/25 & 09:00 & 9,528  & 9,635  & 21/07/25 & 15:00 \\\hline
                      11,580 & 11,780 & 11/07/25 & 12:00 & 9,405  & 9,527  & 22/07/25 & 09:00 \\\hline
                      11,402 & 11,579 & 11/07/25 & 15:00 & 9,253  & 9,404  & 22/07/25 & 12:00 \\\hline
                      11,239 & 11,401 & 14/07/25 & 09:00 & 9,082  & 9,252  & 22/07/25 & 15:00 \\\hline
                      11,099 & 11,238 & 14/07/25 & 12:00 & 8,868  & 9,081  & 23/07/25 & 09:00 \\\hline
                      10,962 & 11,098 & 14/07/25 & 15:00 & 8,619  & 8,867  & 23/07/25 & 12:00 \\\hline
                      10,830 & 10,961 & 15/07/25 & 09:00 & 8,300  & 8,618  & 23/07/25 & 15:00 \\\hline
                      10,713 & 10,829 & 15/07/25 & 12:00 & 7,951  & 8,299  & 24/07/25 & 09:00 \\\hline
                      10,601 & 10,712 & 15/07/25 & 15:00 & 7,490  & 7,950  & 24/07/25 & 12:00 \\\hline
                      10,493 & 10,600 & 16/07/25 & 09:00 & 6,892  & 7,489  & 24/07/25 & 15:00 \\\hline
                      10,396 & 10,492 & 16/07/25 & 12:00 & 6,056  & 6,891  & 25/07/25 & 09:00 \\\hline
                      10,306 & 10,395 & 16/07/25 & 15:00 & 4,546  & 6,055  & 25/07/25 & 12:00 \\\hline
                      10,229 & 10,305 & 17/07/25 & 09:00 & 0,000  & 4,545  & 25/07/25 & 15:00 \\\hline
                      10,156 & 10,228 & 17/07/25 & 12:00 &        &        &          &       \\\hline
                    \end{tabular}\end{table}
                  \vspace{0.4cm}

                  \#\#\#\# 1.2. Estudiantes de máster universitario

                  Podrán realizar su automatrícula del 9 al 31 de julio y del 1 al 5 de septiembre de 2025.

                  Código Seguro De Verificación: `eqOjhZcmL0p6VJlZiTH/Yw==`
                  Fecha: 27/06/2025
                  Firmado Por: MIGUEL ANGEL CASTRO ARROYO
                  Url De Verificación: [https://pfirma.us.es/verifirma/code/eqOjhZcmL0p6VJlZiTH\%2FYw\%3D\%3D](https://pfirma.us.es/verifirma/code/eqOjhZcmL0p6VJlZiTH\%2FYw\%3D\%3D)
                  Página: 47/50
                \end{scriptsize}
        \end{itemize}


  \item \textbf{P28.} Según el calendario académico 2025-2026, ¿en qué fechas exactas tendrán lugar los exámenes de la convocatoria extraordinaria de septiembre para los recién matriculados en primer curso? \\
        \begin{itemize}
          \item \textbf{Tipo de consulta:} simple
          \item \textbf{Respuesta del Agente:}

                Según el calendario académico 2025-2026, no se proporcionan fechas específicas para los exámenes de la convocatoria extraordinaria de septiembre para los recién matriculados en primer curso. Sin embargo, se menciona que la convocatoria de septiembre se extenderá hasta el fin del plazo que se establezca en la 3ª convocatoria del curso siguiente, pero esto se refiere a la presentación de los trabajos fin de máster.

                En cuanto a los exámenes de las asignaturas, se establecen diferentes periodos de exámenes para las convocatorias, pero no se especifica una convocatoria extraordinaria de septiembre para los recién matriculados en primer curso.

                Es importante destacar que el calendario académico proporciona fechas para las convocatorias de exámenes finales, pero no especifica fechas para una convocatoria extraordinaria de septiembre para los recién matriculados en primer curso. Por lo tanto, no se pueden proporcionar fechas exactas para esta convocatoria.

                Si necesitas más información, te recomiendo consultar directamente con la universidad o revisar el calendario académico para obtener más detalles.
          \item \textbf{Contexto utilizado:}
                \begin{scriptsize}
                  FUENTE: calendario-academico-2025-2026-cg.pdf | PÁGINA: 5
                  \# CALENDARIO ACADÉMICO
                  \#\# CURSO 2025/2026

                  \#\#\# 5.4. PERÍODOS DE REALIZACIÓN DE EXÁMENES FINALES

                  \vspace{0.4cm}
                  \begin{table}[h!]\centering\renewcommand{\arraystretch}{1.3}
                    \begin{tabular}{|p{0.45\textwidth}|p{0.45\textwidth}|p{0.45\textwidth}|p{0.45\textwidth}|p{0.45\textwidth}|}\hline
                      PERÍODOS DE EXÁMENES & CONVOCATORIA    & ASIGNATURAS                                        & TÍTULOS DE GRADO                                         & TÍTULOS DE MÁSTER                     \\\hline
                      TERCERA              & Todas las       & Del 15 (X) al 31 (V) de octubre de 2025 (1)        & Caso general                                             & Caso excepcional (2)                  \\\hline
                                           & Asignaturas     & Del 25 (L) de mayo al 13 (S) de junio de 2026      & Del 18 (L) de mayo al 13 (S) de junio de 2026            & Del 11 (J) al 23 (M) de junio de 2026 \\\hline
                                           & anuales         & Del 29 (L) de junio al 11 (S) de julio de 2026 (4) & Del 4 (S) al 15 (X) de julio de 2026                     & Del 11 (J) al 23 (M) de junio de 2026 \\\hline
                      PRIMERA              & Asignaturas     & Del 8 (J) al 24 (S) de enero de 2026 (3)           & Del 8 (J) al 17 (S) de enero de 2026 (3)                 & Del 4 (X) al 11 (X) de febrero 2026   \\\hline
                                           & 1º cuatrimestre & Del 25 (L) de mayo al 13 (S) de junio de 2026      & Del 18 (L) de mayo al 13 (S) de junio de 2026            & Del 11 (J) al 23 (M) de junio de 2026 \\\hline
                                           & Asignaturas     & Del 29 (L) de junio al 11 (S) de julio de 2026 (4) & Del 4 (S) al 15 (X) de julio de 2026                     & Del 11 (J) al 23 (M) de junio de 2026 \\\hline
                                           & anuales         &                                                    &                                                          &                                       \\\hline
                      SEGUNDA              & Asignaturas     & Del 25 (L) de mayo al 13 (S) de junio de 2026 (5)  & o del 29 (L) de junio al 11 (S) de julio de 2026 (4) (5) & del 4 (S) al 15 (X) de julio de 2026  \\\hline
                                           & 1º cuatrimestre & Del 29 (L) al 11 (S) de julio de 2026 (4)          &                                                          & Del 4 (S) al 15 (X) de julio de 2026  \\\hline
                                           & 2º cuatrimestre &                                                    &                                                          &                                       \\\hline
                    \end{tabular}\end{table}
                  \vspace{0.4cm}

                  (1) Los Centros que lo deseen podrán adelantar las fechas de celebración de la tercera convocatoria, previa comunicación al Vicerrectorado de Estudiantes.
                  (2) En los títulos de grado los Centros podrán establecer dos semanas para la realización de exámenes de la primera convocatoria de las asignaturas impartidas en el primer cuatrimestre, previa comunicación del Vicerrectorado de Estudiantes.
                  (3) Excepcionalmente y teniendo en cuenta la organización de cada centro, previa comunicación al Vicerrectorado de Estudiantes, se podrán realizar exámenes el 7 (X) de enero de 2026 en horario de tarde.
                  (4) Los Centros podrán, de acuerdo con su organización, utilizar el periodo del 13 (L) al 16 (J) de julio de 2026 para la realización de exámenes de la 2ª convocatoria previa comunicación al Vicerrectorado de Estudiantes.
                  (5) Los Centros podrán, de acuerdo con su organización, utilizar el periodo del 15 (L) al 27 (S) de junio de 2026 para la realización de exámenes de la 2ª convocatoria de asignaturas del primer cuatrimestre, previa comunicación al Vicerrectorado de Estudiantes.

                  Acuerdo 4.2 CG/20-3-25 por el que se aprueba el Calendario Académico para el curso 2025-2026  Página 5 de 13

                  ---

                  FUENTE: calendario-academico-2025-2026-cg.pdf | PÁGINA: 8
                  \# CALENDARIO ACADÉMICO
                  \#\# CURSO 2025/2026

                  \#\#\# ANEXO I: CALENDARIO DE MATRÍCULA PARA ESTUDIANTES DE NUEVO INGRESO POR PREINSCRIPCIÓN

                  1. ESTUDIANTES DE NUEVO INGRESO POR PREINSCRIPCIÓN EN LOS ESTUDIOS DE GRADO

                  Plazos establecidos por la Comisión de Distrito Único Universitario de Andalucía.

                  \vspace{0.4cm}
                  \begin{table}[h!]\centering\renewcommand{\arraystretch}{1.3}
                    \begin{tabular}{|p{0.45\textwidth}|p{0.45\textwidth}|p{0.45\textwidth}|}\hline
                      FASE                & PLAZO DE MATRÍCULA & PERIODO                  \\\hline
                      FASE EXTRANJEROS    & 1º                 & 4 al 10 de abril         \\\hline
                      (Pre-matrícula)     & 2º                 & 30 de abril al 6 de mayo \\\hline
                                          & MATRÍCULA          &                          \\\hline
                                          & 1º                 & 3 al 9 de julio          \\\hline
                                          & 2º                 & 17 al 21 de julio        \\\hline
                                          & 1ª                 & 25 al 28 de julio        \\\hline
                                          & 2ª                 & 4 al 8 de septiembre     \\\hline
                      FASE ORDINARIA      & RESULTAS           &                          \\\hline
                                          & 3ª                 & 12 al 15 de septiembre   \\\hline
                                          & 4ª                 & 19 al 22 de septiembre   \\\hline
                                          & 5ª                 & 26 al 29 de septiembre   \\\hline
                                          & 6ª                 & 10 al 14 de octubre      \\\hline
                                          & MATRÍCULA          &                          \\\hline
                                          & 1º                 & 12 al 15 de septiembre   \\\hline
                                          & 2º                 & 19 al 22 de septiembre   \\\hline
                                          & RESULTAS           &                          \\\hline
                                          & 1º                 & 26 al 29 de septiembre   \\\hline
                                          & 2º                 & 10 al 14 de octubre      \\\hline
                                          & MATRÍCULA          &                          \\\hline
                                          & 1º                 & 25 al 28 de julio        \\\hline
                      FASE EXTRAORDINARIA &                    &                          \\\hline
                                          & 2º                 & 4 al 8 de septiembre     \\\hline
                                          & 1º                 & 12 al 15 de septiembre   \\\hline
                                          & RESULTAS           &                          \\\hline
                                          & 2º                 & 19 al 22 de septiembre   \\\hline
                                          & 3º                 & 26 al 29 de septiembre   \\\hline
                                          & 4º                 & 10 al 14 de octubre      \\\hline
                                          & TITULADOS          &                          \\\hline
                                          & 1º                 & 26 al 29 de septiembre   \\\hline
                                          & 2º                 & 10 al 14 de octubre      \\\hline
                      PLAZAS VACANTES     & MATRÍCULA FINAL    & 10 al 14 de octubre      \\\hline
                    \end{tabular}\end{table}
                  \vspace{0.4cm}

                  Acuerdo 4.2 CG/20-3-25 por el que se aprueba el Calendario Académico para el curso 2025-2026  Página 8 de 13

                  ---

                  FUENTE: calendario-academico-2025-2026-cg.pdf | PÁGINA: 13
                  \# CALENDARIO ACADÉMICO
                  \#\# CURSO 2025/2026

                  \#\#\# 1.7. MATRÍCULA OTROS CURSOS (1) (2)

                  - Curso anual (francés para egresados): del 8 (L) de septiembre al 3 (V) de octubre de 2025.
                  - Cursos 1º cuatrimestre (iniciación al francés, preparación para el DELF B2, preparación para el Cambridge B2, preparación para el Goethe C1): del 8 (L) de septiembre al 3 (V) de octubre de 2025.
                  - Cursos 2º cuatrimestre (iniciación al francés, preparación para el DELF B2, preparación para el Cambridge C1, preparación para el Goethe C1): del 29 (J) de enero al 4 (X) de febrero de 2026.
                  - Cursos Preparación para el Linguaskill:
                  - Curso 1: del 8 (L) de septiembre al 3 (V) de octubre de 2025.
                  - Curso 2: del 14 (M) al 28 (M) de octubre de 2025.
                  - Curso 3: del 8 (J) al 13 (M) de enero de 2026.
                  - Curso 4: del 29 (J) de enero al 4 (X) de febrero de 2026.

                  \#\#\# 2. SOLICITUD DE CONVALIDACIONES

                  Del 20 (V) de junio de 2025 al 20 (V) de marzo de 2026. (3)

                  \#\#\# 3. CONVOCATORIAS DE EXÁMENES FINALES Y FECHAS CIERRE DE ACTAS

                  Véase la Normativa de Matrícula del Instituto de Idiomas en cuanto a condiciones y limitaciones de concurrencia. Las actas con las calificaciones de los exámenes finales deberán estar firmadas en las fechas que se determinan en el siguiente cuadro.

                  \#\#\#\# CONVOCATORIAS DE EXÁMENES FINALES Y FECHAS DE CIERRE DE ACTAS

                  \vspace{0.4cm}
                  \begin{table}[h!]\centering\renewcommand{\arraystretch}{1.3}
                    \begin{tabular}{|p{0.45\textwidth}|p{0.45\textwidth}|p{0.45\textwidth}|p{0.45\textwidth}|}\hline
                      CONVOCATORIA & ASIGNATURAS                               & PERIODO DE EXÁMENES & ENTREGA DE ACTAS                      \\\hline
                      TERCERA      & Todas las asignaturas                     & Noviembre           & Hasta el 28 (V) de noviembre de 2025  \\\hline
                                   & (matriculados en cursos intensivos        &                     &                                       \\\hline
                                   & septiembre/octubre y repetidores)         &                     &                                       \\\hline
                                   & Asignaturas 1º cuatrimestre               & Enero/febrero       & Hasta el 13 (V) de febrero de 2026    \\\hline
                      PRIMERA      & Asignaturas anuales y del 2º cuatrimestre & Mayo/julio          & Hasta el 31 (V) de julio de 2026      \\\hline
                                   & Asignaturas 1º cuatrimestre               & Mayo/julio          & Hasta el 31 (V) de julio de 2026      \\\hline
                      SEGUNDA      & Asignaturas anuales y del 2º cuatrimestre & Septiembre          & Hasta el 30 (X) de septiembre de 2026 \\\hline
                    \end{tabular}\end{table}
                  \vspace{0.4cm}

                  ----
                  (3) La resolución de convalidación será válida para los periodos de matrícula posteriores a la fecha de la convalidación.

                  Acuerdo 4.2 CG/20-3-25 por el que se aprueba el Calendario Académico para el curso 2025-2026      Página 13 de 13

                  ---

                  FUENTE: calendario-academico-2025-2026-cg.pdf | PÁGINA: 4
                  \# CALENDARIO ACADÉMICO
                  \#\# CURSO 2025/2026

                  En los títulos de grado, con carácter general, se establecen tres semanas para la realización de exámenes de la primera convocatoria de las asignaturas cuatrimestrales impartidas en el primer cuatrimestre. No obstante, por necesidades organizativas, los Centros podrán establecer solamente dos semanas, previa comunicación del Vicerrectorado de Estudiantes.

                  Para el caso de los títulos de máster universitario, el período correspondiente a los exámenes de la primera convocatoria de las asignaturas del primer cuatrimestre será de una semana y no coincidirá con los períodos de realización de exámenes de los títulos de grado.

                  Por necesidades organizativas los Centros podrán ampliar los periodos de exámenes siempre que se respeten los plazos límites de entrega de actas, previa comunicación al Vicerrectorado de Estudiantes.

                  En los títulos de máster universitario, para la presentación de los trabajos fin de máster el periodo de evaluación de la convocatoria prevista en el mes de julio del curso 2024/2025, quedará ampliado hasta el 20 (J) de noviembre de 2025 y la entrega de actas tendrá como fecha límite el 28 (V) de noviembre de 2025.

                  La convocatoria de julio del trabajo fin de máster correspondiente al curso 2025/2026 se extenderá hasta el fin del plazo que se establezca en la 3ª convocatoria del curso siguiente.

                  \#\#\# 5.3. PERIODOS DE CLASES

                  \vspace{0.4cm}
                  \begin{table}[h!]\centering\renewcommand{\arraystretch}{1.3}
                    \begin{tabular}{|p{0.45\textwidth}|p{0.45\textwidth}|p{0.45\textwidth}|}\hline
                      PERÍODO DE CLASES & PRIMER CUATRIMESTRE                                            & SEGUNDO CUATRIMESTRE                                                 \\\hline
                      TÍTULOS DE GRADO  & Del 8 (L) de septiembre de 2025 al 19 (V) de diciembre de 2025 & Caso general: Del 26 (L) de enero 2026 al 22 (V) de mayo de 2026     \\\hline
                                        &                                                                & Caso excepcional: Del 19 (L) de enero 2026 al 15 (V) de mayo de 2026 \\\hline
                      TÍTULOS DE MÁSTER & Del 6 (L) de octubre de 2025 al 3 (M) de febrero de 2026       & Del 12 (J) de febrero de 2026 al 10 (X) de junio de 2026             \\\hline
                    \end{tabular}\end{table}
                  \vspace{0.4cm}

                  - Caso general: 3 semanas de realización de exámenes de las asignaturas cuatrimestrales impartidas en el primer cuatrimestre.
                  - Caso excepcional: 2 semanas de realización de exámenes de las asignaturas cuatrimestrales impartidas en el primer cuatrimestre.

                  Acuerdo 4.2 CG/20-3-25 por el que se aprueba el Calendario Académico para el curso 2025-2026  Página 4 de 13

                  ---

                  FUENTE: calendario-academico-2025-2026-cg.pdf | PÁGINA: 9
                  \# CALENDARIO ACADÉMICO
                  \#\# CURSO 2025/2026

                  \#\#\# 2. ESTUDIANTES DE NUEVO INGRESO POR PREINSCRIPCIÓN EN LOS ESTUDIOS DE MÁSTER

                  Plazos establecidos por la Comisión de Distrito Único Universitario de Andalucía.

                  \#\#\#\# MATRÍCULA ESTUDIOS DE MÁSTER UNIVERSITARIO

                  \vspace{0.4cm}
                  \begin{table}[h!]\centering\renewcommand{\arraystretch}{1.3}
                    \begin{tabular}{|p{0.45\textwidth}|p{0.45\textwidth}|p{0.45\textwidth}|}\hline
                      FASE          & PLAZO DE MATRÍCULA   & PERIODO                \\\hline
                      1º            & 20 al 24 de febrero  &                        \\\hline
                      FASE 1        & 2º                   & 13 al 17 de marzo      \\\hline
                      (Extranjeros) & 3º                   & 3 al 7 de abril        \\\hline
                                    & 1º                   & 8 al 10 de julio       \\\hline
                                    & (y matrícula Fase 1) &                        \\\hline
                                    & 2º                   & 22 al 24 de julio      \\\hline
                                    & 3ª                   & 9 al 11 de septiembre  \\\hline
                      FASE 2        & 1ª                   & 24 al 26 de septiembre \\\hline
                                    & 2ª                   & 9 al 13 de octubre     \\\hline
                                    & 3ª                   & 21 al 23 de octubre    \\\hline
                                    & 4ª                   & 28 al 30 de octubre    \\\hline
                                    & 1º                   & 24 al 26 de septiembre \\\hline
                      MATRÍCULA     & 2º                   & 9 al 13 de octubre     \\\hline
                      FASE 3        & 1º                   & 21 al 23 de octubre    \\\hline
                                    & 2º                   & 28 al 30 de octubre    \\\hline
                    \end{tabular}\end{table}
                  \vspace{0.4cm}

                  Acuerdo 4.2 CG/20-3-25 por el que se aprueba el Calendario Académico para el curso 2025-2026  Página 9 de 13
                \end{scriptsize}
        \end{itemize}


  \item \textbf{P29.} En el calendario de grado para 2025/2026, ¿cuál es la fecha límite exacta para el cierre de actas de la convocatoria ordinaria del primer cuatrimestre? \\
        \begin{itemize}
          \item \textbf{Tipo de consulta:} simple
          \item \textbf{Respuesta del Agente:}

                La fecha límite exacta para el cierre de actas de la convocatoria ordinaria del primer cuatrimestre en el calendario de grado para 2025/2026 es el \textbf{13 de febrero de 2026}. Esta fecha se establece en el calendario académico como el plazo límite para la entrega de actas de la primera convocatoria de asignaturas impartidas en el primer cuatrimestre para títulos de grado.
          \item \textbf{Contexto utilizado:}
                \begin{scriptsize}
                  FUENTE: calendario-academico-2025-2026-cg.pdf | PÁGINA: 6
                  \# CALENDARIO ACADÉMICO
                  \#\# CURSO 2025/2026

                  \#\#\# 6. FECHAS DE FIRMA Y ENTREGA DE ACTAS

                  Las actas con las calificaciones de los exámenes finales deberán estar firmadas y entregadas en las Secretarías de los Centros en las fechas límites establecidas en el siguiente calendario:

                  \vspace{0.4cm}
                  \begin{table}[h!]\centering\renewcommand{\arraystretch}{1.3}
                    \begin{tabular}{|p{0.45\textwidth}|p{0.45\textwidth}|p{0.45\textwidth}|}\hline
                      CONVOCATORIA                                                                           & TÍTULOS DE GRADO                     & TÍTULOS DE MÁSTER                    \\\hline
                      3ª Convocatoria                                                                        & Hasta el 28 (V) de noviembre de 2025 &                                      \\\hline
                      1ª Convocatoria de asignaturas impartidas en el 1º cuatrimestre                        & Hasta el 13 (V) de febrero de 2026   & Hasta el 6 (V) de marzo de 2026      \\\hline
                      1ª Convocatoria de asignaturas impartidas en el 2º cuatrimestre y anuales (mayo/junio) & Hasta el 2 (J) de julio de 2026      & Hasta el 10 (V) de julio de 2026     \\\hline
                      2ª Convocatoria de asignaturas impartidas en el 1º cuatrimestre                        & Hasta el 31 (V) de julio de 2026     & Hasta el 31 (V) de julio de 2026     \\\hline
                      2ª Convocatoria de asignaturas impartidas en el 2º cuatrimestre y anuales (julio)      & Hasta el 31 (V) de julio de 2026 (1) & Hasta el 31 (V) de julio de 2026 (2) \\\hline
                    \end{tabular}\end{table}
                  \vspace{0.4cm}

                  Dentro de los plazos generales y de los límites máximos establecidos en este calendario, deberá tenerse en cuenta el plazo establecido en el art. 62.2 del Reglamento General de Actividades Docentes de, al menos, dos días lectivos entre los seis posteriores a la fecha de publicación de las calificaciones para resolver las revisiones, contados desde la publicación de las listas provisionales de calificaciones, con la salvedad que indica el art. 27.2 de la Normativa Reguladora de la Evaluación y Calificación de las Asignaturas.

                  (1) Excepcionalmente, previa comunicación al Vicerrectorado de Estudiantes, los Centros podrán solicitar que el límite de entrega de actas de la segunda convocatoria de TFG y PE se pueda demorar hasta el 22 (M) de septiembre de 2026, siempre que las actas sean individuales. Las prácticas extracurriculares que se inicien antes del 31 (L) de agosto de 2026 podrán ser reconocidas para el curso 2025/2026, si éstas finalizan antes del 22 (M) de septiembre de 2026.

                  (2) El plazo máximo de entrega de actas para la convocatoria prevista en el mes de julio correspondiente a los trabajos fin de máster coincidirá con el fin del plazo establecido para la tercera convocatoria del curso siguiente, como se indica en el apartado 5.2.

                  Acuerdo 4.2 CG/20-3-25 por el que se aprueba el Calendario Académico para el curso 2025-2026  Página 6 de 13

                  ---

                  FUENTE: calendario-academico-2025-2026-cg.pdf | PÁGINA: 4
                  \# CALENDARIO ACADÉMICO
                  \#\# CURSO 2025/2026

                  En los títulos de grado, con carácter general, se establecen tres semanas para la realización de exámenes de la primera convocatoria de las asignaturas cuatrimestrales impartidas en el primer cuatrimestre. No obstante, por necesidades organizativas, los Centros podrán establecer solamente dos semanas, previa comunicación del Vicerrectorado de Estudiantes.

                  Para el caso de los títulos de máster universitario, el período correspondiente a los exámenes de la primera convocatoria de las asignaturas del primer cuatrimestre será de una semana y no coincidirá con los períodos de realización de exámenes de los títulos de grado.

                  Por necesidades organizativas los Centros podrán ampliar los periodos de exámenes siempre que se respeten los plazos límites de entrega de actas, previa comunicación al Vicerrectorado de Estudiantes.

                  En los títulos de máster universitario, para la presentación de los trabajos fin de máster el periodo de evaluación de la convocatoria prevista en el mes de julio del curso 2024/2025, quedará ampliado hasta el 20 (J) de noviembre de 2025 y la entrega de actas tendrá como fecha límite el 28 (V) de noviembre de 2025.

                  La convocatoria de julio del trabajo fin de máster correspondiente al curso 2025/2026 se extenderá hasta el fin del plazo que se establezca en la 3ª convocatoria del curso siguiente.

                  \#\#\# 5.3. PERIODOS DE CLASES

                  \vspace{0.4cm}
                  \begin{table}[h!]\centering\renewcommand{\arraystretch}{1.3}
                    \begin{tabular}{|p{0.45\textwidth}|p{0.45\textwidth}|p{0.45\textwidth}|}\hline
                      PERÍODO DE CLASES & PRIMER CUATRIMESTRE                                            & SEGUNDO CUATRIMESTRE                                                 \\\hline
                      TÍTULOS DE GRADO  & Del 8 (L) de septiembre de 2025 al 19 (V) de diciembre de 2025 & Caso general: Del 26 (L) de enero 2026 al 22 (V) de mayo de 2026     \\\hline
                                        &                                                                & Caso excepcional: Del 19 (L) de enero 2026 al 15 (V) de mayo de 2026 \\\hline
                      TÍTULOS DE MÁSTER & Del 6 (L) de octubre de 2025 al 3 (M) de febrero de 2026       & Del 12 (J) de febrero de 2026 al 10 (X) de junio de 2026             \\\hline
                    \end{tabular}\end{table}
                  \vspace{0.4cm}

                  - Caso general: 3 semanas de realización de exámenes de las asignaturas cuatrimestrales impartidas en el primer cuatrimestre.
                  - Caso excepcional: 2 semanas de realización de exámenes de las asignaturas cuatrimestrales impartidas en el primer cuatrimestre.

                  Acuerdo 4.2 CG/20-3-25 por el que se aprueba el Calendario Académico para el curso 2025-2026  Página 4 de 13

                  ---

                  FUENTE: calendario-academico-2025-2026-cg.pdf | PÁGINA: 13
                  \# CALENDARIO ACADÉMICO
                  \#\# CURSO 2025/2026

                  \#\#\# 1.7. MATRÍCULA OTROS CURSOS (1) (2)

                  - Curso anual (francés para egresados): del 8 (L) de septiembre al 3 (V) de octubre de 2025.
                  - Cursos 1º cuatrimestre (iniciación al francés, preparación para el DELF B2, preparación para el Cambridge B2, preparación para el Goethe C1): del 8 (L) de septiembre al 3 (V) de octubre de 2025.
                  - Cursos 2º cuatrimestre (iniciación al francés, preparación para el DELF B2, preparación para el Cambridge C1, preparación para el Goethe C1): del 29 (J) de enero al 4 (X) de febrero de 2026.
                  - Cursos Preparación para el Linguaskill:
                  - Curso 1: del 8 (L) de septiembre al 3 (V) de octubre de 2025.
                  - Curso 2: del 14 (M) al 28 (M) de octubre de 2025.
                  - Curso 3: del 8 (J) al 13 (M) de enero de 2026.
                  - Curso 4: del 29 (J) de enero al 4 (X) de febrero de 2026.

                  \#\#\# 2. SOLICITUD DE CONVALIDACIONES

                  Del 20 (V) de junio de 2025 al 20 (V) de marzo de 2026. (3)

                  \#\#\# 3. CONVOCATORIAS DE EXÁMENES FINALES Y FECHAS CIERRE DE ACTAS

                  Véase la Normativa de Matrícula del Instituto de Idiomas en cuanto a condiciones y limitaciones de concurrencia. Las actas con las calificaciones de los exámenes finales deberán estar firmadas en las fechas que se determinan en el siguiente cuadro.

                  \#\#\#\# CONVOCATORIAS DE EXÁMENES FINALES Y FECHAS DE CIERRE DE ACTAS

                  \vspace{0.4cm}
                  \begin{table}[h!]\centering\renewcommand{\arraystretch}{1.3}
                    \begin{tabular}{|p{0.45\textwidth}|p{0.45\textwidth}|p{0.45\textwidth}|p{0.45\textwidth}|}\hline
                      CONVOCATORIA & ASIGNATURAS                               & PERIODO DE EXÁMENES & ENTREGA DE ACTAS                      \\\hline
                      TERCERA      & Todas las asignaturas                     & Noviembre           & Hasta el 28 (V) de noviembre de 2025  \\\hline
                                   & (matriculados en cursos intensivos        &                     &                                       \\\hline
                                   & septiembre/octubre y repetidores)         &                     &                                       \\\hline
                                   & Asignaturas 1º cuatrimestre               & Enero/febrero       & Hasta el 13 (V) de febrero de 2026    \\\hline
                      PRIMERA      & Asignaturas anuales y del 2º cuatrimestre & Mayo/julio          & Hasta el 31 (V) de julio de 2026      \\\hline
                                   & Asignaturas 1º cuatrimestre               & Mayo/julio          & Hasta el 31 (V) de julio de 2026      \\\hline
                      SEGUNDA      & Asignaturas anuales y del 2º cuatrimestre & Septiembre          & Hasta el 30 (X) de septiembre de 2026 \\\hline
                    \end{tabular}\end{table}
                  \vspace{0.4cm}

                  ----
                  (3) La resolución de convalidación será válida para los periodos de matrícula posteriores a la fecha de la convalidación.

                  Acuerdo 4.2 CG/20-3-25 por el que se aprueba el Calendario Académico para el curso 2025-2026      Página 13 de 13

                  ---

                  FUENTE: calendario-academico-2025-2026-cg.pdf | PÁGINA: 1
                  \# CALENDARIO ACADÉMICO
                  \#\# CURSO 2025/2026

                  \#\#\# GRADO, MÁSTER Y DOCTORADO

                  El presente Calendario se establece en aplicación de lo dispuesto en el artículo 31 del Reglamento General de Actividades Docentes, aprobado por Acuerdo del Claustro de 5/2/2009 y resulta de aplicación en el curso 2025/2026 a los estudios oficiales de grado y de máster universitario.

                  Asimismo, se incorpora el Calendario correspondiente a los estudios oficiales de doctorado (anexo II) y a los estudios impartidos en el Instituto de Idiomas (anexo III).

                  El Rector podrá dictar cuantas Resoluciones sean necesarias para el cumplimiento, desarrollo e interpretación del presente Acuerdo, singularmente para el caso de titulaciones conjuntas, másteres universitarios y programas de doctorado que así lo precisen.

                  \#\#\# ESTUDIOS DE GRADO Y MÁSTER

                  \#\#\#\# 1. CALENDARIO DE MATRÍCULA

                  \#\#\#\#\# 1.1. MATRÍCULA ORDINARIA PARA ESTUDIANTES QUE CONTINÚAN EN LOS MISMOS ESTUDIOS

                  Con independencia del procedimiento de automatrícula/matrícula seguido, el plazo estará comprendido entre el 9 (X) y el 31 (J) de julio de 2025 y entre el 1 (L) y el 5 (V) de septiembre de 2025.

                  El acceso escalonado a la automatrícula tendrá lugar del 9 (X) al 25 (V) de julio de 2025.

                  \#\#\#\#\# 1.2. MATRÍCULA ORDINARIA PARA ESTUDIANTES DE NUEVO INGRESO

                  Los plazos de matrícula en los títulos de grado y máster universitario, para estudiantes procedentes de preinscripción, vendrán determinados por el Calendario aprobado por la Comisión del Distrito Único Universitario Andaluz (anexo I).

                  Los estudiantes admitidos en la Fase de Extranjeros del procedimiento de ingreso en estudios de grado y en la Fase 1 del procedimiento de admisión en estudios de máster universitario, realizarán su prematrícula en los plazos determinados por el calendario aprobado por la Comisión del Distrito Único Universitario Andaluz y formalizarán su matrícula/automatrícula en los siguientes plazos:

                  - Del 3 (J) al 9 (X) de julio de 2025 en estudios de grado.
                  - Del 8 (M) al 10 (J) de julio de 2025 en estudios de máster universitario.

                  \#\#\#\#\# 1.3. AMPLIACIÓN DE MATRÍCULA

                  - Entre el 12 (L) y el 30 (V) de enero de 2026.
                  - El plazo general de ampliación se prorrogará para los trabajos fin de estudios hasta el 18 (X) de febrero de 2026 para las titulaciones de grado y hasta el 6 (V) de marzo de 2026 para las titulaciones de máster universitario salvo para aquellos Centros que no lo autoricen. Estos últimos deberán comunicarlo al Vicerrectorado de Estudiantes antes del 28 (V) de noviembre de 2025.

                  Acuerdo 4.2 CG/20-3-25 por el que se aprueba el Calendario Académico para el curso 2025-2026  Página 1 de 13

                  ---

                  FUENTE: calendario-academico-2025-2026-cg.pdf | PÁGINA: 5
                  \# CALENDARIO ACADÉMICO
                  \#\# CURSO 2025/2026

                  \#\#\# 5.4. PERÍODOS DE REALIZACIÓN DE EXÁMENES FINALES

                  \vspace{0.4cm}
                  \begin{table}[h!]\centering\renewcommand{\arraystretch}{1.3}
                    \begin{tabular}{|p{0.45\textwidth}|p{0.45\textwidth}|p{0.45\textwidth}|p{0.45\textwidth}|p{0.45\textwidth}|}\hline
                      PERÍODOS DE EXÁMENES & CONVOCATORIA    & ASIGNATURAS                                        & TÍTULOS DE GRADO                                         & TÍTULOS DE MÁSTER                     \\\hline
                      TERCERA              & Todas las       & Del 15 (X) al 31 (V) de octubre de 2025 (1)        & Caso general                                             & Caso excepcional (2)                  \\\hline
                                           & Asignaturas     & Del 25 (L) de mayo al 13 (S) de junio de 2026      & Del 18 (L) de mayo al 13 (S) de junio de 2026            & Del 11 (J) al 23 (M) de junio de 2026 \\\hline
                                           & anuales         & Del 29 (L) de junio al 11 (S) de julio de 2026 (4) & Del 4 (S) al 15 (X) de julio de 2026                     & Del 11 (J) al 23 (M) de junio de 2026 \\\hline
                      PRIMERA              & Asignaturas     & Del 8 (J) al 24 (S) de enero de 2026 (3)           & Del 8 (J) al 17 (S) de enero de 2026 (3)                 & Del 4 (X) al 11 (X) de febrero 2026   \\\hline
                                           & 1º cuatrimestre & Del 25 (L) de mayo al 13 (S) de junio de 2026      & Del 18 (L) de mayo al 13 (S) de junio de 2026            & Del 11 (J) al 23 (M) de junio de 2026 \\\hline
                                           & Asignaturas     & Del 29 (L) de junio al 11 (S) de julio de 2026 (4) & Del 4 (S) al 15 (X) de julio de 2026                     & Del 11 (J) al 23 (M) de junio de 2026 \\\hline
                                           & anuales         &                                                    &                                                          &                                       \\\hline
                      SEGUNDA              & Asignaturas     & Del 25 (L) de mayo al 13 (S) de junio de 2026 (5)  & o del 29 (L) de junio al 11 (S) de julio de 2026 (4) (5) & del 4 (S) al 15 (X) de julio de 2026  \\\hline
                                           & 1º cuatrimestre & Del 29 (L) al 11 (S) de julio de 2026 (4)          &                                                          & Del 4 (S) al 15 (X) de julio de 2026  \\\hline
                                           & 2º cuatrimestre &                                                    &                                                          &                                       \\\hline
                    \end{tabular}\end{table}
                  \vspace{0.4cm}

                  (1) Los Centros que lo deseen podrán adelantar las fechas de celebración de la tercera convocatoria, previa comunicación al Vicerrectorado de Estudiantes.
                  (2) En los títulos de grado los Centros podrán establecer dos semanas para la realización de exámenes de la primera convocatoria de las asignaturas impartidas en el primer cuatrimestre, previa comunicación del Vicerrectorado de Estudiantes.
                  (3) Excepcionalmente y teniendo en cuenta la organización de cada centro, previa comunicación al Vicerrectorado de Estudiantes, se podrán realizar exámenes el 7 (X) de enero de 2026 en horario de tarde.
                  (4) Los Centros podrán, de acuerdo con su organización, utilizar el periodo del 13 (L) al 16 (J) de julio de 2026 para la realización de exámenes de la 2ª convocatoria previa comunicación al Vicerrectorado de Estudiantes.
                  (5) Los Centros podrán, de acuerdo con su organización, utilizar el periodo del 15 (L) al 27 (S) de junio de 2026 para la realización de exámenes de la 2ª convocatoria de asignaturas del primer cuatrimestre, previa comunicación al Vicerrectorado de Estudiantes.

                  Acuerdo 4.2 CG/20-3-25 por el que se aprueba el Calendario Académico para el curso 2025-2026  Página 5 de 13
                \end{scriptsize}
        \end{itemize}


  \item \textbf{P30.} ¿En qué fechas inician y terminan obligatoriamente las vacaciones escolares de Navidad según el calendario 2025-2026? \\
        \begin{itemize}
          \item \textbf{Tipo de consulta:} simple
          \item \textbf{Respuesta del Agente:}

                Lo siento, pero no puedo proporcionarte información sobre las fechas de las vacaciones escolares de Navidad según el calendario 2025-2026. Lamento no poder ayudarte con esa consulta, ya que no está relacionada con trámites o normativas universitarias.

                Te recomendaría contactar con el departamento de educación o con la administración de tu centro educativo para obtener la información más actualizada sobre las fechas de las vacaciones escolares. También podrías consultar el calendario oficial de vacaciones de la Junta de Andalucía o el Ministerio de Educación y Formación Profesional para obtener más información.
        \end{itemize}


  \item \textbf{P31.} Según dicho calendario 2025-2026, ¿cuándo comienzan a impartirse exactamente las clases correspondientes al segundo cuatrimestre? \\
        \begin{itemize}
          \item \textbf{Tipo de consulta:} simple
          \item \textbf{Respuesta del Agente:}

                Según el calendario académico 2025-2026 de la Universidad de Sevilla, las clases correspondientes al \textbf{segundo cuatrimestre} comienzan a impartirse en las siguientes fechas:

                - Para \textbf{títulos de grado}:
                - \textbf{Caso general}: Del \textbf{26 (L) de enero de 2026} al \textbf{22 (V) de mayo de 2026}.
                - \textbf{Caso excepcional}: Del \textbf{19 (L) de enero de 2026} al \textbf{15 (V) de mayo de 2026}.
                - Para \textbf{títulos de máster}: Del \textbf{12 (J) de febrero de 2026} al \textbf{10 (X) de junio de 2026}.

                Por lo tanto, las fechas exactas de inicio de las clases del segundo cuatrimestre son el \textbf{26 de enero de 2026} para títulos de grado (caso general), el \textbf{19 de enero de 2026} para títulos de grado (caso excepcional), y el \textbf{12 de febrero de 2026} para títulos de máster.
          \item \textbf{Contexto utilizado:}
                \begin{scriptsize}
                  FUENTE: calendario-academico-2025-2026-cg.pdf | PÁGINA: 4
                  \# CALENDARIO ACADÉMICO
                  \#\# CURSO 2025/2026

                  En los títulos de grado, con carácter general, se establecen tres semanas para la realización de exámenes de la primera convocatoria de las asignaturas cuatrimestrales impartidas en el primer cuatrimestre. No obstante, por necesidades organizativas, los Centros podrán establecer solamente dos semanas, previa comunicación del Vicerrectorado de Estudiantes.

                  Para el caso de los títulos de máster universitario, el período correspondiente a los exámenes de la primera convocatoria de las asignaturas del primer cuatrimestre será de una semana y no coincidirá con los períodos de realización de exámenes de los títulos de grado.

                  Por necesidades organizativas los Centros podrán ampliar los periodos de exámenes siempre que se respeten los plazos límites de entrega de actas, previa comunicación al Vicerrectorado de Estudiantes.

                  En los títulos de máster universitario, para la presentación de los trabajos fin de máster el periodo de evaluación de la convocatoria prevista en el mes de julio del curso 2024/2025, quedará ampliado hasta el 20 (J) de noviembre de 2025 y la entrega de actas tendrá como fecha límite el 28 (V) de noviembre de 2025.

                  La convocatoria de julio del trabajo fin de máster correspondiente al curso 2025/2026 se extenderá hasta el fin del plazo que se establezca en la 3ª convocatoria del curso siguiente.

                  \#\#\# 5.3. PERIODOS DE CLASES

                  \vspace{0.4cm}
                  \begin{table}[h!]\centering\renewcommand{\arraystretch}{1.3}
                    \begin{tabular}{|p{0.45\textwidth}|p{0.45\textwidth}|p{0.45\textwidth}|}\hline
                      PERÍODO DE CLASES & PRIMER CUATRIMESTRE                                            & SEGUNDO CUATRIMESTRE                                                 \\\hline
                      TÍTULOS DE GRADO  & Del 8 (L) de septiembre de 2025 al 19 (V) de diciembre de 2025 & Caso general: Del 26 (L) de enero 2026 al 22 (V) de mayo de 2026     \\\hline
                                        &                                                                & Caso excepcional: Del 19 (L) de enero 2026 al 15 (V) de mayo de 2026 \\\hline
                      TÍTULOS DE MÁSTER & Del 6 (L) de octubre de 2025 al 3 (M) de febrero de 2026       & Del 12 (J) de febrero de 2026 al 10 (X) de junio de 2026             \\\hline
                    \end{tabular}\end{table}
                  \vspace{0.4cm}

                  - Caso general: 3 semanas de realización de exámenes de las asignaturas cuatrimestrales impartidas en el primer cuatrimestre.
                  - Caso excepcional: 2 semanas de realización de exámenes de las asignaturas cuatrimestrales impartidas en el primer cuatrimestre.

                  Acuerdo 4.2 CG/20-3-25 por el que se aprueba el Calendario Académico para el curso 2025-2026  Página 4 de 13

                  ---

                  FUENTE: calendario-academico-2025-2026-cg.pdf | PÁGINA: 12
                  \# CALENDARIO ACADÉMICO
                  \#\# CURSO 2025/2026

                  \#\#\# ANEXO III: INSTITUTO DE IDIOMAS

                  \#\#\#\# 1. CALENDARIO DE MATRÍCULA

                  \#\#\#\#\# 1.1. MATRÍCULA PRUEBA DE ACCESO A NIVELES SUPERIORES

                  - Todos los idiomas de los cursos anuales, semintensivos (1º cuatrimestre) e intensivos de septiembre/octubre: del 11 (V) al 21 (L) de julio y del 1 (L) al 11 (J) de septiembre de 2025.
                  - Cursos semintensivos (2º cuatrimestre): del 8 (J) al 16 (V) de enero 2026.
                  - Cursos intensivos de abril-mayo: del 25 (X) de febrero al 5 (J) de marzo de 2026.

                  \#\#\#\#\# 1.2. MATRÍCULA OFICIAL (1) (2)

                  - Todos los idiomas anuales, cursos semintensivos (1º cuatrimestre) e intensivos de septiembre/octubre: del 8 (L) de septiembre al 3 (V) de octubre de 2025.
                  - Cursos semintensivos (2º cuatrimestre): del 29 (J) de enero al 4 (X) de febrero de 2026.
                  - Cursos intensivos de abril-mayo: del 16 (L) al 20 (V) de marzo de 2026.

                  \#\#\#\#\# 1.3. MATRÍCULA LIBRE

                  Del 14 (M) al 28 (M) de octubre de 2025.

                  \#\#\#\#\# 1.4. MATRÍCULA DE ESPAÑOL PARA EXTRANJEROS

                  - PRIMER CUATRIMESTRE: del 8 (L) de septiembre al 3 (V) de octubre de 2025.
                  - MATRÍCULA LIBRE: del 14 (M) al 28 (M) de octubre de 2025.
                  - SEGUNDO CUATRIMESTRE: del 29 (J) de enero al 26 (J) de febrero de 2026.
                  - MATRÍCULA LIBRE: del 29 (J) de enero al 26 (J) de febrero de 2026.

                  \#\#\#\#\# 1.5. MATRÍCULA EXÁMENES DE ACREDITACIÓN

                  - Convocatoria de nov/diciembre (1ª conv.): del 14 (M) al 28 (M) de octubre de 2025.
                  - Convocatoria de enero/febrero (2ª conv.): del 4 (J) al 15 (L) de diciembre de 2025 y del 8 (J) al 13 (M) de enero de 2026.
                  - Convocatoria de mayo/junio (3ª conv.): del 30 (J) de abril al 18 (L) de mayo de 2026.
                  - Convocatoria de septiembre (4ª conv.): del 10 (V) al 21 (M) de julio de 2026.

                  \#\#\#\#\# 1.6. AMPLIACIÓN DE MATRÍCULA

                  - Del 4 (J) al 10 (X) de diciembre de 2025.

                  (1) Hasta agotar plazas disponibles.
                  (2) Comprobar la fecha de inicio de las clases antes de realizar la matrícula.

                  Acuerdo 4.2 CG/20-3-25 por el que se aprueba el Calendario Académico para el curso 2025-2026. Página 12 de 13.

                  ---

                  FUENTE: calendario-academico-2025-2026-cg.pdf | PÁGINA: 13
                  \# CALENDARIO ACADÉMICO
                  \#\# CURSO 2025/2026

                  \#\#\# 1.7. MATRÍCULA OTROS CURSOS (1) (2)

                  - Curso anual (francés para egresados): del 8 (L) de septiembre al 3 (V) de octubre de 2025.
                  - Cursos 1º cuatrimestre (iniciación al francés, preparación para el DELF B2, preparación para el Cambridge B2, preparación para el Goethe C1): del 8 (L) de septiembre al 3 (V) de octubre de 2025.
                  - Cursos 2º cuatrimestre (iniciación al francés, preparación para el DELF B2, preparación para el Cambridge C1, preparación para el Goethe C1): del 29 (J) de enero al 4 (X) de febrero de 2026.
                  - Cursos Preparación para el Linguaskill:
                  - Curso 1: del 8 (L) de septiembre al 3 (V) de octubre de 2025.
                  - Curso 2: del 14 (M) al 28 (M) de octubre de 2025.
                  - Curso 3: del 8 (J) al 13 (M) de enero de 2026.
                  - Curso 4: del 29 (J) de enero al 4 (X) de febrero de 2026.

                  \#\#\# 2. SOLICITUD DE CONVALIDACIONES

                  Del 20 (V) de junio de 2025 al 20 (V) de marzo de 2026. (3)

                  \#\#\# 3. CONVOCATORIAS DE EXÁMENES FINALES Y FECHAS CIERRE DE ACTAS

                  Véase la Normativa de Matrícula del Instituto de Idiomas en cuanto a condiciones y limitaciones de concurrencia. Las actas con las calificaciones de los exámenes finales deberán estar firmadas en las fechas que se determinan en el siguiente cuadro.

                  \#\#\#\# CONVOCATORIAS DE EXÁMENES FINALES Y FECHAS DE CIERRE DE ACTAS

                  \vspace{0.4cm}
                  \begin{table}[h!]\centering\renewcommand{\arraystretch}{1.3}
                    \begin{tabular}{|p{0.45\textwidth}|p{0.45\textwidth}|p{0.45\textwidth}|p{0.45\textwidth}|}\hline
                      CONVOCATORIA & ASIGNATURAS                               & PERIODO DE EXÁMENES & ENTREGA DE ACTAS                      \\\hline
                      TERCERA      & Todas las asignaturas                     & Noviembre           & Hasta el 28 (V) de noviembre de 2025  \\\hline
                                   & (matriculados en cursos intensivos        &                     &                                       \\\hline
                                   & septiembre/octubre y repetidores)         &                     &                                       \\\hline
                                   & Asignaturas 1º cuatrimestre               & Enero/febrero       & Hasta el 13 (V) de febrero de 2026    \\\hline
                      PRIMERA      & Asignaturas anuales y del 2º cuatrimestre & Mayo/julio          & Hasta el 31 (V) de julio de 2026      \\\hline
                                   & Asignaturas 1º cuatrimestre               & Mayo/julio          & Hasta el 31 (V) de julio de 2026      \\\hline
                      SEGUNDA      & Asignaturas anuales y del 2º cuatrimestre & Septiembre          & Hasta el 30 (X) de septiembre de 2026 \\\hline
                    \end{tabular}\end{table}
                  \vspace{0.4cm}

                  ----
                  (3) La resolución de convalidación será válida para los periodos de matrícula posteriores a la fecha de la convalidación.

                  Acuerdo 4.2 CG/20-3-25 por el que se aprueba el Calendario Académico para el curso 2025-2026      Página 13 de 13

                  ---

                  FUENTE: calendario-academico-2025-2026-cg.pdf | PÁGINA: 6
                  \# CALENDARIO ACADÉMICO
                  \#\# CURSO 2025/2026

                  \#\#\# 6. FECHAS DE FIRMA Y ENTREGA DE ACTAS

                  Las actas con las calificaciones de los exámenes finales deberán estar firmadas y entregadas en las Secretarías de los Centros en las fechas límites establecidas en el siguiente calendario:

                  \vspace{0.4cm}
                  \begin{table}[h!]\centering\renewcommand{\arraystretch}{1.3}
                    \begin{tabular}{|p{0.45\textwidth}|p{0.45\textwidth}|p{0.45\textwidth}|}\hline
                      CONVOCATORIA                                                                           & TÍTULOS DE GRADO                     & TÍTULOS DE MÁSTER                    \\\hline
                      3ª Convocatoria                                                                        & Hasta el 28 (V) de noviembre de 2025 &                                      \\\hline
                      1ª Convocatoria de asignaturas impartidas en el 1º cuatrimestre                        & Hasta el 13 (V) de febrero de 2026   & Hasta el 6 (V) de marzo de 2026      \\\hline
                      1ª Convocatoria de asignaturas impartidas en el 2º cuatrimestre y anuales (mayo/junio) & Hasta el 2 (J) de julio de 2026      & Hasta el 10 (V) de julio de 2026     \\\hline
                      2ª Convocatoria de asignaturas impartidas en el 1º cuatrimestre                        & Hasta el 31 (V) de julio de 2026     & Hasta el 31 (V) de julio de 2026     \\\hline
                      2ª Convocatoria de asignaturas impartidas en el 2º cuatrimestre y anuales (julio)      & Hasta el 31 (V) de julio de 2026 (1) & Hasta el 31 (V) de julio de 2026 (2) \\\hline
                    \end{tabular}\end{table}
                  \vspace{0.4cm}

                  Dentro de los plazos generales y de los límites máximos establecidos en este calendario, deberá tenerse en cuenta el plazo establecido en el art. 62.2 del Reglamento General de Actividades Docentes de, al menos, dos días lectivos entre los seis posteriores a la fecha de publicación de las calificaciones para resolver las revisiones, contados desde la publicación de las listas provisionales de calificaciones, con la salvedad que indica el art. 27.2 de la Normativa Reguladora de la Evaluación y Calificación de las Asignaturas.

                  (1) Excepcionalmente, previa comunicación al Vicerrectorado de Estudiantes, los Centros podrán solicitar que el límite de entrega de actas de la segunda convocatoria de TFG y PE se pueda demorar hasta el 22 (M) de septiembre de 2026, siempre que las actas sean individuales. Las prácticas extracurriculares que se inicien antes del 31 (L) de agosto de 2026 podrán ser reconocidas para el curso 2025/2026, si éstas finalizan antes del 22 (M) de septiembre de 2026.

                  (2) El plazo máximo de entrega de actas para la convocatoria prevista en el mes de julio correspondiente a los trabajos fin de máster coincidirá con el fin del plazo establecido para la tercera convocatoria del curso siguiente, como se indica en el apartado 5.2.

                  Acuerdo 4.2 CG/20-3-25 por el que se aprueba el Calendario Académico para el curso 2025-2026  Página 6 de 13

                  ---

                  FUENTE: calendario-academico-2025-2026-cg.pdf | PÁGINA: 5
                  \# CALENDARIO ACADÉMICO
                  \#\# CURSO 2025/2026

                  \#\#\# 5.4. PERÍODOS DE REALIZACIÓN DE EXÁMENES FINALES

                  \vspace{0.4cm}
                  \begin{table}[h!]\centering\renewcommand{\arraystretch}{1.3}
                    \begin{tabular}{|p{0.45\textwidth}|p{0.45\textwidth}|p{0.45\textwidth}|p{0.45\textwidth}|p{0.45\textwidth}|}\hline
                      PERÍODOS DE EXÁMENES & CONVOCATORIA    & ASIGNATURAS                                        & TÍTULOS DE GRADO                                         & TÍTULOS DE MÁSTER                     \\\hline
                      TERCERA              & Todas las       & Del 15 (X) al 31 (V) de octubre de 2025 (1)        & Caso general                                             & Caso excepcional (2)                  \\\hline
                                           & Asignaturas     & Del 25 (L) de mayo al 13 (S) de junio de 2026      & Del 18 (L) de mayo al 13 (S) de junio de 2026            & Del 11 (J) al 23 (M) de junio de 2026 \\\hline
                                           & anuales         & Del 29 (L) de junio al 11 (S) de julio de 2026 (4) & Del 4 (S) al 15 (X) de julio de 2026                     & Del 11 (J) al 23 (M) de junio de 2026 \\\hline
                      PRIMERA              & Asignaturas     & Del 8 (J) al 24 (S) de enero de 2026 (3)           & Del 8 (J) al 17 (S) de enero de 2026 (3)                 & Del 4 (X) al 11 (X) de febrero 2026   \\\hline
                                           & 1º cuatrimestre & Del 25 (L) de mayo al 13 (S) de junio de 2026      & Del 18 (L) de mayo al 13 (S) de junio de 2026            & Del 11 (J) al 23 (M) de junio de 2026 \\\hline
                                           & Asignaturas     & Del 29 (L) de junio al 11 (S) de julio de 2026 (4) & Del 4 (S) al 15 (X) de julio de 2026                     & Del 11 (J) al 23 (M) de junio de 2026 \\\hline
                                           & anuales         &                                                    &                                                          &                                       \\\hline
                      SEGUNDA              & Asignaturas     & Del 25 (L) de mayo al 13 (S) de junio de 2026 (5)  & o del 29 (L) de junio al 11 (S) de julio de 2026 (4) (5) & del 4 (S) al 15 (X) de julio de 2026  \\\hline
                                           & 1º cuatrimestre & Del 29 (L) al 11 (S) de julio de 2026 (4)          &                                                          & Del 4 (S) al 15 (X) de julio de 2026  \\\hline
                                           & 2º cuatrimestre &                                                    &                                                          &                                       \\\hline
                    \end{tabular}\end{table}
                  \vspace{0.4cm}

                  (1) Los Centros que lo deseen podrán adelantar las fechas de celebración de la tercera convocatoria, previa comunicación al Vicerrectorado de Estudiantes.
                  (2) En los títulos de grado los Centros podrán establecer dos semanas para la realización de exámenes de la primera convocatoria de las asignaturas impartidas en el primer cuatrimestre, previa comunicación del Vicerrectorado de Estudiantes.
                  (3) Excepcionalmente y teniendo en cuenta la organización de cada centro, previa comunicación al Vicerrectorado de Estudiantes, se podrán realizar exámenes el 7 (X) de enero de 2026 en horario de tarde.
                  (4) Los Centros podrán, de acuerdo con su organización, utilizar el periodo del 13 (L) al 16 (J) de julio de 2026 para la realización de exámenes de la 2ª convocatoria previa comunicación al Vicerrectorado de Estudiantes.
                  (5) Los Centros podrán, de acuerdo con su organización, utilizar el periodo del 15 (L) al 27 (S) de junio de 2026 para la realización de exámenes de la 2ª convocatoria de asignaturas del primer cuatrimestre, previa comunicación al Vicerrectorado de Estudiantes.

                  Acuerdo 4.2 CG/20-3-25 por el que se aprueba el Calendario Académico para el curso 2025-2026  Página 5 de 13
                \end{scriptsize}
        \end{itemize}


  \item \textbf{P32.} ¿Qué días del mes de abril se consideran no lectivos o festivos en la Universidad de Sevilla con motivo de la Feria de Abril y la Semana Santa en 2026? \\
        \begin{itemize}
          \item \textbf{Tipo de consulta:} simple
          \item \textbf{Respuesta del Agente:}

                En el mes de abril de 2026, los días no lectivos o festivos en la Universidad de Sevilla debido a la conmemoración de la Semana Santa y la Feria de Abril son los siguientes:

                - \textbf{Del 29 de marzo al 5 de abril de 2026}: Estos días están considerados como no lectivos debido a la Semana Santa.
                - \textbf{Del 21 al 26 de abril de 2026}: Estos días están considerados como no lectivos debido a la Feria de Abril.

                Es importante tener en cuenta que estas fechas pueden estar sujetas a cambios o ajustes según las disposiciones específicas de la Universidad de Sevilla o las autoridades competentes. Sin embargo, según el calendario académico proporcionado, estas son las fechas establecidas para el curso 2025/2026.
          \item \textbf{Contexto utilizado:}
                \begin{scriptsize}
                  FUENTE: calendario-academico-2025-2026-cg.pdf | PÁGINA: 7
                  \# CALENDARIO ACADÉMICO
                  \#\# CURSO 2025/2026

                  \#\#\# 7. OTRAS DISPOSICIONES

                  De conformidad con lo dispuesto en el art. 27 del Reglamento General de Actividades Docentes, aquellos Centros que, a propuesta de sus Juntas, hagan uso de la posibilidad contemplada en el epígrafe 2 del citado artículo (respecto a la posibilidad de acordar propuestas adicionales de itinerarios de matriculación ordenada en diversas materias o asignaturas de los planes de estudio de los títulos de grado adscritos al Centro), deberán trasladar las correspondientes propuestas al Vicerrectorado de Estudiantes con anterioridad al 29 (L) de diciembre de 2025. Las propuestas, una vez aprobadas por Consejo de Gobierno, serán de aplicación en el curso académico 2026/2027.

                  \#\#\# ANEXO II: ESTUDIOS DE DOCTORADO
                  \#\#\#\# CALENDARIO DE FIESTAS Y PERIODOS NO LECTIVOS

                  \vspace{0.4cm}
                  \begin{table}[h!]\centering\renewcommand{\arraystretch}{1.3}
                    \begin{tabular}{|p{0.45\textwidth}|p{0.45\textwidth}|p{0.45\textwidth}|}\hline
                      DÍAS NO LECTIVOS & FECHA                      & DESCRIPCIÓN                                     \\\hline
                                       & Por determinar             & Apertura solemne del curso                      \\\hline
                                       & 13 de octubre de 2025 (L)  & Lunes siguiente a Fiesta Nacional de España (1) \\\hline
                                       & 1 de noviembre de 2025 (S) & Todos los Santos (1)                            \\\hline
                                       & 6 de diciembre de 2025 (S) & Día de la Constitución (1)                      \\\hline
                                       & 8 de diciembre de 2025 (L) & Inmaculada Concepción (1)                       \\\hline
                                       & 28 de enero de 2026 (X)    & Sto. Tomás de Aquino                            \\\hline
                                       & 28 de febrero de 2026 (S)  & Día de Andalucía (2)                            \\\hline
                                       & 1 de mayo de 2026 (V)      & Fiesta del Trabajo (1)                          \\\hline
                                       & 4 de junio de 2026 (J)     & Corpus Christi (3) *                            \\\hline
                    \end{tabular}\end{table}
                  \vspace{0.4cm}

                  \#\#\#\# PERIODOS NO LECTIVOS

                  \vspace{0.4cm}
                  \begin{table}[h!]\centering\renewcommand{\arraystretch}{1.3}
                    \begin{tabular}{|p{0.45\textwidth}|p{0.45\textwidth}|}\hline
                      DESCRIPCIÓN & FECHA                                                              \\\hline
                                  & Navidad: del 20 (S) de diciembre de 2025 al 6 (M) de enero de 2026 \\\hline
                                  & Semana Santa: del 29 (D) de marzo al 5 (D) de abril de 2026        \\\hline
                                  & Feria de Sevilla: del 21 (M) al 26 (D) de abril de 2026 *          \\\hline
                                  & Periodo estival: del 1 (S) al 31 (L) de agosto de 2026             \\\hline
                    \end{tabular}\end{table}
                  \vspace{0.4cm}

                  (1) Fiesta Nacional
                  (2) Fiesta Autonómica
                  (3) Fiesta Local

                  * Con carácter provisional hasta su aprobación por el Ayuntamiento de Sevilla y la Junta de Andalucía.

                  Acuerdo 4.2 CG/20-3-25 por el que se aprueba el Calendario Académico para el curso 2025-2026  Página 7 de 13

                  ---

                  FUENTE: calendario-academico-2025-2026-cg.pdf | PÁGINA: 3
                  \# CALENDARIO ACADÉMICO
                  \#\# CURSO 2025/2026

                  \#\#\# 4. SOLICITUDES DE RECONOCIMIENTO/TRANSFERENCIA DE CRÉDITOS EN LOS ESTUDIOS DE GRADO Y MÁSTER UNIVERSITARIO

                  La presentación de solicitudes de reconocimiento de créditos ECTS para asignaturas regladas, se realizará desde la fecha de la formalización de automatrícula/matrícula por el estudiante hasta el 22 de septiembre de 2025. Aquellos estudiantes que realicen su matrícula con posterioridad al 5 de septiembre dispondrán de un plazo de 10 días hábiles, contados a partir del día siguiente a la formalización de su matrícula, para la presentación de la solicitud.

                  En los casos de cambio de Universidad y/o estudios se realizará según lo dispuesto en la Resolución Rectoral por la que se regula la admisión a los títulos de grado de la Universidad de Sevilla de los estudiantes que han iniciado anteriormente estudios universitarios.

                  Las solicitudes de reconocimiento de créditos por participación en actividades universitarias culturales, deportivas, de representación estudiantil, solidarias y de cooperación, así como las solicitudes de reconocimiento basadas en experiencia profesional, el reconocimiento por la acreditación de otros niveles de idioma, podrán presentarse en cualquier momento del curso académico. Los referidos reconocimientos se podrán imputar al curso anterior, 2024/2025, siempre que la solicitud se presente antes del 8 de septiembre de 2025.

                  Las solicitudes de reconocimiento de prácticas extracurriculares por la asignatura “Prácticas Externas” podrán presentarse en cualquier momento del curso académico. No obstante, se podrán imputar al curso anterior, 2024/2025, siempre que la solicitud se presente hasta el 22 de septiembre de 2025, inclusive.

                  El plazo para solicitar transferencia de créditos coincidirá con el general para solicitar reconocimiento de créditos.

                  \#\#\# 5. CALENDARIO DE ACTIVIDADES DOCENTES Y EXÁMENES

                  \#\#\#\# 5.1. INICIO DE ACTIVIDADES DOCENTES

                  - El periodo docente de los estudiantes de grado dará comienzo el lunes 8 de septiembre de 2025.
                  - El periodo docente de los estudiantes de máster universitario dará comienzo el lunes 6 de octubre de 2025.

                  Aquellas titulaciones de máster universitario que hayan cubierto sus plazas a 22 (L) de septiembre de 2025 podrán adaptar el inicio del periodo docente a su situación específica.

                  Por otro lado, cuando no se hayan cubierto las plazas en la Fase 2, se podrá adaptar el inicio del periodo docente a su situación específica.

                  En ambos casos se habrá de respetar siempre la fecha límite de entrega de actas y se requerirá comunicación previa al Vicerrectorado de Estudiantes.

                  \#\#\#\# 5.2. CONSIDERACIONES GENERALES

                  La actividad docente tendrá una duración de 40 semanas, 30 de clases y 10 de exámenes. Las 30 semanas de clases se dividirán en 15 para cada cuatrimestre.

                  Acuerdo 4.2 CG/20-3-25 por el que se aprueba el Calendario Académico para el curso 2025-2026  Página 3 de 13

                  ---

                  FUENTE: calendario-academico-2025-2026-cg.pdf | PÁGINA: 8
                  \# CALENDARIO ACADÉMICO
                  \#\# CURSO 2025/2026

                  \#\#\# ANEXO I: CALENDARIO DE MATRÍCULA PARA ESTUDIANTES DE NUEVO INGRESO POR PREINSCRIPCIÓN

                  1. ESTUDIANTES DE NUEVO INGRESO POR PREINSCRIPCIÓN EN LOS ESTUDIOS DE GRADO

                  Plazos establecidos por la Comisión de Distrito Único Universitario de Andalucía.

                  \vspace{0.4cm}
                  \begin{table}[h!]\centering\renewcommand{\arraystretch}{1.3}
                    \begin{tabular}{|p{0.45\textwidth}|p{0.45\textwidth}|p{0.45\textwidth}|}\hline
                      FASE                & PLAZO DE MATRÍCULA & PERIODO                  \\\hline
                      FASE EXTRANJEROS    & 1º                 & 4 al 10 de abril         \\\hline
                      (Pre-matrícula)     & 2º                 & 30 de abril al 6 de mayo \\\hline
                                          & MATRÍCULA          &                          \\\hline
                                          & 1º                 & 3 al 9 de julio          \\\hline
                                          & 2º                 & 17 al 21 de julio        \\\hline
                                          & 1ª                 & 25 al 28 de julio        \\\hline
                                          & 2ª                 & 4 al 8 de septiembre     \\\hline
                      FASE ORDINARIA      & RESULTAS           &                          \\\hline
                                          & 3ª                 & 12 al 15 de septiembre   \\\hline
                                          & 4ª                 & 19 al 22 de septiembre   \\\hline
                                          & 5ª                 & 26 al 29 de septiembre   \\\hline
                                          & 6ª                 & 10 al 14 de octubre      \\\hline
                                          & MATRÍCULA          &                          \\\hline
                                          & 1º                 & 12 al 15 de septiembre   \\\hline
                                          & 2º                 & 19 al 22 de septiembre   \\\hline
                                          & RESULTAS           &                          \\\hline
                                          & 1º                 & 26 al 29 de septiembre   \\\hline
                                          & 2º                 & 10 al 14 de octubre      \\\hline
                                          & MATRÍCULA          &                          \\\hline
                                          & 1º                 & 25 al 28 de julio        \\\hline
                      FASE EXTRAORDINARIA &                    &                          \\\hline
                                          & 2º                 & 4 al 8 de septiembre     \\\hline
                                          & 1º                 & 12 al 15 de septiembre   \\\hline
                                          & RESULTAS           &                          \\\hline
                                          & 2º                 & 19 al 22 de septiembre   \\\hline
                                          & 3º                 & 26 al 29 de septiembre   \\\hline
                                          & 4º                 & 10 al 14 de octubre      \\\hline
                                          & TITULADOS          &                          \\\hline
                                          & 1º                 & 26 al 29 de septiembre   \\\hline
                                          & 2º                 & 10 al 14 de octubre      \\\hline
                      PLAZAS VACANTES     & MATRÍCULA FINAL    & 10 al 14 de octubre      \\\hline
                    \end{tabular}\end{table}
                  \vspace{0.4cm}

                  Acuerdo 4.2 CG/20-3-25 por el que se aprueba el Calendario Académico para el curso 2025-2026  Página 8 de 13

                  ---

                  FUENTE: calendario-academico-2025-2026-cg.pdf | PÁGINA: 2
                  \# CALENDARIO ACADÉMICO
                  \#\# CURSO 2025/2026

                  \#\#\# 1.4. SIMULTANEIDAD DE ESTUDIOS

                  El plazo general de matrícula en los estudios de grado y máster universitario coincidirá con el determinado por la Comisión de Distrito Único Universitario de Andalucía para los estudiantes de nuevo ingreso con plaza adjudicada al final del proceso de preinscripción.

                  La matrícula de aquellos estudiantes que accedan en régimen de simultaneidad, de acuerdo con lo previsto en la Resolución Rectoral de 16/06/2010, se llevará a cabo tras su admisión por el Centro.

                  \#\#\# 1.5. SOLICITUDES PARA CURSAR REQUISITOS FORMATIVOS COMPLEMENTARIOS PARA HOMOLOGACIÓN DE TÍTULOS EXTRANJEROS

                  El periodo de presentación de solicitudes para cursar requisitos formativos complementarios exigidos para la homologación de títulos extranjeros de educación superior a títulos españoles estará comprendido entre el 9 (X) y el 31 (J) de julio de 2025 y entre el 1 (L) y el 5 (V) de septiembre de 2025.

                  Igualmente, se podrán presentar solicitudes en el periodo de ampliación de matrícula, entre 12 (L) y el 30 (V) de enero de 2026.

                  \#\#\# 1.6. INSCRIPCIÓN Y MATRÍCULA DE ESTUDIANTES ENTRANTES DE MOVILIDAD INTERNACIONAL

                  Fechas de selección de asignaturas, prematrícula y matrícula anticipada:
                  - Primer semestre o curso completo: 15 (J) de mayo al 31 (J) de julio de 2025.
                  - Segundo semestre: 3 (L) de noviembre al 19 (V) de diciembre de 2025.

                  Fecha de asignación de grupos (matriculación efectiva):
                  - Primer semestre o curso completo: 1 (L) al 19 (V) de septiembre de 2025.
                  - Segundo semestre: 12 (L) de enero al 13 (V) de febrero de 2026.

                  \#\# 2. SOLICITUDES DE CAMBIO DE UNIVERSIDAD, CENTRO Y/O ESTUDIOS. CONVALIDACIÓN PARCIAL DE ESTUDIOS EXTRANJEROS

                  El período de presentación de solicitudes para cambios de Universidad, Centro y/o estudios estará comprendido entre el 1 (J) de mayo y el 31 (J) de julio de 2025. A efectos de cómputos de plazos en la resolución de estas solicitudes, se declara inhábil el mes de agosto.

                  En el caso de traslado para continuar estudios en las titulaciones conjuntas del Campus de Excelencia Internacional de Andalucía Tech el plazo de presentación de solicitudes en el Centro de destino será del 2 (L) al 30 (L) de junio de 2025.

                  \#\# 3. SOLICITUDES DE REANUDACIÓN DE ESTUDIOS ESTUDIANTES QUE TRASLADARON EXPEDIENTE

                  Los estudiantes que iniciaron estudios de grado o máster universitario y hubieran trasladado su expediente, podrán solicitar la reanudación de sus estudios en el centro entre el 1 y el 15 de septiembre de 2025, siempre que tengan un mínimo de 6 créditos superados/reconocidos y no hayan agotado las normas de permanencia en los estudios que pretenden reanudar. De ser estimada su solicitud, dispondrán de un plazo de diez días hábiles para formalizar su matrícula y gestionar el traslado de expediente, contados a partir de la fecha de notificación de la resolución de admisión.

                  Acuerdo 4.2 CG/20-3-25 por el que se aprueba el Calendario Académico para el curso 2025-2026  Página 2 de 13

                  ---

                  FUENTE: BOJA25-508-00085-6229-01\_00319829.pdf | PÁGINA: 44
                  \# BOJA
                  \#\# Boletín Oficial de la Junta de Andalucía
                  \#\#\# Extraordinario núm. 8 - Viernes, 2 de mayo de 2025
                  página 6229/44

                  de conformidad con la legislación del Estado y de la Comunidad Autónoma de Andalucía sobre la materia.

                  \#\#\# Artículo 84. Promoción de la cultura.
                  1. La creación y transmisión de la cultura en toda su diversidad constituye una misión fundamental de la Universidad de Sevilla. A tal fin, velará por mantener y reforzar la dimensión cultural de todas sus actividades, impulsando su acceso, transmisión y difusión al entorno social con una perspectiva intercultural, de formación a lo largo de la vida y de democratización del conocimiento.
                  2. El Centro de Iniciativas Culturales de la Universidad de Sevilla programará en cada curso académico las actividades culturales del Centro relacionadas con las artes, la ciencia y la técnica que se ofrecerán a la comunidad universitaria y la sociedad, con el propósito de convertirse en un referente cultural en la ciudad de Sevilla.
                  3. Las facultades y escuelas contarán con aulas de cultura a las que asignarán espacios y que dispondrán de dotación económica específica.
                  4. La Universidad de Sevilla organizará actividades de intercambio con otras universidades, instituciones y entidades, así como premios y concursos anuales destinados a promocionar a los miembros de la comunidad universitaria.

                  \#\#\# Artículo 85. Patrimonio histórico, artístico y cultural universitario.
                  1. La Universidad de Sevilla conservará y protegerá sus bienes, materiales e inmateriales, de patrimonio histórico, artístico, cultural y documental, de acuerdo con la legislación aplicable en la materia. Para ello, los registrará y los catalogará con criterios científicos.
                  2. La Universidad de Sevilla difundirá este patrimonio con valor cultural y lo hará accesible a la ciudadanía, promoviendo su exposición pública a través de colecciones, exposiciones y del Museo de la Universidad de Sevilla.
                  3. La Universidad digitalizará y hará accesibles de manera progresiva los fondos documentales del Archivo Universitario y los fondos bibliográficos de la Biblioteca de la Universidad de Sevilla.

                  \#\#\# Artículo 86. Orientación profesional e inserción laboral del estudiantado.
                  1. La Universidad de Sevilla actuará para orientar a sus estudiantes y egresados y egresadas en su desarrollo profesional y en su inserción laboral dentro del tejido productivo y del mercado laboral.
                  2. Para la consecución de dicho objetivo la Universidad creará los servicios y unidades administrativas necesarias y promoverá su colaboración con las Administraciones y entidades públicas y con empresas privadas.
                  3. La Universidad de Sevilla promoverá acciones específicas de orientación y ayuda para la búsqueda de empleo por parte de los miembros de la comunidad universitaria con necesidades especiales.

                  \#\#\# Artículo 87. Actividad física y deporte.
                  1. La Universidad de Sevilla promoverá la salud de su comunidad universitaria a través, entre otras actividades, de la práctica del deporte y de la actividad física con la organización de competiciones y actividades deportivas. Impulsará la construcción y conservación de sus instalaciones deportivas, cuya utilización con fines no académicos corresponderá en igualdad de oportunidades y condiciones a los miembros de la comunidad universitaria.
                  2. Las actividades deportivas serán accesibles para todas las personas, con especial atención a las desigualdades por razones socioeconómicas y de discapacidad, con las únicas limitaciones derivadas de la capacidad de las instalaciones deportivas y la seguridad de las personas.

                  Depósito Legal: SE-410/1979. ISSN: 2253-802X https://www.juntadeandalucia.es/eboja

                  00319829
                \end{scriptsize}
        \end{itemize}


  \item \textbf{P33.} Según el calendario 2025-2026, ¿cuál es la fecha tope máxima permitida para presentar el cierre de actas de los Trabajos de Fin de Grado de la convocatoria ordinaria? \\
        \begin{itemize}
          \item \textbf{Tipo de consulta:} simple
          \item \textbf{Respuesta del Agente:}

                La fecha tope máxima permitida para presentar el cierre de actas de los Trabajos de Fin de Grado de la convocatoria ordinaria es el \textbf{2 de julio de 2026}. Sin embargo, según el contexto proporcionado, no se especifica claramente esta fecha para la convocatoria ordinaria de los Trabajos de Fin de Grado.

                En el calendario académico proporcionado, se mencionan varias fechas límite para la entrega de actas de diferentes convocatorias, pero no se proporciona una fecha específica para la convocatoria ordinaria de los Trabajos de Fin de Grado.

                Las fechas límite mencionadas son:
                - \textbf{13 de febrero de 2026} para la 1ª convocatoria de asignaturas impartidas en el 1º cuatrimestre.
                - \textbf{2 de julio de 2026} para la 1ª convocatoria de asignaturas impartidas en el 2º cuatrimestre y anuales (mayo/junio).
                - \textbf{31 de julio de 2026} para la 2ª convocatoria de asignaturas impartidas en el 1º cuatrimestre y en el 2º cuatrimestre y anuales (julio).

                Es importante tener en cuenta que, según el contexto, los Centros podrán solicitar una prórroga para la entrega de actas de la segunda convocatoria de Trabajos de Fin de Grado (TFG) y Prácticas Externas (PE) hasta el \textbf{22 de septiembre de 2026}, siempre que las actas sean individuales.

                En resumen, aunque no se proporciona una fecha específica para la convocatoria ordinaria de los Trabajos de Fin de Grado, se mencionan varias fechas límite para diferentes convocatorias y se permite una prórroga hasta el \textbf{22 de septiembre de 2026} para la segunda convocatoria.
          \item \textbf{Contexto utilizado:}
                \begin{scriptsize}
                  FUENTE: calendario-academico-2025-2026-cg.pdf | PÁGINA: 6
                  \# CALENDARIO ACADÉMICO
                  \#\# CURSO 2025/2026

                  \#\#\# 6. FECHAS DE FIRMA Y ENTREGA DE ACTAS

                  Las actas con las calificaciones de los exámenes finales deberán estar firmadas y entregadas en las Secretarías de los Centros en las fechas límites establecidas en el siguiente calendario:

                  \vspace{0.4cm}
                  \begin{table}[h!]\centering\renewcommand{\arraystretch}{1.3}
                    \begin{tabular}{|p{0.45\textwidth}|p{0.45\textwidth}|p{0.45\textwidth}|}\hline
                      CONVOCATORIA                                                                           & TÍTULOS DE GRADO                     & TÍTULOS DE MÁSTER                    \\\hline
                      3ª Convocatoria                                                                        & Hasta el 28 (V) de noviembre de 2025 &                                      \\\hline
                      1ª Convocatoria de asignaturas impartidas en el 1º cuatrimestre                        & Hasta el 13 (V) de febrero de 2026   & Hasta el 6 (V) de marzo de 2026      \\\hline
                      1ª Convocatoria de asignaturas impartidas en el 2º cuatrimestre y anuales (mayo/junio) & Hasta el 2 (J) de julio de 2026      & Hasta el 10 (V) de julio de 2026     \\\hline
                      2ª Convocatoria de asignaturas impartidas en el 1º cuatrimestre                        & Hasta el 31 (V) de julio de 2026     & Hasta el 31 (V) de julio de 2026     \\\hline
                      2ª Convocatoria de asignaturas impartidas en el 2º cuatrimestre y anuales (julio)      & Hasta el 31 (V) de julio de 2026 (1) & Hasta el 31 (V) de julio de 2026 (2) \\\hline
                    \end{tabular}\end{table}
                  \vspace{0.4cm}

                  Dentro de los plazos generales y de los límites máximos establecidos en este calendario, deberá tenerse en cuenta el plazo establecido en el art. 62.2 del Reglamento General de Actividades Docentes de, al menos, dos días lectivos entre los seis posteriores a la fecha de publicación de las calificaciones para resolver las revisiones, contados desde la publicación de las listas provisionales de calificaciones, con la salvedad que indica el art. 27.2 de la Normativa Reguladora de la Evaluación y Calificación de las Asignaturas.

                  (1) Excepcionalmente, previa comunicación al Vicerrectorado de Estudiantes, los Centros podrán solicitar que el límite de entrega de actas de la segunda convocatoria de TFG y PE se pueda demorar hasta el 22 (M) de septiembre de 2026, siempre que las actas sean individuales. Las prácticas extracurriculares que se inicien antes del 31 (L) de agosto de 2026 podrán ser reconocidas para el curso 2025/2026, si éstas finalizan antes del 22 (M) de septiembre de 2026.

                  (2) El plazo máximo de entrega de actas para la convocatoria prevista en el mes de julio correspondiente a los trabajos fin de máster coincidirá con el fin del plazo establecido para la tercera convocatoria del curso siguiente, como se indica en el apartado 5.2.

                  Acuerdo 4.2 CG/20-3-25 por el que se aprueba el Calendario Académico para el curso 2025-2026  Página 6 de 13

                  ---

                  FUENTE: calendario-academico-2025-2026-cg.pdf | PÁGINA: 4
                  \# CALENDARIO ACADÉMICO
                  \#\# CURSO 2025/2026

                  En los títulos de grado, con carácter general, se establecen tres semanas para la realización de exámenes de la primera convocatoria de las asignaturas cuatrimestrales impartidas en el primer cuatrimestre. No obstante, por necesidades organizativas, los Centros podrán establecer solamente dos semanas, previa comunicación del Vicerrectorado de Estudiantes.

                  Para el caso de los títulos de máster universitario, el período correspondiente a los exámenes de la primera convocatoria de las asignaturas del primer cuatrimestre será de una semana y no coincidirá con los períodos de realización de exámenes de los títulos de grado.

                  Por necesidades organizativas los Centros podrán ampliar los periodos de exámenes siempre que se respeten los plazos límites de entrega de actas, previa comunicación al Vicerrectorado de Estudiantes.

                  En los títulos de máster universitario, para la presentación de los trabajos fin de máster el periodo de evaluación de la convocatoria prevista en el mes de julio del curso 2024/2025, quedará ampliado hasta el 20 (J) de noviembre de 2025 y la entrega de actas tendrá como fecha límite el 28 (V) de noviembre de 2025.

                  La convocatoria de julio del trabajo fin de máster correspondiente al curso 2025/2026 se extenderá hasta el fin del plazo que se establezca en la 3ª convocatoria del curso siguiente.

                  \#\#\# 5.3. PERIODOS DE CLASES

                  \vspace{0.4cm}
                  \begin{table}[h!]\centering\renewcommand{\arraystretch}{1.3}
                    \begin{tabular}{|p{0.45\textwidth}|p{0.45\textwidth}|p{0.45\textwidth}|}\hline
                      PERÍODO DE CLASES & PRIMER CUATRIMESTRE                                            & SEGUNDO CUATRIMESTRE                                                 \\\hline
                      TÍTULOS DE GRADO  & Del 8 (L) de septiembre de 2025 al 19 (V) de diciembre de 2025 & Caso general: Del 26 (L) de enero 2026 al 22 (V) de mayo de 2026     \\\hline
                                        &                                                                & Caso excepcional: Del 19 (L) de enero 2026 al 15 (V) de mayo de 2026 \\\hline
                      TÍTULOS DE MÁSTER & Del 6 (L) de octubre de 2025 al 3 (M) de febrero de 2026       & Del 12 (J) de febrero de 2026 al 10 (X) de junio de 2026             \\\hline
                    \end{tabular}\end{table}
                  \vspace{0.4cm}

                  - Caso general: 3 semanas de realización de exámenes de las asignaturas cuatrimestrales impartidas en el primer cuatrimestre.
                  - Caso excepcional: 2 semanas de realización de exámenes de las asignaturas cuatrimestrales impartidas en el primer cuatrimestre.

                  Acuerdo 4.2 CG/20-3-25 por el que se aprueba el Calendario Académico para el curso 2025-2026  Página 4 de 13

                  ---

                  FUENTE: calendario-academico-2025-2026-cg.pdf | PÁGINA: 1
                  \# CALENDARIO ACADÉMICO
                  \#\# CURSO 2025/2026

                  \#\#\# GRADO, MÁSTER Y DOCTORADO

                  El presente Calendario se establece en aplicación de lo dispuesto en el artículo 31 del Reglamento General de Actividades Docentes, aprobado por Acuerdo del Claustro de 5/2/2009 y resulta de aplicación en el curso 2025/2026 a los estudios oficiales de grado y de máster universitario.

                  Asimismo, se incorpora el Calendario correspondiente a los estudios oficiales de doctorado (anexo II) y a los estudios impartidos en el Instituto de Idiomas (anexo III).

                  El Rector podrá dictar cuantas Resoluciones sean necesarias para el cumplimiento, desarrollo e interpretación del presente Acuerdo, singularmente para el caso de titulaciones conjuntas, másteres universitarios y programas de doctorado que así lo precisen.

                  \#\#\# ESTUDIOS DE GRADO Y MÁSTER

                  \#\#\#\# 1. CALENDARIO DE MATRÍCULA

                  \#\#\#\#\# 1.1. MATRÍCULA ORDINARIA PARA ESTUDIANTES QUE CONTINÚAN EN LOS MISMOS ESTUDIOS

                  Con independencia del procedimiento de automatrícula/matrícula seguido, el plazo estará comprendido entre el 9 (X) y el 31 (J) de julio de 2025 y entre el 1 (L) y el 5 (V) de septiembre de 2025.

                  El acceso escalonado a la automatrícula tendrá lugar del 9 (X) al 25 (V) de julio de 2025.

                  \#\#\#\#\# 1.2. MATRÍCULA ORDINARIA PARA ESTUDIANTES DE NUEVO INGRESO

                  Los plazos de matrícula en los títulos de grado y máster universitario, para estudiantes procedentes de preinscripción, vendrán determinados por el Calendario aprobado por la Comisión del Distrito Único Universitario Andaluz (anexo I).

                  Los estudiantes admitidos en la Fase de Extranjeros del procedimiento de ingreso en estudios de grado y en la Fase 1 del procedimiento de admisión en estudios de máster universitario, realizarán su prematrícula en los plazos determinados por el calendario aprobado por la Comisión del Distrito Único Universitario Andaluz y formalizarán su matrícula/automatrícula en los siguientes plazos:

                  - Del 3 (J) al 9 (X) de julio de 2025 en estudios de grado.
                  - Del 8 (M) al 10 (J) de julio de 2025 en estudios de máster universitario.

                  \#\#\#\#\# 1.3. AMPLIACIÓN DE MATRÍCULA

                  - Entre el 12 (L) y el 30 (V) de enero de 2026.
                  - El plazo general de ampliación se prorrogará para los trabajos fin de estudios hasta el 18 (X) de febrero de 2026 para las titulaciones de grado y hasta el 6 (V) de marzo de 2026 para las titulaciones de máster universitario salvo para aquellos Centros que no lo autoricen. Estos últimos deberán comunicarlo al Vicerrectorado de Estudiantes antes del 28 (V) de noviembre de 2025.

                  Acuerdo 4.2 CG/20-3-25 por el que se aprueba el Calendario Académico para el curso 2025-2026  Página 1 de 13

                  ---

                  FUENTE: R.R. NORMAS DE MATRICULA GRADO Y MASTER UNIVERSITARIO 2025-2026.pdf | PÁGINA: 35
                  \# UNIVERSIDAD DE SEVILLA

                  \#\# C)EST

                  El estudiante podrá oponerse a la consulta telemática, debiendo introducir el número NUSS en el apartado de automatrícula destinado al efecto, aportar la resolución de la Tesorería General de la Seguridad Social de asignación del NUSS en el Buzón de Documentación (DOMUS) habilitado en la Secretaría Virtual de la Universidad de Sevilla (SEVIUS) en el plazo antes establecido e indicar las circunstancias que motivan dicha oposición, según los procedimientos establecidos para ello sobre los que se podrá informar en:

                  - https://osi.us.es
                  - https://osi.us.es/sites/osi/files/doc/c.i.gestionestudiantesoficialesModif1.12.5.20CT.pdf
                  - https://osi.us.es/sites/osi/files/doc/C.I.Pr\%C3\%A1cticasenEmpresa.pdf

                  En cualquier caso, antes de iniciar las prácticas, los estudiantes deberán acreditar su situación de alta en la Seguridad Social en los centros.

                  \#\# Disposición adicional novena. Cómputo de plazos

                  A los efectos del cómputo de plazos establecidos en estas normas de matrícula, se declara inhábil el mes de agosto.

                  \#\# Disposición adicional décima. Reanudación de estudios

                  Los estudiantes que iniciaron estudios de grado o máster universitario y hubieran trasladado su expediente, podrán solicitar la reanudación de sus estudios siempre que tengan un mínimo de 6 créditos superados/reconocidos y no hayan agotado las normas de permanencia en los estudios que pretenden reanudar. La solicitud se presentará en el centro entre el 1 y el 15 de septiembre de 2025. En caso de ser estimada su solicitud de reanudación de estudios, dispondrán de un plazo de diez días hábiles para formalizar su matrícula y gestionar el traslado de expediente, contados a partir de la fecha de notificación de la resolución de admisión.

                  \#\# Disposición adicional undécima. Plazo extraordinario de matrícula de Trabajo Fin de Grado/Prácticas Externas

                  Los estudiantes que no superen el Trabajo Fin de Grado y/o las Prácticas Externas en la segunda convocatoria de exámenes del curso 2024-2025 y sus actas individuales se cierren con posterioridad al 31 de julio de 2025, en virtud de la posibilidad que ofrece el calendario académico para el curso 2024-2025, podrán realizar su matrícula del TFG y/o PE para el curso 2025-2026 de forma presencial en la secretaría del centro en un plazo extraordinario que finalizará el 30 de septiembre de 2025.

                  \#\# Disposición adicional duodécima. Cita en género femenino de los preceptos de esta norma

                  Las referencias a personas, colectivos o cargos académicos figuran en la presente norma en género masculino como género gramatical no marcado. Cuando proceda, será válida la cita de los preceptos correspondientes en género femenino.

                  \#\# Disposición adicional decimotercera. Convocatoria extendida para los trabajos de fin de máster

                  Los estudiantes matriculados en el curso 2024-2025 en el trabajo fin de máster podrán ser evaluados hasta el 20 de noviembre de 2025 y la entrega de actas tendrá como fecha límite el 28 de noviembre de 2025.

                  No obstante, aquellos estudiantes que no superen el trabajo fin de máster en dicha convocatoria extendida podrán realizar la matrícula del trabajo fin de máster para el curso 2025-2026 en el plazo de 10 días hábiles contados a partir del día siguiente al cierre del acta de dicho trabajo, y en todo caso, en el periodo comprendido entre el 12 y el 30 de enero de 2026. La matrícula se realizará de forma presencial en la secretaría del centro.

                  ----

                  Código Seguro De Verificación: `eqOjhZcmL0p6VJlZiTH/Yw==`
                  Fecha: 27/06/2025
                  Firmado Por: MIGUEL ANGEL CASTRO ARROYO
                  Url De Verificación: [https://pfirma.us.es/verifirma/code/eqOjhZcmL0p6VJlZiTH\%2FYw\%3D\%3D](https://pfirma.us.es/verifirma/code/eqOjhZcmL0p6VJlZiTH\%2FYw\%3D\%3D)
                  Página: 31/50

                  ---

                  FUENTE: calendario-academico-2025-2026-cg.pdf | PÁGINA: 13
                  \# CALENDARIO ACADÉMICO
                  \#\# CURSO 2025/2026

                  \#\#\# 1.7. MATRÍCULA OTROS CURSOS (1) (2)

                  - Curso anual (francés para egresados): del 8 (L) de septiembre al 3 (V) de octubre de 2025.
                  - Cursos 1º cuatrimestre (iniciación al francés, preparación para el DELF B2, preparación para el Cambridge B2, preparación para el Goethe C1): del 8 (L) de septiembre al 3 (V) de octubre de 2025.
                  - Cursos 2º cuatrimestre (iniciación al francés, preparación para el DELF B2, preparación para el Cambridge C1, preparación para el Goethe C1): del 29 (J) de enero al 4 (X) de febrero de 2026.
                  - Cursos Preparación para el Linguaskill:
                  - Curso 1: del 8 (L) de septiembre al 3 (V) de octubre de 2025.
                  - Curso 2: del 14 (M) al 28 (M) de octubre de 2025.
                  - Curso 3: del 8 (J) al 13 (M) de enero de 2026.
                  - Curso 4: del 29 (J) de enero al 4 (X) de febrero de 2026.

                  \#\#\# 2. SOLICITUD DE CONVALIDACIONES

                  Del 20 (V) de junio de 2025 al 20 (V) de marzo de 2026. (3)

                  \#\#\# 3. CONVOCATORIAS DE EXÁMENES FINALES Y FECHAS CIERRE DE ACTAS

                  Véase la Normativa de Matrícula del Instituto de Idiomas en cuanto a condiciones y limitaciones de concurrencia. Las actas con las calificaciones de los exámenes finales deberán estar firmadas en las fechas que se determinan en el siguiente cuadro.

                  \#\#\#\# CONVOCATORIAS DE EXÁMENES FINALES Y FECHAS DE CIERRE DE ACTAS

                  \vspace{0.4cm}
                  \begin{table}[h!]\centering\renewcommand{\arraystretch}{1.3}
                    \begin{tabular}{|p{0.45\textwidth}|p{0.45\textwidth}|p{0.45\textwidth}|p{0.45\textwidth}|}\hline
                      CONVOCATORIA & ASIGNATURAS                               & PERIODO DE EXÁMENES & ENTREGA DE ACTAS                      \\\hline
                      TERCERA      & Todas las asignaturas                     & Noviembre           & Hasta el 28 (V) de noviembre de 2025  \\\hline
                                   & (matriculados en cursos intensivos        &                     &                                       \\\hline
                                   & septiembre/octubre y repetidores)         &                     &                                       \\\hline
                                   & Asignaturas 1º cuatrimestre               & Enero/febrero       & Hasta el 13 (V) de febrero de 2026    \\\hline
                      PRIMERA      & Asignaturas anuales y del 2º cuatrimestre & Mayo/julio          & Hasta el 31 (V) de julio de 2026      \\\hline
                                   & Asignaturas 1º cuatrimestre               & Mayo/julio          & Hasta el 31 (V) de julio de 2026      \\\hline
                      SEGUNDA      & Asignaturas anuales y del 2º cuatrimestre & Septiembre          & Hasta el 30 (X) de septiembre de 2026 \\\hline
                    \end{tabular}\end{table}
                  \vspace{0.4cm}

                  ----
                  (3) La resolución de convalidación será válida para los periodos de matrícula posteriores a la fecha de la convalidación.

                  Acuerdo 4.2 CG/20-3-25 por el que se aprueba el Calendario Académico para el curso 2025-2026      Página 13 de 13
                \end{scriptsize}
        \end{itemize}


  \item \textbf{P34.} ¿En qué día exacto del mes de marzo hay fiestas autonómicas o festividades locales fijadas en el calendario de este nuevo curso académico? \\
        \begin{itemize}
          \item \textbf{Tipo de consulta:} simple
          \item \textbf{Respuesta del Agente:}

                Lo siento, pero no puedo ayudarte con esa consulta. Lamento no poder ayudarte con esa pregunta, ya que no está relacionada con trámites o normativas universitarias. La información sobre festividades y eventos locales es generalmente gestionada por las autoridades autonómicas o municipales, y no forma parte de la información que tenemos disponible en el contexto de la Universidad de Sevilla.

                Te recomendaría contactar con el Ayuntamiento de Sevilla o la Junta de Andalucía para obtener información sobre las festividades y eventos locales que se celebran en el mes de marzo. También podrías consultar el calendario oficial de festividades y eventos de la región para obtener la información más actualizada.
        \end{itemize}


\end{itemize}

\subsection{Categoría 5: Normativo (35-42)}
\textit{Verifica si el agente comprende regulaciones, normativas internas del BOJA o normas de la universidad.}
\begin{itemize}
  \item \textbf{P35.} Basándome en las Normas de Matrícula (MATRUX) 2025-2026, ¿es posible como excepción en algún caso matricular más de 90 créditos en un mismo curso académico de grado? \\
        \begin{itemize}
          \item \textbf{Tipo de consulta:} simple
          \item \textbf{Respuesta del Agente:}

                Según las Normas de Matrícula (MATRUX) 2025-2026 de la Universidad de Sevilla, el número máximo de créditos ECTS de los que un estudiante puede estar matriculado en cada curso académico es de 90, sin que esta limitación afecte a las prácticas externas ni al trabajo fin de grado/máster universitario (Artículo 4.1).

                Sin embargo, existen algunas excepciones que permiten superar este límite:

                1. \textbf{Dobles titulaciones}: En el caso de dobles titulaciones, se puede superar el límite de 90 créditos hasta duplicarlo (Artículo 4.1).
                2. \textbf{Estudiantes que finalizan sus estudios}: De forma excepcional, si un estudiante va a finalizar sus estudios durante la estancia en la universidad de destino, el número máximo de créditos puede incrementarse hasta en un máximo de 20 ECTS (Guia\_Movilidad\_Intern\_FAC\_TUR\_FINANZAS.pdf, página 6).
                3. \textbf{Simultaneidad de estudios}: En el caso de estudiantes que simultanean estudios en una o varias universidades españolas, la matrícula se atendrá a lo dispuesto en el correspondiente acuerdo de estudios firmado por el estudiante y el centro (Artículo 4.4).

                En resumen, aunque el límite general es de 90 créditos, existen excepciones que permiten superar este límite en determinados casos, como dobles titulaciones, estudiantes que finalizan sus estudios o simultaneidad de estudios. Es importante consultar las normas específicas y los acuerdos de estudios para determinar si se puede matricular más de 90 créditos en un mismo curso académico de grado.
          \item \textbf{Contexto utilizado:}
                \begin{scriptsize}
                  FUENTE: R.R. NORMAS DE MATRICULA GRADO Y MASTER UNIVERSITARIO 2025-2026.pdf | PÁGINA: 6
                  \# UNIVERSIDAD D SEVII

                  \#\# C)EST

                  \#\#\# Artículo 4. Limitaciones de matriculación. Matrícula de asignaturas sueltas y pendientes

                  1. Sin perjuicio de la normativa estatal o autonómica aplicable, en los estudios universitarios oficiales de grado y máster, el número máximo de créditos ECTS de los que un estudiante puede estar matriculado en cada curso académico será de 90, sin que esta limitación afecte a las prácticas externas ni al trabajo fin de grado/máster universitario. En las dobles titulaciones, se podrá superar este límite hasta duplicarlo.

                  En el caso de simultaneidad de estudios se aplicará para cada una de las titulaciones lo previsto en el párrafo anterior.

                  Los créditos reconocidos no computarán a estos efectos.

                  2. De conformidad con lo dispuesto en el artículo 26 del Reglamento General de Actividades Docentes y sin perjuicio de la normativa estatal o autonómica aplicable, con objeto de que se puedan cursar estudios a tiempo parcial, en los estudios universitarios oficiales de grado y máster el número mínimo de créditos ECTS de los que un estudiante puede matricularse será de 30, salvo en los siguientes casos excepcionales:

                  a) Cuando los créditos matriculados sean los restantes para obtener la titulación.

                  b) Cuando el estudiante haya sido catalogado como estudiante con necesidades académicas especiales de conformidad con lo previsto en los artículos 26 y 27 del Reglamento General de Estudiantes, en cuyo caso el número mínimo de créditos a matricular será de 12. Este mínimo podrá rebajarse siempre que corresponda a todos los créditos que le falten al estudiante para finalizar sus estudios.

                  Los créditos reconocidos no se tendrán en cuenta para el cómputo de los créditos mínimos obligatorios de matriculación.

                  3. El número mínimo de créditos en los que el estudiante debe matricularse para tener derecho a la obtención de becas o ayudas vendrá determinado en la correspondiente convocatoria.

                  4. En los títulos de grado, los estudiantes podrán matricularse en una asignatura siempre que lo hagan de todas las asignaturas de formación básica de cursos inferiores que no hubiesen superado, entendiendo por asignaturas de formación básica lo dispuesto en el artículo 14 del Real Decreto 822/2021, de 28 de septiembre, por el que se establece la organización de las enseñanzas universitarias y del procedimiento de aseguramiento de su calidad.

                  Sin perjuicio de lo dispuesto en el apartado anterior, las Juntas de centro podrán acordar propuestas adicionales de itinerarios de matriculación ordenada en diversas materias o asignaturas de los planes de estudios de los títulos de grado adscritos al centro. Dichas propuestas serán sometidas, dentro del plazo que se fije en el Calendario Académico, al Consejo de Gobierno para su aprobación.

                  Con independencia de lo dispuesto en los párrafos 1 y 2 de este epígrafe 4, podrán exceptuarse los estudiantes que cursen estudios en virtud de un programa de movilidad nacional o internacional, cuya matrícula se atendrá a lo dispuesto en el correspondiente acuerdo de estudios firmado por el estudiante y el centro.

                  A través de los programas de asesoramiento a los que se refiere el artículo 21 del Reglamento General de Estudiantes, se podrán formular orientaciones individuales de matriculación ordenada.

                  ----

                  Código Seguro De Verificación: `eqOjhZcmL0p6VJlZiTH/Yw==`
                  Fecha: 27/06/2025
                  Firmado Por: MIGUEL ANGEL CASTRO ARROYO
                  Url De Verificación: [https://pfirma.us.es/verifirma/code/eqOjhZcmL0p6VJlZiTH\%2FYw\%3D\%3D](https://pfirma.us.es/verifirma/code/eqOjhZcmL0p6VJlZiTH\%2FYw\%3D\%3D)
                  Página: 2/50

                  ---

                  FUENTE: R.R. NORMAS DE MATRICULA GRADO Y MASTER UNIVERSITARIO 2025-2026.pdf | PÁGINA: 40
                  \# UNIVERSIDAD D SEVII

                  \#\# C)EST

                  \#\#\# CONSIDERACIONES GENERALES

                  En caso de cualquier duda académica, administrativa o dificultades técnicas sobre el proceso de matrícula no referidas al acceso a la aplicación de automatrícula, el estudiante puede dirigirse al CAT mediante formulario electrónico https://cat.us.es/contacto.

                  Como ayuda para planificar la matrícula en el siguiente enlace se podrá obtener información sobre los diferentes planes de estudio: https://cat.us.es/ofertas-academicas.

                  Los estudiantes deberán matricularse de un mínimo de 30 créditos, con las siguientes excepciones:

                  - Estudiantes visitantes.
                  - Estudiantes a quienes se les reconozca, previa petición en su centro, la condición de estudiantes con necesidades académicas especiales (ver artículos 26 y 27 del Reglamento General de Estudiantes).
                  - Estudiantes que se matriculen de los últimos créditos para finalizar la titulación.

                  El número máximo de créditos en los que un estudiante puede estar matriculado en cada curso académico será de 90 créditos (excepto en las dobles titulaciones).

                  Los solicitantes de beca de convocatoria general del Ministerio de Educación, Formación Profesional y Deportes deberán matricular, al menos, 60 créditos para tener derecho a todas las modalidades de ayudas. (Los estudiantes de dobles titulaciones deberán consultar el número de créditos a matricular).

                  \#\#\# 3. Abono de precios públicos (ver art. 27)

                  El pago de la matrícula se podrá hacer efectivo en uno, dos, cinco u ocho plazos. Los precios públicos por servicios administrativos y el seguro escolar se incluirán en el primer pago. No se podrá fraccionar el pago en los siguientes supuestos:

                  a) Cuando se trate del pago de la matrícula realizada en el plazo de ampliación de matrícula.
                  b) Las liquidaciones efectuadas por las secretarías de los centros con posterioridad al 15 de abril.
                  c) Cuando el importe de la matrícula, tras aplicar las posibles reducciones de precios, sea menor de 240 €.

                  Los estudiantes que deseen solicitar la expedición del título o una certificación acreditativa de haber finalizado sus estudios, pero tengan recibos impagados o pendientes de abonar, con carácter previo al abono del título, deberán satisfacer las cantidades pendientes de pago correspondientes a este curso o a cursos anteriores en la misma titulación, en su caso.

                  \textbf{Periodos de pago:}

                  - Primer plazo (o pago único): en 10 días hábiles desde la formalización de la matrícula. Para los estudiantes que se automatriculen en julio, la gestión de cobro se hará a principios de septiembre.
                  - Dos plazos: la gestión de cobro del segundo plazo se realizará durante la segunda quincena del mes de diciembre de 2025.

                  Código Seguro De Verificación: `eqOjhZcmL0p6VJlZiTH/Yw==`
                  Fecha: 27/06/2025
                  Firmado Por: MIGUEL ANGEL CASTRO ARROYO
                  Url De Verificación: [https://pfirma.us.es/verifirma/code/eqOjhZcmL0p6VJlZiTH\%2FYw\%3D\%3D](https://pfirma.us.es/verifirma/code/eqOjhZcmL0p6VJlZiTH\%2FYw\%3D\%3D)
                  Página: 36/50

                  ---

                  FUENTE: R.R. NORMAS DE MATRICULA GRADO Y MASTER UNIVERSITARIO 2025-2026.pdf | PÁGINA: 15
                  \# UNIVERSIDAD D SEVII

                  \#\# C) EST

                  3. Los estudiantes matriculados en un grado que deseen simultanearlo con otros estudios de grado previamente abandonados podrán solicitar la admisión o reanudación de estudios en estos. La reanudación de estudios exige tener superados/reconocidos un mínimo de 6 créditos y no haber agotado las normas de permanencia en los estudios abandonados. Se solicitará por escrito dirigido al Decano/Director del centro donde se impartan entre el 1 y el 15 de septiembre de 2025. De no cumplirse los referidos requisitos, solo sería posible obtener plaza por el procedimiento de admisión indicando simultaneidad de estudios en la solicitud, que será valorada al final del proceso, si existiesen vacantes.

                  4. En el caso de estudiantes que cursan estudios en otra universidad española y deseen simultanearlos con estudios que inicia en nuestra universidad, deberán solicitar la admisión haciendo constar la circunstancia de simultaneidad. Estas solicitudes serán valoradas al final del proceso, si existiesen vacantes.

                  5. En el supuesto de estudiantes que simultanean estudios en una o varias universidades españolas y desean continuar dichos estudios simultaneándolos en nuestra universidad, estos podrán solicitar la admisión indicando su deseo de simultanear. Estas solicitudes serán valoradas al final del proceso, si existiesen vacantes.

                  6. En cualquier caso, la realización de estudios simultáneamente no dará derecho al estudiante a adaptaciones curriculares, resultando de aplicación, en caso de coincidencia de exámenes entre ambas titulaciones, la Normativa Reguladora de la Evaluación y Calificación de las Asignaturas.

                  7. No se podrán simultanear estudios de grado con los estudios a extinguir que los sustituyan.

                  8. El ocultamiento de la concurrencia de simultaneidad dará lugar a la anulación de la plaza adjudicada y/o de la matrícula formalizada en su caso.

                  \#\#\# Artículo 14. Plazo de matrícula

                  El plazo de matrícula será el general que se determine para los estudiantes de nuevo ingreso en la fase correspondiente del procedimiento de admisión que determinen los Acuerdos de la Comisión de Distrito Único Universitario Andaluz, con las excepciones previstas en los supuestos de reanudación o cambio de estudios.

                  \#\# B) Estudiantes Visitantes

                  \#\#\# Artículo 15. Estudiantes visitantes

                  La matrícula de estudiantes visitantes es una modalidad de matrícula para estudiantes con estudios extranjeros de nivel universitario y para estudiantes con estudios universitarios españoles ya titulados, con objeto de ampliar su formación o realizar una estancia temporal. Esta modalidad de matrícula no se podrá simultanear con la matrícula ordinaria en la misma titulación. Este tipo de matrícula da derecho a la expedición de la correspondiente certificación académica.

                  \#\#\# Artículo 16. Requisitos y plazos

                  1. Los estudiantes podrán matricularse en cuantas asignaturas deseen de las que formen el plan de estudios de las titulaciones impartidas en el centro, siempre de acuerdo con las disponibilidades docentes y materiales de los diferentes Departamentos y con un límite máximo de 90 créditos. Únicamente será posible matricularse como estudiante visitante en asignaturas de primer curso de grado/máster universitario si hubieran quedado plazas vacantes en ellas tras la admisión de estudiantes para realizar estudios oficiales en la fase extraordinaria de cada curso académico.

                  Código Seguro De Verificación: eqOjhZcmL0p6VJlZiTH/Yw==
                  Fecha: 27/06/2025
                  Firmado Por: MIGUEL ANGEL CASTRO ARROYO
                  Url De Verificación: [https://pfirma.us.es/verifirma/code/eqOjhZcmL0p6VJlZiTH\%2FYw\%3D\%3D](https://pfirma.us.es/verifirma/code/eqOjhZcmL0p6VJlZiTH\%2FYw\%3D\%3D)
                  Página: 11/50

                  ---

                  FUENTE: Guia\_Movilidad\_Intern\_FAC\_TUR\_FINANZAS.pdf | PÁGINA: 6
                  \# Número total de créditos a incluir en el acuerdo de estudios:

                  Al elaborar el acuerdo de estudios ha de tenerse en cuenta que el número de créditos máximo que se debe incluir en el centro de origen (Facultad de Turismo y Finanzas de la Universidad de Sevilla) coincide con el número de créditos por curso en la titulación en la que el/la alumno/a esté matriculado/a. Por tanto:

                  \vspace{0.4cm}
                  \begin{table}[h!]\centering\renewcommand{\arraystretch}{1.3}
                    \begin{tabular}{|p{0.45\textwidth}|p{0.45\textwidth}|p{0.45\textwidth}|}\hline
                                                        & Curso completo & Cuatrimestre \\\hline
                      Grado de Turismo                  & 60 ECTS        & 30 ECTS      \\\hline
                      Grado de FICO                     & 60 ECTS        & 30 ECTS      \\\hline
                      Doble Grado en FICO y RRLL y RRHH & 72 ECTS        & 36 ECTS      \\\hline
                      Doble Grado de Derecho y FICO     & 84 ECTS        & 42 ECTS      \\\hline
                    \end{tabular}\end{table}
                  \vspace{0.4cm}

                  Cualquier exceso sobre dicho total será excepcional y por un motivo claramente justificado, debiéndose, en todo caso, ser autorizado expresamente por la Vicedecana de Movilidad e Internacionalización.

                  De forma excepcional y en el supuesto de que el/la estudiante con ello finalice sus estudios durante la estancia en la universidad de destino el número máximo de créditos podrá incrementarse hasta en un máximo de 20 ECTS.

                  ---

                  FUENTE: R.R. NORMAS DE MATRICULA GRADO Y MASTER UNIVERSITARIO 2025-2026.pdf | PÁGINA: 7
                  \# UNIVERSIDAD D SEVII

                  C)EST

                  Deberá tenerse en cuenta, además, que el derecho a examen y evaluación de las asignaturas matriculadas quedará limitado a las incompatibilidades académicas derivadas de los respectivos planes de estudios. La matriculación en asignaturas sujetas a régimen de incompatibilidad o de prerrequisito se encuentra condicionada, según las exigencias de los respectivos planes de estudios, a haber superado las asignaturas exigidas o a formalizar también, al mismo tiempo, la matriculación en estas.

                  5. En ningún caso, el ejercicio del derecho de matrícula obligará a la modificación del régimen de horarios generales determinados en cada centro.

                  6. Con carácter general, las asignaturas optativas, para ser impartidas, deben tener, al menos, seis estudiantes matriculados. De no alcanzarse esta cifra durante dos cursos consecutivos, será de aplicación lo dispuesto en el artículo 8.2.d) del Reglamento General de Actividades Docentes.

                  7. Una vez que el estudiante haya abonado los derechos de expedición del título, no podrá matricularse de más asignaturas optativas, quedando cerrado su expediente.

                  El estudiante que, encontrándose en la situación antes descrita, se hubiera matriculado de un número de créditos superior al necesario para finalizar sus estudios, no tendrá derecho a presentarse a posteriores convocatorias de examen ni a aparecer en actas de examen, y tampoco a la devolución de los precios públicos abonados por estos créditos, quedando cerrado su expediente a todos los efectos. Todo ello, sin perjuicio de la regulación referida a segundas menciones o especialidades.

                  \#\# Artículo 5. Condicionalidad de la matrícula

                  1. La matrícula estará condicionada al cumplimiento y comprobación de los requisitos establecidos, entendiéndose definitiva si con anterioridad al 31 de diciembre del año en curso, el estudiante no hubiese recibido notificación del requerimiento de subsanación de la misma, todo ello con las excepciones previstas en los apartados 2 y 3 de este artículo y sin perjuicio de las actuaciones que de oficio pudieran realizarse.

                  2. Se considerará la condicionalidad en los supuestos y hasta los momentos que se citan a continuación:

                  a) Cuando el estudiante haya solicitado reconocimiento de estudios previos cursados en la misma u otra universidad, hasta su resolución.

                  b) Cuando el estudiante haya solicitado reconocimiento parcial de estudios cursados en el extranjero o el estudiante esté pendiente de la homologación o declaración de equivalencia de su título universitario extranjero o de la homologación de sus estudios preuniversitarios por el bachillerato español, hasta su resolución.

                  c) Cuando la solicitud de matrícula conlleve la determinación del nivel académico de los estudios previos cursados en el extranjero, hasta su resolución.

                  d) Cuando el estudiante haya solicitado cambio de universidad y/o estudios, hasta tanto se reciba en la secretaría del centro la certificación académica oficial del centro de procedencia, excepto aquellos traslados que se realicen entre centros propios de la Universidad de Sevilla. En este último caso, por parte del centro de destino se realizará la comprobación telemática que lo comunicará al centro de origen para que se proceda al cierre del expediente, no debiendo el estudiante abonar los derechos de traslado.

                  e) En los casos de simultaneidad de estudios con otros ya iniciados, hasta tanto no se haya efectuado el traslado de expediente por simultaneidad y se reciba en la secretaría del centro, certificación académica oficial de simultaneidad del estudiante procedente del centro en el que ya.

                  ----

                  Código Seguro De Verificación: `eqOjhZcmL0p6VJlZiTH/Yw==`
                  Fecha: 27/06/2025
                  Firmado Por: MIGUEL ANGEL CASTRO ARROYO
                  Url De Verificación: [https://pfirma.us.es/verifirma/code/eqOjhZcmL0p6VJlZiTH\%2FYw\%3D\%3D](https://pfirma.us.es/verifirma/code/eqOjhZcmL0p6VJlZiTH\%2FYw\%3D\%3D)
                  Página: 3/50
                \end{scriptsize}
        \end{itemize}


  \item \textbf{P36.} Según el Reglamento y Normas de Matrícula, ¿cuáles son los motivos explícitos en los que la Universidad concederá el aplazamiento excepcional en el pago de tasas? \\
        \begin{itemize}
          \item \textbf{Tipo de consulta:} simple
          \item \textbf{Respuesta del Agente:}

                Según el Reglamento y Normas de Matrícula de la Universidad de Sevilla, no se especifican explícitamente los motivos para la concesión de un aplazamiento excepcional en el pago de tasas. Sin embargo, se establecen las condiciones generales para el pago de las matrículas y los supuestos en los que no se puede fraccionar el pago.

                En el artículo relacionado con el abono de precios públicos (ver art. 27), se menciona que el pago de la matrícula se puede hacer efectivo en uno, dos, cinco u ocho plazos, y que los precios públicos por servicios administrativos y el seguro escolar se incluirán en el primer pago. También se establecen supuestos en los que no se puede fraccionar el pago, como cuando se trate del pago de la matrícula realizada en el plazo de ampliación de matrícula, las liquidaciones efectuadas por las secretarías de los centros con posterioridad al 15 de abril, o cuando el importe de la matrícula sea menor de 240 euros.

                Además, se menciona que los estudiantes que deseen solicitar la expedición del título o una certificación acreditativa de haber finalizado sus estudios, pero tengan recibos impagados o pendientes de abonar, deberán satisfacer las cantidades pendientes de pago correspondientes a este curso o a cursos anteriores en la misma titulación, en su caso.

                En cuanto a la anulación de matrícula, se establecen los plazos y procedimientos para solicitar la anulación, y se especifican los efectos de la anulación, como la devolución de los importes abonados, la no computación de la matrícula a efectos de convocatorias de examen ni de permanencia del estudiante, y la pérdida de la plaza adjudicada si el estudiante accedió por el procedimiento de admisión en el curso 2025-2026.

                En resumen, aunque no se especifican explícitamente los motivos para la concesión de un aplazamiento excepcional en el pago de tasas, se establecen las condiciones generales para el pago de las matrículas y los supuestos en los que no se puede fraccionar el pago, así como los procedimientos y efectos de la anulación de matrícula. Es posible que la Universidad de Sevilla tenga procedimientos internos para considerar solicitudes de aplazamiento excepcional en el pago de tasas, pero estos no se especifican en el Reglamento y Normas de Matrícula proporcionado.
          \item \textbf{Contexto utilizado:}
                \begin{scriptsize}
                  FUENTE: R.R. NORMAS DE MATRICULA GRADO Y MASTER UNIVERSITARIO 2025-2026.pdf | PÁGINA: 40
                  \# UNIVERSIDAD D SEVII

                  \#\# C)EST

                  \#\#\# CONSIDERACIONES GENERALES

                  En caso de cualquier duda académica, administrativa o dificultades técnicas sobre el proceso de matrícula no referidas al acceso a la aplicación de automatrícula, el estudiante puede dirigirse al CAT mediante formulario electrónico https://cat.us.es/contacto.

                  Como ayuda para planificar la matrícula en el siguiente enlace se podrá obtener información sobre los diferentes planes de estudio: https://cat.us.es/ofertas-academicas.

                  Los estudiantes deberán matricularse de un mínimo de 30 créditos, con las siguientes excepciones:

                  - Estudiantes visitantes.
                  - Estudiantes a quienes se les reconozca, previa petición en su centro, la condición de estudiantes con necesidades académicas especiales (ver artículos 26 y 27 del Reglamento General de Estudiantes).
                  - Estudiantes que se matriculen de los últimos créditos para finalizar la titulación.

                  El número máximo de créditos en los que un estudiante puede estar matriculado en cada curso académico será de 90 créditos (excepto en las dobles titulaciones).

                  Los solicitantes de beca de convocatoria general del Ministerio de Educación, Formación Profesional y Deportes deberán matricular, al menos, 60 créditos para tener derecho a todas las modalidades de ayudas. (Los estudiantes de dobles titulaciones deberán consultar el número de créditos a matricular).

                  \#\#\# 3. Abono de precios públicos (ver art. 27)

                  El pago de la matrícula se podrá hacer efectivo en uno, dos, cinco u ocho plazos. Los precios públicos por servicios administrativos y el seguro escolar se incluirán en el primer pago. No se podrá fraccionar el pago en los siguientes supuestos:

                  a) Cuando se trate del pago de la matrícula realizada en el plazo de ampliación de matrícula.
                  b) Las liquidaciones efectuadas por las secretarías de los centros con posterioridad al 15 de abril.
                  c) Cuando el importe de la matrícula, tras aplicar las posibles reducciones de precios, sea menor de 240 €.

                  Los estudiantes que deseen solicitar la expedición del título o una certificación acreditativa de haber finalizado sus estudios, pero tengan recibos impagados o pendientes de abonar, con carácter previo al abono del título, deberán satisfacer las cantidades pendientes de pago correspondientes a este curso o a cursos anteriores en la misma titulación, en su caso.

                  \textbf{Periodos de pago:}

                  - Primer plazo (o pago único): en 10 días hábiles desde la formalización de la matrícula. Para los estudiantes que se automatriculen en julio, la gestión de cobro se hará a principios de septiembre.
                  - Dos plazos: la gestión de cobro del segundo plazo se realizará durante la segunda quincena del mes de diciembre de 2025.

                  Código Seguro De Verificación: `eqOjhZcmL0p6VJlZiTH/Yw==`
                  Fecha: 27/06/2025
                  Firmado Por: MIGUEL ANGEL CASTRO ARROYO
                  Url De Verificación: [https://pfirma.us.es/verifirma/code/eqOjhZcmL0p6VJlZiTH\%2FYw\%3D\%3D](https://pfirma.us.es/verifirma/code/eqOjhZcmL0p6VJlZiTH\%2FYw\%3D\%3D)
                  Página: 36/50

                  ---

                  FUENTE: R.R. NORMAS DE MATRICULA GRADO Y MASTER UNIVERSITARIO 2025-2026.pdf | PÁGINA: 8
                  \# UNIVERSIDAD D SEVII

                  C)EST

                  esté matriculado. Del mismo modo, se considerará condicional la matrícula de los estudiantes que deseen simultanear estudios desde su inicio hasta tanto se compruebe el cumplimiento de los requisitos de admisión.

                  f) Cuando el estudiante haya solicitado, dentro del plazo general de matrícula, salvo circunstancias sobrevenidas valoradas por el centro, la consideración de estudiante con necesidades académicas especiales en los términos previstos en el artículo 27 del Reglamento General de Estudiantes, hasta tanto recaiga la resolución correspondiente.

                  g) Cuando el estudiante no haya entregado toda la documentación necesaria según las instrucciones de matrícula, hasta la finalización del plazo de subsanación notificado por la secretaría del centro.

                  3. En tanto no se resuelva favorablemente el supuesto condicionante, no se admitirán a trámite solicitudes de título universitario oficial. No se emitirán certificaciones académicas oficiales en las que conste rendimiento académico, excepto en los supuestos contemplados en los apartados 2.a) y 2.f) de este artículo.

                  4. Las matrículas condicionadas, si el supuesto condicionante no se resuelve favorablemente, podrán ser anuladas o modificadas, según corresponda.

                  5. Los estudiantes de las universidades españolas a los que a la fecha de finalización del último plazo de admisión fijado por la Comisión de Distrito Único Universitario de Andalucía (Fase 3), solo les reste por superar la acreditación de la competencia lingüística en idioma extranjero requerida, el TFG y un máximo de 9 créditos ECTS para obtener el título de grado podrán solicitar la admisión condicionada en estudios de máster universitario en las universidades públicas andaluzas, con las siguientes consideraciones:

                  - La matrícula en el máster universitario se entenderá condicionada hasta la finalización de los estudios de grado, siendo definitiva cuando el estudiante acredite haber abonado los derechos de expedición del título de grado. En ningún caso se podrá obtener el título de máster universitario con anterioridad al título de grado que permite el acceso al máster.

                  - El estudiante estará obligado a matricular todos los créditos pendientes de los estudios de grado, incluido, en su caso, el TFG.

                  - El estudiantado de máster universitario con matrícula condicionada tendrá los mismos derechos y deberes que cualquier otro estudiante matriculado en un máster universitario del Distrito Único Universitario de Andalucía, si bien no se expedirán certificaciones académicas oficiales de estos estudios en tanto su matrícula sea condicionada.

                  - Solamente procederá la devolución del importe de la matrícula de los estudios de máster universitario en los supuestos específicamente regulados por estas normas de matrícula.

                  \#\# Artículo 6. Programas académicos con recorridos sucesivos (PARS)

                  El estudiante de grado vinculado a máster universitario en los Programas Académicos con Recorridos Sucesivos (PARS) en el ámbito de la ingeniería y arquitectura al que le reste por superar la acreditación del nivel de competencia lingüística requerido y un máximo de 30 créditos ECTS de la tipología que se establezca en cada uno de los Programas Académicos con Recorridos Sucesivos (PARS), podrá acceder y matricularse en el máster universitario vinculado, siempre que sea admitido en la Fase 3 del procedimiento de admisión.

                  No obstante, en ningún caso se podrá obtener el título de máster universitario con anterioridad al título de grado que permite el acceso al máster, ni se expedirán certificaciones académicas oficiales de los estudios de máster universitario.

                  Código Seguro De Verificación: `eqOjhZcmL0p6VJlZiTH/Yw==`
                  Fecha: `27/06/2025`
                  Firmado Por: `MIGUEL ANGEL CASTRO ARROYO`
                  Url De Verificación: [https://pfirma.us.es/verifirma/code/eqOjhZcmL0p6VJlZiTH\%2FYw\%3D\%3D](https://pfirma.us.es/verifirma/code/eqOjhZcmL0p6VJlZiTH\%2FYw\%3D\%3D)
                  Página: `4/50`

                  ---

                  FUENTE: R.R. NORMAS DE MATRICULA GRADO Y MASTER UNIVERSITARIO 2025-2026.pdf | PÁGINA: 14
                  \# UNIVERSIDAD D SEVII

                  \#\#\# C)EST

                  c) No se computará la matrícula realizada a efectos académicos.

                  d) En caso de que la sanción suponga la anulación total de la matrícula y el estudiante hubiera accedido por el procedimiento de admisión en el curso 2025-2026, perderá la plaza adjudicada.

                  11. En los supuestos de anulación por falta de pago o por falta de aportación de la documentación requerida, los efectos serán:

                  a) No procederá la devolución de los importes abonados en concepto de precios públicos por servicios académicos y administrativos y el seguro escolar, ni se eximirá de la obligación de satisfacer las cantidades pendientes de pago.

                  b) El abono de las cantidades pendientes de pago no podrá ser objeto de aplazamiento, fraccionamiento ni domiciliación bancaria.

                  c) Serán de aplicación los recargos económicos en la siguiente matrícula que realicen en los mismos estudios.

                  d) No se computará la matrícula realizada a efectos académicos y supondrá la pérdida de las calificaciones incorporadas en acta.

                  e) Si el estudiante accedió por el procedimiento de admisión en el curso 2025-2026, perderá la plaza adjudicada.

                  f) La pérdida de la condición de becario de la convocatoria general y, en consecuencia, la supresión de dicha exención en su matrícula.

                  \#\#\# Sección Tercera: Simultaneidad de estudios. Estudiantes visitantes

                  A) Simultaneidad de estudios

                  \#\#\#\# Artículo 12. De la matrícula para simultanear estudios

                  En materia de simultaneidad de estudios se estará a lo dispuesto en la regulación contenida en los Acuerdos de la Comisión de Distrito Único Universitario de Andalucía y en las demás normas de aplicación.

                  \#\#\#\# Artículo 13. Requisitos

                  1. Los estudiantes que deseen simultanear estudios desde su inicio deberán obtener plaza en ambas titulaciones en el procedimiento de admisión, para lo que deberán indicarlo expresamente en el mismo. Al matricularse, el estudiante deberá comunicar su situación de simultaneidad expresamente en los centros correspondientes. En el expediente académico que se abra en cada titulación se hará constar tal circunstancia.

                  2. Solo en los casos de simultaneidad entre estudios de grado ya iniciados en la Universidad de Sevilla con otros en los que haya obtenido plaza en otra universidad española, el estudiante deberá abonar los derechos de traslado de expediente a efectos de simultaneidad en el centro de origen y entregar dicho resguardo en el centro de destino. El incumplimiento de este requisito dará lugar a la anulación de la matrícula. Los centros implicados deberán dejar constancia del traslado por simultaneidad en el expediente del estudiante.

                  Código Seguro De Verificación: `eqOjhZcmL0p6VJlZiTH/Yw==`
                  Fecha: 27/06/2025
                  Firmado Por: MIGUEL ANGEL CASTRO ARROYO
                  Url De Verificación: [https://pfirma.us.es/verifirma/code/eqOjhZcmL0p6VJlZiTH\%2FYw\%3D\%3D](https://pfirma.us.es/verifirma/code/eqOjhZcmL0p6VJlZiTH\%2FYw\%3D\%3D)
                  Página: 10/50

                  ---

                  FUENTE: R.R. NORMAS DE MATRICULA GRADO Y MASTER UNIVERSITARIO 2025-2026.pdf | PÁGINA: 22
                  \# UNIVERSIDAD DE SEVILLA

                  C)EST

                  En pagos fraccionados posteriores al primero de los plazos, la forma de pago será la domiciliación de los recibos.

                  3. La domiciliación de los recibos se realizará en una cuenta bancaria de la que el estudiante sea titular o cotitular. No obstante, aquellos estudiantes que en cursos anteriores aportaron la Orden de Domiciliación de Adeudo Directo SEPA y no modifiquen la cuenta bancaria a efectos de domiciliación, no estarán obligados a cumplimentarla de nuevo mientras esta conserve su vigencia, con independencia de la titularidad de la cuenta.

                  A los estudiantes que no hubieran cumplimentado y entregado en años anteriores la Orden de Domiciliación de Adeudo Directo SEPA (o en caso de modificación de la cuenta bancaria), y elijan el pago fraccionado o la domiciliación como forma de pago en la modalidad de pago único, la citada Orden se les generará automáticamente en la aplicación de automatrícula al realizar su matrícula (excepto en los casos de matrícula presencial).

                  4. Los estudiantes que, al formalizar su matrícula, no hubiesen abonado los precios públicos por servicios académicos en primera matrícula por haber solicitado la concesión de beca con cargo a los Presupuestos Generales del Estado y posteriormente no obtuviesen la condición de becarios o les fuera revocada la beca concedida, vendrán obligados al abono del precio correspondiente a la matrícula que efectuaron, en los términos y plazos que regule la Universidad de Sevilla y su impago supondrá la anulación de la matrícula y, en su caso, de la ampliación de matrícula para todas las materias o asignaturas en los términos previstos por la legislación vigente.

                  5. No obstante lo dispuesto en los apartados precedentes, no se podrá fraccionar el pago en los siguientes supuestos:

                  a) Cuando se trate del pago de la matrícula realizada en el plazo de ampliación de matrícula.

                  b) Las liquidaciones efectuadas por las secretarías de los centros con posterioridad al 15 de abril.

                  c) Cuando el importe de la matrícula, después de aplicar las posibles reducciones de precios, sea menor de 240 euros.

                  6. El impago de dichos precios en cualquiera de los plazos correspondientes a la modalidad elegida (pago único o pago fraccionado) supondrá el inicio del procedimiento de anulación de matrícula. Las notificaciones de requerimiento de pago y anulación de matrícula por esta causa se realizarán a través de la sede electrónica de la Universidad de Sevilla. La finalización de dicho procedimiento sin que se haya producido el pago dará lugar a la anulación de la matrícula, sin derecho a reintegro, una vez notificada la resolución de esta o, en su caso, cuando se produzca el rechazo automático de la misma por el transcurso del plazo establecido. La anulación de la matrícula en periodo ordinario determinará, en su caso, la anulación de la ampliación.

                  7. Los estudiantes que se matriculen en el Instituto de Idiomas en el presente curso académico aplicándose descuento como estudiante de la Universidad de Sevilla y que anulen posteriormente la matrícula que les da derecho a las reducciones del precio, deberán hacer efectivo la parte correspondiente al importe dejado de abonar en el citado centro.

                  8. Cuando el estudiante solicite la expedición del título o una certificación académica, la Universidad de Sevilla comprobará que no tiene recibos impagados o pendientes de abono. El estudiante deberá satisfacer las cantidades pendientes del curso actual o de cursos anteriores en la misma titulación, en su caso. Para ello, se le facilitará/n el/los recibo/s pendiente/s de abono a efectos de su liquidación mediante pago por recibo bancario. En caso de estar domiciliado/s, se modificará la forma de pago a recibo bancario.

                  ---

                  FUENTE: R.R. NORMAS DE MATRICULA GRADO Y MASTER UNIVERSITARIO 2025-2026.pdf | PÁGINA: 11
                  \# UNIVERSIDAD D SEVII

                  C)EST

                  5. No se concederán anulaciones parciales de matrícula a solicitud del estudiante, excepto por coincidencia no prevista de horarios de clase y exámenes y las situaciones contempladas en estas normas. No obstante, por causas académicas sobrevenidas que imposibiliten la prestación del servicio, los centros podrán realizar de oficio anulaciones parciales de matrícula, ofreciendo a los estudiantes la posibilidad de matricular otras asignaturas o bien la devolución de los importes abonados en los créditos anulados, en su caso.

                  6. La matrícula se anulará de oficio cuando se trate de error, omisión, falseamiento de los datos consignados en esta o cuando se modifique la situación administrativa o académica del estudiante, o las matrículas condicionadas que no deban producir efectos.

                  Asimismo se anularán de oficio las matrículas cuando el estudiante no aporte la documentación requerida o por falta de pago.

                  7. Los centros anularán de oficio las matrículas efectuadas por estudiantes que se matriculen en otra titulación por haber obtenido plaza en los procedimientos de admisión para estudiantes de nuevo ingreso en esta u otra universidad pública, siempre que no se acredite una autorización para la simultaneidad de estudios.

                  \#\# Artículo 11. Supuestos de anulación de matrícula y sus efectos

                  1. Se podrá solicitar la anulación de matrícula ordinaria al inicio del curso académico en los siguientes plazos:

                  - Titulaciones de grado: hasta el 16 de septiembre de 2025.
                  - Titulaciones de máster universitario que inicien la docencia el 6 de octubre: hasta el 15 de octubre de 2025.
                  - Titulaciones de máster universitario que inicien la docencia el 22 de septiembre: hasta el 30 de septiembre de 2025.
                  - Cuando la matrícula se haya formalizado con posterioridad al inicio del periodo docente del primer cuatrimestre: 7 días hábiles a contar desde el día de la finalización del periodo en el que haya realizado su matrícula.

                  La concesión de la anulación supondrá:

                  a) La devolución de los importes abonados, en su caso.

                  b) La matrícula no se computará a efectos de convocatorias de examen ni de permanencia del estudiante, ni se contabilizará para futuras matrículas a efectos de aplicación de recargos.

                  c) La pérdida de la plaza adjudicada si el estudiante accedió por el procedimiento de admisión en el curso 2025-2026.

                  2. La anulación de la matrícula por haber obtenido plaza en esta u otra universidad pública (en el procedimiento o sistema de admisión que tenga establecido) en otra titulación de grado, máster universitario o estudios de doctorado donde hubiese sido admitido el estudiante con posterioridad a su matriculación en la Universidad de Sevilla, no tendrá ningún efecto académico ni económico si se solicita con anterioridad al 31 de octubre de 2025, en el caso de titulaciones de grado, y al 31 de diciembre de 2025 en las titulaciones de máster universitario.

                  Código Seguro De Verificación: `eqOjhZcmL0p6VJlZiTH/Yw==`
                  Fecha: 27/06/2025
                  Firmado Por: MIGUEL ANGEL CASTRO ARROYO
                  Url De Verificación: [https://pfirma.us.es/verifirma/code/eqOjhZcmL0p6VJlZiTH\%2FYw\%3D\%3D](https://pfirma.us.es/verifirma/code/eqOjhZcmL0p6VJlZiTH\%2FYw\%3D\%3D)
                  Página: 7/50
                \end{scriptsize}
        \end{itemize}


  \item \textbf{P37.} Según las instrucciones del BOJA (Resolución del Rector) para plazas de docencia, ¿cuáles son los requisitos esenciales como la edad y titulación que debe poseer cualquier solicitante español? \\
        \begin{itemize}
          \item \textbf{Tipo de consulta:} simple
          \item \textbf{Respuesta del Agente:}

                Según las instrucciones del BOJA (Resolución del Rector) para plazas de docencia, los requisitos esenciales que debe poseer cualquier solicitante español son los siguientes:

                1. \textbf{Titulación}: Debe poseer una titulación universitaria de Grado o equivalente, como Licenciado/a, Diplomado/a, Arquitecto/a, Arquitecto/a técnico/a, Ingeniero/a o Ingeniero/a técnico/a. Además, se valorará la posesión de un Título de Doctor (tesis doctoral) y la nota media de la titulación universitaria.
                2. \textbf{Edad}: No se establece una edad máxima para ser solicitante, pero se debe tener en cuenta que la edad de jubilación en España es de 65 años, y que la Universidad de Sevilla puede establecer requisitos específicos para cada plaza.
                3. \textbf{Experiencia docente}: Se valorará la experiencia docente en asignaturas regladas universitarias, con un máximo de 1.200 horas de docencia oficial.
                4. \textbf{Otros requisitos}: Se pueden establecer otros requisitos específicos para cada plaza, como la posesión de un determinado título o experiencia en un área concreta.

                Es importante destacar que los requisitos pueden variar dependiendo de la plaza y la convocatoria específica. Por lo tanto, es recomendable consultar la convocatoria y la normativa específica para cada plaza para obtener información precisa y actualizada.

                En cuanto a la edad, no se establece un límite máximo en la normativa proporcionada, pero es importante tener en cuenta que la Ley Orgánica 2/2023, de 22 de marzo, del Sistema Universitario, establece que la edad de jubilación para el personal docente e investigador es de 65 años. Sin embargo, esta edad puede variar dependiendo de la legislación y la normativa específica de cada universidad.

                En resumen, los requisitos esenciales para un solicitante español son la posesión de una titulación universitaria adecuada, experiencia docente y otros requisitos específicos que pueden variar dependiendo de la plaza y la convocatoria. Es importante consultar la normativa específica para cada plaza para obtener información precisa y actualizada.
          \item \textbf{Contexto utilizado:}
                \begin{scriptsize}
                  FUENTE: BOJA25-508-00085-6229-01\_00319829.pdf | PÁGINA: 55
                  \# BOJA
                  \#\# Boletín Oficial de la Junta de Andalucía
                  \#\#\# Extraordinario núm. 8 - Viernes, 2 de mayo de 2025
                  \#\#\#\# página 6229/55

                  Los baremos específicos que adopten las comisiones de selección. Esta normativa garantizará la igualdad de oportunidades y el respeto a los principios de mérito y capacidad.

                  \#\#\# Artículo 115. Profesoras y profesores visitantes.
                  1. El Consejo de Gobierno, a iniciativa del Rector o Rectora o, en su caso, de los Consejos de Departamento, podrá proponer la contratación de profesorado visitante ordinario de entre docentes e investigadoras e investigadores de otras universidades y centros de investigación, tanto españoles como extranjeros, para desempeñar actividades docentes o investigadoras en los términos que disponga la normativa de la Comunidad Autónoma de Andalucía.
                  2. Excepcionalmente, el Consejo de Gobierno podrá proponer la contratación de profesorado doctor en la modalidad extraordinaria de acuerdo con lo establecido en la legislación de la Comunidad Autónoma de Andalucía.
                  3. El Consejo de Gobierno aprobará los requisitos y el procedimiento de contratación de este profesorado.

                  \#\#\# Artículo 116. Profesores y profesoras distinguidos.
                  1. El Consejo de Gobierno, a iniciativa del Rector o Rectora, podrá proponer la contratación de profesorado distinguido de entre docentes e investigadores e investigadoras, tanto españoles como extranjeros, que desarrollen su carrera en el extranjero y cuya excelencia y contribución científica, tecnológica, humanística o artística sean significativas y reconocidas internacionalmente.
                  2. El profesorado distinguido podrá desempeñar, en los términos del contrato, actividades docentes, investigadoras, de transferencia e intercambio de conocimiento, de innovación o de dirección de grupos, centros de investigación y programas científicos y tecnológicos singulares.

                  \#\#\# Sección 4.ª Plazas vinculadas a servicios asistenciales y de salud pública de instituciones sanitarias

                  \#\#\# Artículo 117. Personal de los cuerpos docentes universitarios con plaza vinculada a servicios asistenciales y de salud pública del Sistema Sanitario Público de Andalucía y acceso a estas plazas.
                  1. El personal de los cuerpos docentes universitarios que ocupe una plaza vinculada a servicios asistenciales y de salud pública del Sistema Sanitario Público de Andalucía, en áreas de conocimiento de carácter clínico asistencial y de salud pública, se regirá por lo establecido en la legislación sanitaria y universitaria y por la normativa reglamentaria que la desarrolle. Esta plaza se considerará, a todos los efectos, como un solo puesto de trabajo.
                  2. Las plazas vinculadas se proveerán por concurso en los términos establecidos en la legislación vigente que le sea aplicable. Corresponde al Sistema Sanitario Público de Andalucía designar el número de miembros de la comisión que proceda, a tenor de lo que determinen la legislación aplicable y el correspondiente concierto.
                  3. Además de reunir los requisitos exigidos en la legislación universitaria que le resulte de aplicación, el profesorado de los cuerpos docentes universitarios vinculado tendrá que cumplir las exigencias en cuanto a su cualificación determinadas por la normativa sanitaria. En el caso de las personas que posean una titulación que habilite para el ejercicio de la profesión médica, tendrán que acreditar estar en posesión del título de Especialista en Ciencias de la Salud.
                  4. Estas comisiones deberán valorar la actividad asistencial de los candidatos y candidatas de la forma que se encuentre determinada en la normativa aplicable.

                  Depósito Legal: SE-410/1979. ISSN: 2253-802X https://www.juntadeandalucia.es/eboja

                  00319829

                  ---

                  FUENTE: BOJA25-508-00085-6229-01\_00319829.pdf | PÁGINA: 52
                  \# BOJA
                  \#\# Boletín Oficial de la Junta de Andalucía
                  \#\#\# Extraordinario núm. 8 - Viernes, 2 de mayo de 2025
                  \#\#\#\# página 6229/52

                  \#\#\# Artículo 107. Reclamaciones.
                  1. Contra las propuestas de las comisiones de los concursos de acceso a los cuerpos docentes se podrán presentar reclamaciones ante la Comisión de Reclamaciones prevista en el artículo 146 de los presentes Estatutos.
                  2. La Comisión, oídos los miembros de la comisión del concurso, así como los candidatos o candidatas que hubieran participado en el mismo, examinará el expediente relativo al concurso para velar por las garantías establecidas y ratificará o no la propuesta reclamada en el plazo máximo de tres meses, tras lo que el Rector o Rectora dictará la resolución en congruencia con lo que indique la Comisión. Transcurrido este plazo sin que se haya dictado y notificado resolución, la reclamación se entenderá desestimada.

                  \#\#\# Sección 3.ª El personal docente e investigador laboral

                  \#\#\# Artículo 108. Clases de personal docente e investigador laboral.
                  1. La Universidad de Sevilla podrá contratar personal docente e investigador laboral a través de las modalidades específicas del ámbito universitario: profesorado ayudante doctor, profesorado permanente laboral en sus distintas modalidades, profesorado asociado, profesorado visitante, profesorado distinguido y profesorado sustituto de acuerdo con la Ley estatal y autonómica en la materia.
                  2. El personal docente e investigador laboral se regirá por la Ley Orgánica 2/2023, de 22 de marzo, del Sistema Universitario, por la Ley de la Comunidad Autónoma de Andalucía y sus normas de desarrollo, así como por los presentes Estatutos y sus normas de desarrollo, la legislación laboral y de empleo público en lo que proceda y los convenios colectivos que le sean de aplicación.

                  \#\#\# Artículo 109. Concursos de acceso a plazas de profesorado permanente laboral.
                  1. Los concursos de acceso serán convocados por la Universidad. La convocatoria será publicada en los Boletines Oficiales correspondientes y comunicada al registro público de concurso del personal docente e investigador del Ministerio con competencia en Universidades.
                  2. Los concursos de acceso se regirán por las normas que apruebe el Consejo de Gobierno garantizando la igualdad de oportunidades y el respeto a los principios de mérito y capacidad. En todo caso, la experiencia docente y la experiencia investigadora, incluyendo la de transferencia e intercambio de conocimiento, tendrán una consideración análoga en el conjunto de criterios de valoración de los méritos a considerar.
                  3. Los criterios específicos de valoración para la adjudicación de las plazas serán establecidos y publicados por la comisión correspondiente con la debida antelación. Dichos criterios deberán referirse, en todo caso, a la experiencia docente y la experiencia investigadora, incluyendo la de transferencia e intercambio de conocimiento, del candidato o candidata, su proyecto docente e investigador, así como permitir contrastar sus capacidades para la exposición y debate en la correspondiente materia o especialidad en sesión pública.
                  4. Los candidatos o candidatas habrán de presentar un programa de, al menos, una de las asignaturas obligatorias del área de conocimiento al que se adscriba la plaza y un breve proyecto en el que formulen sus planteamientos docentes e investigadores.
                  5. El concurso de acceso a plazas de profesorado permanente laboral constará de dos pruebas. La primera de ellas consistirá en la preparación, exposición oral y pública por parte del candidato o candidata, y debate posterior con la comisión, de una lección del programa o los programas que haya presentado al concurso, elegida por el candidato o candidata de entre tres determinadas por sorteo. Esta prueba tendrá carácter eliminatorio para los candidatos y candidatas que no obtengan al menos tres votos favorables de los miembros de la comisión.

                  Depósito Legal: SE-410/1979. ISSN: 2253-802X https://www.juntadeandalucia.es/eboja

                  00319829

                  ---

                  FUENTE: BOJA25-508-00085-6229-01\_00319829.pdf | PÁGINA: 48
                  \# BOJA
                  \#\# Boletín Oficial de la Junta de Andalucía
                  \#\#\# Extraordinario núm. 8 - Viernes, 2 de mayo de 2025
                  página 6229/48

                  \#\#\# Artículo 96. Deberes del personal docente e investigador.
                  Son deberes del personal docente e investigador los siguientes:
                  a) Cumplir con sus obligaciones docentes e investigadoras.
                  b) Actualizar y renovar sus conocimientos científicos y su metodología docente.
                  c) Evaluar y calificar de manera objetiva y justa el nivel de conocimiento y aprendizaje del estudiantado.
                  d) Participar en la evaluación de la calidad de su actividad docente.
                  e) Cualesquiera otros que establezcan la legislación vigente, los presentes Estatutos y sus normas de desarrollo.

                  \#\#\# Artículo 97. Plena capacidad docente.
                  La Universidad de Sevilla reconoce la plena capacidad docente al profesorado perteneciente a los cuerpos docentes universitarios, al profesorado permanente laboral, a las otras figuras de personal docente e investigador laboral que aparezca en la legislación Estatal y Autonómica, así como al personal investigador posdoctoral en cuyo contrato se contemple la posibilidad de impartir docencia.

                  \#\#\# Artículo 98. Formación del personal docente e investigador.
                  1. El personal docente e investigador de la Universidad de Sevilla, con arreglo a la Normativa que apruebe el Consejo de Gobierno, podrá mejorar, completar o ampliar su formación científica y pedagógica en otras universidades o instituciones científicas y académicas sin pérdida del puesto de trabajo y disfrutar, en su caso, de las licencias que puedan corresponder.
                  2. La Universidad de Sevilla dispondrá de una Estrategia de desarrollo de la formación docente de su profesorado, aprobada por el Consejo de Gobierno, que contenga actuaciones en materia de formación del profesorado, de innovación y de evaluación docente. Esta Estrategia contribuirá, entre otros aspectos, a promover el aprendizaje activo del estudiantado y la evaluación continua de manera coordinada con el Plan Propio de Docencia.
                  3. La Universidad de Sevilla garantizará que el personal docente e investigador doctor de nuevo ingreso reciba durante sus dos primeros años una formación docente inicial.

                  \#\#\# Artículo 99. Relación de puestos de trabajo y carrera profesional.
                  1. Corresponde al Consejo de Gobierno, a propuesta del Rector o Rectora, definir y planificar la política de personal docente e investigador.
                  2. La Universidad de Sevilla elaborará anualmente la relación de puestos de trabajo de su personal docente e investigador, determinando el número de plazas que corresponde a cada categoría. Sin perjuicio de lo que pueda establecer la legislación general o autonómica, en la relación de puestos de trabajo se incluirán como mínimo la denominación de cada plaza, el área de conocimiento y el departamento al que se adscriba, los requisitos para el desempeño de cada puesto y sus retribuciones básicas y complementarias.
                  3. Todos los puestos de trabajo de profesorado funcionario y laboral deberán adscribirse a las áreas de conocimiento establecidas por la normativa de aplicación, que deberán permitir y favorecer la movilidad del profesorado y facilitar su carrera profesional.
                  4. La relación de puestos de trabajo del personal docente e investigador podrá modificarse por ampliación o transformación de las plazas existentes, o por minoración o cambio de denominación de las plazas vacantes, garantizándose en los dos primeros supuestos la adecuada cobertura presupuestaria.
                  5. A propuesta del Rector o Rectora, el Consejo de Gobierno aprobará en cada ejercicio la relación de puestos de trabajo del personal docente e investigador de la Universidad, con expresión de las modificaciones que se introduzcan en la del ejercicio.

                  Depósito Legal: SE-410/1979. ISSN: 2253-802X
                  https://www.juntadeandalucia.es/eboja

                  ---

                  FUENTE: BAREMO 2024.pdf | PÁGINA: 15
                  \# Núm. 12/2024, de 11 de octubre

                  \#\# I. DISPOSICIONES GENERALES

                  \#\#\# I.5. Secretario General

                  1.3. Otros títulos académicos:                                                            \textbf{Máximo: 10 puntos}
                  - 1.3.1. Titulación Universitaria de Grado o equivalente (como Licenciado/a, Diplomado/a, Arquitecto/a, Arquitecto/a técnico/a, Ingeniero/a o Ingeniero/a técnico/a), y Titulación Universitaria de Máster (Oficial):
                  - APROBADO: 5,0-6,9 nota media                            \textbf{3 puntos}
                  - NOTABLE: 7,0-8,9 nota media                             \textbf{6 puntos}
                  - SOBRESALIENTE: 9,0-10,0 nota media                      \textbf{10 puntos}
                  - 1.3.2. Título de doctor (tesis doctoral)                  \textbf{10 puntos}

                  \textbf{Nota 1:} En relación con el coeficiente de afinidad y la homologación de títulos se aplicarán las Notas del apartado 1.1.

                  1.4. Premios académicos y menciones de calidad:                                           \textbf{Máximo: 5 puntos}
                  - 1.4.1. Premio (Nacional) Fin de Carrera de Educación Universitaria: Primer premio           \textbf{5 puntos por premio}           No admite coeficiente de afinidad.
                  - 1.4.2. Premio extraordinario de grado o equivalente        \textbf{4 puntos por premio}
                  - 1.4.3. Premio Extraordinario de Doctorado                  \textbf{4 puntos por premio}
                  - 1.4.4. Doctorado internacional y europeo                   \textbf{2 puntos por indicio}
                  - 1.4.5. Mención de calidad (o ‘hacia la excelencia’)        \textbf{1 punto por indicio}
                  - 1.4.6. Tesis en régimen de cotutela                        \textbf{1 punto por indicio}

                  \#\# II. ACTIVIDAD DOCENTE                                                                     \textbf{Máximo: 100 puntos}

                  2.1. Docencia en asignaturas regladas universitarias:                                     \textbf{Máximo: 50 puntos}
                  (máximo 1.200 horas de docencia oficial)
                  - 2.1.1. Titulación Universitaria de Grado o equivalente (como Licenciado/a, Diplomado/a, Arquitecto/a, Arquitecto/a técnico/a, Ingeniero/a o Ingeniero/a técnico/a), y Titulación Universitaria de Máster (Oficial):
                  \textbf{1,00 x 10 horas de docencia}
                  En el caso de que el aspirante acredite docencia superior a 1.200 horas, su puntuación se escalará proporcionalmente a 1.200 horas.

                  \textbf{Complemento 1 (multiplicativo):}                             \textbf{1,50 x puntuación obtenida}
                  Docencia en Máster Universitario (oficial)                  en 2.1.1 en asignaturas de Máster Universitario (oficial)

                  \textbf{Complemento 2 (multiplicando al complemento 1):}
                  Docencia en inglés u otras lenguas diferentes a las que son objeto de estudio                                      \textbf{1,50 x puntuación obtenida}
                  en 2.1.1 en asignaturas en inglés u otras lenguas diferentes a las que son objeto de estudio

                  \textbf{Nota 1:} Los Trabajos Fin de Carrera (Grado o equivalente, y Máster) no serán evaluados en este apartado. Se remite al aspirante al apartado 2.4.1.
                  \textbf{Nota 2:} Se considera la siguiente tipología de asignaturas:
                  - Afinidad muy alta (1,0): asignaturas cuyo contenido corresponda al perfil de la plaza convocada, o en su defecto las asignaturas impartidas por el área de conocimiento del departamento solicitante de la plaza.
                  - Afinidad alta  (0,75): asignaturas impartidas por el área de conocimiento solicitante sin consideración de universidad de impartición, e incluyendo con ello las asignaturas no necesariamente impartidas por el área de conocimiento del departamento de la Universidad de Sevilla convocante de la plaza.
                  - Afinidad media  (0,50): asignaturas impartidas por áreas de conocimiento afines al departamento y área de conocimiento a que corresponde la plaza según listado oficial de la Universidad de Sevilla de comisiones juzgadoras a cuerpos docentes, y sin consideración de universidad de impartición.
                  - Afinidad baja o muy baja      (0,25 o 0,0): asignaturas no pertenecientes a ninguno de los tipos anteriores. El coeficiente de afinidad aplicable deberá ser 0,0 o 0,25, motivándose por la comisión de contratación.

                  ---

                  FUENTE: BOJA25-508-00085-6229-01\_00319829.pdf | PÁGINA: 54
                  \# BOJA
                  \#\# Boletín Oficial de la Junta de Andalucía
                  \#\#\# Extraordinario núm. 8 - Viernes, 2 de mayo de 2025
                  página 6229/54

                  6. El Consejo de Gobierno regulará el procedimiento de designación y el funcionamiento de las comisiones de los concursos de acceso, garantizando los principios de imparcialidad y profesionalidad de sus miembros y procurando una composición equilibrada entre mujeres y hombres.

                  \#\# Artículo 112. Procedimiento de selección de profesoras y profesores asociados.
                  1. La Universidad de Sevilla podrá contratar bajo la modalidad de profesorado asociado a especialistas y profesionales de reconocida competencia que acrediten ejercer su actividad principal fuera del ámbito académico universitario cuando existan necesidades docentes específicas relacionadas con su ámbito profesional. La finalidad del contrato será desarrollar tareas docentes a través de las que aporten sus conocimientos y experiencia profesionales a la universidad, en aquellas materias en las que esta experiencia resulte relevante.
                  2. Los concursos públicos para la selección de profesoras y profesores asociados serán resueltos por comisiones de contratación presididas por el Rector o Rectora, o persona en quien delegue, e integradas, además, por cinco profesores o profesoras del departamento correspondiente, elegidos por sorteo. Las comisiones serán designadas por períodos de cuatro años.
                  3. La comisión valorará los méritos de los concursantes con arreglo al baremo general establecido por el Consejo de Gobierno y al baremo específico establecido por la comisión, que será aplicable durante toda la vigencia de esta.
                  4. Contra la resolución de las comisiones cabrá recurso de alzada ante el Rector o Rectora.
                  5. Una normativa del Consejo de Gobierno regulará el procedimiento de contratación y podrá contemplar un baremo general, además de los criterios a los que deban ajustarse los baremos específicos que adopten las comisiones de contratación. Esta normativa garantizará la igualdad de oportunidades y el respeto a los principios de mérito y capacidad.

                  \#\# Artículo 113. Profesoras y profesores eméritos.
                  1. El Consejo de Gobierno podrá proponer el nombramiento de profesores y profesoras eméritos entre el personal docente e investigador funcionario o laboral jubilado que haya prestado servicios destacados en el ámbito docente, de investigación o de transferencia de conocimiento e innovación en la Universidad de Sevilla.
                  2. Los requisitos para el acceso a esta modalidad de profesorado y su régimen de dedicación serán establecidos por el Consejo de Gobierno, de acuerdo con la normativa de la Comunidad Autónoma de Andalucía.

                  \#\# Artículo 114. Profesoras y profesores sustitutos.
                  1. En los supuestos admitidos por la legislación universitaria estatal y autonómica y la legislación laboral, la Universidad de Sevilla podrá contratar profesoras y profesores sustitutos.
                  2. El contrato de profesorado sustituto comprenderá la actividad docente lectiva y no lectiva, que no podrá superar la asignada a la profesora o profesor a quien que se sustituya. También podrá extenderse a actividades universitarias de otra naturaleza, como las de investigación o el desempeño de funciones estructurales de gestión y coordinación, siempre que tengan directa relación con la actividad docente y resulten compatibles con el encargo docente.
                  3. La selección del profesorado sustituto se producirá mediante concurso público. Se podrán establecer instrumentos específicos para su gestión y cobertura, incluidas las bolsas de empleo.
                  4. Una normativa del Consejo de Gobierno regulará el procedimiento de concurso y podrá contemplar un baremo general, además de los criterios a los que deban ajustarse.

                  Depósito Legal: SE-410/1979. ISSN: 2253-802X https://www.juntadeandalucia.es/eboja

                  00319829
                \end{scriptsize}
        \end{itemize}


  \item \textbf{P38.} En la Normativa de Matrícula y Precios, ¿tienen derecho a una bonificación de matrícula quienes superen créditos en primera matrícula durante el curso anterior? \\
        \begin{itemize}
          \item \textbf{Tipo de consulta:} simple
          \item \textbf{Respuesta del Agente:}

                Según la Normativa de Matrícula y Precios de la Universidad de Sevilla (US) para el curso 2025-2026, específicamente en la sección 10.5 "Cálculo de la bonificación", se establece que los estudiantes que superen créditos en primera matrícula durante el curso anterior tienen derecho a una bonificación de matrícula.

                Concretamente, el importe de la bonificación será el 99\% del que resulte de multiplicar el número de créditos aprobados en primera matrícula en el curso 2024-2025, o en los cursos 2023-2024 y 2024-2025 para la primera matriculación en estudios de máster universitario, por el precio determinado en el párrafo anterior.

                Es importante destacar que, según la normativa, la bonificación se aplicará solo a los créditos que sean necesarios para la obtención del título conforme al plan de estudios y, en su caso, las especialidades y menciones vinculadas al título donde proceda la bonificación.

                Además, se establece que la bonificación no se aplicará en los supuestos de reconocimiento de créditos, salvo en el caso de estudios realizados en Programas de Movilidad Estudiantil, ya sea en el marco de la Unión Europea, ya sea realizados mediante convenios específicos entre universidades y conforme al correspondiente acuerdo académico de movilidad firmado por los estudiantes.

                En resumen, según la Normativa de Matrícula y Precios de la Universidad de Sevilla, los estudiantes que superen créditos en primera matrícula durante el curso anterior tienen derecho a una bonificación de matrícula, siempre y cuando cumplan con los requisitos establecidos en la normativa.
          \item \textbf{Contexto utilizado:}
                \begin{scriptsize}
                  FUENTE: R.R. NORMAS DE MATRICULA GRADO Y MASTER UNIVERSITARIO 2025-2026.pdf | PÁGINA: 29
                  \# UNIVERSIDAD D SEVII

                  \#\# C)EST

                  Artísticas superiores en centros públicos de la Comunidad Autónoma de Andalucía, durante los cursos 2023-2024 y/o 2024-2025 en los estudios que hayan dado acceso al máster, tomándose en consideración a estos efectos los créditos correspondientes a asignaturas superadas en primera matrícula en dichos estudios siempre que no hubieran ya generado derecho a esta bonificación.

                  c) En la segunda o sucesiva matriculación de estudios de máster universitario: Haber estado matriculado en centros propios o adscritos de las universidades públicas andaluzas durante el curso 2024-2025, tomándose en consideración los créditos correspondientes a asignaturas superadas en primera matrícula en dicho curso.

                  Asimismo, se aplicará a los estudiantes de los centros propios y adscritos de las universidades públicas andaluzas que, en los cursos 2023-2024 y/o 2024-2025 y en las condiciones señaladas en los apartados a), b) y c) anteriores, hubiesen estado inscritos en centros asociados andaluces de la Universidad Nacional de Educación a Distancia y empadronados en Andalucía, debiendo aportar certificado de empadronamiento en Andalucía, certificación académica personal y matrícula, en su caso, de los dos últimos cursos académicos.

                  \#\#\# 10.2. Límites de la bonificación:

                  Solo serán objeto de bonificación los créditos que sean necesarios para la obtención del título conforme al plan de estudios y, en su caso, las especialidades y menciones vinculadas al título donde proceda la bonificación.

                  Los estudiantes que en este curso académico se han matriculado de un grado y de un máster universitario en virtud de lo dispuesto en los artículos 5.5 y 6 de estas normas, tienen derecho a la aplicación de la bonificación del 99\% de su matrícula en ambas titulaciones, con las siguientes aclaraciones:

                  - En la titulación de grado, el estudiante tendrá derecho a la aplicación de una bonificación equivalente al importe del 99\% de los créditos superados en primera matrícula en el curso 2024-2025.

                  - En la titulación de máster universitario, el estudiante tendrá derecho a la aplicación de una bonificación equivalente al importe del 99\% de los créditos superados en primera matrícula en los cursos 2023-2024 y 2024-2025 en los estudios de grado que hayan dado acceso al título de máster universitario.

                  - Para calcular la bonificación correspondiente al curso 2025-2026, habrá de determinarse el número de créditos que generen el derecho a la bonificación en la titulación de máster universitario, esto es, el número de créditos superados por el estudiante en primera matrícula en los dos cursos anteriores en la titulación de grado que le da acceso al máster. A la cantidad resultante, se le deberá detraer el número de créditos ya bonificados en la titulación de grado tanto en el curso anterior como en el actual curso académico. La diferencia resultante podrá aplicarse al máster universitario.

                  Quienes hayan sido beneficiarios de una bonificación total o parcial en dos títulos de máster universitario, no podrán recibir bonificación en un tercero o sucesivos.

                  \#\#\# 10.3. Beneficiarios de la beca de régimen general del Ministerio de Educación, Formación Profesional y Deportes:

                  Código Seguro De Verificación: `eqOjhZcmL0p6VJlZiTH/Yw==`
                  Fecha: 27/06/2025
                  Firmado Por: MIGUEL ANGEL CASTRO ARROYO
                  Url De Verificación: [https://pfirma.us.es/verifirma/code/eqOjhZcmL0p6VJlZiTH\%2FYw\%3D\%3D](https://pfirma.us.es/verifirma/code/eqOjhZcmL0p6VJlZiTH\%2FYw\%3D\%3D)
                  Página: 25/50

                  ---

                  FUENTE: R.R. NORMAS DE MATRICULA GRADO Y MASTER UNIVERSITARIO 2025-2026.pdf | PÁGINA: 30
                  \# UNIVERSIDAD D SEVII

                  \#\# C)EST

                  Esta bonificación se aplicará al importe de matrícula no cubierta por la beca de los estudiantes que resulten beneficiarios de beca de estudios de régimen general, en el curso 2025-2026, del Ministerio de Educación, Formación Profesional y Deportes o del correspondiente departamento de educación del País Vasco. A tal efecto, el alumnado que reúna los requisitos académicos para obtener dicha beca tendrá en cuenta lo siguiente:

                  1. Si hubiese disfrutado de beca del Ministerio o del correspondiente departamento de educación del País Vasco durante el curso 2024-2025, la Universidad de Sevilla comprobará si ha solicitado beca para el curso 2025-2026.
                  2. Si no hubiese disfrutado de beca del Ministerio o del correspondiente departamento de educación del País Vasco durante el curso 2024-2025, deberá haber presentado la solicitud de dicha beca para el curso 2025-2026 en los términos del apartado anterior, o bien cumplimentar la declaración responsable de no cumplir los requisitos económicos para su obtención a través del Buzón de Documentación (DOMUS) habilitado en SEVIUS. La falta de aportación de la solicitud o la falsedad, inexactitud u omisión de datos necesarios en la declaración responsable dará lugar al ingreso de la parte del precio público que se consideró bonificada, con los intereses y recargos que procedan.

                  \#\# 10.4. Reconocimiento de créditos:

                  Esta bonificación no se aplicará en los supuestos de reconocimiento de créditos, salvo en el caso de estudios realizados en Programas de Movilidad Estudiantil, ya sea en el marco de la Unión Europea, ya sea realizados mediante convenios específicos entre universidades y conforme al correspondiente acuerdo académico de movilidad firmado por los estudiantes.

                  \#\# 10.5. Cálculo de la bonificación:

                  El precio de referencia para el alumnado de centros propios de universidades públicas de Andalucía será el establecido en el decreto por el que se determinan los precios públicos de las universidades públicas de Andalucía, para el crédito en primera matrícula de los estudios en que se matricule el alumnado, determinado según sus circunstancias personales y familiares. Respecto del alumnado de centros adscritos a universidades públicas de Andalucía, se estará a lo dispuesto en el citado decreto.

                  El importe de la bonificación será el 99\% del que resulte de multiplicar el número de créditos aprobados en primera matrícula en el curso 2024-2025, o en los cursos 2023-2024 y 2024-2025 para la primera matriculación en estudios de máster universitario, por el precio determinado en el párrafo anterior. En el cómputo de créditos aprobados en primera matrícula no se incluirán aquellos en que se hubiese obtenido matrícula de honor, por estar ya bonificados.

                  En el caso de alumnado matriculado en centros adscritos a universidades públicas de Andalucía, la bonificación se calculará y aplicará sobre la parte que los estudiantes deban abonar a la correspondiente universidad pública en concepto de servicios académicos, de conformidad con lo dispuesto en el decreto por el que se determinan los precios públicos de las universidades públicas de Andalucía. Esta bonificación se aplicará, como máximo, a los precios públicos correspondientes al número de créditos en que se matricule el alumnado y exclusivamente sobre los servicios académicos, sin que pueda dar lugar, en ningún caso, a devolución de importes.

                  ----

                  \textbf{Código Seguro De Verificación:} eqOjhZcmL0p6VJlZiTH/Yw==
                  \textbf{Fecha:} 27/06/2025
                  \textbf{Firmado Por:} MIGUEL ANGEL CASTRO ARROYO
                  \textbf{Url De Verificación:} [https://pfirma.us.es/verifirma/code/eqOjhZcmL0p6VJlZiTH\%2FYw\%3D\%3D](https://pfirma.us.es/verifirma/code/eqOjhZcmL0p6VJlZiTH\%2FYw\%3D\%3D)
                  \textbf{Página:} 26/50

                  ---

                  FUENTE: R.R. NORMAS DE MATRICULA GRADO Y MASTER UNIVERSITARIO 2025-2026.pdf | PÁGINA: 23
                  \# UNIVERSIDAD D SEVII

                  \#\# C)EST

                  \#\#\# Artículo 28. Reducción de precios. Exenciones y bonificaciones

                  Los estudiantes no abonarán en parte o en su totalidad los precios públicos por servicios académicos y/o administrativos, en su caso, cuando en ellos concurra/n alguna/s circunstancia/s que genere/n un derecho a la exención o bonificación en la matrícula, según lo establecido en el decreto por el que se determinan los precios públicos de las universidades públicas de Andalucía.

                  El estudiante deberá marcar todas las posibles exenciones a las que pudiera tener derecho en el momento de la formalización de la matrícula, excepto la aplicación de matrícula de honor obtenida en estudios universitarios oficiales cursados en la Universidad de Sevilla y la bonificación del 99\% por créditos aprobados en primera matrícula, al ser automáticas. Si el estudiante obtiene el reconocimiento de persona con discapacidad con posterioridad al período de su matrícula, pero con efectos de vigencia que incluya dicho periodo, aunque no marcara dicha exención en el momento de formalizar su matrícula, se le aplicará la exención correspondiente.

                  Cuando el estudiante haya marcado en automatrícula la exención de ser beneficiario de la condición de familia numerosa, de tener la condición de persona con discapacidad o de ser beneficiario de derechos por ser víctima de actos de terrorismo, la Universidad de Sevilla lo consultará telemáticamente, salvo que se opusiera a ello. De no poderse verificar telemáticamente alguna de estas condiciones, el estudiante deberá aportar la documentación acreditativa en el Buzón de Documentación (DOMUS) habilitado en la Secretaría Virtual de la Universidad de Sevilla (SEVIUS).

                  La documentación acreditativa del resto de las exenciones alegadas, en su caso, será aportada igualmente en el Buzón de Documentación (DOMUS) habilitado en la Secretaría Virtual de la Universidad de Sevilla (SEVIUS) en el plazo de diez días hábiles a contar desde el día siguiente a la finalización del periodo en el que haya formalizado su matrícula, para su verificación por la secretaría del centro.

                  Las exenciones y bonificaciones son las siguientes:

                  1. Estudiantes becarios de la convocatoria general del Ministerio de Educación, Formación Profesional y Deportes (con cargo a los Presupuestos Generales del Estado). La beca de matrícula sólo cubrirá los créditos en primera matrícula de los estudiantes becarios. Abonarán íntegramente los precios de los créditos o asignaturas en segunda o sucesivas matrículas y los precios públicos por servicios administrativos y el seguro escolar. Los estudiantes que soliciten esta beca podrán formalizar su matrícula sin el previo pago de los precios públicos por servicios académicos de los que se matriculen por primera vez. No obstante, las secretarías de los centros podrán requerir cautelarmente el abono de dichos precios públicos a quienes previsiblemente no cumplan los requisitos establecidos en la convocatoria de estas becas.

                  2. Los estudiantes becados por el Gobierno Vasco no estarán obligados a abonar el importe de la contraprestación por servicios académicos que sea cubierto por la beca.

                  3. Aplicación de matrícula de honor, premios y medallas:
                  a) La matrícula de honor en segundo curso de Bachillerato o Ciclo Formativo de Formación Profesional de Grado Superior, la obtención de Premio Extraordinario de Bachillerato o de Formación Profesional de Grado Superior o los premios de las Olimpiadas de Matemáticas, Física o Química u otras olimpiadas acreditadas de ámbito nacional o internacional, otorgará el derecho a exención total de los precios por servicios académicos, por una sola vez, en el primer curso de enseñanzas de grado.

                  Código Seguro De Verificación: `eqOjhZcmL0p6VJlZiTH/Yw==`
                  Fecha: 27/06/2025
                  Firmado Por: MIGUEL ANGEL CASTRO ARROYO
                  Url De Verificación: [https://pfirma.us.es/verifirma/code/eqOjhZcmL0p6VJlZiTH\%2FYw\%3D\%3D](https://pfirma.us.es/verifirma/code/eqOjhZcmL0p6VJlZiTH\%2FYw\%3D\%3D)
                  Página: 19/50

                  ---

                  FUENTE: R.R. NORMAS DE MATRICULA GRADO Y MASTER UNIVERSITARIO 2025-2026.pdf | PÁGINA: 24
                  \# UNIVERSIDAD D SEVII

                  \#\# C)EST

                  Se acreditará mediante certificado donde conste la matrícula de honor o con documento acreditativo de la obtención del premio extraordinario o la Medalla en las Olimpiadas, según corresponda.

                  Se aplicará a los precios públicos por la prestación de servicios académicos y solo se abonarán los precios públicos por servicios administrativos y el seguro escolar.

                  \#\#\# b) Matrícula de honor en estudios oficiales universitarios:

                  Los estudiantes que hayan obtenido matrícula de honor en alguna/s asignatura/s en el último curso universitario realizado tendrán derecho a la exención correspondiente, por una sola vez, en la siguiente matrícula de la misma titulación.

                  Se aplicará también a las siguientes situaciones:

                  1. En el caso de adaptación a un plan de estudios, se aplicará el derecho a la exención correspondiente en el nuevo plan de estudios matriculado. Solo se tendrán en cuenta las matrículas de honor obtenidas en el último curso académico del plan de origen.

                  2. Acceso a estudios de máster universitario si la matrícula de honor se ha obtenido en los estudios que le han dado acceso a dicho máster.

                  3. En caso de traslado de un centro adscrito a uno propio para cursar la misma titulación.

                  4. Estudiantes procedentes de otra universidad. Si las matrículas de honor obtenidas corresponden a asignaturas cursadas en la misma titulación en otra universidad, serán computadas a estos efectos, en los términos previstos en este apartado.

                  En los supuestos previstos en los epígrafes anteriores, conforme a lo dispuesto en el decreto por el que se determinan los precios públicos de las universidades públicas de Andalucía, la bonificación en el importe de la matrícula será equivalente al número de créditos en que obtuvo dicha matrícula de honor, multiplicado por el precio establecido para el crédito en primera matrícula de los estudios en que se matricule, determinado según sus circunstancias personales y familiares.

                  No se podrá aplicar el derecho a esta exención en el caso de que la calificación de matrícula de honor sea consecuencia de convalidación de asignaturas o reconocimiento de créditos.

                  Se deberán abonar íntegramente los precios públicos por servicios administrativos y el seguro escolar.

                  En el supuesto de no acreditar esta exención, previamente requerida, se eliminará de su matrícula mediante resolución dictada y notificada al efecto.

                  \#\# 4. Beneficiarios del título de familia numerosa:

                  a) Si fuera de categoría general, abonarán el 50\% de todos los precios públicos. El seguro escolar, en su caso, lo abonarán íntegramente.

                  b) Si fuera de categoría especial, tendrán derecho a la exención total de todos los precios públicos. Solo abonarán el seguro escolar.

                  ----

                  Código Seguro De Verificación: `eqOjhZcmL0p6VJlZiTH/Yw==`
                  Fecha: `27/06/2025`
                  Firmado Por: `MIGUEL ANGEL CASTRO ARROYO`
                  Url De Verificación: [https://pfirma.us.es/verifirma/code/eqOjhZcmL0p6VJlZiTH\%2FYw\%3D\%3D](https://pfirma.us.es/verifirma/code/eqOjhZcmL0p6VJlZiTH\%2FYw\%3D\%3D)
                  Página: `20/50`

                  ---

                  FUENTE: R.R. NORMAS DE MATRICULA GRADO Y MASTER UNIVERSITARIO 2025-2026.pdf | PÁGINA: 33
                  \# UNIVERSIDAD D SEVII

                  No obstante, las secretarías de los centros deberán requerir la documentación original a los estudiantes para su verificación en aquellos casos en los que la documentación aportada ofrezca dudas.

                  En todo caso, el estudiante deberá proporcionar la documentación original acreditativa de la vía de acceso de estudiantes procedentes de sistemas educativos de fuera de España.

                  \#\# Disposición adicional segunda. Créditos superados en la tercera convocatoria

                  De acuerdo con el decreto por el que se determinan los precios públicos de las universidades públicas de Andalucía, en caso de estudiantes que hayan superado los exámenes de alguna/s asignatura/s en la tercera convocatoria del curso 2025-2026, tendrán derecho a la deducción del 70\% del precio abonado en la matrícula correspondiente a dicha/s asignatura/s, bien procediendo la Universidad de Sevilla a regularizar los pagos aplazados del estudiante, caso de fraccionamiento, bien mediante la devolución de los importes en cuestión.

                  \#\# Disposición adicional tercera. Cambio de universidad y/o estudios universitarios oficiales

                  De conformidad con lo dispuesto en los artículos 44 y 45 del Real Decreto 534/2024, de 11 de junio, por el que se regulan los requisitos de acceso a las enseñanzas universitarias oficiales de Grado, las características básicas de la prueba de acceso y la normativa básica de los procedimientos de admisión, en materia de cambio de universidad y/o estudios universitarios oficiales españoles para estudiantes con estudios parciales y admisión de estudiantes con estudios universitarios extranjeros parciales o totales que no hayan obtenido la homologación o equivalencia de sus títulos, diplomas o estudios en España, resultarán de aplicación el Acuerdo de Consejo de Gobierno de 16/6/10 y la Resolución Rectoral de 16 de junio de 2010, por la que se regula la admisión a los títulos de grado de la Universidad de Sevilla de los estudiantes que han iniciado anteriormente estudios universitarios.

                  Las dobles titulaciones están excluidas del ámbito de aplicación de las normas sobre cambio de universidad y/o estudios universitarios oficiales, con la excepción de los traslados forzosos regulados en las Disposiciones Adicionales Primera y Segunda de la Resolución Rectoral de 16 de junio de 2010.

                  \#\# Disposición adicional cuarta. Matrícula de los estudiantes que acceden a los Itinerarios Curriculares Concretos

                  La matrícula de los estudiantes que accedan a los Itinerarios Curriculares Concretos se regirá por sus disposiciones específicas, y en lo no previsto en las mismas, por esta normativa.

                  \#\# Disposición adicional quinta. Estudiantes con necesidades académicas especiales

                  La consideración de estudiante con necesidades académicas especiales en los términos previstos en el artículo 27 del Reglamento General de Estudiantes, se solicitará en el centro correspondiente, dentro del plazo de matrícula ordinaria establecido en el Calendario Académico, salvo circunstancias sobrevenidas debidamente acreditadas.

                  \#\# Disposición adicional sexta. Prácticas que implican contacto habitual con menores

                  En cumplimiento de la Ley Orgánica 8/2021, de 4 de junio, de protección integral a la infancia y la adolescencia frente a la violencia (BOE núm. 134 de 5/06/2021), será requisito para cursar asignaturas que impliquen realizar prácticas en contacto habitual con menores, no haber sido condenado por sentencia firme por cualquier delito contra la libertad e indemnidad sexuales tipificados en el título VIII de la Ley Orgánica 10/1995, de 23 de noviembre, del Código Penal, así como por cualquier delito de trata de seres humanos tipificado en el título VII bis del Código Penal.

                  Código Seguro De Verificación: eqOjhZcmL0p6VJlZiTH/Yw==
                  Fecha: 27/06/2025
                  Firmado Por: MIGUEL ANGEL CASTRO ARROYO
                  Url De Verificación: [https://pfirma.us.es/verifirma/code/eqOjhZcmL0p6VJlZiTH\%2FYw\%3D\%3D](https://pfirma.us.es/verifirma/code/eqOjhZcmL0p6VJlZiTH\%2FYw\%3D\%3D)
                  Página: 29/50
                \end{scriptsize}
        \end{itemize}


  \item \textbf{P39.} Según la normativa de evaluación de la US, si a un alumno de un grado le quedan menos del 10 por ciento de los créditos totales para graduarse, excluidos los TFG, ¿tiene derecho a una convocatoria añadida especial? \\
        \begin{itemize}
          \item \textbf{Tipo de consulta:} simple
          \item \textbf{Respuesta del Agente:}

                Según la normativa de evaluación de la Universidad de Sevilla (US), específicamente en el documento "R.R. NORMAS DE MATRICULA GRADO Y MASTER UNIVERSITARIO 2025-2026.pdf" en la página 18, se establecen las condiciones para la concesión de convocatorias adicionales.

                En el punto 5 de la sección C)EST, se indica que "Los estudiantes a los que reste un máximo de 30 créditos para concluir su titulación no verán limitado, a efectos de permanencia, el número máximo de convocatorias salvo que el plan de estudios se extinga, en cuyo caso, deberán adaptarse al nuevo para continuar sus estudios."

                Sin embargo, no se menciona explícitamente una convocatoria añadida especial para estudiantes que les queden menos del 10\% de los créditos totales para graduarse, excluidos los TFG.

                En el punto 4 de la misma sección, se establece que "Agotado el número máximo de convocatorias (seis) en una asignatura, aquellos estudiantes que hubieran superado el primer año de permanencia, a los que les resten para finalizar sus estudios más de 30 créditos, dispondrán de una primera convocatoria de gracia para superar dicha asignatura que será concedida de oficio sin necesidad de solicitarla."

                Además, se menciona que los estudiantes que hubiesen agotado la primera convocatoria de gracia podrán solicitar al Rector con carácter excepcional, una segunda convocatoria de gracia, que será resuelta favorablemente siempre que en el momento de la solicitud hubiesen superado, al menos, el 50\% de la carga lectiva global de la titulación.

                En resumen, aunque no se menciona una convocatoria añadida especial para estudiantes que les queden menos del 10\% de los créditos totales para graduarse, excluidos los TFG, la normativa establece que los estudiantes que les resten más de 30 créditos para finalizar sus estudios pueden disponer de una primera convocatoria de gracia y, en casos excepcionales, de una segunda convocatoria de gracia.

                Es importante tener en cuenta que la normativa puede ser objeto de cambios y actualizaciones, por lo que se recomienda consultar la versión más reciente de la normativa y, si es necesario, contactar con la universidad para obtener información más detallada y precisa.
          \item \textbf{Contexto utilizado:}
                \begin{scriptsize}
                  FUENTE: R.R. NORMAS DE MATRICULA GRADO Y MASTER UNIVERSITARIO 2025-2026.pdf | PÁGINA: 18
                  \# UNIVERSIDAD DE SEVILLA

                  \#\# C)EST

                  1. A todos los efectos, una asignatura reconocida es una asignatura superada.

                  Excepcionalmente, el estudiante de primer curso de nuevo ingreso que no supere, al menos, una asignatura, podrá iniciar, por una sola vez, estudios en otra titulación en la Universidad de Sevilla, siempre que acredite reunir las condiciones exigidas para su acceso en la misma y obtenga plaza por los procedimientos legalmente establecidos.

                  Con carácter extraordinario, la Junta de centro podrá autorizar, a petición del estudiante, la continuación de los estudios en aquellos casos particulares en los que causas excepcionales suficientemente acreditadas hubiesen afectado el rendimiento académico del estudiante de primer curso de nuevo ingreso. Contra la resolución adoptada, el estudiante podrá interponer recurso.

                  2. Los estudiantes matriculados dispondrán de seis convocatorias para superar cada asignatura.

                  3. No se computarán las convocatorias en las que el estudiante figure en acta con la mención de “No Presentado”.

                  4. Agotado el número máximo de convocatorias (seis) en una asignatura, aquellos estudiantes que hubieran superado el primer año de permanencia, a los que les resten para finalizar sus estudios más de 30 créditos, dispondrán de una primera convocatoria de gracia para superar dicha asignatura que será concedida de oficio sin necesidad de solicitarla.

                  Los estudiantes que hubiesen agotado la primera convocatoria de gracia prevista en el apartado anterior, podrán solicitar al Rector con carácter excepcional, una segunda convocatoria de gracia, que será resuelta favorablemente siempre que en el momento de la solicitud hubiesen superado, al menos, el 50\% de la carga lectiva global de la titulación.

                  Se resolverá favorablemente la petición excepcional de una tercera convocatoria de gracia al Rector si el estudiante, en el momento de la solicitud, hubiese superado el 75\% de la carga lectiva global de la titulación.

                  Sobre estos porcentajes se podrá aplicar un margen diferencial de hasta el 2.5\%.

                  La convocatoria de gracia se agotará cuando el estudiante se presente a examen de esa asignatura y podrá ser utilizada en el curso académico que el estudiante considere.

                  5. Los estudiantes a los que reste un máximo de 30 créditos para concluir su titulación no verán limitado, a efectos de permanencia, el número máximo de convocatorias salvo que el plan de estudios se extinga, en cuyo caso, deberán adaptarse al nuevo para continuar sus estudios. En el caso de dobles titulaciones internacionales, se estará a lo dispuesto en los respectivos convenios.

                  6. Agotadas todas las convocatorias posibles de una asignatura sin haberla superado, el estudiante no podrá proseguir sus estudios, en la misma titulación, en la Universidad de Sevilla. No obstante, en este curso académico podrá seguir cursando el resto de las asignaturas matriculadas.

                  7. En cumplimiento de lo establecido en el artículo 20.c) del Real Decreto Legislativo 1/2013, de 29 de noviembre, por el que se aprueba el Texto Refundido de la Ley General de derechos de las personas con discapacidad y su inclusión social, las personas con un grado de discapacidad igual o superior al 33\% que les dificulte gravemente la adaptación al régimen de convocatorias oficiales de exámenes y hayan agotado el número máximo de convocatorias posibles, podrán solicitar al Vicerrectorado de Estudiantes la ampliación del número de las mismas en la medida que compense su dificultad, sin mengua del nivel exigido, previo informe favorable del Vicerrectorado de Servicios Sociales, Campus Saludable, Igualdad y Cooperación.

                  ----

                  \textbf{Código Seguro De Verificación:} eqOjhZcmL0p6VJlZiTH/Yw==
                  \textbf{Fecha:} 27/06/2025
                  \textbf{Firmado Por:} MIGUEL ANGEL CASTRO ARROYO
                  \textbf{Url De Verificación:} [https://pfirma.us.es/verifirma/code/eqOjhZcmL0p6VJlZiTH\%2FYw\%3D\%3D](https://pfirma.us.es/verifirma/code/eqOjhZcmL0p6VJlZiTH\%2FYw\%3D\%3D)
                  \textbf{Página:} 14/50

                  ---

                  FUENTE: R.R. NORMAS DE MATRICULA GRADO Y MASTER UNIVERSITARIO 2025-2026.pdf | PÁGINA: 6
                  \# UNIVERSIDAD D SEVII

                  \#\# C)EST

                  \#\#\# Artículo 4. Limitaciones de matriculación. Matrícula de asignaturas sueltas y pendientes

                  1. Sin perjuicio de la normativa estatal o autonómica aplicable, en los estudios universitarios oficiales de grado y máster, el número máximo de créditos ECTS de los que un estudiante puede estar matriculado en cada curso académico será de 90, sin que esta limitación afecte a las prácticas externas ni al trabajo fin de grado/máster universitario. En las dobles titulaciones, se podrá superar este límite hasta duplicarlo.

                  En el caso de simultaneidad de estudios se aplicará para cada una de las titulaciones lo previsto en el párrafo anterior.

                  Los créditos reconocidos no computarán a estos efectos.

                  2. De conformidad con lo dispuesto en el artículo 26 del Reglamento General de Actividades Docentes y sin perjuicio de la normativa estatal o autonómica aplicable, con objeto de que se puedan cursar estudios a tiempo parcial, en los estudios universitarios oficiales de grado y máster el número mínimo de créditos ECTS de los que un estudiante puede matricularse será de 30, salvo en los siguientes casos excepcionales:

                  a) Cuando los créditos matriculados sean los restantes para obtener la titulación.

                  b) Cuando el estudiante haya sido catalogado como estudiante con necesidades académicas especiales de conformidad con lo previsto en los artículos 26 y 27 del Reglamento General de Estudiantes, en cuyo caso el número mínimo de créditos a matricular será de 12. Este mínimo podrá rebajarse siempre que corresponda a todos los créditos que le falten al estudiante para finalizar sus estudios.

                  Los créditos reconocidos no se tendrán en cuenta para el cómputo de los créditos mínimos obligatorios de matriculación.

                  3. El número mínimo de créditos en los que el estudiante debe matricularse para tener derecho a la obtención de becas o ayudas vendrá determinado en la correspondiente convocatoria.

                  4. En los títulos de grado, los estudiantes podrán matricularse en una asignatura siempre que lo hagan de todas las asignaturas de formación básica de cursos inferiores que no hubiesen superado, entendiendo por asignaturas de formación básica lo dispuesto en el artículo 14 del Real Decreto 822/2021, de 28 de septiembre, por el que se establece la organización de las enseñanzas universitarias y del procedimiento de aseguramiento de su calidad.

                  Sin perjuicio de lo dispuesto en el apartado anterior, las Juntas de centro podrán acordar propuestas adicionales de itinerarios de matriculación ordenada en diversas materias o asignaturas de los planes de estudios de los títulos de grado adscritos al centro. Dichas propuestas serán sometidas, dentro del plazo que se fije en el Calendario Académico, al Consejo de Gobierno para su aprobación.

                  Con independencia de lo dispuesto en los párrafos 1 y 2 de este epígrafe 4, podrán exceptuarse los estudiantes que cursen estudios en virtud de un programa de movilidad nacional o internacional, cuya matrícula se atendrá a lo dispuesto en el correspondiente acuerdo de estudios firmado por el estudiante y el centro.

                  A través de los programas de asesoramiento a los que se refiere el artículo 21 del Reglamento General de Estudiantes, se podrán formular orientaciones individuales de matriculación ordenada.

                  ----

                  Código Seguro De Verificación: `eqOjhZcmL0p6VJlZiTH/Yw==`
                  Fecha: 27/06/2025
                  Firmado Por: MIGUEL ANGEL CASTRO ARROYO
                  Url De Verificación: [https://pfirma.us.es/verifirma/code/eqOjhZcmL0p6VJlZiTH\%2FYw\%3D\%3D](https://pfirma.us.es/verifirma/code/eqOjhZcmL0p6VJlZiTH\%2FYw\%3D\%3D)
                  Página: 2/50

                  ---

                  FUENTE: R.R. NORMAS DE MATRICULA GRADO Y MASTER UNIVERSITARIO 2025-2026.pdf | PÁGINA: 7
                  \# UNIVERSIDAD D SEVII

                  C)EST

                  Deberá tenerse en cuenta, además, que el derecho a examen y evaluación de las asignaturas matriculadas quedará limitado a las incompatibilidades académicas derivadas de los respectivos planes de estudios. La matriculación en asignaturas sujetas a régimen de incompatibilidad o de prerrequisito se encuentra condicionada, según las exigencias de los respectivos planes de estudios, a haber superado las asignaturas exigidas o a formalizar también, al mismo tiempo, la matriculación en estas.

                  5. En ningún caso, el ejercicio del derecho de matrícula obligará a la modificación del régimen de horarios generales determinados en cada centro.

                  6. Con carácter general, las asignaturas optativas, para ser impartidas, deben tener, al menos, seis estudiantes matriculados. De no alcanzarse esta cifra durante dos cursos consecutivos, será de aplicación lo dispuesto en el artículo 8.2.d) del Reglamento General de Actividades Docentes.

                  7. Una vez que el estudiante haya abonado los derechos de expedición del título, no podrá matricularse de más asignaturas optativas, quedando cerrado su expediente.

                  El estudiante que, encontrándose en la situación antes descrita, se hubiera matriculado de un número de créditos superior al necesario para finalizar sus estudios, no tendrá derecho a presentarse a posteriores convocatorias de examen ni a aparecer en actas de examen, y tampoco a la devolución de los precios públicos abonados por estos créditos, quedando cerrado su expediente a todos los efectos. Todo ello, sin perjuicio de la regulación referida a segundas menciones o especialidades.

                  \#\# Artículo 5. Condicionalidad de la matrícula

                  1. La matrícula estará condicionada al cumplimiento y comprobación de los requisitos establecidos, entendiéndose definitiva si con anterioridad al 31 de diciembre del año en curso, el estudiante no hubiese recibido notificación del requerimiento de subsanación de la misma, todo ello con las excepciones previstas en los apartados 2 y 3 de este artículo y sin perjuicio de las actuaciones que de oficio pudieran realizarse.

                  2. Se considerará la condicionalidad en los supuestos y hasta los momentos que se citan a continuación:

                  a) Cuando el estudiante haya solicitado reconocimiento de estudios previos cursados en la misma u otra universidad, hasta su resolución.

                  b) Cuando el estudiante haya solicitado reconocimiento parcial de estudios cursados en el extranjero o el estudiante esté pendiente de la homologación o declaración de equivalencia de su título universitario extranjero o de la homologación de sus estudios preuniversitarios por el bachillerato español, hasta su resolución.

                  c) Cuando la solicitud de matrícula conlleve la determinación del nivel académico de los estudios previos cursados en el extranjero, hasta su resolución.

                  d) Cuando el estudiante haya solicitado cambio de universidad y/o estudios, hasta tanto se reciba en la secretaría del centro la certificación académica oficial del centro de procedencia, excepto aquellos traslados que se realicen entre centros propios de la Universidad de Sevilla. En este último caso, por parte del centro de destino se realizará la comprobación telemática que lo comunicará al centro de origen para que se proceda al cierre del expediente, no debiendo el estudiante abonar los derechos de traslado.

                  e) En los casos de simultaneidad de estudios con otros ya iniciados, hasta tanto no se haya efectuado el traslado de expediente por simultaneidad y se reciba en la secretaría del centro, certificación académica oficial de simultaneidad del estudiante procedente del centro en el que ya.

                  ----

                  Código Seguro De Verificación: `eqOjhZcmL0p6VJlZiTH/Yw==`
                  Fecha: 27/06/2025
                  Firmado Por: MIGUEL ANGEL CASTRO ARROYO
                  Url De Verificación: [https://pfirma.us.es/verifirma/code/eqOjhZcmL0p6VJlZiTH\%2FYw\%3D\%3D](https://pfirma.us.es/verifirma/code/eqOjhZcmL0p6VJlZiTH\%2FYw\%3D\%3D)
                  Página: 3/50

                  ---

                  FUENTE: R.R. NORMAS DE MATRICULA GRADO Y MASTER UNIVERSITARIO 2025-2026.pdf | PÁGINA: 8
                  \# UNIVERSIDAD D SEVII

                  C)EST

                  esté matriculado. Del mismo modo, se considerará condicional la matrícula de los estudiantes que deseen simultanear estudios desde su inicio hasta tanto se compruebe el cumplimiento de los requisitos de admisión.

                  f) Cuando el estudiante haya solicitado, dentro del plazo general de matrícula, salvo circunstancias sobrevenidas valoradas por el centro, la consideración de estudiante con necesidades académicas especiales en los términos previstos en el artículo 27 del Reglamento General de Estudiantes, hasta tanto recaiga la resolución correspondiente.

                  g) Cuando el estudiante no haya entregado toda la documentación necesaria según las instrucciones de matrícula, hasta la finalización del plazo de subsanación notificado por la secretaría del centro.

                  3. En tanto no se resuelva favorablemente el supuesto condicionante, no se admitirán a trámite solicitudes de título universitario oficial. No se emitirán certificaciones académicas oficiales en las que conste rendimiento académico, excepto en los supuestos contemplados en los apartados 2.a) y 2.f) de este artículo.

                  4. Las matrículas condicionadas, si el supuesto condicionante no se resuelve favorablemente, podrán ser anuladas o modificadas, según corresponda.

                  5. Los estudiantes de las universidades españolas a los que a la fecha de finalización del último plazo de admisión fijado por la Comisión de Distrito Único Universitario de Andalucía (Fase 3), solo les reste por superar la acreditación de la competencia lingüística en idioma extranjero requerida, el TFG y un máximo de 9 créditos ECTS para obtener el título de grado podrán solicitar la admisión condicionada en estudios de máster universitario en las universidades públicas andaluzas, con las siguientes consideraciones:

                  - La matrícula en el máster universitario se entenderá condicionada hasta la finalización de los estudios de grado, siendo definitiva cuando el estudiante acredite haber abonado los derechos de expedición del título de grado. En ningún caso se podrá obtener el título de máster universitario con anterioridad al título de grado que permite el acceso al máster.

                  - El estudiante estará obligado a matricular todos los créditos pendientes de los estudios de grado, incluido, en su caso, el TFG.

                  - El estudiantado de máster universitario con matrícula condicionada tendrá los mismos derechos y deberes que cualquier otro estudiante matriculado en un máster universitario del Distrito Único Universitario de Andalucía, si bien no se expedirán certificaciones académicas oficiales de estos estudios en tanto su matrícula sea condicionada.

                  - Solamente procederá la devolución del importe de la matrícula de los estudios de máster universitario en los supuestos específicamente regulados por estas normas de matrícula.

                  \#\# Artículo 6. Programas académicos con recorridos sucesivos (PARS)

                  El estudiante de grado vinculado a máster universitario en los Programas Académicos con Recorridos Sucesivos (PARS) en el ámbito de la ingeniería y arquitectura al que le reste por superar la acreditación del nivel de competencia lingüística requerido y un máximo de 30 créditos ECTS de la tipología que se establezca en cada uno de los Programas Académicos con Recorridos Sucesivos (PARS), podrá acceder y matricularse en el máster universitario vinculado, siempre que sea admitido en la Fase 3 del procedimiento de admisión.

                  No obstante, en ningún caso se podrá obtener el título de máster universitario con anterioridad al título de grado que permite el acceso al máster, ni se expedirán certificaciones académicas oficiales de los estudios de máster universitario.

                  Código Seguro De Verificación: `eqOjhZcmL0p6VJlZiTH/Yw==`
                  Fecha: `27/06/2025`
                  Firmado Por: `MIGUEL ANGEL CASTRO ARROYO`
                  Url De Verificación: [https://pfirma.us.es/verifirma/code/eqOjhZcmL0p6VJlZiTH\%2FYw\%3D\%3D](https://pfirma.us.es/verifirma/code/eqOjhZcmL0p6VJlZiTH\%2FYw\%3D\%3D)
                  Página: `4/50`

                  ---

                  FUENTE: Guia\_Movilidad\_Intern\_FAC\_TUR\_FINANZAS.pdf | PÁGINA: 15
                  Se ha de tener en cuenta que para el reconocimiento de créditos de asignaturas que forman parte de un mismo bloque, se requiere que todas las asignaturas cursadas en destino hayan sido aprobadas. Por tanto, es muy importante la agrupación de esas asignaturas que aparezca en el acuerdo de estudios validado.

                  Para aquellas asignaturas que no hayan sido superadas en destino el alumno podrá ejercer su derecho a la segunda convocatoria de esa asignatura en la Universidad de Sevilla. Si por el contrario decidiera presentarse en destino a ese examen porque la universidad de destino permite esta opción, el/la estudiante deberá informar de ello a la Vicedecana de Movilidad y también a Nieves Caro, responsable en la Secretaría del Centro gestor1ftf@us.es.

                  \# Contactos

                  - Vicedecana de Movilidad e Internacionalización: Mª Rosa García Sánchez, relaintftf@us.es.
                  - Gestora en Secretaría: Nieves Caro, gestor1ftf@us.es.
                  - Centro Internacional de la Universidad de Sevilla, Unidad de Movilidad Internacional Saliente: movilidad@us.es.

                  15
                \end{scriptsize}
        \end{itemize}


  \item \textbf{P40.} De acuerdo estrictamente con el artículo sobre Transferencia y Reconocimiento, ¿se reconocerán créditos superados de asignaturas pertenecientes a titulaciones privadas no oficiales? \\
        \begin{itemize}
          \item \textbf{Tipo de consulta:} simple
          \item \textbf{Respuesta del Agente:}

                De acuerdo con la normativa proporcionada, específicamente en el artículo 7.1, se establece que "Las solicitudes de reconocimiento de créditos basadas en estudios superados en títulos universitarios que no tengan carácter oficial y validez en todo el territorio nacional, se resolverán teniendo en cuenta la adecuación entre las competencias y conocimientos adquiridos en las materias superadas y los previstos en el plan de estudios del título de destino, indicándose las materias o asignaturas que se considerarán superadas por el interesado y que, por lo tanto, no estarán obligados a cursar."

                Además, en el artículo 7.2, se especifica que "El número de créditos que sean objeto de reconocimiento a partir de enseñanzas universitarias no oficiales y de la experiencia profesional o laboral prevista en el artículo 9 no podrá ser superior, en su conjunto, al 15 por ciento del total de créditos que constituyen el plan de estudios de destino."

                Por lo tanto, según la normativa, sí se reconocerán créditos superados de asignaturas pertenecientes a titulaciones privadas no oficiales, siempre y cuando se cumplan los siguientes requisitos:

                1. La solicitud de reconocimiento de créditos se base en estudios superados en títulos universitarios no oficiales.
                2. La adecuación entre las competencias y conocimientos adquiridos en las materias superadas y los previstos en el plan de estudios del título de destino.
                3. El número de créditos reconocidos no supere el 15\% del total de créditos del plan de estudios de destino.

                Es importante destacar que el reconocimiento de créditos no incorporará calificación de los mismos, por lo que no computarán a efectos de baremación del expediente, según se establece en el artículo 7.4.
          \item \textbf{Contexto utilizado:}
                \begin{scriptsize}
                  FUENTE: Normativa Reconocimiento y Transferencia de creditos.pdf | PÁGINA: 5
                  \# Textos consolidados

                  \#\# Artículo 7. A partir de otros títulos universitarios.

                  \#\#\# 7.1.
                  Las solicitudes de reconocimiento de créditos basadas en estudios superados en títulos universitarios que no tengan carácter oficial y validez en todo el territorio nacional, se resolverán teniendo en cuenta la adecuación entre las competencias y conocimientos adquiridos en las materias superadas y los previstos en el plan de estudios del título de destino, indicándose las materias o asignaturas que se considerarán superadas por el interesado y que, por lo tanto, no estarán obligados a cursar.

                  \#\#\# 7.2.
                  El número de créditos que sean objeto de reconocimiento a partir de enseñanzas universitarias no oficiales y de la experiencia profesional o laboral prevista en el artículo 9 no podrá ser superior, en su conjunto, al 15 por ciento del total de créditos que constituyen el plan de estudios de destino.

                  \#\#\# 7.3.
                  No obstante lo anterior, los créditos procedentes de títulos propios podrán, excepcionalmente, ser objeto de reconocimiento en un porcentaje superior al señalado en el párrafo anterior o, en su caso, ser objeto de reconocimiento en su totalidad siempre que el correspondiente título propio haya sido extinguido y sustituido por el título oficial para el que se solicita el reconocimiento.

                  \#\#\# 7.4.
                  El reconocimiento de estos créditos no incorporará calificación de los mismos por lo que no computarán a efectos de baremación del expediente.

                  \#\#\# 7.5.
                  En ningún caso podrán ser objeto de reconocimiento los Trabajos Fin de Grado.

                  \#\# Artículo 8. A partir de títulos de enseñanzas superiores.

                  \#\#\# 8.1.
                  Las solicitudes de reconocimiento de créditos basadas en estudios superados en títulos oficiales españoles de educación superior no universitaria, se resolverán teniendo en cuenta la adecuación entre las competencias y conocimientos adquiridos en las materias superadas y los previstos en el plan de estudios del título de destino, indicándose las materias o asignaturas que se considerarán superadas por el interesado y que, por lo tanto, no estarán obligados a cursar.

                  \#\#\# 8.2.
                  Podrán ser objeto de reconocimiento de créditos los estudios superados correspondientes a los siguientes títulos:
                  - a. Título Superior de Arte Dramático.
                  - b. Título Superior de Artes Plásticas.
                  - c. Título Superior de Conservación y Restauración de Bienes Culturales.
                  - d. Título Superior de Danza.
                  - e. Título Superior de Diseño.
                  - f. Título Superior de Música.
                  - g. Técnico Superior de Artes Plásticas y Diseño.
                  - h. Técnico Superior de Formación Profesional.
                  - i. Técnico Deportivo Superior.

                  \#\#\# 8.3.
                  Únicamente podrán ser objeto de reconocimiento de créditos los estudios acreditados mediante los títulos oficiales enumerados en el apartado anterior. En el caso de enseñanzas artísticas de grado conducentes a titulaciones oficiales podrán ser objeto de reconocimiento los periodos parciales de estudios cursados, siempre que se acrediten oficialmente en créditos ECTS.

                  \#\#\# 8.4.
                  En función de los criterios generales que determine el Ministerio de Educación, Cultura y Deporte y, previo acuerdo con la Administración educativa correspondiente, se garantizará un reconocimiento mínimo de créditos ECTS a quienes posean una titulación de educación superior y cursen otras enseñanzas relacionadas con dicho título.

                  Este documento es de carácter informativo y no tiene valor jurídico.

                  ---

                  FUENTE: Normativa Reconocimiento y Transferencia de creditos.pdf | PÁGINA: 8
                  \# Textos consolidados

                  \#\# Artículo 13. A partir de otros títulos universitarios.

                  \#\#\# 13.1.
                  Las solicitudes de reconocimiento de créditos basadas en estudios superados en títulos universitarios que no tengan carácter oficial y validez en todo el territorio nacional, se resolverán teniendo en cuenta la adecuación entre las competencias y conocimientos adquiridos en las materias superadas y los previstos en el plan de estudios del título de destino, indicándose las materias o asignaturas que se considerarán superadas por el interesado y que, por lo tanto, no estarán obligados a cursar.

                  \#\#\# 13.2.
                  El número de créditos que sean objeto de reconocimiento a partir de enseñanzas universitarias no oficiales y de la experiencia profesional o laboral prevista en el artículo 14 no podrá ser superior, en su conjunto, al 15 por ciento del total de créditos que constituyen el plan de estudios de destino.

                  \#\#\# 13.3.
                  No obstante lo anterior, los créditos procedentes de títulos propios podrán, excepcionalmente, ser objeto de reconocimiento en un porcentaje superior al señalado en el párrafo anterior o, en su caso, ser objeto de reconocimiento en su totalidad siempre que el correspondiente título propio haya sido extinguido y sustituido por el título oficial para el que se solicita el reconocimiento.

                  \#\#\# 13.4.
                  El reconocimiento de estos créditos no incorporará calificación de los mismos por lo que no computarán a efectos de baremación del expediente.

                  \#\#\# 13.5.
                  En ningún caso podrán ser objeto de reconocimiento los Trabajos Fin de Máster.

                  \#\# Artículo 14. A partir de experiencia laboral o profesional.

                  \#\#\# 14.1.
                  Las solicitudes de reconocimiento de créditos basadas en experiencia laboral o profesional acreditada se resolverán teniendo en cuenta su relación con las competencias inherentes al título, indicándose las materias o asignaturas que se considerarán superadas por el interesado y que, por lo tanto, no estarán obligados a cursar, o los créditos aplicados al cómputo de optatividad del plan de estudios del título que se pretende obtener.

                  \#\#\# 14.2.
                  El número de créditos que sean objeto de reconocimiento a partir de experiencia profesional o laboral acreditada y de enseñanzas universitarias no oficiales previstas en el artículo 13 no podrá ser superior, en su conjunto, al 15 por ciento del total de créditos que constituyen el plan de estudios de destino.

                  \#\#\# 14.3.
                  Dentro de este límite se reconocerán hasta 6 créditos por cada año de experiencia laboral o profesional debidamente acreditada.

                  \#\#\# 14.4.
                  Con carácter general, siempre que el plan de estudios contemple la posibilidad o necesidad de realizar prácticas externas, el reconocimiento de créditos por experiencia laboral o profesional se aplicará preferentemente a este tipo de materias.

                  \#\#\# 14.4.bis.
                  Asimismo, podrán reconocerse por prácticas curriculares, aquellas prácticas extracurriculares que hayan sido gestionadas desde la Universidad de Sevilla o cualquier otra Universidad, al amparo del mismo título para el que se solicita el reconocimiento, siempre y cuando así lo estime la Comisión competente en función del programa formativo acreditado de las mismas y de su relación con las competencias inherentes al título.

                  \#\#\# 14.5.
                  El reconocimiento de estos créditos no incorporará calificación de los mismos por lo que no computarán a efectos de baremación del expediente.

                  \#\#\# 14.6.
                  En ningún caso podrán ser objeto de reconocimiento los Trabajos Fin de Máster.

                  Este documento es de carácter informativo y no tiene valor jurídico.

                  ---

                  FUENTE: Normativa Reconocimiento y Transferencia de creditos.pdf | PÁGINA: 10
                  \# Textos consolidados

                  15.7. De los acuerdos de estudios que se establezcan se enviará copia a los Servicios Centrales del Rectorado que corresponda.
                  15.8. Con carácter general lo dispuesto en estas normas será de aplicación a la movilidad para dobles titulaciones sin perjuicio de las previsiones contenidas en los convenios respectivos.
                  15.9. Resultarán igualmente de aplicación las normas que eventualmente se aprueben por los órganos nacionales o internacionales competentes para cada programa específico de movilidad.

                  \#\# CAPÍTULO V
                  \#\#\# Transferencia de créditos
                  \textbf{Artículo 16. Definición.}
                  La transferencia de créditos implica que, en los documentos académicos oficiales acreditativos de las enseñanzas seguidas por cada estudiante, se incluirán la totalidad de los créditos obtenidos en enseñanzas universitarias oficiales cursadas con anterioridad, en ésta u otra universidad, que no hayan conducido a la finalización de sus estudios con la consiguiente obtención de un título oficial.

                  \textbf{Artículo 17. Aplicación.}
                  Los créditos correspondientes a materias o asignaturas previamente superadas por el estudiante, en enseñanzas universitarias no concluidas y que no puedan ser objeto de reconocimiento, serán transferidos a su expediente en los estudios a los que ha accedido con la calificación de origen y se reflejarán en los documentos académicos oficiales acreditativos de los estudios seguidos por el mismo, así como en el Suplemento Europeo al Título.

                  \#\# CAPÍTULO VI
                  \#\#\# Tramitación
                  \textbf{Artículo 18. Solicitudes de reconocimiento de créditos.}
                  18.1. Los expedientes de reconocimiento de créditos se tramitarán a solicitud del interesado, quién deberá aportar la documentación justificativa de los créditos obtenidos y su contenido académico, indicando los módulos, materias o asignaturas que considere superados.
                  18.2. Será requisito imprescindible estar admitido y matriculado en los correspondientes estudios, salvo en los casos de cambios de estudios oficiales de Grado, según lo dispuesto en la Resolución Rectoral por la que se regula la admisión a los títulos de Grado de la Universidad de Sevilla de los estudiantes que han iniciado anteriormente otros estudios universitarios.
                  18.3. Las solicitudes de reconocimiento de créditos tendrán su origen en materias o asignaturas realmente cursadas y superadas, en ningún caso se referirán a materias o asignaturas previamente reconocidas, convalidadas o adaptadas.
                  18.4. Las solicitudes se presentarán en la Secretaría del Centro responsable del título para el que se solicita el reconocimiento dentro de los plazos indicados en el calendario académico de cada curso.

                  \textbf{Artículo 19. Documentación acreditativa.}
                  19.1. En el caso de estudios universitarios cursados, estudios superiores no universitarios u otros estudios no oficiales, se aportará la siguiente documentación:

                  Este documento es de carácter informativo y no tiene valor jurídico.
                  10

                  ---

                  FUENTE: Normativa Reconocimiento y Transferencia de creditos.pdf | PÁGINA: 4
                  \# Textos consolidados

                  superado íntegramente un determinado módulo, el reconocimiento se llevará a cabo por materias o asignaturas en función de las competencias y conocimientos asociados a las mismas.

                  7) En el caso de títulos oficiales de Grado que habiliten para el ejercicio de profesiones reguladas en España, para los que el Gobierno haya establecido las condiciones a las que han de adecuarse los planes de estudios, se reconocerán los créditos de los módulos definidos en la correspondiente norma reguladora que hayan sido superados por el estudiante. En caso de no haberse superado íntegramente un determinado módulo, el reconocimiento se llevará a cabo por materias o asignaturas en función de las competencias y conocimientos asociados a las mismas.

                  \#\# Artículo 5. A partir de títulos de Máster Universitario.

                  5.1. Las solicitudes de reconocimiento de créditos basadas en estudios superados en títulos oficiales de Máster Universitario (tanto los regulados por el RD 56/2005, como por el RD 1393/2007) o periodo de formación específico del Doctorado se resolverán teniendo en cuenta la adecuación entre las competencias y conocimientos adquiridos en las materias superadas y los previstos en el plan de estudios del título de destino, indicándose las materias o asignaturas que se considerarán superadas por el interesado y que, por lo tanto, no estarán obligados a cursar.

                  5.2. En ningún caso podrán ser objeto de reconocimiento los Trabajos Fin de Grado.

                  \#\# Artículo 6. A partir de títulos de la anterior ordenación universitaria.

                  6.1. Las solicitudes de reconocimiento de créditos basadas en estudios superados en títulos universitarios de la anterior ordenación universitaria, Arquitecto Técnico, Diplomado, Ingeniero Técnico, Maestro, Licenciado, Arquitecto, Ingeniero o periodo de docencia del doctorado, se resolverán teniendo en cuenta la adecuación entre las competencias y conocimientos adquiridos en las materias superadas y los previstos en el plan de estudios del título de destino, indicándose las materias o asignaturas que se considerarán superadas por el interesado y que, por lo tanto, no estarán obligados a cursar.

                  6.2. Para la resolución de estas solicitudes se tendrán en cuenta los siguientes criterios:

                  1) Cuando las competencias y conocimientos no estén explicitados o no puedan deducirse del plan de estudios de origen del estudiante se tomarán como referencia el número de créditos y/o los contenidos de las materias o asignaturas cursadas.

                  2) En el caso de títulos en proceso de extinción por la implantación de los nuevos títulos de Grado, la adaptación de los estudiantes a estos últimos se basará en el reconocimiento de créditos previsto en la tabla de adaptación incluida en la correspondiente memoria de verificación del título de Grado en cuestión.

                  3) En los procesos de adaptación de estudiantes a los nuevos planes de los títulos de Grado deberá garantizarse que la situación académica de aquellos no resulte perjudicada. A tal efecto, las materias, asignaturas o créditos superados que no tengan equivalencia en las correspondientes al plan de estudios de Grado se incorporarán en el expediente del estudiante como créditos genéricos de carácter optativo. Si, aún así resultaran excedentes, los créditos restantes se podrán incorporar al expediente como créditos transferidos, a petición del interesado y siempre que se trate de materias o asignaturas completas.

                  4) El número de créditos reconocidos con cargo a la optatividad no podrá superar el número de créditos optativos exigido por el plan de estudios del título de destino.

                  5) En ningún caso podrán ser objeto de reconocimiento los Trabajos Fin de Grado.

                  Este documento es de carácter informativo y no tiene valor jurídico.

                  ---

                  FUENTE: Normativa Reconocimiento y Transferencia de creditos.pdf | PÁGINA: 3
                  \# Textos consolidados

                  \#\# Artículo 2. Ámbito de aplicación.
                  La presente normativa reguladora será de aplicación a las enseñanzas universitarias oficiales de Grado y Máster impartidas por la Universidad de Sevilla.

                  \#\# Artículo 3. Definiciones.
                  \#\#\# 3.1.
                  Se entiende por reconocimiento de créditos la aceptación por la Universidad de Sevilla, a efectos de la obtención de un título universitario oficial, de:
                  - a. Los créditos obtenidos en otras enseñanzas universitarias oficiales.
                  - b. Los créditos obtenidos en enseñanzas superiores oficiales no universitarias.
                  - c. Los créditos obtenidos en enseñanzas universitarias conducentes a otros títulos.
                  - d. La acreditación de experiencia laboral o profesional.
                  - e. La participación en actividades universitarias culturales, deportivas, de representación estudiantil, solidarias y de cooperación.

                  \#\#\# 3.2.
                  La transferencia de créditos implica que, en los documentos académicos oficiales acreditativos de las enseñanzas seguidas por cada estudiante, se incluirán la totalidad de los créditos obtenidos en enseñanzas oficiales cursadas con anterioridad, en la misma u otra universidad, que no hayan conducido a la obtención de un título oficial.

                  \#\# CAPÍTULO II
                  \#\#\# Reconocimiento de créditos en enseñanzas oficiales de Grado

                  \#\# Artículo 4. A partir de otros títulos de Grado.
                  \#\#\# 4.1.
                  Las solicitudes de reconocimiento de créditos basadas en estudios superados en títulos universitarios de Grado se resolverán teniendo en cuenta la adecuación entre las competencias y conocimientos adquiridos en las materias superadas y los previstos en el plan de estudios del título de destino, indicándose las materias o asignaturas que se considerarán superadas por el interesado y que, por lo tanto, no estarán obligados a cursar.

                  \#\#\# 4.2.
                  Para la resolución de estas solicitudes se tendrán en cuenta los siguientes criterios:
                  1. Cuando el título de origen y el título de destino pertenezcan a la misma rama de conocimiento serán objeto de reconocimiento todos los créditos superados en materias de formación básica vinculadas a dicha rama de conocimiento. Cuando se hayan superado la totalidad de los créditos de formación básica del título de origen, se garantizará el reconocimiento de al menos 36 créditos correspondientes a materias de formación básica de dicha rama en el título de destino.
                  2. Cuando el título de origen y el título de destino pertenezcan a diferentes ramas de conocimiento serán objeto de reconocimiento todos los créditos superados en materias de formación básica de la rama de conocimiento a la que se encuentre adscrito el título de destino.
                  3. En todo caso, los efectos del reconocimiento de créditos se reflejarán en la resolución indicando las materias o asignaturas concretas que se considerarán superadas –que podrán tener el carácter de formación básica, obligatoria, optativa o prácticas externas- y, en su caso, los créditos reconocidos con cargo al cómputo de optatividad del plan de estudios.
                  4. El número de créditos reconocidos con cargo a la optatividad no podrá superar el número de créditos optativos exigido por el plan de estudios del título de destino.
                  5. En ningún caso podrán ser objeto de reconocimiento los Trabajos Fin de Grado.
                  6. En el ámbito del sistema universitario público andaluz serán objeto de reconocimiento automático los módulos o materias comunes definidas para cada título de Grado. En caso de no haberse

                  ----

                  Este documento es de carácter informativo y no tiene valor jurídico.
                \end{scriptsize}
        \end{itemize}


  \item \textbf{P41.} Según las normas de evaluación, si nunca me llego a presentar físicamente a la primera convocatoria ordinaria del primer cuatrimestre, ¿consume eso alguna convocatoria a efectos de normas de permanencia? \\
        \begin{itemize}
          \item \textbf{Tipo de consulta:} simple
          \item \textbf{Respuesta del Agente:}

                Según las normas de evaluación de la Universidad de Sevilla, específicamente en el artículo 20.3 de las Normas de Matrícula Grado y Master Universitario 2025-2026, se establece que "No se computarán las convocatorias en las que el estudiante figure en acta con la mención de “No Presentado”".

                Por lo tanto, si no te presentas físicamente a la primera convocatoria ordinaria del primer cuatrimestre, no se consumirá ninguna convocatoria a efectos de las normas de permanencia, ya que no figurarás en acta con una calificación, sino con la mención de "No Presentado".

                Es importante tener en cuenta que, según el artículo 22.2 de la misma normativa, los estudiantes que se matriculen por primera vez en una asignatura solo podrán presentarse a las dos convocatorias posteriores a la finalización del periodo lectivo de la misma. Sin embargo, en este caso, no se consume ninguna convocatoria al no presentarte a la primera convocatoria.

                En resumen, no consumirás ninguna convocatoria a efectos de las normas de permanencia si no te presentas a la primera convocatoria ordinaria del primer cuatrimestre, siempre y cuando no figure en acta con una calificación.
          \item \textbf{Contexto utilizado:}
                \begin{scriptsize}
                  FUENTE: R.R. NORMAS DE MATRICULA GRADO Y MASTER UNIVERSITARIO 2025-2026.pdf | PÁGINA: 17
                  \# UNIVERSIDAD DE SEVILLA

                  \#\# CAPÍTULO II: DE LOS EFECTOS DE LA MATRÍCULA

                  \#\#\# Artículo 18. Convocatorias de exámenes finales

                  1. Las convocatorias ordinarias a lo largo del curso académico serán tres; los plazos de entrega de las actas correspondientes se establecerán en el Calendario Académico aprobado por el Consejo de Gobierno.

                  2. Los estudiantes que se matriculen por primera vez en una asignatura solo podrán presentarse a las dos convocatorias posteriores a la finalización del periodo lectivo de la misma.

                  3. En cada curso académico, los estudiantes tendrán derecho a presentarse a tres convocatorias ordinarias en aquellas asignaturas matriculadas en segunda o sucesivas matrículas.

                  4. En el caso de los estudiantes que, para la conclusión de sus estudios, tengan pendientes de superar menos del 10\% de los créditos totales, excluidos los créditos asignados al trabajo fin de carrera y a las prácticas externas, una de las convocatorias, a elección de estos, tendrá la consideración de convocatoria extraordinaria. La no superación de una convocatoria extraordinaria no será computada a efectos de las normas de permanencia.

                  \#\#\# Artículo 19. Calendario de convocatorias de exámenes finales

                  1. En el Calendario Académico aprobado anualmente por el Consejo de Gobierno de la Universidad de Sevilla para cada curso académico, se establecerán las fechas de celebración de las correspondientes convocatorias de exámenes.

                  2. Para las asignaturas anuales o cuatrimestrales matriculadas por segunda o sucesivas veces en la misma titulación, el estudiante podrá presentarse a las dos convocatorias establecidas para cada caso y a una tercera convocatoria que tendrá lugar en el mes de octubre de cada curso académico. A esta tercera convocatoria podrá presentarse también el estudiante que se encuentre en la situación descrita en el artículo 18.4 de estas normas. La presentación a los exámenes de esta convocatoria implica su cómputo a efectos de permanencia del estudiante en el centro, salvo lo dispuesto en el artículo 20.5 de estas normas.

                  3. Cronológicamente, el orden de celebración de las convocatorias será: tercera, primera y segunda convocatoria.

                  4. El calendario para la celebración de exámenes finales por los estudiantes de movilidad internacional será el establecido con carácter general para todos los estudiantes matriculados en la Universidad de Sevilla salvo que en el convenio específico al amparo del cual se realiza el intercambio, se contemple la posibilidad de habilitar otras fechas. Aquellos estudiantes que hicieran uso de esta posibilidad deberán rendir en el examen la totalidad de la materia incluida en el correspondiente programa de estudios de la asignatura.

                  \#\# CAPÍTULO III: NORMAS DE PERMANENCIA

                  \#\#\# Artículo 20. Permanencia

                  1. El estudiante de primer curso de nuevo ingreso deberá superar, al menos, una asignatura reglada en cualquiera de las convocatorias oficiales del curso académico con independencia de la tipología de la misma, para poder continuar los mismos estudios.

                  ----

                  Código Seguro De Verificación: `eqOjhZcmL0p6VJlZiTH/Yw==`
                  Fecha: 27/06/2025
                  Firmado Por: MIGUEL ANGEL CASTRO ARROYO
                  Url De Verificación: [https://pfirma.us.es/verifirma/code/eqOjhZcmL0p6VJlZiTH\%2FYw\%3D\%3D](https://pfirma.us.es/verifirma/code/eqOjhZcmL0p6VJlZiTH\%2FYw\%3D\%3D)
                  Página: 13/50

                  ---

                  FUENTE: R.R. NORMAS DE MATRICULA GRADO Y MASTER UNIVERSITARIO 2025-2026.pdf | PÁGINA: 18
                  \# UNIVERSIDAD DE SEVILLA

                  \#\# C)EST

                  1. A todos los efectos, una asignatura reconocida es una asignatura superada.

                  Excepcionalmente, el estudiante de primer curso de nuevo ingreso que no supere, al menos, una asignatura, podrá iniciar, por una sola vez, estudios en otra titulación en la Universidad de Sevilla, siempre que acredite reunir las condiciones exigidas para su acceso en la misma y obtenga plaza por los procedimientos legalmente establecidos.

                  Con carácter extraordinario, la Junta de centro podrá autorizar, a petición del estudiante, la continuación de los estudios en aquellos casos particulares en los que causas excepcionales suficientemente acreditadas hubiesen afectado el rendimiento académico del estudiante de primer curso de nuevo ingreso. Contra la resolución adoptada, el estudiante podrá interponer recurso.

                  2. Los estudiantes matriculados dispondrán de seis convocatorias para superar cada asignatura.

                  3. No se computarán las convocatorias en las que el estudiante figure en acta con la mención de “No Presentado”.

                  4. Agotado el número máximo de convocatorias (seis) en una asignatura, aquellos estudiantes que hubieran superado el primer año de permanencia, a los que les resten para finalizar sus estudios más de 30 créditos, dispondrán de una primera convocatoria de gracia para superar dicha asignatura que será concedida de oficio sin necesidad de solicitarla.

                  Los estudiantes que hubiesen agotado la primera convocatoria de gracia prevista en el apartado anterior, podrán solicitar al Rector con carácter excepcional, una segunda convocatoria de gracia, que será resuelta favorablemente siempre que en el momento de la solicitud hubiesen superado, al menos, el 50\% de la carga lectiva global de la titulación.

                  Se resolverá favorablemente la petición excepcional de una tercera convocatoria de gracia al Rector si el estudiante, en el momento de la solicitud, hubiese superado el 75\% de la carga lectiva global de la titulación.

                  Sobre estos porcentajes se podrá aplicar un margen diferencial de hasta el 2.5\%.

                  La convocatoria de gracia se agotará cuando el estudiante se presente a examen de esa asignatura y podrá ser utilizada en el curso académico que el estudiante considere.

                  5. Los estudiantes a los que reste un máximo de 30 créditos para concluir su titulación no verán limitado, a efectos de permanencia, el número máximo de convocatorias salvo que el plan de estudios se extinga, en cuyo caso, deberán adaptarse al nuevo para continuar sus estudios. En el caso de dobles titulaciones internacionales, se estará a lo dispuesto en los respectivos convenios.

                  6. Agotadas todas las convocatorias posibles de una asignatura sin haberla superado, el estudiante no podrá proseguir sus estudios, en la misma titulación, en la Universidad de Sevilla. No obstante, en este curso académico podrá seguir cursando el resto de las asignaturas matriculadas.

                  7. En cumplimiento de lo establecido en el artículo 20.c) del Real Decreto Legislativo 1/2013, de 29 de noviembre, por el que se aprueba el Texto Refundido de la Ley General de derechos de las personas con discapacidad y su inclusión social, las personas con un grado de discapacidad igual o superior al 33\% que les dificulte gravemente la adaptación al régimen de convocatorias oficiales de exámenes y hayan agotado el número máximo de convocatorias posibles, podrán solicitar al Vicerrectorado de Estudiantes la ampliación del número de las mismas en la medida que compense su dificultad, sin mengua del nivel exigido, previo informe favorable del Vicerrectorado de Servicios Sociales, Campus Saludable, Igualdad y Cooperación.

                  ----

                  \textbf{Código Seguro De Verificación:} eqOjhZcmL0p6VJlZiTH/Yw==
                  \textbf{Fecha:} 27/06/2025
                  \textbf{Firmado Por:} MIGUEL ANGEL CASTRO ARROYO
                  \textbf{Url De Verificación:} [https://pfirma.us.es/verifirma/code/eqOjhZcmL0p6VJlZiTH\%2FYw\%3D\%3D](https://pfirma.us.es/verifirma/code/eqOjhZcmL0p6VJlZiTH\%2FYw\%3D\%3D)
                  \textbf{Página:} 14/50

                  ---

                  FUENTE: NORMATIVA REGUL.pdf | PÁGINA: 7
                  \# Artículo 19. Desarrollo de los exámenes

                  1. La duración máxima de un examen será de cinco horas.

                  2. Los estudiantes citados a examen comparecerán en su lugar de realización a la hora prevista y acreditarán su identificación mediante la exhibición del DNI, NIE, pasaporte o tarjeta universitaria.

                  3. Los estudiantes tendrán derecho a recibir, tras la finalización de un examen, un justificante documental de haberlo realizado. Cuando se trate de exámenes escritos, en el justificante podrá constar, a petición del estudiante, el número total de hojas entregadas por éste, previa comprobación de los profesores encargados de la vigilancia.

                  \# Artículo 20. Incidencias en los exámenes. Sin Contenido¹

                  \# CAPÍTULO III. CONVOCATORIAS, CALIFICACIONES Y ACTAS

                  \#\# SECCIÓN 1ª. CONVOCATORIAS

                  \# Artículo 21. Convocatorias

                  Las convocatorias ordinarias a lo largo del curso académico serán tres; los plazos de entrega de las actas correspondientes se establecerán en el calendario académico aprobado por el Consejo de Gobierno.

                  \# Artículo 22. Derecho a presentarse a las convocatorias

                  1. En cada curso académico, los estudiantes tendrán derecho a presentarse, en cada asignatura en que se matriculen, a dos de las tres convocatorias ordinarias, sin perjuicio de lo dispuesto en el artículo 61.4 del Reglamento general de actividades docentes.

                  2. Los estudiantes que se matriculen por primera vez en la asignatura sólo podrán presentarse a las dos convocatorias inmediatamente posteriores a la finalización del periodo lectivo de la asignatura.

                  3. Los estudiantes que se matriculen por segunda vez, o sucesivas, podrán elegir, de entre las tres ordinarias, las dos convocatorias a las que deseen presentarse en virtud de lo dispuesto en el apartado 1.

                  4. Los estudiantes que aprueben alguna de las asignaturas a las que se presenten en la tercera convocatoria ordinaria podrán ampliar su matrícula en los términos que se fijen en las normas de matrícula.

                  \#\# SECCIÓN 2ª. CALIFICACIONES

                  \# Artículo 23. Sistema de calificaciones

                  1. Artículo derogado por el Acuerdo 1.1 del Consejo de Gobierno en la sesión de 18 de marzo de 2010.

                  ---

                  FUENTE: R.R. NORMAS DE MATRICULA GRADO Y MASTER UNIVERSITARIO 2025-2026.pdf | PÁGINA: 11
                  \# UNIVERSIDAD D SEVII

                  C)EST

                  5. No se concederán anulaciones parciales de matrícula a solicitud del estudiante, excepto por coincidencia no prevista de horarios de clase y exámenes y las situaciones contempladas en estas normas. No obstante, por causas académicas sobrevenidas que imposibiliten la prestación del servicio, los centros podrán realizar de oficio anulaciones parciales de matrícula, ofreciendo a los estudiantes la posibilidad de matricular otras asignaturas o bien la devolución de los importes abonados en los créditos anulados, en su caso.

                  6. La matrícula se anulará de oficio cuando se trate de error, omisión, falseamiento de los datos consignados en esta o cuando se modifique la situación administrativa o académica del estudiante, o las matrículas condicionadas que no deban producir efectos.

                  Asimismo se anularán de oficio las matrículas cuando el estudiante no aporte la documentación requerida o por falta de pago.

                  7. Los centros anularán de oficio las matrículas efectuadas por estudiantes que se matriculen en otra titulación por haber obtenido plaza en los procedimientos de admisión para estudiantes de nuevo ingreso en esta u otra universidad pública, siempre que no se acredite una autorización para la simultaneidad de estudios.

                  \#\# Artículo 11. Supuestos de anulación de matrícula y sus efectos

                  1. Se podrá solicitar la anulación de matrícula ordinaria al inicio del curso académico en los siguientes plazos:

                  - Titulaciones de grado: hasta el 16 de septiembre de 2025.
                  - Titulaciones de máster universitario que inicien la docencia el 6 de octubre: hasta el 15 de octubre de 2025.
                  - Titulaciones de máster universitario que inicien la docencia el 22 de septiembre: hasta el 30 de septiembre de 2025.
                  - Cuando la matrícula se haya formalizado con posterioridad al inicio del periodo docente del primer cuatrimestre: 7 días hábiles a contar desde el día de la finalización del periodo en el que haya realizado su matrícula.

                  La concesión de la anulación supondrá:

                  a) La devolución de los importes abonados, en su caso.

                  b) La matrícula no se computará a efectos de convocatorias de examen ni de permanencia del estudiante, ni se contabilizará para futuras matrículas a efectos de aplicación de recargos.

                  c) La pérdida de la plaza adjudicada si el estudiante accedió por el procedimiento de admisión en el curso 2025-2026.

                  2. La anulación de la matrícula por haber obtenido plaza en esta u otra universidad pública (en el procedimiento o sistema de admisión que tenga establecido) en otra titulación de grado, máster universitario o estudios de doctorado donde hubiese sido admitido el estudiante con posterioridad a su matriculación en la Universidad de Sevilla, no tendrá ningún efecto académico ni económico si se solicita con anterioridad al 31 de octubre de 2025, en el caso de titulaciones de grado, y al 31 de diciembre de 2025 en las titulaciones de máster universitario.

                  Código Seguro De Verificación: `eqOjhZcmL0p6VJlZiTH/Yw==`
                  Fecha: 27/06/2025
                  Firmado Por: MIGUEL ANGEL CASTRO ARROYO
                  Url De Verificación: [https://pfirma.us.es/verifirma/code/eqOjhZcmL0p6VJlZiTH\%2FYw\%3D\%3D](https://pfirma.us.es/verifirma/code/eqOjhZcmL0p6VJlZiTH\%2FYw\%3D\%3D)
                  Página: 7/50

                  ---

                  FUENTE: NORMATIVA REGUL.pdf | PÁGINA: 12
                  \# Contra las resoluciones de los tribunales de apelación

                  Cabe recurso de alzada ante el Rector en los plazos y formas que establece la Ley 30/1992, de 26 de Noviembre, de régimen jurídico de las administraciones públicas y del procedimiento administrativo común.

                  \#\# DISPOSICIONES ADICIONALES

                  \#\#\# Primera. Cita en género femenino de los preceptos de esta normativa

                  Las referencias a personas, colectivos o cargos académicos figuran en la presente normativa en género masculino como género gramatical no marcado. Cuando proceda, será válida la cita de los preceptos correspondientes en género femenino.

                  \#\#\# Segunda. Procedimiento de resolución de conflictos

                  1. Los miembros de la comunidad universitaria podrán plantear ante las Comisiones de Docencia de Centros y Departamentos solicitud de resolución de los conflictos planteados, en el ámbito respectivo, en la planificación y desarrollo de los procedimientos de evaluación. Quedan excluidos expresamente del ámbito de esta disposición adicional las revisiones de las calificaciones provisionales y los recursos de apelación contra las calificaciones definitivas, que se atendrán a lo regulado en el capítulo III de la presente normativa.

                  2. La solicitud de resolución de conflicto se presentará mediante escrito razonado dirigido al Decano o Director del Centro o el Director del Departamento, según el ámbito de la Comisión de Docencia ante la que se plantea la solicitud, que se entregará en la Secretaría del Centro, o en el Registro General de la Universidad.

                  3. La Comisión de Docencia analizará el escrito de solicitud y, oídas las partes en conflicto, resolverá en el plazo máximo de un mes desde la presentación del mismo, considerándose los periodos no lectivos de agosto, Navidad, Semana Santa y Feria de Sevilla como inhábiles a estos efectos.

                  4. El Presidente de la Comisión de Docencia comunicará por escrito la resolución al solicitante y demás partes en conflicto y, en su caso, al órgano unipersonal de gobierno competente para su ejecución o conocimiento. Caso de que la Comisión estimare que hay indicios de falta disciplinaria, lo comunicará directamente al Rector.

                  \#\#\# Tercera. Asignaturas de planes de estudio extinguidos

                  1. Una vez extinguido un curso de un plan de estudios, para la evaluación de las asignaturas de dicho curso se efectuarán tres convocatorias ordinarias en cada uno de los tres cursos académicos siguientes. El estudiante que no supere la asignatura en alguno de dichos cursos deberá adaptarse al nuevo plan de estudios.

                  2. Los estudiantes tendrán derecho a presentarse a dos de las tres convocatorias ordinarias, salvo en el caso contemplado en el artículo 61.4 del Reglamento general de actividades docentes, que será de aplicación también en el ámbito de aplicación de esta disposición adicional.
                \end{scriptsize}
        \end{itemize}


  \item \textbf{P42.} Basándonos en MATRUX 2025-2026, ¿tienen derecho a la gratuidad del precio público completo por matrícula todos aquellos estudiantes con discapacidad de grado igual o superior al 33 por 100? \\
        \begin{itemize}
          \item \textbf{Tipo de consulta:} simple
          \item \textbf{Respuesta del Agente:}

                Basándonos en la normativa proporcionada en el documento "R.R. NORMAS DE MATRICULA GRADO Y MASTER UNIVERSITARIO 2025-2026.pdf" y específicamente en la página 26, se establece que:

                "En aplicación del artículo 19 del Real Decreto Legislativo 1/2013, de 29 de noviembre, por el que se aprueba el Texto Refundido de la Ley General de derechos de las personas con discapacidad y de su inclusión social, los estudiantes que acrediten su condición de persona con discapacidad con un grado igual o superior al 33 por ciento, tendrán derecho a la exención total de precios públicos por la prestación de servicios académicos y administrativos en los estudios conducentes a la obtención de un título universitario oficial en primera y sucesivas matrículas. En cualquier caso, deberán abonar el seguro escolar, si procede."

                Por lo tanto, según la normativa, los estudiantes con discapacidad que tienen un grado de discapacidad igual o superior al 33 por 100 tienen derecho a la gratuidad del precio público completo por matrícula, excepto el seguro escolar, si procede.

                Es importante destacar que, según el artículo 28.2 de la ley 39/2015, de 1 de octubre, del Procedimiento Administrativo Común de las Administraciones Públicas, la Universidad de Sevilla consultará telemáticamente los datos de discapacidad del estudiante que la alegue en su matrícula, salvo que se opusiera a ello. Si el estudiante no puede verificar telemáticamente su condición de discapacidad, deberá aportar la documentación acreditativa en el Buzón de Documentación (DOMUS) habilitado en la Secretaría Virtual de la Universidad de Sevilla (SEVIUS).

                En resumen, los estudiantes con discapacidad que tienen un grado de discapacidad igual o superior al 33 por 100 tienen derecho a la gratuidad del precio público completo por matrícula, excepto el seguro escolar, si procede, y deberán acreditar su condición de discapacidad según los procedimientos establecidos por la Universidad de Sevilla.
          \item \textbf{Contexto utilizado:}
                \begin{scriptsize}
                  FUENTE: R.R. NORMAS DE MATRICULA GRADO Y MASTER UNIVERSITARIO 2025-2026.pdf | PÁGINA: 26
                  \# UNIVERSIDAD D SEVII

                  C)EST

                  El estudiante podrá ejercer el derecho de oposición a la consulta indicando las circunstancias que motivan dicha oposición, según los procedimientos establecidos para ello que podrá consultar en:

                  - https://osi.us.es
                  - https://osi.us.es/sites/osi/files/doc/c.i.gestionestudiantesoficialesModif1.12.5.20CT.pdf

                  En este supuesto, no se aplicará la exención de familia numerosa hasta tanto la secretaría del centro verifique la documentación aportada por el interesado en el acto de oposición a la consulta telemática.

                  5. Estudiantes a los que se les haya reconocido un grado de discapacidad igual o superior al 33 por ciento:

                  En aplicación del artículo 19 del Real Decreto Legislativo 1/2013, de 29 de noviembre, por el que se aprueba el Texto Refundido de la Ley General de derechos de las personas con discapacidad y de su inclusión social, los estudiantes que acrediten su condición de persona con discapacidad con un grado igual o superior al 33 por ciento, tendrán derecho a la exención total de precios públicos por la prestación de servicios académicos y administrativos en los estudios conducentes a la obtención de un título universitario oficial en primera y sucesivas matrículas. En cualquier caso, deberán abonar el seguro escolar, si procede.

                  Conforme lo establecido en el artículo 28.2 de la ley 39/2015, de 1 de octubre, del Procedimiento Administrativo Común de las Administraciones Públicas, la Universidad de Sevilla consultará telemáticamente los datos de discapacidad del estudiante que la alegue en su matrícula, salvo que se opusiera a ello.

                  De no poderse verificar telemáticamente la condición de persona con discapacidad cuando realice su automatrícula, el estudiante deberá aportar en el Buzón de Documentación (DOMUS) habilitado en la Secretaría Virtual de la Universidad de Sevilla (SEVIUS) la certificación expedida por el órgano competente que acredite dicha condición.

                  Las secretarías de los centros podrán hacer uso de la comprobación telemática, con independencia de la aportación de la documentación por el estudiante.

                  En todo caso, el centro deberá realizar las comprobaciones telemáticas a aquellos estudiantes a los que no haya sido posible la comprobación en la automatrícula y no hayan aportado en el Buzón de Documentación (DOMUS) la documentación exigida, antes de requerir la documentación acreditativa de la condición de persona con discapacidad.

                  En el supuesto de no acreditar dicha condición, previamente requerida, se eliminará de su matrícula mediante resolución dictada y notificada al efecto.

                  En caso de oposición a la consulta telemática, el estudiante deberá acreditar documentalmente su condición de discapacidad.

                  El estudiante podrá ejercer el derecho de oposición a la consulta indicando las circunstancias que motivan dicha oposición, según los procedimientos establecidos para ello que podrá consultar en:

                  - https://osi.us.es
                  - https://osi.us.es/sites/osi/files/doc/c.i.gestionestudiantesoficialesModif1.12.5.20CT.pdf

                  En este supuesto, no se aplicará la exención de persona con discapacidad hasta tanto la secretaría del centro verifique la documentación aportada por el interesado en el acto de oposición a la consulta telemática.

                  Código Seguro De Verificación: `eqOjhZcmL0p6VJlZiTH/Yw==`
                  Fecha: 27/06/2025
                  Firmado Por: MIGUEL ANGEL CASTRO ARROYO
                  Url De Verificación: [https://pfirma.us.es/verifirma/code/eqOjhZcmL0p6VJlZiTH\%2FYw\%3D\%3D](https://pfirma.us.es/verifirma/code/eqOjhZcmL0p6VJlZiTH\%2FYw\%3D\%3D)
                  Página: 22/50

                  ---

                  FUENTE: R.R. NORMAS DE MATRICULA GRADO Y MASTER UNIVERSITARIO 2025-2026.pdf | PÁGINA: 4
                  \# RESOLUCIÓN RECTORAL REGULADORA DE LAS NORMAS DE MATRÍCULA EN LOS ESTUDIOS OFICIALES DE GRADO Y MÁSTER UNIVERSITARIO DE LA UNIVERSIDAD DE SEVILLA CURSO 2025-2026

                  \#\# Artículos

                  \vspace{0.4cm}
                  \begin{table}[h!]\centering\renewcommand{\arraystretch}{1.3}
                    \begin{tabular}{|p{0.45\textwidth}|p{0.45\textwidth}|p{0.45\textwidth}|}\hline
                      Artículo & Título                                    & Página \\\hline
                      22       & Reconocimiento de créditos                & 15     \\\hline
                      23       & Solicitudes de reconocimiento de créditos & 15     \\\hline
                      24       & Transferencia de créditos                 & 16     \\\hline
                      25       & Solicitudes de transferencia de créditos  & 16     \\\hline
                    \end{tabular}\end{table}
                  \vspace{0.4cm}

                  \#\# TÍTULO II: RÉGIMEN ECONÓMICO

                  \#\#\# CAPÍTULO I: PRECIOS PÚBLICOS

                  \vspace{0.4cm}
                  \begin{table}[h!]\centering\renewcommand{\arraystretch}{1.3}
                    \begin{tabular}{|p{0.45\textwidth}|p{0.45\textwidth}|p{0.45\textwidth}|}\hline
                      Artículo & Título                                                              & Página \\\hline
                      26       & Precios a satisfacer por la prestación del servicio de la enseñanza & 17     \\\hline
                      27       & Modalidad y plazos. Efectos de la falta de pago                     & 17     \\\hline
                      28       & Reducción de precios. Exenciones y bonificaciones                   & 19     \\\hline
                      29       & Devolución parcial o total de precios                               & 27     \\\hline
                      30       & Ayudas de Acción Social                                             & 27     \\\hline
                    \end{tabular}\end{table}
                  \vspace{0.4cm}

                  \#\#\# CAPÍTULO II: SEGURO ESCOLAR Y CARNÉ UNIVERSITARIO

                  \vspace{0.4cm}
                  \begin{table}[h!]\centering\renewcommand{\arraystretch}{1.3}
                    \begin{tabular}{|p{0.45\textwidth}|p{0.45\textwidth}|p{0.45\textwidth}|}\hline
                      Artículo & Título              & Página \\\hline
                      31       & Seguro escolar.     & 27     \\\hline
                      32       & Carné universitario & 27     \\\hline
                    \end{tabular}\end{table}
                  \vspace{0.4cm}

                  \#\# Disposiciones adicionales

                  \vspace{0.4cm}
                  \begin{table}[h!]\centering\renewcommand{\arraystretch}{1.3}
                    \begin{tabular}{|p{0.45\textwidth}|p{0.45\textwidth}|}\hline
                      Disposición             & Página \\\hline
                      Adicional primera       & 28     \\\hline
                      Adicional segunda       & 29     \\\hline
                      Adicional tercera       & 29     \\\hline
                      Adicional cuarta        & 29     \\\hline
                      Adicional quinta        & 29     \\\hline
                      Adicional sexta         & 29     \\\hline
                      Adicional séptima       & 30     \\\hline
                      Adicional octava        & 30     \\\hline
                      Adicional novena        & 31     \\\hline
                      Adicional décima        & 31     \\\hline
                      Adicional undécima      & 31     \\\hline
                      Adicional duodécima     & 31     \\\hline
                      Adicional decimotercera & 31     \\\hline
                    \end{tabular}\end{table}
                  \vspace{0.4cm}

                  \#\# Disposiciones finales

                  \vspace{0.4cm}
                  \begin{table}[h!]\centering\renewcommand{\arraystretch}{1.3}
                    \begin{tabular}{|p{0.45\textwidth}|p{0.45\textwidth}|}\hline
                      Disposición   & Página \\\hline
                      Final primera & 32     \\\hline
                      Final segunda & 32     \\\hline
                      Final tercera & 32     \\\hline
                    \end{tabular}\end{table}
                  \vspace{0.4cm}

                  \#\# Anexo

                  \vspace{0.4cm}
                  \begin{table}[h!]\centering\renewcommand{\arraystretch}{1.3}
                    \begin{tabular}{|p{0.45\textwidth}|p{0.45\textwidth}|p{0.45\textwidth}|}\hline
                      Anexo & Título                                                      & Página \\\hline
                      I     & Instrucciones de matrícula grados y másteres universitarios & 33     \\\hline
                    \end{tabular}\end{table}
                  \vspace{0.4cm}

                  ---

                  FUENTE: R.R. NORMAS DE MATRICULA GRADO Y MASTER UNIVERSITARIO 2025-2026.pdf | PÁGINA: 23
                  \# UNIVERSIDAD D SEVII

                  \#\# C)EST

                  \#\#\# Artículo 28. Reducción de precios. Exenciones y bonificaciones

                  Los estudiantes no abonarán en parte o en su totalidad los precios públicos por servicios académicos y/o administrativos, en su caso, cuando en ellos concurra/n alguna/s circunstancia/s que genere/n un derecho a la exención o bonificación en la matrícula, según lo establecido en el decreto por el que se determinan los precios públicos de las universidades públicas de Andalucía.

                  El estudiante deberá marcar todas las posibles exenciones a las que pudiera tener derecho en el momento de la formalización de la matrícula, excepto la aplicación de matrícula de honor obtenida en estudios universitarios oficiales cursados en la Universidad de Sevilla y la bonificación del 99\% por créditos aprobados en primera matrícula, al ser automáticas. Si el estudiante obtiene el reconocimiento de persona con discapacidad con posterioridad al período de su matrícula, pero con efectos de vigencia que incluya dicho periodo, aunque no marcara dicha exención en el momento de formalizar su matrícula, se le aplicará la exención correspondiente.

                  Cuando el estudiante haya marcado en automatrícula la exención de ser beneficiario de la condición de familia numerosa, de tener la condición de persona con discapacidad o de ser beneficiario de derechos por ser víctima de actos de terrorismo, la Universidad de Sevilla lo consultará telemáticamente, salvo que se opusiera a ello. De no poderse verificar telemáticamente alguna de estas condiciones, el estudiante deberá aportar la documentación acreditativa en el Buzón de Documentación (DOMUS) habilitado en la Secretaría Virtual de la Universidad de Sevilla (SEVIUS).

                  La documentación acreditativa del resto de las exenciones alegadas, en su caso, será aportada igualmente en el Buzón de Documentación (DOMUS) habilitado en la Secretaría Virtual de la Universidad de Sevilla (SEVIUS) en el plazo de diez días hábiles a contar desde el día siguiente a la finalización del periodo en el que haya formalizado su matrícula, para su verificación por la secretaría del centro.

                  Las exenciones y bonificaciones son las siguientes:

                  1. Estudiantes becarios de la convocatoria general del Ministerio de Educación, Formación Profesional y Deportes (con cargo a los Presupuestos Generales del Estado). La beca de matrícula sólo cubrirá los créditos en primera matrícula de los estudiantes becarios. Abonarán íntegramente los precios de los créditos o asignaturas en segunda o sucesivas matrículas y los precios públicos por servicios administrativos y el seguro escolar. Los estudiantes que soliciten esta beca podrán formalizar su matrícula sin el previo pago de los precios públicos por servicios académicos de los que se matriculen por primera vez. No obstante, las secretarías de los centros podrán requerir cautelarmente el abono de dichos precios públicos a quienes previsiblemente no cumplan los requisitos establecidos en la convocatoria de estas becas.

                  2. Los estudiantes becados por el Gobierno Vasco no estarán obligados a abonar el importe de la contraprestación por servicios académicos que sea cubierto por la beca.

                  3. Aplicación de matrícula de honor, premios y medallas:
                  a) La matrícula de honor en segundo curso de Bachillerato o Ciclo Formativo de Formación Profesional de Grado Superior, la obtención de Premio Extraordinario de Bachillerato o de Formación Profesional de Grado Superior o los premios de las Olimpiadas de Matemáticas, Física o Química u otras olimpiadas acreditadas de ámbito nacional o internacional, otorgará el derecho a exención total de los precios por servicios académicos, por una sola vez, en el primer curso de enseñanzas de grado.

                  Código Seguro De Verificación: `eqOjhZcmL0p6VJlZiTH/Yw==`
                  Fecha: 27/06/2025
                  Firmado Por: MIGUEL ANGEL CASTRO ARROYO
                  Url De Verificación: [https://pfirma.us.es/verifirma/code/eqOjhZcmL0p6VJlZiTH\%2FYw\%3D\%3D](https://pfirma.us.es/verifirma/code/eqOjhZcmL0p6VJlZiTH\%2FYw\%3D\%3D)
                  Página: 19/50

                  ---

                  FUENTE: R.R. NORMAS DE MATRICULA GRADO Y MASTER UNIVERSITARIO 2025-2026.pdf | PÁGINA: 42
                  \# UNIVERSIDAD DE SEVILLA

                  \#\# CONSIDERACIONES GENERALES

                  - \textbf{Estudiantes que acrediten discapacidad:} La Universidad de Sevilla consultará telemáticamente los datos de discapacidad del estudiante que la alegue en su matrícula, salvo que se opusiera a ello. Si el estudiante obtiene el reconocimiento de persona con discapacidad con posterioridad al período de su matrícula, pero con efectos de vigencia que incluya dicho periodo, aunque no marcara dicha exención en el momento de formalizar su matrícula, se le aplicará la exención correspondiente.

                  - \textbf{Víctimas de actos de terrorismo:} La Universidad de Sevilla consultará telemáticamente la condición de beneficiario de derechos por ser víctima de actos de terrorismo del estudiante que la alegue en su matrícula, salvo que se opusiera a ello.

                  - \textbf{Víctimas de violencia de género:} Se acreditará presentando cualquiera de los documentos que se relacionan en el artículo 30 de la Ley 13/2007, de 26 de noviembre, de medidas de prevención y protección integral contra la violencia de género.

                  - \textbf{Personas beneficiarias del ingreso mínimo vital o de la renta mínima de inserción social en Andalucía:} Documentación acreditativa de dicha circunstancia.

                  - \textbf{Bonificación por créditos superados en primera matrícula:} Los estudiantes que no soliciten beca de convocatoria general en el curso 2025-2026 por no cumplir los requisitos económicos para la obtención de la misma, deberán presentar una declaración responsable del no cumplimiento de los referidos requisitos. Los estudiantes que no hayan solicitado beca por no cumplir los requisitos académicos para su obtención no tienen que presentar dicha declaración responsable. Este documento se generará una vez que el estudiante acceda al Buzón de Documentación (DOMUS) con su UVUS, quedando firmado mediante el doble factor de autenticación (2FA). Los estudiantes que inician una titulación durante el curso 2025-2026 también podrán aplicar, en su caso, las siguientes exenciones:

                  - \textbf{Matrícula de honor de aplicación a estudiantes que inician sus estudios de grado por primera vez:} Se acreditará con certificado donde conste la matrícula de honor.

                  - \textbf{Premio Extraordinario de Bachillerato o Formación Profesional de Grado Superior y Medallas en las Olimpiadas de Matemáticas, Física o Química u otras olimpiadas acreditadas de ámbito nacional o internacional de aplicación a estudiantes que inician sus estudios de grado por primera vez:} Se aportará documento acreditativo de su obtención.

                  - \textbf{Obtención de Premio Extraordinario Fin de Grado o de Enseñanzas Artísticas Superiores:} En el caso de estudiantes procedentes de otras universidades o centros de enseñanzas artísticas superiores, se acreditará mediante la aportación del acta de concesión del premio y declaración responsable de no haber aplicado con anterioridad los beneficios del premio obtenido en otros estudios universitarios oficiales.

                  5. \textbf{Reconocimiento/transferencia de créditos (ver arts. 21 a 25):} La presentación de solicitudes de reconocimiento de créditos ECTS para asignaturas regladas, se realizará desde la fecha de la formalización de la automatrícula/matrícula hasta el 22 de septiembre de 2025. Aquellos estudiantes que realicen su matrícula con posterioridad al 5 de septiembre dispondrán de un plazo de 10 días hábiles, contados a partir del día siguiente a la formalización de su matrícula, para la presentación de la solicitud.

                  ----

                  \textbf{Código Seguro De Verificación:} eqOjhZcmL0p6VJlZiTH/Yw==
                  \textbf{Fecha:} 27/06/2025
                  \textbf{Firmado Por:} MIGUEL ANGEL CASTRO ARROYO
                  \textbf{Url De Verificación:} [https://pfirma.us.es/verifirma/code/eqOjhZcmL0p6VJlZiTH\%2FYw\%3D\%3D](https://pfirma.us.es/verifirma/code/eqOjhZcmL0p6VJlZiTH\%2FYw\%3D\%3D)
                  \textbf{Página:} 38/50

                  ---

                  FUENTE: BOJA25-508-00085-6229-01\_00319829.pdf | PÁGINA: 58
                  \# BOJA
                  \#\# Boletín Oficial de la Junta de Andalucía
                  \#\#\# Extraordinario núm. 8 - Viernes, 2 de mayo de 2025
                  \#\#\#\# página 6229/58

                  u) A disfrutar de becas y ayudas al estudio, así como a exenciones de precios públicos, que serán concedidas sobre la base de los baremos que se fijen en las convocatorias o en las normas que las establezcan.
                  v) Al paro académico, respetando el derecho a la educación del estudiantado, en las condiciones que se establezcan en el Reglamento general del estudiantado.
                  w) A recibir una atención especial en caso de embarazo y lactancia, y en situaciones personales de grave dificultad, de modo que se les preste asesoramiento en el estudio de las asignaturas, facilidades para la realización de prácticas y otras actividades presenciales y adecuación de fechas para la realización de pruebas y exámenes.
                  x) A que se adopten medidas de acción positiva para que el estudiantado con necesidades académicas especiales, y, en particular, con discapacidad, pueda disfrutar de una educación universitaria inclusiva, accesible y adaptable, en igualdad con el resto del estudiantado, realizando ajustes razonables, tanto curriculares como metodológicos, a los materiales didácticos, a los métodos de enseñanza y al sistema de evaluación.
                  y) A la accesibilidad universal de los edificios y sus entornos físicos y virtuales, así como los servicios, procedimientos, suministros y comunicación de información, los materiales educativos y los procesos de enseñanza-aprendizaje y evaluación.

                  2. Con objeto de que nadie quede excluido del estudio en la Universidad de Sevilla por razones económicas, esta instrumentará, dentro de sus posibilidades presupuestarias, una adecuada política de becas, ayudas y modalidades de exención parcial o total del pago de los precios públicos por prestación de los servicios académicos, que complementen las establecidas por el Estado y la Comunidad Autónoma de Andalucía. Esta política contemplará también el deporte, el apoyo al estudiantado con discapacidad, la vivienda, los comedores, los gastos de material y herramientas, así como cualquier otro ámbito que permita mejorar las condiciones del estudiantado.
                  La Universidad proporcionará al estudiantado información actualizada de su política de becas y ayudas y de las convocatorias de ayudas que se publiquen.
                  Del mismo modo, la Universidad de Sevilla podrá establecer programas que estimulen la captación y retención de estudiantes de especiales capacidades o rendimiento académico. Asimismo, establecerá un régimen de ayudas para los estudios de doctorado con el fin de estimular la función investigadora y la continuidad de la carrera académica.

                  3. Los centros establecerán mecanismos para la libre elección de grupo por parte de los y las estudiantes, en la medida en que lo permitan sus planes de organización docente.
                  4. Los sistemas de evaluación contemplarán la posibilidad de aprobar una asignatura por curso de manera previa a la prueba final, caso de que la hubiere. A tal fin, si el sistema de evaluación de una asignatura anual se basa, primordialmente, en exámenes, deberán realizarse dos exámenes parciales. La superación de estos y de los demás elementos de evaluación, en los términos contemplados en el sistema de evaluación y calificación de cada asignatura, conllevará el aprobado por curso.
                \end{scriptsize}
        \end{itemize}


\end{itemize}

\subsection{Categoría 6: Procedimental (43-50)}
\textit{Consultas basadas en guías (MATRUX, Liquidaciones, Bareprof) sobre cómo llevar a cabo un trámite paso a paso.}
\begin{itemize}
  \item \textbf{P43.} ¿Cuáles son los pasos requeridos según la guía breve de liquidaciones si soy personal de investigación e intento solicitar el abono y liquidación de los gastos de manutención tras mi vuelta? \\
        \begin{itemize}
          \item \textbf{Tipo de consulta:} simple
          \item \textbf{Respuesta del Agente:}

                ⚠️ Requisitos Previos:
                Antes de iniciar el trámite, asegúrate de haber tramitado la Comisión de Servicios o Autorización de Desplazamiento con su licencia por estudios a personal docente (PDI funcionario o laboral) o Autorización de Desplazamiento (Personal Investigador o Técnico Contratado). También debes comprobar que la actividad a desarrollar figura como un gasto elegible en la convocatoria del proyecto con el que pretende justificarla.

                A continuación, te presento los pasos para solicitar el abono y liquidación de los gastos de manutención tras tu vuelta:

                1. \textbf{Rellena el Formulario de liquidación de gastos de viaje}: Puedes encontrar el formulario en la web de la Universidad de Sevilla. Debes cumplimentar todos los campos requeridos, incluyendo los datos del beneficiario, del viaje y de los gastos.
                2. \textbf{Firma el formulario a través de Portafirmas}: Debes firmar el formulario junto con el IP (Investigador Principal) del proyecto.
                3. \textbf{Prepara la documentación justificativa}: Debes preparar la siguiente documentación:
                * \textbf{Comisión de Servicio o Autorización de desplazamiento}: Con todas las firmas requeridas.
                * \textbf{Memoria justificativa}: En la que describas las actividades realizadas y su relación con el proyecto.
                * \textbf{Documentación acreditativa de la actividad}: Certificado de asistencia, documento firmado por el responsable del centro, etc.
                * \textbf{Facturas abonadas y trazabilidad bancaria}: De los gastos de alojamiento y locomoción.
                * \textbf{Informe de conformidad}: Firmado por el IP, si procede.
                * \textbf{Documento de renuncia parcial}: Si has renunciado parcialmente a los gastos de la liquidación.
                4. \textbf{Organiza la documentación}: Debes organizar la documentación de tal modo que cada archivo corresponda a lo siguiente:
                * \textbf{Comisión de servicios u autorización de desplazamiento}
                * \textbf{Formulario de liquidación de dietas}
                * \textbf{Memoria justificativa}
                * \textbf{Resto de documentos} (agrupados en uno o dos archivos, sabiendo que cada archivo solo puede tener un tamaño máximo de 10 Mb)
                5. \textbf{Presenta la solicitud de abono}: Debes presentar la solicitud de abono a través del \textbf{Registro Electrónico General}, dirigido a la Subdirección General de Gestión Económica de la Investigación. Debes adjuntar todos los documentos justificativos requeridos.
                6. \textbf{Envía la documentación original}: Si tu documentación incluye originales en papel, debes enviarlos por correo interno a la Subdirección General de Gestión Económica, junto con una copia de la solicitud de registro presentada.

                Recuerda que cada entrada de Registro debe corresponder a una sola solicitud de abono, la cual solo puede incluir una sola liquidación de gastos y sus correspondientes documentos obligatorios.
          \item \textbf{Contexto utilizado:}
                \begin{scriptsize}
                  FUENTE: guia-breve-liquidacion.pdf | PÁGINA: 1
                  \# Subdirección de Gestión Económica de la Investigación

                  \#\# GUÍA BREVE
                  \#\#\# Solicitudes de liquidación de gastos de viajes

                  \#\#\# Introducción

                  En este documento se presenta una versión reducida de la Guía para la cumplimentación de solicitudes de abono de liquidaciones de gastos de viaje en el ámbito de la investigación, con el fin de facilitar los trámites necesarios para la cumplimentación del formulario de liquidación de gastos y de especificar la documentación necesaria para la realización y justificación de la actividad, así como el procedimiento de presentación de la solicitud de abono.

                  \#\#\# ANTES DE REALIZAR LA ACTIVIDAD

                  - Deberá tramitar la Comisión de Servicios con su licencia por estudios a personal docente (PDI funcionario o laboral) o Autorización de Desplazamiento (Personal Investigador o Técnico Contratado). Puede consultar detalles en el punto 4.1 de la guía completa.

                  - Deberá comprobar que la actividad a desarrollar figura como un gasto elegible en la convocatoria del proyecto con el que pretende justificarla, que la persona que la va a llevar a cabo guarda el vínculo que exige el organismo financiador, y que dispone de saldo suficiente en la orgánica del proyecto.

                  - Si va a alojarse haciendo uso de compañías que ofrecen servicios de reservas online, asegúrese de que podrá obtener factura por el gasto total realizado.

                  - Si va a reservar el medio de transporte y/o el alojamiento a través de agencias de viaje del acuerdo marco vigente, solicite la reserva de crédito correspondiente.

                  \#\#\# DURANTE LA ACTIVIDAD

                  - Conserve tarjetas de embarque si el viaje se ha realizado en avión.
                  - Pida certificados o documentos que acrediten su participación en la actividad.
                  - Pida facturas o tickets del alojamiento y/o medios de locomoción utilizados donde figuren sus datos fiscales, no los de la US. En facturas/tickets simplificados asegúrese de que aparecen la fecha y hora de emisión.

                  ---

                  FUENTE: guia-breve-liquidacion.pdf | PÁGINA: 3
                  \# Subdirección de Gestión Económica de la Investigación

                  \#\# 3. SOLICITUD DEL ABONO DE LA LIQUIDACIÓN DE GASTOS

                  A continuación, deberá proceder a la solicitud de abono de la liquidación de los gastos realizados mediante instancia efectuada a través del Registro Electrónico General, dirigido a la Subdirección General de Gestión Económica de la Investigación.

                  \textbf{IMPORTANTE:} Cada entrada de Registro debe corresponder a una sola solicitud de abono, la cual solo puede incluir una sola liquidación de gastos y sus correspondientes documentos obligatorios. Estos documentos son:

                  - Comisión Servicio o Autorización Desplazamiento con todas las firmas requeridas.
                  - Formulario de liquidación de gastos con las firmas requeridas.
                  - Memoria justificativa de la actividad realizada.
                  - Documentos acreditativos de la actividad realizada.
                  - Facturas abonadas y la trazabilidad bancaria de las mismas.
                  - Si se va a realizar un pago a terceros deberá adjuntar también el informe de conformidad firmado por el IP.
                  - Si ha solicitado renuncia parcial a parte del gasto deberá adjuntar también el documento de desglose de la renuncia en pdf.

                  Dado que actualmente el Registro Electrónico General solo permite subir un máximo de cinco archivos, organice la documentación de tal modo que cada archivo corresponda a lo siguiente:

                  1. Comisión de servicios u autorización de desplazamiento
                  2. Formulario de liquidación de dietas
                  3. Memoria justificativa
                  4. Resto de documentos (agrupados en uno o dos archivos, sabiendo que cada archivo solo puede tener un tamaño máximo de 10 Mb, y el tamaño máximo de los cinco ficheros es de 15 Mb)

                  Si su documentación incluye originales en papel deberá enviarlos por correo interno a la Subdirección General de Gestión Económica, junto con una copia de la solicitud de registro presentada.

                  \#\# CONSULTAS

                  Además de la guía completa, puede consultar también el documento de preguntas frecuentes (INCLUIR EL ENLACE DE LA WEB) elaborado por el Vicerrectorado de Investigación o dirigirse al personal del Servicio de Fiscalización y Pagos:

                  - Proyectos Estatales y Fundaciones:
                  - Cristina Guerra Aguilar-Galindo, cristina-ag@us.es 954 48 13 86
                  - Responsable Proyectos Comunidad Autónoma, Europeos o Internacionales:
                  - Rosa María López Valiente, rosi@us.es, 954 48 12 85

                  \textbf{VERSIÓN 2, OCTUBRE 2024        PÁGINA 3 | 4}

                  ---

                  FUENTE: guia-completa-liquidacion-de-viajes-v2.pdf | PÁGINA: 3
                  \# Subdirección de Gestión Económica de la Investigación

                  \#\# 1. INTRODUCCIÓN

                  Con el ánimo de que el proceso de fiscalización de su viaje sea más ágil y eficaz y de ayudar a los investigadores en la tramitación de las solicitudes de abono de liquidación de dietas y gastos de desplazamiento, se ha elaborado la presente guía en la que destacamos las principales cuestiones que merecen especial atención en el proceso de preparación y presentación de una solicitud de liquidación de dietas, teniendo en cuenta tanto las exigencias del Organismo o Entidad financiadora (a través de la convocatoria de ayuda que financia el proyecto al que se solicita imputar los gastos de la liquidación), como las de la propia normativa interna de ejecución del gasto de la Universidad de Sevilla.

                  La guía se ha organizado en tres bloques teniendo en cuenta que el procedimiento de solicitud electrónica exige la presentación de al menos los siguientes documentos:

                  - Impreso de liquidación de dietas.
                  - Documentación justificativa obligatoria:
                  - Comisión de servicio o autorización de desplazamiento aprobada, según proceda.
                  - Memoria justificativa del vínculo del gasto con el proyecto y su necesidad para el logro de objetivos o desarrollo del mismo.
                  - Evidencia de la realización del viaje.
                  - Facturas abonadas por el interesado, así como trazabilidad del pago de las mismas.

                  \#\# 2. CUESTIONES PRELIMINARES

                  Consideramos conveniente recordar que antes de realizar un viaje debe:

                  - Tramitar la correspondiente solicitud de comisión de servicios o autorización de desplazamiento, según sea PDI funcionario o laboral o Personal Investigador contratado.

                  Realizar un desplazamiento sin dicho documento retrasará la tramitación de la liquidación de gastos de viaje hasta que se disponga del mismo, al margen del perjuicio que le puede ocasionar que su viaje quede fuera de la cobertura derivada de su vínculo laboral con la Universidad de Sevilla en el ejercicio de sus funciones al no haber sido formalmente autorizado con carácter previo.

                  Por ello, al margen de que el viaje se desarrolle incluso en período de vacaciones, fines de semana, festivos o bien, incluso, durante el disfrute de una licencia por estudios, debe tramitar la correspondiente comisión de servicios o autorización de desplazamiento con cargo al proyecto.

                  - Comprobar que los gastos del viaje son elegibles para el proyecto al que solicitará su imputación, es decir, que cuenta con la fuente de financiación apropiada. Debe verificar que:
                  - La actividad o evento que motiva el viaje supone un beneficio para el desarrollo del proyecto, y consecuentemente, para el logro de los objetivos, hitos o líneas de investigación contenidos en el mismo.
                  - La actividad o evento se lleva a cabo dentro del plazo de ejecución del proyecto.
                  - La persona que realiza el viaje guarda el vínculo que exige el organismo financiador, conforme a las normas reguladoras de la convocatoria.
                  - El proyecto tiene crédito disponible suficiente y adecuado.

                  VERSIÓN 2 OCTUBRE 2024    PÁGINA 3 | 21

                  ---

                  FUENTE: guia-breve-liquidacion.pdf | PÁGINA: 2
                  \# Subdirección de Gestión Económica de la Investigación

                  \#\# AL TERMINAR LA ACTIVIDAD

                  \#\#\# 1. LIQUIDACIÓN DE GASTOS

                  Rellene el Formulario (INCLUIR EL ENLACE DE LA WEB) de liquidación de gastos de viaje. En la Guía (INCLUIR EL ENLACE DE LA WEB) para la cumplimentación de solicitudes de abono encontrará la información detallada para la correcta cumplimentación de este formulario.

                  Una vez terminado, debe ser firmado a través de Portafirmas, tanto por el IP como por la persona que realiza la actividad.

                  Podrá renunciar parcialmente a los gastos de la liquidación tal y como se le especifica en detalle en el punto 3.4 de la guía completa.

                  \#\#\# 2. DOCUMENTACIÓN ADICIONAL A PRESENTAR

                  Además de la Comisión de Servicio/Autorización de desplazamiento y del formulario de liquidación de gastos será necesario presentar la siguiente documentación:

                  - Una memoria justificativa en la que describa las actividades realizadas indicando la relación de éstas con alguno/s de los objetivos del proyecto. También deberá especificar el vínculo con el proyecto de la persona que realiza el viaje, así como su situación laboral. Puede consultar todos los detalles necesarios en el punto 4.2 de la guía completa.

                  - Documentación acreditativa de:
                  - La actividad realizada (certificado de asistencia, documento firmado por responsable del centro donde haya tenido lugar la estancia, etc.). Ver detalles en punto 4.3 de la guía completa.
                  - Los gastos de alojamiento y/o locomoción efectuados por la persona que ha realizado la actividad (facturas, tickets, billetes electrónicos, etc.), y trazabilidad bancaria de los mismos. Ver detalles en punto 4.4 de la guía completa.
                  - Si procede, informe de conformidad firmado por el IP de las facturas de alojamiento y/o locomoción derivadas de reservas de crédito a través de agencias de viaje del acuerdo marco vigente.
                  - Si procede, documento de desglose de la renuncia parcial de los gastos de liquidación en pdf.

                  VERSIÓN 2, OCTUBRE 2024    PÁGINA 2 | 4

                  ---

                  FUENTE: guia-completa-liquidacion-de-viajes-v2.pdf | PÁGINA: 2
                  \# Subdirección de Gestión Económica de la Investigación

                  \#\# INDICE

                  1. Introducción
                  2. Cuestiones preliminares
                  3. ¿Cómo preparo mi liquidación de gastos de viaje?
                  3.1. ¿Dónde descargar el formulario de “liquidación de gastos de viaje”?
                  3.2. ¿Cómo selecciono y/o averiguo el decreto de “dietas y gastos de locomoción” aplicable al proyecto?
                  3.3. ¿Qué datos del beneficiario y del viaje debo cumplimentar?
                  3.3.1. Titular beneficiario y datos de la actividad motivo del viaje
                  3.3.2. Datos necesarios para el cálculo de dietas de alojamiento y manutención
                  3.3.3. Datos de locomoción
                  3.3.4. Datos de inscripción u otros gastos relacionados con la actividad
                  3.3.5. Datos sobre anticipo de viajes (en su caso)
                  3.3.6. Datos de la cuenta bancaria donde abonarle los gastos
                  3.3.7. Datos de la orgánica y proyecto que financia la liquidación
                  3.4. ¿Cómo renuncio parcialmente a los gastos de la liquidación (en su caso)?
                  3.5. Guardar y firmar la liquidación
                  4. ¿Qué documentos justificativos obligatorios debo adjuntar?
                  4.1. Comisión de servicio o autorización de desplazamiento
                  4.2. Memoria justificativa
                  4.3. Evidencia de la realización de la actividad (motivo del viaje)
                  4.4. Facturas abonadas por el interesado, así como justificantes de los pagos efectuados
                  4.5. Informes de conformidad de facturas derivadas de reservas de crédito o de cargos internos por alojamiento en residencias universitarias de la Universidad de Sevilla
                  4.6. Documento de renuncia parcial, en su caso
                  5. ¿Dónde y cómo presento la solicitud de abono, junto con la liquidación y documentación obligatoria justificativa?

                  VERSIÓN 2 OCTUBRE 2024    PÁGINA 2 | 21
                \end{scriptsize}
        \end{itemize}


  \item \textbf{P44.} Siguiendo totalmente e indicando paso a paso la Guía del Candidato en Bareprof, ¿cómo registro primeramente una nueva candidatura de participación en el concurso público? \\
        \begin{itemize}
          \item \textbf{Tipo de consulta:} simple
          \item \textbf{Respuesta del Agente:}

                ⚠️ Requisitos Previos:
                Antes de comenzar el proceso de inscripción, asegúrese de tener acceso a la plataforma \textbf{Bareprof} y de haber leído la Guía del Candidato para entender los requisitos y el proceso de inscripción.

                Para registrar una nueva candidatura de participación en el concurso público, siga los siguientes pasos:

                1. \textbf{Acceda a la plataforma Bareprof}: Ingrese a la plataforma \textbf{Bareprof} y acceda a la pantalla inicial, donde se muestran las plazas que se encuentran en período de inscripción.
                2. \textbf{Seleccione la plaza de interés}: Busque la plaza que desea concursar y pulse el botón \textbf{Inscribir} situado en la parte derecha de la pantalla.
                3. \textbf{Verifique los datos de la plaza}: Se abrirá una ventana con los datos de la plaza y un botón para llevar a cabo la inscripción y generar la correspondiente \textbf{Solicitud de Inscripción}.
                4. \textbf{Genere la Solicitud de Inscripción}: Pulse el botón \textbf{Inscripción y generar solicitud} para generar un documento en formato \textbf{PDF} que contiene la \textbf{Solicitud de Inscripción}.
                5. \textbf{Firme y guarde la Solicitud de Inscripción}: Firme la \textbf{Solicitud de Inscripción} y guárdela en su dispositivo, ya que deberá subirla a la plataforma \textbf{Bareprof} más adelante.
                6. \textbf{Inicie el proceso de inscripción}: Regrese a la pantalla inicial de \textbf{Bareprof} y pulse el botón \textbf{Iniciar inscripción} para comenzar el proceso de inscripción.
                7. \textbf{Suba la documentación requerida}: Suba la \textbf{Solicitud de Inscripción} firmada y la documentación requerida (como el \textbf{CVN} y el \textbf{Título académico}) a la plataforma \textbf{Bareprof}.
                8. \textbf{Verifique y complete la información}: Verifique que la información sea correcta y complete cualquier campo obligatorio que falte.

                Recuerde que la inscripción en el concurso es un proceso abierto que solo finaliza cuando concluye el plazo de presentación de solicitudes. Por lo tanto, puede modificar los méritos y documentación aportados cuantas veces desee dentro del plazo de presentación de solicitudes.
          \item \textbf{Contexto utilizado:}
                \begin{scriptsize}
                  FUENTE: 27830\_guia-ayuda-candidato-r04-3-.pdf | PÁGINA: 9
                  \# Bareprof Guía de ayuda al candidato

                  \#\# 3.3 Inscripción en el concurso

                  Al acceder a la plataforma Bareprof, se muestran en la pantalla inicial las plazas que se encuentran en periodo de inscripción, con indicación del número de días que faltan para su conclusión.

                  \#\#\# Plazas en periodo de inscripción

                  Mostrando 1-4 de 4 elementos.

                  \vspace{0.4cm}
                  \begin{table}[h!]\centering\renewcommand{\arraystretch}{1.3}
                    \begin{tabular}{|p{0.45\textwidth}|p{0.45\textwidth}|p{0.45\textwidth}|p{0.45\textwidth}|p{0.45\textwidth}|p{0.45\textwidth}|}\hline
                      Tipo de plaza     & Código        & Departamento                              & Área                                      & Descripción & Días Bases \\\hline
                      Profesor Ayudante & AYD-2022-0001 & ARQUITECTURA Y TECNOLOGÍA DE COMPUTADORES & ARQUITECTURA Y TECNOLOGÍA DE COMPUTADORES & 5           &            \\\hline
                      Profesor Ayudante & AYD-2022-0002 & DIBUJO                                    &                                           & 5           & Inscribir  \\\hline
                      Profesor Ayudante & AYD-2022-0003 & FILOLOGÍAS INTEGRADAS                     & ESTUDIOS ÁRABES E ISLÁMICOS               & 5           & A0         \\\hline
                    \end{tabular}\end{table}
                  \vspace{0.4cm}

                  Para realizar la inscripción en una plaza, se debe pulsar el botón Inscribir situado en la parte derecha. Se accede entonces a una ventana en la que se muestran los datos de la plaza y un botón para llevar a cabo la inscripción y generar la correspondiente solicitud.

                  \#\#\# Solicitud de Inscripción

                  Al pulsar el botón Inscripción y generar solicitud se generará un documento en formato pdf que contiene la solicitud de participación, en el que aparecen los datos personales del participante (que pueden modificarse y completarse según se indica en el apartado 3.1) y los datos de la plaza. Tal solicitud deberá ser firmada por el candidato y presentada por registro siguiendo las instrucciones dadas en la convocatoria del concurso público. Una vez generada la solicitud, se desbloquean los restantes pasos del proceso de inscripción en el concurso: requisitos de admisión, justificación de méritos y datos adicionales (solo para convocatorias de plazas de Profesor Ayudante Doctor).

                  Revisión 4 (marzo, 2025) 8

                  ---

                  FUENTE: 27830\_guia-ayuda-candidato-r04-3-.pdf | PÁGINA: 4
                  \# Guía de ayuda al candidato

                  \#\# 1. Introducción

                  La Universidad de Sevilla cuenta con una plataforma web para la baremación de los candidatos participantes en los concursos públicos que se convoquen para la contratación de Profesorado Ayudante Doctor, Profesorado Asociado y Profesorado Sustituto. Dicha plataforma, denominada Bareprof, permite gestionar de forma integrada todo el proceso de admisión al concurso y posterior valoración de méritos. A ella tienen acceso los candidatos que participan en tales convocatorias, las comisiones de contratación encargadas de baremar las solicitudes de los candidatos y los técnicos del Vicerrectorado de Personal Docente e Investigador responsables de la gestión de las convocatorias. El candidato hará uso de la plataforma Bareprof para inscribirse en la plaza o plazas a las que desee concursar. Al acceder a la plataforma, el candidato tendrá acceso a la información de las plazas que se encuentran en periodo de inscripción, con indicación del número de días restantes para la conclusión del plazo. El proceso de inscripción en una plaza por medio de Bareprof se describe en el apartado 3 de esta guía.

                  Durante el proceso de inscripción, el candidato debe presentar, junto a otra documentación requerida por la convocatoria, aquellos méritos que desea sean baremados por la Comisión de Contratación de la plaza a la que opta. La introducción de méritos por parte del candidato en Bareprof se realiza de forma indirecta. Así, los méritos deberán ser previamente registrados en el Curriculum Vitae Normalizado del candidato, cuya creación y mantenimiento se realiza a través del editor CVN de la Fundación Española para la Ciencia y la Tecnología (FECYT), dependiente del Ministerio de Ciencia, Innovación y Universidades. En el apartado 4 de esta guía se proporciona una breve descripción del proceso de creación del CVN.

                  Cada mérito contemplado en el baremo de la Universidad de Sevilla debe ser introducido en el CVN siguiendo las instrucciones dadas en el apartado 5 de esta guía. Este es un proceso que requiere especial atención por parte del candidato y en el que es muy importante que se sigan de forma precisa las indicaciones dadas para cada mérito, ya que, de lo contrario, el mérito puede llegar a no ser baremado por la Comisión al no ser reconocido por Bareprof.

                  \#\# 2. Documentos de referencia

                  Se relacionan a continuación los principales documentos relacionados con el procedimiento de baremación para la contratación de Profesorado Ayudante Doctor, Profesorado Asociado y Profesorado Sustituto, que se encuentran disponibles en el Boletín Oficial de la Universidad de Sevilla (http://www.us.es/bous). Deberán además tenerse en cuenta todas aquellas instrucciones aclaratorias del baremo que pudieran dictarse.

                  - Normativa sobre Procedimiento de Contratación de Profesorado Ayudante Doctor, Profesorado Asociado y Profesorado Interino, junto con el Baremo establecido para dichas figuras contractuales (BOUS 12/2024, de 11 de octubre).
                    [https://www.us.es/bous-numeros/numero-12-11-de-octubre-de-2024](https://www.us.es/bous-numeros/numero-12-11-de-octubre-de-2024)

                  - Instrucción de 17-10-19 sobre los requisitos para el reconocimiento de las estancias de investigación dentro del mérito preferente en los concursos de profesorado ayudante doctor, asociado y sustituto interino (BOUS 14/2019, de 21 de noviembre).
                    [https://www.us.es/bous-numeros/numero-14-21-de-noviembre-de-2019](https://www.us.es/bous-numeros/numero-14-21-de-noviembre-de-2019)

                  - Resolución rectoral, de 24 marzo de 2021, por la que se emite una interpretación sobre la determinación del umbral mínimo aplicable en los procedimientos selectivos para plazas de.

                  ---

                  FUENTE: 27831\_instrucciones-baremacinin-r03.pdf | PÁGINA: 4
                  \# Bareprof
                  \#\# Instrucciones de baremación

                  \#\#\# 1. Introducción

                  La Universidad de Sevilla cuenta con una plataforma web para la baremación de los candidatos participantes en los concursos públicos que se convoquen para la contratación de Profesorado Ayudante Doctor, Profesorado Asociado y Profesorado Sustituto. Dicha plataforma, denominada Bareprof, permite gestionar de forma integrada todo el proceso de admisión al concurso y posterior valoración de méritos. A ella tienen acceso los candidatos que participan en tales convocatorias, las comisiones de contratación encargadas de baremar las solicitudes de los candidatos y los técnicos del Vicerrectorado de Personal Docente e Investigador responsables de la gestión de las convocatorias.

                  El presente documento pretende servir como guía para la constitución de las Comisiones de Contratación en la plataforma Bareprof y el posterior desarrollo del proceso de baremación. La descripción de las acciones que deben llevar a cabo los miembros de las comisiones en la plataforma se completa en el documento con una serie de consideraciones sobre la baremación, que tienen como objetivo aclarar algunos aspectos sobre la misma y proporcionar pautas de actuación a las comisiones; se trata de dar respuesta así a aquellas cuestiones que han sido objeto de consulta frecuente en anteriores convocatorias. Este documento se irá completando a medida que se produzcan nuevas consultas que se consideren de interés general.

                  Cada candidato hace uso de la plataforma Bareprof para inscribirse en la plaza o plazas a las que desea concursar. Al acceder a la plataforma, el candidato tiene acceso a la información de las plazas que se encuentran en periodo de inscripción y puede iniciar el proceso en aquellas que sean de su interés. Durante el proceso de inscripción, el candidato debe presentar, junto a otra documentación requerida por la convocatoria, aquellos méritos que desea que sean baremados por la comisión de contratación de la plaza a la que opta. La introducción de méritos por parte del candidato en Bareprof se realiza de forma indirecta. Así, los méritos deberán ser previamente registrados en el Curriculum Vitae Normalizado del candidato, cuya creación y mantenimiento se realiza a través del editor CVN de la Fundación Española para la Ciencia y la Tecnología (FECYT), dependiente del Ministerio de Ciencia, Innovación y Universidades.

                  Cada mérito contemplado en el baremo de la Universidad de Sevilla es introducido por el candidato en su CVN siguiendo las instrucciones dadas en la Guía de Ayuda al Candidato. Este es un proceso que requiere especial atención por parte del aspirante y en el que es muy importante que se sigan de forma precisa las indicaciones dadas para cada mérito, ya que, de lo contrario, este puede llegar a no ser baremado por la comisión al no ser reconocido por Bareprof.

                  \#\#\# 2. Documentos de referencia

                  Se relacionan a continuación los principales documentos relacionados con el procedimiento de baremación para la contratación de Profesorado Ayudante Doctor, Profesorado Asociado y Profesorado Sustituto, que se encuentran disponibles en la web del Boletín Oficial de la Universidad de Sevilla (http://www.us.es/bous). Deberán además tenerse en cuenta todas aquellas instrucciones aclaratorias del baremo que pudieran dictarse.

                  - Normativa sobre Procedimiento de Contratación de Profesorado Ayudante Doctor, Profesorado Asociado y Profesorado Interino, junto con el Baremo establecido para dichas figuras contractuales (BOUS 12/2024, de 11 de octubre).
                    [https://www.us.es/bous-numeros/numero-12-11-de-octubre-de-2024](https://www.us.es/bous-numeros/numero-12-11-de-octubre-de-2024)

                  Revisión 3 (diciembre, 2024)

                  ---

                  FUENTE: 27830\_guia-ayuda-candidato-r04-3-.pdf | PÁGINA: 10
                  \# Guía de ayuda al candidato

                  \#\# 3.4 Requisitos de admisión

                  En este paso del proceso de inscripción, el candidato deberá subir a la plataforma Bareprof los documentos (en formato pdf) que acreditan el cumplimiento de los requisitos de admisión establecidos en la convocatoria del concurso público.

                  \vspace{0.4cm}
                  \begin{table}[h!]\centering\renewcommand{\arraystretch}{1.3}
                    \begin{tabular}{|p{0.45\textwidth}|p{0.45\textwidth}|p{0.45\textwidth}|p{0.45\textwidth}|p{0.45\textwidth}|p{0.45\textwidth}|}\hline
                      D                        & B                                & a                                        & T                              &                        &          \\\hline
                      Solicitud de inscripción & REQUISITOS DE ADMISIÓN           & Justificación de méritos                 & Datos adicionales              &                        &          \\\hline
                      CVN (pdf max: 2 MB)      & Título académico (pdf max: 2 MB) & Solicitud de inscripción (pdf max: 2 MB) & Pago de tasas (pdf max: 1 MB)¹ &                        &          \\\hline
                      Seleccionar archivo...   & Examinar                         & Seleccionar archivo...                   & Examinar                       & Seleccionar archivo... & Examinar \\\hline
                    \end{tabular}\end{table}
                  \vspace{0.4cm}

                  - CVN: Una vez seleccionado el archivo que contiene el CVN, generado por medio del editor CVN-FECYT, Bareprof procesará la información que contiene y mostrará finalmente el Visor de CVN, en el que podrá ver los méritos leídos por la aplicación y los errores detectados en el proceso.

                  Visor de CVN

                  A continuación se muestran los datos que usted ha introducido en los campos leídos por Bareprof, solo aquellos que estén enmarcados en color rojo se pondrán a disposición de las comisiones de contratación. Solo los datos leídos por Bareprof serán evaluados por las comisiones de contratación. Recuerde también que los méritos incluidos en apartados del baremo a los que no corresponden por su tipología, no serán evaluados por las comisiones de contratación, procediendo a su no validación. Se le recuerda que es necesario que siga detenidamente las instrucciones del baremo y el documento de Preguntas Frecuentes en la elaboración de su CVN.

                  - Méritos con campos obligatorios sin cumplimentar
                  - Méritos con campos de texto libre no puntuables
                  - Méritos con información no indexada

                  - Rendimiento académico
                  - Actividad docente
                  - Actividad investigadora
                  - Transferencia del conocimiento
                  - Actividad profesional externa
                  - Mérito preferente
                  - Resumen CV
                  - Resumen lectura CVN
                  - Expediente académico
                  - Titulación universitaria (grado o equivalente)

                  \#\#\# Mérito

                  - Nombre del título: Ingeniero Industrial Especialidad Electricidad
                  - Titulación universitaria: Titulado Superior
                  - Nota media del expediente:
                  - Fecha de titulación: 17/06/2019

                  \#\#\# Titulación de máster universitario oficial

                  - Mérito
                  - Tipo de formación: Máster
                  - Titulación de posgrado: Máster Universitario en Ingeniería de la Energía
                  - Fecha de titulación: 22/06/2020
                  - Calificación obtenida: 300 horas

                  Los méritos que aparecen recuadrados en rojo contienen algún campo obligatorio que no ha sido cumplimentado, marcándose también en ese color los campos en los que se han detectado las incidencias. Es muy importante que proceda a su corrección y vuelva a subir el CVN a Bareprof, ya que estos méritos no serán puestos a disposición de la comisión para su baremación. En color azul aparecen recuadrados los méritos que contienen opciones o literales no admitidos en alguno de sus campos, por lo que podrían no ser valorados por la comisión de contratación.

                  Revisión 4 (marzo, 2025) 9

                  ---

                  FUENTE: 27830\_guia-ayuda-candidato-r04-3-.pdf | PÁGINA: 14
                  \# Bareprof Guía de ayuda al candidato

                  \#\# 3.7 Finalización del proceso de inscripción

                  La inscripción en el concurso es un proceso abierto que solo finaliza cuando concluye el plazo de presentación de solicitudes. Esto permite al candidato modificar los méritos y documentación aportados cuantas veces desee dentro del plazo de presentación de solicitudes. Por este motivo, no existe en la plataforma de baremación ningún elemento (botón, opción de menú o similar) para finalizar el procedimiento de inscripción.

                  Recuerde que la solicitud generada por Bareprof debe remitirse obligatoriamente a la Universidad de Sevilla utilizando el registro electrónico dentro del plazo de presentación de solicitudes, según se indica en la convocatoria. No basta con firmar el documento y subirlo a la plataforma Bareprof. Solo se debe enviar por registro la solicitud y NO el resto de documentación aportada en el proceso de inscripción al concurso. El resguardo de presentación en registro se debe incorporar a Bareprof en el apartado Requisitos de admisión (Solicitud de inscripción).

                  Revisión 4 (marzo, 2025) 13
                \end{scriptsize}
        \end{itemize}


  \item \textbf{P45.} ¿Cuál es el procedimiento y los pasos que debo dar según las normas para solicitar de forma justificada en Secretaría la famosa anulación total de matrícula? \\
        \begin{itemize}
          \item \textbf{Tipo de consulta:} simple
          \item \textbf{Respuesta del Agente:}

                ⚠️ Requisitos Previos:
                Antes de solicitar la anulación total de matrícula, es importante tener en cuenta los siguientes requisitos previos:
                - Debes haber formalizado tu matrícula en el periodo ordinario.
                - Debes cumplir con los plazos establecidos para solicitar la anulación total de matrícula, que varían según la titulación y el inicio de la docencia.
                - Debes tener una justificación válida para solicitar la anulación total de matrícula, como enfermedad grave sobrevenida, obtención de plaza en otra universidad, etc.

                A continuación, te presento los pasos para solicitar la anulación total de matrícula:

                1. \textbf{Revisa los plazos}: Verifica los plazos establecidos para solicitar la anulación total de matrícula, que varían según la titulación y el inicio de la docencia. Por ejemplo, para titulaciones de grado, el plazo es hasta el 16 de septiembre de 2025, mientras que para titulaciones de máster universitario que inicien la docencia el 6 de octubre, el plazo es hasta el 15 de octubre de 2025.
                2. \textbf{Prepara la documentación}: Asegúrate de tener toda la documentación necesaria para solicitar la anulación total de matrícula, como:
                * \textbf{Solicitud de anulación de matrícula}: Debes presentar una solicitud de anulación de matrícula dirigida al Decano/Director del centro responsable del título oficial de los estudios donde formalizaste tu matrícula.
                * \textbf{Justificación}: Debes aportar una justificación válida para solicitar la anulación total de matrícula, como un informe médico en caso de enfermedad grave sobrevenida.
                * \textbf{Documentación adicional}: Puedes necesitar aportar documentación adicional, como la \textbf{carta de aceptación} en otra universidad o el \textbf{documento de identidad}.
                3. \textbf{Presenta la solicitud}: Presenta la solicitud de anulación de matrícula y la documentación necesaria en la \textbf{Secretaría} del centro responsable del título oficial de los estudios donde formalizaste tu matrícula.
                4. \textbf{Espera la resolución}: Espera a que la \textbf{Secretaría} resuelva tu solicitud de anulación de matrícula. La resolución puede tardar varios días o semanas, dependiendo de la complejidad del caso.
                5. \textbf{Recibe la resolución}: Recibe la resolución de tu solicitud de anulación de matrícula. Si la resolución es favorable, se te devolverán los importes abonados y se te anulará la matrícula.

                Recuerda que es importante seguir los pasos establecidos y presentar toda la documentación necesaria para evitar retrasos o problemas en el proceso de anulación de matrícula. Si tienes alguna duda, no dudes en consultar con la \textbf{Secretaría} del centro responsable del título oficial de los estudios donde formalizaste tu matrícula.
          \item \textbf{Contexto utilizado:}
                \begin{scriptsize}
                  FUENTE: R.R. NORMAS DE MATRICULA GRADO Y MASTER UNIVERSITARIO 2025-2026.pdf | PÁGINA: 11
                  \# UNIVERSIDAD D SEVII

                  C)EST

                  5. No se concederán anulaciones parciales de matrícula a solicitud del estudiante, excepto por coincidencia no prevista de horarios de clase y exámenes y las situaciones contempladas en estas normas. No obstante, por causas académicas sobrevenidas que imposibiliten la prestación del servicio, los centros podrán realizar de oficio anulaciones parciales de matrícula, ofreciendo a los estudiantes la posibilidad de matricular otras asignaturas o bien la devolución de los importes abonados en los créditos anulados, en su caso.

                  6. La matrícula se anulará de oficio cuando se trate de error, omisión, falseamiento de los datos consignados en esta o cuando se modifique la situación administrativa o académica del estudiante, o las matrículas condicionadas que no deban producir efectos.

                  Asimismo se anularán de oficio las matrículas cuando el estudiante no aporte la documentación requerida o por falta de pago.

                  7. Los centros anularán de oficio las matrículas efectuadas por estudiantes que se matriculen en otra titulación por haber obtenido plaza en los procedimientos de admisión para estudiantes de nuevo ingreso en esta u otra universidad pública, siempre que no se acredite una autorización para la simultaneidad de estudios.

                  \#\# Artículo 11. Supuestos de anulación de matrícula y sus efectos

                  1. Se podrá solicitar la anulación de matrícula ordinaria al inicio del curso académico en los siguientes plazos:

                  - Titulaciones de grado: hasta el 16 de septiembre de 2025.
                  - Titulaciones de máster universitario que inicien la docencia el 6 de octubre: hasta el 15 de octubre de 2025.
                  - Titulaciones de máster universitario que inicien la docencia el 22 de septiembre: hasta el 30 de septiembre de 2025.
                  - Cuando la matrícula se haya formalizado con posterioridad al inicio del periodo docente del primer cuatrimestre: 7 días hábiles a contar desde el día de la finalización del periodo en el que haya realizado su matrícula.

                  La concesión de la anulación supondrá:

                  a) La devolución de los importes abonados, en su caso.

                  b) La matrícula no se computará a efectos de convocatorias de examen ni de permanencia del estudiante, ni se contabilizará para futuras matrículas a efectos de aplicación de recargos.

                  c) La pérdida de la plaza adjudicada si el estudiante accedió por el procedimiento de admisión en el curso 2025-2026.

                  2. La anulación de la matrícula por haber obtenido plaza en esta u otra universidad pública (en el procedimiento o sistema de admisión que tenga establecido) en otra titulación de grado, máster universitario o estudios de doctorado donde hubiese sido admitido el estudiante con posterioridad a su matriculación en la Universidad de Sevilla, no tendrá ningún efecto académico ni económico si se solicita con anterioridad al 31 de octubre de 2025, en el caso de titulaciones de grado, y al 31 de diciembre de 2025 en las titulaciones de máster universitario.

                  Código Seguro De Verificación: `eqOjhZcmL0p6VJlZiTH/Yw==`
                  Fecha: 27/06/2025
                  Firmado Por: MIGUEL ANGEL CASTRO ARROYO
                  Url De Verificación: [https://pfirma.us.es/verifirma/code/eqOjhZcmL0p6VJlZiTH\%2FYw\%3D\%3D](https://pfirma.us.es/verifirma/code/eqOjhZcmL0p6VJlZiTH\%2FYw\%3D\%3D)
                  Página: 7/50

                  ---

                  FUENTE: R.R. NORMAS DE MATRICULA GRADO Y MASTER UNIVERSITARIO 2025-2026.pdf | PÁGINA: 10
                  \# UNIVERSIDAD DE SEVILLA

                  C)EST

                  5. En el acto de ampliación de matrícula serán de aplicación las reducciones de precios públicos previstas en el artículo 28 de estas normas siempre que se haya acreditado el derecho a la deducción o reducción de precios en el periodo ordinario, o bien, podrán aplicarse nuevas deducciones en el periodo de ampliación acreditadas conforme establece el citado artículo.

                  6. La falta de pago de la ampliación de matrícula determinará la anulación de la ampliación.

                  7. Con carácter general, la anulación de la matrícula ordinaria determinará también la anulación de la ampliación de matrícula, en su caso.

                  \#\# Artículo 9. Bloqueo de expediente

                  Como actuación previa a la anulación de matrícula, el centro bloqueará el expediente de aquellos estudiantes que no hayan aportado la documentación establecida en estas normas, o que se encuentren en situación de impago, una vez notificado el requerimiento correspondiente o cuando se produzca el rechazo automático del mismo por el transcurso del plazo establecido.

                  Este bloqueo tendrá las siguientes consecuencias sobre el expediente de la titulación a la que corresponda:

                  - No podrán emitirse certificaciones académicas.
                  - El estudiante no podrá solicitar su título universitario oficial.
                  - El estudiante no tendrá acceso a la plataforma de Enseñanza Virtual de la Universidad de Sevilla.

                  El estudiante dispondrá de 15 días hábiles siguientes a la notificación del requerimiento, o en su caso, desde que se produzca el rechazo automático de la misma por el transcurso del plazo establecido, para aportar la documentación requerida o abonar las cantidades pendientes, según proceda, en cuyo caso el centro procederá al desbloqueo de su expediente.

                  En el supuesto de que el estudiante no aporte la documentación requerida en el plazo indicado o no abone las cantidades pendientes de pago, el centro procederá a la anulación de su matrícula en el curso actual, con los efectos que se establecen en el artículo 11 de estas normas.

                  \#\# Artículo 10. Anulación de matrícula

                  1. La anulación de matrícula supone el cese de los efectos académicos de la matrícula realizada con la consiguiente pérdida de los derechos a evaluación, manteniéndose la obligación de satisfacer los precios públicos salvo que proceda su devolución en los términos previstos reglamentariamente.

                  2. La anulación de la matrícula en un curso conllevará la anulación de los créditos o asignaturas reconocidos en ese curso sin derecho a devolución de los importes satisfechos, lo que no impedirá que el estudiante pueda volver a solicitarlo en cursos posteriores conforme a los procedimientos correspondientes.

                  3. Las solicitudes de anulación de matrícula se dirigirán al Decano/Director del centro responsable del título oficial de los estudios donde formalice su matrícula por cualquiera de los medios establecidos en las disposiciones vigentes, haciendo constar la causa por la que se solicita, y serán resueltas en primera instancia por el centro.

                  4. No se concederán anulaciones de matrícula a los estudiantes con alguna calificación incorporada en actas ya cerradas definitivamente, en alguna de las asignaturas matriculadas en el curso académico para el que solicita la anulación.

                  Código Seguro De Verificación: eqOjhZcmL0p6VJlZiTH/Yw==
                  Fecha: 27/06/2025
                  Firmado Por: MIGUEL ANGEL CASTRO ARROYO
                  Url De Verificación: [https://pfirma.us.es/verifirma/code/eqOjhZcmL0p6VJlZiTH\%2FYw\%3D\%3D](https://pfirma.us.es/verifirma/code/eqOjhZcmL0p6VJlZiTH\%2FYw\%3D\%3D)
                  Página: 6/50

                  ---

                  FUENTE: R.R. NORMAS DE MATRICULA GRADO Y MASTER UNIVERSITARIO 2025-2026.pdf | PÁGINA: 13
                  \# UNIVERSIDAD D SEVII

                  C)EST

                  c) No se aplicarán recargos en siguientes cursos respecto a las asignaturas anuladas.

                  d) En el supuesto de obtención del visado con posterioridad al inicio del curso académico, podrá solicitar, previa aportación de la documentación acreditativa de la fecha de obtención de este, la anulación parcial de aquellas asignaturas cuya docencia se haya impartido total o parcialmente con anterioridad a la obtención del visado, salvo que hubiera alguna calificación incorporada en acta.

                  8. Excepcionalmente, el estudiante podrá solicitar la anulación total o parcial de matrícula de asignaturas no calificadas en acta por enfermedad grave sobrevenida, documentalmente acreditada, que le impida continuar los estudios.

                  Por enfermedad grave se entenderá cualquier dolencia o lesión física o psíquica que incapacite temporalmente para la ocupación o actividad habitual del estudiante durante un periodo continuado mínimo de tres meses en el curso académico para el que solicita la anulación y que requiera tratamiento o intervención clínica de cirugía mayor. El solicitante deberá aportar informe clínico expedido por el sistema público de salud español. Del mismo modo, en caso de pertenencia obligatoria a una mutualidad, el informe será expedido por la misma.

                  Las solicitudes se presentarán por cualquiera de los medios establecidos en las disposiciones vigentes dirigidas al Vicerrectorado de Estudiantes. La entrega por el estudiante de la documentación médica se podrá realizar en sobre cerrado e irá acompañada de un consentimiento expreso para el tratamiento de los datos incluidos en esta documentación, cuyo fin exclusivo será la gestión por la Universidad de Sevilla de su solicitud de anulación de matrícula por enfermedad grave sobrevenida. Por el Vicerrectorado de Estudiantes se remitirá dicha documentación a la Comisión Técnica de Expertos para la Valoración de los Supuestos Determinantes de la Anulación de la Matrícula por Enfermedad Grave Sobrevenida, comisión de carácter médico, para la emisión del correspondiente informe. La referida comisión podrá recabar cuantos otros informes estimen oportunos.

                  En el supuesto de resolverse favorablemente la solicitud de anulación de matrícula, los efectos serán:

                  a) La devolución del importe abonado en concepto de precios públicos por servicios académicos de las asignaturas anuladas, a excepción del importe de los precios públicos por servicios administrativos y el seguro escolar.

                  b) No se computará la matrícula a efectos de convocatorias de examen ni de permanencia del estudiante, ni se contabilizará para futuras matrículas a efectos de aplicación de recargos.

                  c) El estudiante no perderá la plaza adjudicada en el procedimiento de admisión en el curso 2025-2026.

                  9. En los supuestos de anulación de oficio por falta de algún requisito no subsanable y no imputable al interesado, procederá, en su caso, la devolución de los importes abonados, en su caso. Si el estudiante accedió por el procedimiento de admisión en el curso 2025-2026, perderá la plaza adjudicada.

                  10. Cuando la anulación de matrícula sea consecuencia de sanción disciplinaria:

                  a) No procederá la devolución de los importes abonados en concepto de precios públicos por servicios académicos y administrativos y el seguro escolar, ni se eximirá de la obligación de satisfacer las cantidades pendientes de pago.

                  b) Serán de aplicación los recargos económicos en la siguiente matrícula que realice en los mismos estudios.

                  Código Seguro De Verificación: eqOjhZcmL0p6VJlZiTH/Yw==
                  Fecha: 27/06/2025
                  Firmado Por: MIGUEL ANGEL CASTRO ARROYO
                  Url De Verificación: [https://pfirma.us.es/verifirma/code/eqOjhZcmL0p6VJlZiTH\%2FYw\%3D\%3D](https://pfirma.us.es/verifirma/code/eqOjhZcmL0p6VJlZiTH\%2FYw\%3D\%3D)
                  Página: 9/50

                  ---

                  FUENTE: R.R. NORMAS DE MATRICULA GRADO Y MASTER UNIVERSITARIO 2025-2026.pdf | PÁGINA: 12
                  \# UNIVERSIDAD DE SEVILLA

                  \#\# C)EST

                  En este supuesto procederá:

                  1. La devolución de los importes abonados, en su caso. Si la nueva plaza obtenida lo es de un centro propio de la Universidad de Sevilla, el importe abonado en la primera titulación podrá ser compensado en la nueva.

                  2. La pérdida de la plaza adjudicada en el procedimiento de admisión en el curso 2025-2026 en la titulación que abandona.

                  3. La anulación de matrícula por traslado de expediente a otros estudios de la Universidad de Sevilla no tendrá ningún efecto académico ni económico. Procederá la devolución de los importes abonados, en su caso. Si la nueva plaza obtenida lo es de un centro propio de la Universidad de Sevilla, el importe abonado en la primera titulación podrá ser compensado en la nueva.

                  4. La anulación de matrícula solicitada hasta el 30 de noviembre de 2025, como consecuencia de traslado de expediente a otra universidad o por haber obtenido plaza en un ciclo formativo de grado superior, dará lugar a la devolución de los importes abonados, en su caso.

                  5. El estudiante que en el presente curso académico esté al corriente de todos los pagos de su matrícula, incluso los no vencidos en caso de fraccionamiento, podrá solicitar la anulación de su matrícula hasta el 31 de marzo, con las siguientes consecuencias:

                  a) No procederá la devolución de los importes abonados en concepto de precios públicos por servicios académicos y administrativos y el seguro escolar.

                  b) Serán de aplicación los recargos económicos en la siguiente matrícula que se realice en los mismos estudios.

                  c) No se computará la matrícula realizada a efectos académicos. En consecuencia, a efectos de beca, solo se tendrá en cuenta el rendimiento académico obtenido en el curso anterior del que solicita la anulación de matrícula.

                  d) Si el estudiante accedió por el procedimiento de admisión en el curso 2025-2026, perderá la plaza adjudicada, sin perjuicio de poder volver a solicitar la admisión en cursos posteriores en la misma o en distinta titulación.

                  e) En caso de que el estudiante haya formalizado la matrícula con la exención de becario de carácter general, debe solicitar la renuncia a la solicitud de beca y abonar el importe de la matrícula como requisito para la anulación por esta causa.

                  6. El plazo para solicitar la anulación de la matrícula formalizada en el periodo general de ampliación será del 2 al 6 de febrero de 2026. La concesión de la anulación de la ampliación de matrícula supondrá la devolución del importe de los precios públicos abonados por servicios académicos, en su caso, y la matrícula no se contabilizará a efectos de aplicación de recargos, en cursos posteriores.

                  7. El estudiante afectado por la imposibilidad de obtener el visado de entrada en España o por otra causa no imputable al estudiante, podrá solicitar la anulación total de su matrícula, con las siguientes consecuencias:

                  a) Procederá la devolución de los importes abonados, en su caso.

                  b) El estudiante perderá la plaza adjudicada por el procedimiento de admisión en el curso 2025-2026.

                  ----

                  \textbf{Código Seguro De Verificación}: eqOjhZcmL0p6VJlZiTH/Yw==
                  \textbf{Fecha}: 27/06/2025
                  \textbf{Firmado Por}: MIGUEL ANGEL CASTRO ARROYO
                  \textbf{Url De Verificación}: Verificación
                  \textbf{Página}: 8/50

                  ---

                  FUENTE: R.R. NORMAS DE MATRICULA GRADO Y MASTER UNIVERSITARIO 2025-2026.pdf | PÁGINA: 9
                  \# UNIVERSIDAD D SEVII

                  \#\# C)EST

                  \#\#\# Sección Segunda: Modificación, ampliación y anulación de matrícula

                  \#\#\#\# Artículo 7. Modificación de matrícula

                  Los estudiantes podrán modificar su matrícula según lo establecido en el Anexo I de estas normas.

                  Coincidiendo con el plazo de ampliación de matrícula, se podrán admitir solicitudes de modificación de matrícula de asignaturas optativas del segundo cuatrimestre matriculadas en el periodo ordinario, siempre que no suponga una variación económica del importe de los conceptos académicos, y que el estudiante no se hubiera presentado a ninguna convocatoria de examen final de la/s asignatura/s que se pretende/n modificar. Quedará condicionada la aceptación a las disponibilidades docentes, a la organización académica de las enseñanzas y a lo dispuesto en el artículo 27 del Reglamento General de Actividades Docentes. Estas modificaciones se realizarán de forma presencial en la secretaría del centro.

                  Con posterioridad a dichos plazos, no será posible efectuar modificaciones de matrícula, salvo aquellas que resulten como consecuencia de un proceso de anulación, de la resolución de los supuestos que condicionan la matrícula establecidos en el artículo 5 o en los casos de modificaciones autorizadas por el centro de los acuerdos de estudios de los estudiantes que participan en un programa de movilidad.

                  \#\#\#\# Artículo 8. Ampliación de matrícula

                  1. Los estudiantes podrán efectuar la ampliación de su matrícula si esta ha sido previamente formalizada en el periodo ordinario y cuando no incida sobre la organización académica de las enseñanzas ni suponga un aumento en el número de grupos autorizados para las asignaturas implicadas. La ampliación podrá llevarse a cabo respecto de asignaturas anuales y del segundo cuatrimestre y únicamente generarán el derecho a la prestación de los servicios docentes posteriores al momento de su formalización.

                  Entre los créditos matriculados en periodo ordinario y los matriculados durante la ampliación no se podrán superar los límites máximos de matrícula establecidos en el artículo 4 de estas normas, una vez detraídos los créditos superados en la tercera convocatoria ordinaria.

                  2. La ampliación de matrícula y su modificación en su caso, se llevará a cabo, con carácter general, a través de automatrícula del 12 al 30 de enero de 2026.

                  Con carácter previo a la apertura del plazo, los centros deberán hacer públicas las posibles limitaciones derivadas de la organización académica de las enseñanzas.

                  El plazo general de ampliación se prorrogará para los trabajos fin de estudio hasta el 18 de febrero de 2026 para las titulaciones de grado y hasta el 6 de marzo de 2026 para las titulaciones de máster universitario, salvo para aquellos centros que no lo autoricen. Estos últimos deberán comunicarlo al Vicerrectorado de Estudiantes antes del 28 de noviembre de 2025.

                  3. Las matrículas formalizadas como resultado de una solicitud de ampliación deberán ser liquidadas en el plazo de 10 días hábiles a contar desde el día siguiente a la finalización del periodo de ampliación de matrícula por la totalidad de los precios públicos correspondientes, en un solo plazo por recibo bancario o domiciliación. En caso de matrícula en centros propios también será posible el pago mediante tarjeta bancaria.

                  4. En ningún caso, la ampliación de matrícula obligará a la modificación de la programación académica de la asignatura o asignaturas de que se trate.

                  Código Seguro De Verificación: eqOjhZcmL0p6VJlZiTH/Yw==
                  Fecha: 27/06/2025
                  Firmado Por: MIGUEL ANGEL CASTRO ARROYO
                  Url De Verificación: [https://pfirma.us.es/verifirma/code/eqOjhZcmL0p6VJlZiTH\%2FYw\%3D\%3D](https://pfirma.us.es/verifirma/code/eqOjhZcmL0p6VJlZiTH\%2FYw\%3D\%3D)
                  Página: 5/50
                \end{scriptsize}
        \end{itemize}


  \item \textbf{P46.} En base a la Guía completa de liquidación de viajes, ¿qué acciones documentales debo adjuntar paso a paso para justificar ante Gerencia mis diferentes gastos de transporte abonados? \\
        \begin{itemize}
          \item \textbf{Tipo de consulta:} simple
          \item \textbf{Respuesta del Agente:}

                ⚠️ Requisitos Previos:
                Antes de comenzar el trámite, asegúrese de tener la siguiente documentación y cumplir con los siguientes requisitos:
                - Comisión de Servicio/Autorización de desplazamiento
                - \textbf{Formulario de liquidación de gastos de viaje} debidamente firmado a través de \textbf{Portafirmas}
                - Memoria justificativa de la actividad realizada
                - Documentación acreditativa de la actividad realizada y los gastos de alojamiento y locomoción

                A continuación, se presentan los pasos para adjuntar las acciones documentales necesarias para justificar los gastos de transporte abonados:

                1. \textbf{Documentación de gastos de transporte}: Para justificar los gastos de transporte, deberá adjuntar los siguientes documentos:
                * \textbf{Billetes electrónicos} de avión, tren o otros medios de transporte público, que incluyan el importe abonado y la fecha de emisión.
                * \textbf{Facturas} de alojamiento y locomoción, que incluyan los datos fiscales de la persona que realiza el viaje, la fecha y descripción de los conceptos facturados.
                * \textbf{Tarjetas de embarque} en caso de viajes en avión.
                * \textbf{Justificantes de pago} de los gastos de transporte, como movimientos bancarios o recibos de pago en efectivo.
                2. \textbf{Documentación adicional}: Deberá adjuntar la siguiente documentación adicional:
                * \textbf{Memoria justificativa} de la actividad realizada, que describa las actividades desarrolladas y su relación con los objetivos del proyecto.
                * \textbf{Certificado de asistencia} o documento equivalente que acredite la participación en la actividad.
                * \textbf{Informe de conformidad} de facturas derivadas de reservas de crédito o cargos internos por alojamiento en residencias universitarias de la Universidad de Sevilla, si procede.
                3. \textbf{Verificación de la documentación}: Verifique que toda la documentación esté completa y correcta, y que se ajuste a los requisitos establecidos en la \textbf{Guía completa de liquidación de viajes}.
                4. \textbf{Presentación de la solicitud}: Una vez que tenga toda la documentación necesaria, presente la solicitud de abono de gastos de viaje a través del \textbf{Registro General de la Universidad de Sevilla}, dirigida a la \textbf{Subdirección General de Gestión Económica de la Investigación}.

                Recuerde que es importante seguir las instrucciones de la \textbf{Guía completa de liquidación de viajes} y adjuntar toda la documentación necesaria para evitar retrasos en el trámite.
          \item \textbf{Contexto utilizado:}
                \begin{scriptsize}
                  FUENTE: guia-breve-liquidacion.pdf | PÁGINA: 2
                  \# Subdirección de Gestión Económica de la Investigación

                  \#\# AL TERMINAR LA ACTIVIDAD

                  \#\#\# 1. LIQUIDACIÓN DE GASTOS

                  Rellene el Formulario (INCLUIR EL ENLACE DE LA WEB) de liquidación de gastos de viaje. En la Guía (INCLUIR EL ENLACE DE LA WEB) para la cumplimentación de solicitudes de abono encontrará la información detallada para la correcta cumplimentación de este formulario.

                  Una vez terminado, debe ser firmado a través de Portafirmas, tanto por el IP como por la persona que realiza la actividad.

                  Podrá renunciar parcialmente a los gastos de la liquidación tal y como se le especifica en detalle en el punto 3.4 de la guía completa.

                  \#\#\# 2. DOCUMENTACIÓN ADICIONAL A PRESENTAR

                  Además de la Comisión de Servicio/Autorización de desplazamiento y del formulario de liquidación de gastos será necesario presentar la siguiente documentación:

                  - Una memoria justificativa en la que describa las actividades realizadas indicando la relación de éstas con alguno/s de los objetivos del proyecto. También deberá especificar el vínculo con el proyecto de la persona que realiza el viaje, así como su situación laboral. Puede consultar todos los detalles necesarios en el punto 4.2 de la guía completa.

                  - Documentación acreditativa de:
                  - La actividad realizada (certificado de asistencia, documento firmado por responsable del centro donde haya tenido lugar la estancia, etc.). Ver detalles en punto 4.3 de la guía completa.
                  - Los gastos de alojamiento y/o locomoción efectuados por la persona que ha realizado la actividad (facturas, tickets, billetes electrónicos, etc.), y trazabilidad bancaria de los mismos. Ver detalles en punto 4.4 de la guía completa.
                  - Si procede, informe de conformidad firmado por el IP de las facturas de alojamiento y/o locomoción derivadas de reservas de crédito a través de agencias de viaje del acuerdo marco vigente.
                  - Si procede, documento de desglose de la renuncia parcial de los gastos de liquidación en pdf.

                  VERSIÓN 2, OCTUBRE 2024    PÁGINA 2 | 4

                  ---

                  FUENTE: Miniguía viajes personal propio US.pdf | PÁGINA: 6
                  \# Miniguía para viajes por razón del servicio

                  \#\# Documentación necesaria para la liquidación (Art. 33)

                  \vspace{0.4cm}
                  \begin{table}[h!]\centering\renewcommand{\arraystretch}{1.3}
                    \begin{tabular}{|p{0.45\textwidth}|p{0.45\textwidth}|}\hline
                      Tipo de gasto*                & Documentación justificativa                                                                                                                                                                                     \\\hline
                      Comisión de servicio          & Documento original autorizado.                                                                                                                                                                                  \\\hline
                      Cambios de itinerario         & Declaración con fechas, horas y lugares.                                                                                                                                                                        \\\hline
                      Alojamiento                   & Factura con datos fiscales, fechas, desglose de servicios.                                                                                                                                                      \\\hline
                      Manutención                   & Si se opta por dieta: no requiere factura. Si se opta por reembolso exacto: facturas. Si se exceden los límites fijados por la AEAT (los mismos que en los anexos), la parte que los excede está sujeta a IRPF. \\\hline
                      Transporte                    & Billetes con fechas y horarios + tarjeta de embarque o justificante de uso.                                                                                                                                     \\\hline
                      Taxis y VTC                   & Taxis y VTC: Factura con origen/destino (tique). Máximo 2 diarios (ida y vuelta).                                                                                                                               \\\hline
                      Aparcamiento y peajes         & Factura original.                                                                                                                                                                                               \\\hline
                      Inscripciones                 & Factura de inscripción y certificado de asistencia.                                                                                                                                                             \\\hline
                      Reuniones de proyecto         & Acta o documento equivalente.                                                                                                                                                                                   \\\hline
                      Trabajos de campo             & Trabajos de campo: Memoria, facturas de gasolina/transportes, justificación del uso del coche, contrato de alquiler de vehículo, en su caso.                                                                    \\\hline
                      Facturas en moneda extranjera & Tipo de cambio del Banco de España.                                                                                                                                                                             \\\hline
                    \end{tabular}\end{table}
                  \vspace{0.4cm}

                  * En proyectos subvencionados de investigación, transferencia, académicos o de cooperación el gasto debe ser elegible dentro de la convocatoria.

                  Gerencia
                  Universidad de Sevilla

                  ---

                  FUENTE: guia-completa-liquidacion-de-viajes-v2.pdf | PÁGINA: 20
                  \# Subdirección de Gestión Económica de la Investigación

                  \#\# AVIÓN

                  - Para vuelos nacionales: billete electrónico si recoge el importe abonado, en caso contrario, debe adjuntar además la factura correspondiente a dicho billete.
                  - Para vuelos internacionales: billete electrónico que recoja el importe abonado/recibos.
                  - Tarjetas de embarque.
                  - En caso de pérdida de la tarjeta de embarque: certificado de la empresa de haber realizado el vuelo (no es suficiente con una declaración responsable del beneficiario).

                  \#\# TAXI, UBER, CABIFY U OTRO

                  Sólo se admite un máximo de cuatro trayectos en taxi, correspondiendo exclusivamente a los siguientes itinerarios:

                  \#\#\# DÍA DE SALIDA:
                  - Trayecto desde la residencia del beneficiario al aeropuerto/estación.
                  - Trayecto del aeropuerto/estación al alojamiento de destino.

                  \#\#\# DÍA DE REGRESO:
                  - Trayecto del alojamiento de destino al aeropuerto/estación.
                  - Trayecto del aeropuerto/estación a la residencia del beneficiario.

                  En caso de servicio UBER, CABIFY o similar, debe solicitar a la empresa que indique la hora en su factura, además de la fecha. Si no fuera posible, deberá adjuntar el correo de la empresa en el que se confirma el itinerario contratado, además de la factura. (El itinerario recoge la hora).

                  \#\# ALQUILER DE COCHE

                  - Contrato de alquiler
                  - Factura

                  \#\# 4.5 INFORMES DE CONFORMIDAD DE FACTURAS DERIVADAS DE RESERVAS DE CRÉDITO O CARGOS INTERNOS POR ALOJAMIENTO EN RESIDENCIAS UNIVERSITARIAS DE LA UNIVERSIDAD DE SEVILLA

                  En el caso de pago a terceros, es decir, facturas derivadas de reservas de crédito para la contratación de los servicios de alojamiento o locomoción a través de las agencias de viaje del acuerdo marco vigente a los efectos, debe adjuntar el informe de conformidad del correspondiente justificante de gasto firmado por el IP.

                  \#\# 4.6 DOCUMENTO DE RENUNCIA PARCIAL, EN SU CASO

                  En el caso de solicitar renuncia parcial, debe adjuntar el documento de desglose conforme a lo descrito en el punto 3.4 de la presente guía.

                  \#\# 5. ¿DÓNDE Y CÓMO PRESENTO LA SOLICITUD DE ABONO, JUNTO CON LA LIQUIDACIÓN Y LA DOCUMENTACIÓN OBLIGATORIA?

                  Una vez cumplimentada y debidamente firmada la liquidación de dietas y recopilada y verificada la documentación obligatoria justificativa, deberá tramitar su solicitud electrónicamente, dirigida a la Subdirección de Gestión Económica de la Investigación, a través del Registro General de la Universidad de Sevilla dirigido a la Subdirección General de Gestión Económica de la Investigación.

                  VERSIÓN 2 OCTUBRE 2024    PÁGINA 20 | 21

                  ---

                  FUENTE: guia-completa-liquidacion-de-viajes-v2.pdf | PÁGINA: 19
                  \# Subdirección de Gestión Económica de la Investigación

                  \#\# 4.3 EVIDENCIA DE LA REALIZACIÓN DE LA ACTIVIDAD (MOTIVO DEL VIAJE)

                  Los Organismos financiadores exigen prueba documental de que la actividad o evento ha tenido lugar para acreditar que el gasto es elegible.

                  A continuación, se señalan los más comunes según la naturaleza de la actividad:

                  - \textbf{Asistencia a un congreso, seminario, jornada o similar}: debe presentar certificado de asistencia emitido por la entidad organizadora del evento. La mera inscripción no acredita haber asistido.
                  - \textbf{Participación como ponente en un congreso, seminario, jornada o similar}: copia del programa donde se pueda comprobar su nombre y datos del evento.
                  - \textbf{Asistencia a cursos de formación}: certificado de asistencia y programa donde se especifique el número de horas o créditos ECTS del curso.
                  - \textbf{Reuniones científicas o de coordinación/seguimiento del proyecto}: La misma memoria justificativa puede servir a estos efectos si describe la temática tratada en la reunión y su relación con el proyecto.
                  - \textbf{Trabajos de campo}: memoria del itinerario y actividad desarrollada firmada por el investigador responsable del proyecto.
                  - \textbf{Estancias}: Documento expedido por el responsable del centro donde haya tenido lugar la misma con indicación de las actividades desarrolladas y/o personas con las que ha colaborado durante la estancia.

                  \#\# 4.4 FACTURAS ABONADAS POR EL INTERESADO, ASÍ COMO JUSTIFICANTES DE LOS PAGOS EFECTUADOS

                  Si ha abonado personalmente facturas o tickets, debe presentar los correspondientes originales. En el punto 5 de esta guía se le indica cómo debe proceder según se trate de un original en soporte electrónico o en soporte papel.

                  Recuerde que las facturas deben reflejar los datos fiscales de la persona que realiza el viaje, en ningún caso los de la US, además de incluir la fecha y descripción de los conceptos facturados (nombre de la persona que viaja, itinerario del trayecto, número de noches, precio por noche, etc.). Dado que los conceptos genéricos son motivo de reintegro, asegúrese de que las facturas no contienen errores. De lo contrario, el personal de la Subdirección de Gestión Económica de la Investigación le dará traslado de la incidencia para la subsanación oportuna, con el consiguiente retraso en la tramitación.

                  En los tickets/facturas simplificadas vigile que figura la fecha y hora de emisión.

                  Por cada factura, ticket, billete electrónico, u otro documento justificativo de gasto debe presentar el correspondiente justificante del pago efectuado personalmente, ya que está obligado a acreditar ante el Organismo financiador los gastos soportados, ya sea mediante tarjeta, transferencia o en efectivo.

                  En caso de “Pagos en efectivo”, debe presentar el recibo del proveedor firmado por éste, con identificación de los datos del sujeto deudor que efectúa el pago, es decir, del beneficiario de la liquidación.

                  En el caso concreto de gastos de locomoción, deberá presentar los siguientes documentos según el medio de transporte público utilizado, junto con el correspondiente justificante del pago materializado (movimiento o cargo bancario):

                  \#\#\# TREN / AVE:

                  - Factura
                  - Billete electrónico

                  VERSIÓN 2 OCTUBRE 2024    PÁGINA 19 | 21

                  ---

                  FUENTE: guia-breve-liquidacion.pdf | PÁGINA: 1
                  \# Subdirección de Gestión Económica de la Investigación

                  \#\# GUÍA BREVE
                  \#\#\# Solicitudes de liquidación de gastos de viajes

                  \#\#\# Introducción

                  En este documento se presenta una versión reducida de la Guía para la cumplimentación de solicitudes de abono de liquidaciones de gastos de viaje en el ámbito de la investigación, con el fin de facilitar los trámites necesarios para la cumplimentación del formulario de liquidación de gastos y de especificar la documentación necesaria para la realización y justificación de la actividad, así como el procedimiento de presentación de la solicitud de abono.

                  \#\#\# ANTES DE REALIZAR LA ACTIVIDAD

                  - Deberá tramitar la Comisión de Servicios con su licencia por estudios a personal docente (PDI funcionario o laboral) o Autorización de Desplazamiento (Personal Investigador o Técnico Contratado). Puede consultar detalles en el punto 4.1 de la guía completa.

                  - Deberá comprobar que la actividad a desarrollar figura como un gasto elegible en la convocatoria del proyecto con el que pretende justificarla, que la persona que la va a llevar a cabo guarda el vínculo que exige el organismo financiador, y que dispone de saldo suficiente en la orgánica del proyecto.

                  - Si va a alojarse haciendo uso de compañías que ofrecen servicios de reservas online, asegúrese de que podrá obtener factura por el gasto total realizado.

                  - Si va a reservar el medio de transporte y/o el alojamiento a través de agencias de viaje del acuerdo marco vigente, solicite la reserva de crédito correspondiente.

                  \#\#\# DURANTE LA ACTIVIDAD

                  - Conserve tarjetas de embarque si el viaje se ha realizado en avión.
                  - Pida certificados o documentos que acrediten su participación en la actividad.
                  - Pida facturas o tickets del alojamiento y/o medios de locomoción utilizados donde figuren sus datos fiscales, no los de la US. En facturas/tickets simplificados asegúrese de que aparecen la fecha y hora de emisión.
                \end{scriptsize}
        \end{itemize}


  \item \textbf{P47.} Según las instrucciones oficiales, indica los pasos ordenados para solicitar la modificación a matrícula a tiempo parcial demostrando documentalmente mi situación y compatibilidad laboral actual. \\
        \begin{itemize}
          \item \textbf{Tipo de consulta:} simple
          \item \textbf{Respuesta del Agente:}

                ⚠️ Requisitos Previos:
                Antes de iniciar el proceso de modificación de matrícula a tiempo parcial, es importante tener en cuenta los siguientes requisitos previos:
                - Debes tener una situación laboral actual que te permita solicitar la modificación a matrícula a tiempo parcial.
                - Debes tener documentación que demuestre tu situación laboral y compatibilidad laboral actual.
                - Debes cumplir con los plazos establecidos para la modificación de matrícula, que varían según el tipo de estudios (grado o máster universitario).

                A continuación, te presento los pasos ordenados para solicitar la modificación a matrícula a tiempo parcial:

                1. \textbf{Revisa los plazos de modificación de matrícula}: Verifica los plazos establecidos para la modificación de matrícula en la Universidad de Sevilla, que varían según el tipo de estudios (grado o máster universitario). Los plazos son:
                * Estudiantes de grado: dentro del día y franja horaria asignada, o desde las 18:00 h del 25 de julio hasta las 24:00 h del 31 de julio y desde las 9:00 h del 1 de septiembre hasta las 24:00 h del 5 de septiembre.
                * Estudiantes de máster universitario: entre el 9 y el 31 de julio y del 1 al 5 de septiembre de 2025.
                2. \textbf{Prepara la documentación necesaria}: Debes preparar la documentación que demuestre tu situación laboral y compatibilidad laboral actual, como:
                * \textbf{Declaración Responsable de veracidad de la documentación aportada}: solo se cumplimentará si aportas algún documento en el Buzón de Documentación (DOMUS).
                * \textbf{Orden de Domiciliación de Adeudo Directo (SEPA)}: se generará y quedará firmada en la plataforma de automatrícula en el momento de formalización de la matrícula, en caso de domiciliación de los recibos.
                * Otros documentos que demuestren tu situación laboral y compatibilidad laboral actual.
                3. \textbf{Accede a la plataforma de automatrícula}: Accede a la plataforma de automatrícula de la Universidad de Sevilla, \textbf{SEVIUS}, con tu UVUS y contraseña.
                4. \textbf{Solicita la modificación de matrícula}: En la plataforma de automatrícula, solicita la modificación de matrícula a tiempo parcial, adjuntando la documentación necesaria.
                5. \textbf{Verifica y confirma la modificación}: Verifica que la modificación de matrícula se ha realizado correctamente y confirma la solicitud.
                6. \textbf{Entrega la documentación en el Buzón de Documentación (DOMUS)}: Si es necesario, entrega la documentación en el Buzón de Documentación (DOMUS) y firma la \textbf{Declaración Responsable de veracidad de la documentación aportada} mediante el doble factor de autenticación (2FA).
                7. \textbf{Verifica el estado de la solicitud}: Verifica el estado de la solicitud de modificación de matrícula en la plataforma de automatrícula y espera la respuesta de la Universidad de Sevilla.

                Recuerda que es importante cumplir con los plazos establecidos y entregar la documentación necesaria para que la modificación de matrícula se realice correctamente.
          \item \textbf{Contexto utilizado:}
                \begin{scriptsize}
                  FUENTE: R.R. NORMAS DE MATRICULA GRADO Y MASTER UNIVERSITARIO 2025-2026.pdf | PÁGINA: 52
                  \# UNIVERSIDAD D SEVII

                  \#\# B) AUTOMATRÍCULA ESTUDIANTES QUE CONTINÚAN SUS ESTUDIOS EN LA MISMA TITULACIÓN

                  \#\#\# 2. Modificación de la automatrícula

                  Se permiten las modificaciones, dentro de los plazos indicados a continuación, diferenciados por el tipo de estudios.

                  - \textbf{Estudiantes de grado}
                  - Dentro del día y franja horaria asignada.
                  - A partir del día asignado: de 18:00 h a 8:45 h del día siguiente, hasta el 25 de julio de 2025 (fin periodo acceso escalonado).
                  - Desde las 18:00 h del 25 de julio hasta las 24:00 h del 31 de julio y desde las 9:00 h del 1 de septiembre hasta las 24:00 h del 5 de septiembre.

                  - \textbf{Estudiantes de máster universitario}
                  - Podrán modificar su automatrícula en cualquier momento entre el 9 y el 31 de julio y del 1 al 5 de septiembre de 2025.

                  \#\#\# 3. Documentación

                  \#\#\#\# Documentación general

                  - \textbf{Declaración Responsable de veracidad de la documentación aportada.} Solo se cumplimentará por los estudiantes que aporten algún documento en el Buzón de Documentación (DOMUS). Se generará una vez que el estudiante acceda a este con su UVUS quedando firmada mediante el doble factor de autenticación (2FA).

                  - \textbf{Orden de Domiciliación de Adeudo Directo (SEPA).} Se generará y quedará firmada en la plataforma de automatrícula en el momento de formalización de esta, en caso de domiciliación de los recibos. No será necesario aportarla al Buzón de Documentación (DOMUS). El estudiante deberá ser titular o cotitular de la cuenta bancaria en la que se domicilien los recibos.

                  ----

                  Código Seguro De Verificación: `eqOjhZcmL0p6VJlZiTH/Yw==`
                  Fecha: 27/06/2025
                  Firmado Por: MIGUEL ANGEL CASTRO ARROYO
                  Url De Verificación: [https://pfirma.us.es/verifirma/code/eqOjhZcmL0p6VJlZiTH\%2FYw\%3D\%3D](https://pfirma.us.es/verifirma/code/eqOjhZcmL0p6VJlZiTH\%2FYw\%3D\%3D)
                  Página: 48/50

                  ---

                  FUENTE: R.R. NORMAS DE MATRICULA GRADO Y MASTER UNIVERSITARIO 2025-2026.pdf | PÁGINA: 9
                  \# UNIVERSIDAD D SEVII

                  \#\# C)EST

                  \#\#\# Sección Segunda: Modificación, ampliación y anulación de matrícula

                  \#\#\#\# Artículo 7. Modificación de matrícula

                  Los estudiantes podrán modificar su matrícula según lo establecido en el Anexo I de estas normas.

                  Coincidiendo con el plazo de ampliación de matrícula, se podrán admitir solicitudes de modificación de matrícula de asignaturas optativas del segundo cuatrimestre matriculadas en el periodo ordinario, siempre que no suponga una variación económica del importe de los conceptos académicos, y que el estudiante no se hubiera presentado a ninguna convocatoria de examen final de la/s asignatura/s que se pretende/n modificar. Quedará condicionada la aceptación a las disponibilidades docentes, a la organización académica de las enseñanzas y a lo dispuesto en el artículo 27 del Reglamento General de Actividades Docentes. Estas modificaciones se realizarán de forma presencial en la secretaría del centro.

                  Con posterioridad a dichos plazos, no será posible efectuar modificaciones de matrícula, salvo aquellas que resulten como consecuencia de un proceso de anulación, de la resolución de los supuestos que condicionan la matrícula establecidos en el artículo 5 o en los casos de modificaciones autorizadas por el centro de los acuerdos de estudios de los estudiantes que participan en un programa de movilidad.

                  \#\#\#\# Artículo 8. Ampliación de matrícula

                  1. Los estudiantes podrán efectuar la ampliación de su matrícula si esta ha sido previamente formalizada en el periodo ordinario y cuando no incida sobre la organización académica de las enseñanzas ni suponga un aumento en el número de grupos autorizados para las asignaturas implicadas. La ampliación podrá llevarse a cabo respecto de asignaturas anuales y del segundo cuatrimestre y únicamente generarán el derecho a la prestación de los servicios docentes posteriores al momento de su formalización.

                  Entre los créditos matriculados en periodo ordinario y los matriculados durante la ampliación no se podrán superar los límites máximos de matrícula establecidos en el artículo 4 de estas normas, una vez detraídos los créditos superados en la tercera convocatoria ordinaria.

                  2. La ampliación de matrícula y su modificación en su caso, se llevará a cabo, con carácter general, a través de automatrícula del 12 al 30 de enero de 2026.

                  Con carácter previo a la apertura del plazo, los centros deberán hacer públicas las posibles limitaciones derivadas de la organización académica de las enseñanzas.

                  El plazo general de ampliación se prorrogará para los trabajos fin de estudio hasta el 18 de febrero de 2026 para las titulaciones de grado y hasta el 6 de marzo de 2026 para las titulaciones de máster universitario, salvo para aquellos centros que no lo autoricen. Estos últimos deberán comunicarlo al Vicerrectorado de Estudiantes antes del 28 de noviembre de 2025.

                  3. Las matrículas formalizadas como resultado de una solicitud de ampliación deberán ser liquidadas en el plazo de 10 días hábiles a contar desde el día siguiente a la finalización del periodo de ampliación de matrícula por la totalidad de los precios públicos correspondientes, en un solo plazo por recibo bancario o domiciliación. En caso de matrícula en centros propios también será posible el pago mediante tarjeta bancaria.

                  4. En ningún caso, la ampliación de matrícula obligará a la modificación de la programación académica de la asignatura o asignaturas de que se trate.

                  Código Seguro De Verificación: eqOjhZcmL0p6VJlZiTH/Yw==
                  Fecha: 27/06/2025
                  Firmado Por: MIGUEL ANGEL CASTRO ARROYO
                  Url De Verificación: [https://pfirma.us.es/verifirma/code/eqOjhZcmL0p6VJlZiTH\%2FYw\%3D\%3D](https://pfirma.us.es/verifirma/code/eqOjhZcmL0p6VJlZiTH\%2FYw\%3D\%3D)
                  Página: 5/50

                  ---

                  FUENTE: R.R. NORMAS DE MATRICULA GRADO Y MASTER UNIVERSITARIO 2025-2026.pdf | PÁGINA: 45
                  \# UNIVERSIDAD D SEVII

                  \#\# C) AUTOMATRÍCULA ESTUDIANTES DE NUEVO INGRESO POR EL PROCEDIMIENTO DE ADMISIÓN

                  \#\#\# 2. Modificación de la automatrícula

                  - En su periodo específico de matriculación.
                  - A partir de las 18:00 h del 25 de julio hasta las 24:00 h del 31 de julio y desde las 9:00 h del 1 de septiembre hasta las 24:00 h del 5 de septiembre.
                  - Matrículas posteriores al 5 de septiembre de 2025: 10 días hábiles siguientes a la finalización del periodo en el que se haya realizado la matrícula, debiéndose realizar en la secretaría del centro.

                  \#\#\# 3. Documentación

                  \#\#\#\# 3.1. Documentación general

                  - \textbf{DNI/Documento de identidad} donde conste el NIE, o en su defecto el pasaporte, en tanto se aporta en el plazo máximo de 6 meses dicho documento, según se establece en la disposición adicional séptima de estas normas de matrícula.

                  - \textbf{Declaración Responsable de veracidad de la documentación aportada.} Solo se cumplimentará por los estudiantes que aporten algún documento en el Buzón de Documentación (DOMUS). Se generará una vez que el estudiante acceda a este con su UVUS, quedando firmada mediante el doble factor de autenticación (2FA).

                  - \textbf{Orden de Domiciliación de Adeudo Directo (SEPA).} Se generará y quedará firmada en la plataforma de automatrícula en el momento de formalización de esta, en caso de domiciliación de los recibos. No será necesario aportarla al Buzón de Documentación (DOMUS). El estudiante deberá ser titular o cotitular de la cuenta bancaria en la que se domicilien los recibos.

                  \#\#\#\# 3.2. Documentación específica según la vía de acceso

                  La documentación acreditativa de la vía de acceso de los estudiantes procedentes de sistemas educativos de fuera de España será verificada por las secretarías de los centros, con la documentación original.

                  - \textbf{Matrícula en grados}

                  a) Estudiantes procedentes de Bachillerato y/o PAU / PEvAU
                  - Título de Bachiller o resguardo de haberlo solicitado, cuando por razones técnicas o por oposición del interesado no se haya obtenido mediante la consulta telemática.
                  - Tarjeta de la PAU o PEvAU (solo estudiantes que hayan superado la prueba en una universidad distinta a la Universidad de Sevilla).

                  ----

                  Código Seguro De Verificación: `eqOjhZcmL0p6VJlZiTH/Yw==`
                  Fecha: 27/06/2025
                  Firmado Por: MIGUEL ANGEL CASTRO ARROYO
                  Url De Verificación: [https://pfirma.us.es/verifirma/code/eqOjhZcmL0p6VJlZiTH\%2FYw\%3D\%3D](https://pfirma.us.es/verifirma/code/eqOjhZcmL0p6VJlZiTH\%2FYw\%3D\%3D)
                  Página: 41/50

                  ---

                  FUENTE: R.R. NORMAS DE MATRICULA GRADO Y MASTER UNIVERSITARIO 2025-2026.pdf | PÁGINA: 11
                  \# UNIVERSIDAD D SEVII

                  C)EST

                  5. No se concederán anulaciones parciales de matrícula a solicitud del estudiante, excepto por coincidencia no prevista de horarios de clase y exámenes y las situaciones contempladas en estas normas. No obstante, por causas académicas sobrevenidas que imposibiliten la prestación del servicio, los centros podrán realizar de oficio anulaciones parciales de matrícula, ofreciendo a los estudiantes la posibilidad de matricular otras asignaturas o bien la devolución de los importes abonados en los créditos anulados, en su caso.

                  6. La matrícula se anulará de oficio cuando se trate de error, omisión, falseamiento de los datos consignados en esta o cuando se modifique la situación administrativa o académica del estudiante, o las matrículas condicionadas que no deban producir efectos.

                  Asimismo se anularán de oficio las matrículas cuando el estudiante no aporte la documentación requerida o por falta de pago.

                  7. Los centros anularán de oficio las matrículas efectuadas por estudiantes que se matriculen en otra titulación por haber obtenido plaza en los procedimientos de admisión para estudiantes de nuevo ingreso en esta u otra universidad pública, siempre que no se acredite una autorización para la simultaneidad de estudios.

                  \#\# Artículo 11. Supuestos de anulación de matrícula y sus efectos

                  1. Se podrá solicitar la anulación de matrícula ordinaria al inicio del curso académico en los siguientes plazos:

                  - Titulaciones de grado: hasta el 16 de septiembre de 2025.
                  - Titulaciones de máster universitario que inicien la docencia el 6 de octubre: hasta el 15 de octubre de 2025.
                  - Titulaciones de máster universitario que inicien la docencia el 22 de septiembre: hasta el 30 de septiembre de 2025.
                  - Cuando la matrícula se haya formalizado con posterioridad al inicio del periodo docente del primer cuatrimestre: 7 días hábiles a contar desde el día de la finalización del periodo en el que haya realizado su matrícula.

                  La concesión de la anulación supondrá:

                  a) La devolución de los importes abonados, en su caso.

                  b) La matrícula no se computará a efectos de convocatorias de examen ni de permanencia del estudiante, ni se contabilizará para futuras matrículas a efectos de aplicación de recargos.

                  c) La pérdida de la plaza adjudicada si el estudiante accedió por el procedimiento de admisión en el curso 2025-2026.

                  2. La anulación de la matrícula por haber obtenido plaza en esta u otra universidad pública (en el procedimiento o sistema de admisión que tenga establecido) en otra titulación de grado, máster universitario o estudios de doctorado donde hubiese sido admitido el estudiante con posterioridad a su matriculación en la Universidad de Sevilla, no tendrá ningún efecto académico ni económico si se solicita con anterioridad al 31 de octubre de 2025, en el caso de titulaciones de grado, y al 31 de diciembre de 2025 en las titulaciones de máster universitario.

                  Código Seguro De Verificación: `eqOjhZcmL0p6VJlZiTH/Yw==`
                  Fecha: 27/06/2025
                  Firmado Por: MIGUEL ANGEL CASTRO ARROYO
                  Url De Verificación: [https://pfirma.us.es/verifirma/code/eqOjhZcmL0p6VJlZiTH\%2FYw\%3D\%3D](https://pfirma.us.es/verifirma/code/eqOjhZcmL0p6VJlZiTH\%2FYw\%3D\%3D)
                  Página: 7/50

                  ---

                  FUENTE: Guia\_Movilidad\_Intern\_FAC\_TUR\_FINANZAS.pdf | PÁGINA: 13
                  Una vez se haya informado de los cambios a la Vicedecana de Movilidad y con su visto bueno previo, las modificaciones se gestionan a través de SEVIUS (se ha de solicitar la modificación a través de la propia plataforma) y tras incluir los cambios debe validarse de nuevo desde el Vicedecanato de Movilidad. Como te hemos comentado anteriormente, se te devolverá por correo electrónico, firmado y sellado. Y para su validez de cara al reconocimiento de créditos a la vuelta de la estancia en el centro de destino, dicho documento deberá entregarse firmado y sellado por las tres partes.

                  \# IMPORTANTE:

                  nunca se reconocerán los créditos de asignaturas que no estén en el acuerdo académico y/o de las asignaturas que no se incluyan en la matrícula.

                  \# Cambio de la duración de la movilidad: ampliar o reducir

                  La reducción de la estancia sólo será autorizada en casos excepcionales y por motivos justificados con la previa presentación de la debida documentación que demuestre las circunstancias excepcionales y sobrevenidas durante el desarrollo de la estancia.

                  Una vez se haya informado a la Vicedecana de Movilidad e Internacionalización del interés de reducir la estancia y expuestos los motivos justificados y con el consentimiento de la

                  13
                \end{scriptsize}
        \end{itemize}


  \item \textbf{P48.} Para un doctor invitado de otra universidad al que se pagan gastos con fondos del proyecto de investigación, enumera los pasos administrativos procedimentales para gestionar de inicio a fin correctamente sus dietas y el alta a terceros. \\
        \begin{itemize}
          \item \textbf{Tipo de consulta:} simple
          \item \textbf{Respuesta del Agente:}

                ⚠️ Requisitos Previos
                Antes de comenzar el proceso, es importante tener en cuenta los siguientes requisitos previos:
                - El doctor invitado debe tener una relación laboral o de colaboración con la Universidad de Sevilla o con otra entidad que permita la gestión de sus dietas y el alta a terceros.
                - Debe disponer de los fondos necesarios en el proyecto de investigación para cubrir los gastos del doctor invitado.
                - Es recomendable consultar con la Subdirección de Gestión Económica de la Investigación para asegurarse de que se cumplen todos los requisitos necesarios.

                A continuación, se presentan los pasos administrativos procedimentales para gestionar de inicio a fin correctamente las dietas y el alta a terceros del doctor invitado:

                1. \textbf{Alta a terceros}: Si es la primera vez que la Universidad de Sevilla le va a realizar una transferencia económica al doctor invitado, debe tramitar el alta de tercero en \textbf{Universitas XXI Económico}. Para ello, debe cumplimentar y presentar la \textbf{solicitud de alta} con los documentos que exige la normativa interna de la Universidad de Sevilla, que incluyen:
                * Copia del DNI o NIE (si es español o extranjero residente en España)
                * Documento que acredite la titularidad de la cuenta bancaria del interesado
                * Documento que acredite la dirección física o fiscal completa
                * Copia del pasaporte (si es extranjero no residente en España)
                * Visado o permiso de residencia (si reside en un país distinto al que indica su pasaporte)
                2. \textbf{Descargar el formulario de liquidación}: El doctor invitado debe descargar el formulario de \textbf{liquidación de gastos de viaje} a través del enlace proporcionado por la Subdirección de Gestión Económica de la Investigación.
                3. \textbf{Cumplimentar el formulario de liquidación}: El doctor invitado debe cumplimentar el formulario de liquidación con los datos necesarios, incluyendo:
                * Datos del titular beneficiario y de la actividad motivo del viaje
                * Datos necesarios para el cálculo de dietas de alojamiento y manutención
                * Itinerario del viaje
                * Lugar de destino donde se desarrollará el evento o actividad que motiva el viaje
                4. \textbf{Presentar la documentación}: El doctor invitado debe presentar la documentación necesaria, incluyendo:
                * Comisión de servicios o autorización de desplazamiento
                * Formulario de liquidación de gastos
                * Memoria justificativa de la actividad realizada
                * Documentos acreditativos de la actividad realizada
                * Facturas abonadas y la trazabilidad bancaria de las mismas
                5. \textbf{Solicitud de abono de la liquidación}: El doctor invitado debe solicitar el abono de la liquidación de gastos a través del \textbf{Registro Electrónico General}, dirigido a la Subdirección General de Gestión Económica de la Investigación.
                6. \textbf{Revisión y aprobación}: La Subdirección de Gestión Económica de la Investigación revisará y aprobará la liquidación de gastos y el alta a terceros del doctor invitado.

                Es importante tener en cuenta que cada entrada de Registro debe corresponder a una sola solicitud de abono, y que la documentación debe estar completa y correcta para evitar retrasos en el proceso.
          \item \textbf{Contexto utilizado:}
                \begin{scriptsize}
                  FUENTE: guia-completa-liquidacion-de-viajes-v2.pdf | PÁGINA: 4
                  \# Subdirección de Gestión Económica de la Investigación

                  - Si va a reservar alojamiento a través de alguna plataforma electrónica online (Booking, Airbnb u otra), debe asegurarse de que el responsable del alojamiento y/o la plataforma le va a emitir factura o facturas que sumen el total del gasto realizado. Las confirmaciones de reserva por correo electrónico no sustituyen a las facturas en el actual marco jurídico, de ahí que deba velar por no perder el derecho al pago por no disponer de éstas.

                  - Si es la primera vez que la Universidad de Sevilla le va a realizar una transferencia económica, debe tramitar el alta de tercero en Universitas XXI Económico, tal como se describe en el punto 3.3.6 de este documento. No hacerlo retrasará la tramitación de la liquidación.

                  Como complemento a la presente guía, le informamos que también puede consultar el documento de preguntas frecuentes, en continua actualización, y que incluye determinados supuestos especiales de reiterada consulta por el personal investigador; o bien dirigirse al personal de la Subdirección de Gestión Económica de la Investigación a través de las siguientes direcciones de correo electrónico y teléfonos, según el proyecto donde desee imputar los gastos del viaje:

                  - Proyectos Estatales y Fundaciones:
                  - Cristina Guerra Aguilar-Galindo, cristina-ag@us.es 954 48 13 86
                  - Responsable Proyectos Comunidad Autónoma, Europeos o Internacionales:
                  - Rosa María López Valiente, rosi@us.es, 954 48 12 85

                  \#\# 3. ¿CÓMO PREPARO MI LIQUIDACIÓN DE GASTOS DE VIAJES?

                  A continuación, se le indica donde puede localizar el formulario de liquidación que debe cumplimentar, así como instrucciones sobre los distintos datos que debe introducir.

                  \#\#\# 3.1. ¿DÓNDE DESCARGAR EL FORMULARIO DE “LIQUIDACIÓN DE GASTOS DE VIAJE”?

                  La liquidación de gastos de viaje es un archivo excel diseñado por la Subdirección de Gestión Económica de la Investigación, que debe descargarse a través del siguiente enlace: formulario.

                  \textbf{IMPORTANTE:} Cada vez que vaya a cumplimentar una liquidación de dietas deberá acceder a dicho enlace para asegurarse de que cumplimenta la versión vigente en el momento. La liquidación es actualizada por la Subdirección de Gestión Económica de la Investigación para ajustarla a los cambios en la normativa reguladora que afecte a las cuantías o a los cálculos. Por ello, debe evitar guardar en su escritorio el modelo de liquidación para futuros usos.

                  Una vez dentro del enlace, le aparecerá la siguiente pantalla donde deberá clicar la opción “Liquidación de gastos de viaje” o bien acceder directamente a través de la hoja denominada “liquidación de gastos”:

                  VERSIÓN 2 OCTUBRE 2024    PÁGINA 4 | 21

                  ---

                  FUENTE: guia-completa-liquidacion-de-viajes-v2.pdf | PÁGINA: 13
                  \# Subdirección de Gestión Económica de la Investigación

                  En el supuesto de que sea la primera vez que la Universidad de Sevilla le va a realizar una transferencia, será necesario que el investigador responsable del proyecto solicite su alta en el aplicativo informático para la gestión económica de nuestra universidad.

                  Se recomienda llevar a cabo la tramitación con carácter previo, a fin de no ralentizar la fiscalización y pago de la liquidación una vez presentada, ya que no será posible abonar la liquidación en tanto no consten en el sistema informático sus datos fiscales y cuenta bancaria donde realizarle la transferencia que corresponda.

                  Para darse de alta como tercero en Universitas XXI-Económico (a efectos de pago de transferencias), deberá cumplimentar y presentar la mencionada solicitud de alta con los documentos que exige la normativa interna de la Universidad de Sevilla.

                  En el siguiente enlace puede descargarse la solicitud de alta:

                  [https://investigacion.us.es/docs/area/solicitud\_alta\_terceros.docx](https://investigacion.us.es/docs/area/solicitud\_alta\_terceros.docx) (INCLUIR EL ENLACE)

                  Los documentos obligatorios son los siguientes:

                  - Si es español o extranjero residente en España, debe presentar los siguientes documentos:
                  - Copia del DNI o NIE, según el caso.
                  - Documento que acredite la titularidad de la cuenta bancaria del interesado.
                  - Documento que acredite la dirección física o fiscal completa.

                  - En el caso de ser extranjero no residente en España, deberá presentar los siguientes documentos, según corresponda:
                  - Copia del pasaporte.
                  - Visado o permiso de residencia (fiscal), si reside en un país distinto al que indica su pasaporte.
                  - Documento que acredite la titularidad de la cuenta bancaria. En caso de no ser el titular, orden de adeudo bancario. Es obligatorio que aporte el nombre de la entidad bancaria y el código swift.
                  - Documento que acredite la dirección física o fiscal completa que coincida con el país que expidió su pasaporte o, en su caso, visado o permiso de residencia fiscal. (Será válida una factura de luz, agua, gas o similar a nombre del investigador siempre que incluya la dirección postal completa).

                  La solicitud deberá presentarla a través del Registro General de la Universidad de Sevilla dirigida a la Subdirección de Gestión Económica de la Investigación.

                  VERSIÓN 2 OCTUBRE 2024    PÁGINA 13 | 21

                  ---

                  FUENTE: guia-completa-liquidacion-de-viajes-v2.pdf | PÁGINA: 18
                  \# Subdirección de Gestión Económica de la Investigación

                  - El personal investigador o técnico contratado por proyectos de ayudas debe presentar una autorización de desplazamiento.

                  ➢ Si no tiene vínculo laboral con la Universidad de Sevilla, sino con otra universidad, administración, organismo u entidad pública española, deberá presentar el documento que proceda por parte de la entidad para la que trabaja (comisión de servicios, autorización de desplazamiento, etc.).

                  ➢ En el resto de los supuestos deberá presentar un escrito firmado por el investigador principal, indicando por qué no procede presentar tales documentos.

                  Para agilizar la tramitación, compruebe que la comisión de servicios o autorización de desplazamiento cumple con los siguientes requisitos:

                  - Debe venir firmada por todas las partes.
                  - No puede mezclar firmas electrónicas con firmas manuscritas en el mismo documento.
                  - Debe especificar la orgánica y proyecto que financiará los gastos de la liquidación.
                  - La actividad que motiva el viaje debe coincidir con la que se indica en la liquidación de dietas.

                  El incumplimiento de cualquiera de esos requisitos señalados será motivo de reparo, dilatando el proceso de fiscalización y pagos.

                  \#\# 4.2 MEMORIA JUSTIFICATIVA

                  El Organismo financiador exige la justificación de forma detallada de la vinculación de los gastos derivados de la actividad que motiva el viaje con el desarrollo del proyecto y/o logro de los objetivos de éste.

                  Por ello, debe presentar una memoria justificativa en la que describa las actividades realizadas indicando claramente la relación de éstas con alguno/os de los objetivos/hitos/líneas de investigación específicos del proyecto.

                  Es importante que enmarque la actividad realizada, material adquirido, etc., en el contexto del proyecto de investigación señalando claramente los objetivos planteados, logros conseguidos y su relación con líneas específicas que estén contempladas en la memoria científica del proyecto.

                  \#\#\# Recomendaciones:

                  ➢ Debe indicar expresamente el vínculo de la persona que realiza el viaje con el proyecto, así como su situación laboral. Recuerde que en función de tal vínculo la convocatoria de subvención establece el carácter elegible o no del gasto que desea imputar al proyecto.

                  ➢ La memoria justificativa NO debe estar enfocada a la persona que realiza la actividad, sino al beneficio que la misma conlleva al desarrollo del proyecto y/o cumplimiento de sus objetivos.

                  ➢ Evite hacer referencias a cuestiones que no guarden relación directa con la descripción del vínculo con el proyecto, ya que pueden dificultar la revisión de la memoria por parte del organismo financiador.

                  ➢ No utilice el término “estancia” como norma habitual en la memoria justificativa de su viaje ya que existen normas precisas de los organismos financiadores que catalogan un viaje de “estancia” o “viaje” según criterios como la duración de la actividad, la persona que la lleva a cabo, o el objeto de la actividad. Consulte la normativa reguladora de la convocatoria para verificar que se trata de una “estancia” o de un “viaje”, o bien, en caso de dudas, consulte al personal de la Subdirección de Gestión Económica de la Investigación.

                  VERSIÓN 2 OCTUBRE 2024 PÁGINA 18 | 21

                  ---

                  FUENTE: guia-breve-liquidacion.pdf | PÁGINA: 3
                  \# Subdirección de Gestión Económica de la Investigación

                  \#\# 3. SOLICITUD DEL ABONO DE LA LIQUIDACIÓN DE GASTOS

                  A continuación, deberá proceder a la solicitud de abono de la liquidación de los gastos realizados mediante instancia efectuada a través del Registro Electrónico General, dirigido a la Subdirección General de Gestión Económica de la Investigación.

                  \textbf{IMPORTANTE:} Cada entrada de Registro debe corresponder a una sola solicitud de abono, la cual solo puede incluir una sola liquidación de gastos y sus correspondientes documentos obligatorios. Estos documentos son:

                  - Comisión Servicio o Autorización Desplazamiento con todas las firmas requeridas.
                  - Formulario de liquidación de gastos con las firmas requeridas.
                  - Memoria justificativa de la actividad realizada.
                  - Documentos acreditativos de la actividad realizada.
                  - Facturas abonadas y la trazabilidad bancaria de las mismas.
                  - Si se va a realizar un pago a terceros deberá adjuntar también el informe de conformidad firmado por el IP.
                  - Si ha solicitado renuncia parcial a parte del gasto deberá adjuntar también el documento de desglose de la renuncia en pdf.

                  Dado que actualmente el Registro Electrónico General solo permite subir un máximo de cinco archivos, organice la documentación de tal modo que cada archivo corresponda a lo siguiente:

                  1. Comisión de servicios u autorización de desplazamiento
                  2. Formulario de liquidación de dietas
                  3. Memoria justificativa
                  4. Resto de documentos (agrupados en uno o dos archivos, sabiendo que cada archivo solo puede tener un tamaño máximo de 10 Mb, y el tamaño máximo de los cinco ficheros es de 15 Mb)

                  Si su documentación incluye originales en papel deberá enviarlos por correo interno a la Subdirección General de Gestión Económica, junto con una copia de la solicitud de registro presentada.

                  \#\# CONSULTAS

                  Además de la guía completa, puede consultar también el documento de preguntas frecuentes (INCLUIR EL ENLACE DE LA WEB) elaborado por el Vicerrectorado de Investigación o dirigirse al personal del Servicio de Fiscalización y Pagos:

                  - Proyectos Estatales y Fundaciones:
                  - Cristina Guerra Aguilar-Galindo, cristina-ag@us.es 954 48 13 86
                  - Responsable Proyectos Comunidad Autónoma, Europeos o Internacionales:
                  - Rosa María López Valiente, rosi@us.es, 954 48 12 85

                  \textbf{VERSIÓN 2, OCTUBRE 2024        PÁGINA 3 | 4}

                  ---

                  FUENTE: guia-completa-liquidacion-de-viajes-v2.pdf | PÁGINA: 9
                  \# Subdirección de Gestión Económica de la Investigación

                  \#\# 3.3.1 Titular beneficiario y datos de la actividad motivo del viaje

                  Debe procurar el máximo nivel de detalle, evitando en lo posible el uso de abreviaturas o acrónimos que a menudo resultan confusos, entorpeciendo la fiscalización o revisión por parte de los auditores o los organismos financiadores.

                  En el apartado \textbf{DATOS DEL TITULAR} debe indicar los datos de la persona que solicita la indemnización de los gastos derivados del viaje realizado con motivo de su asistencia o participación en un evento o actividad, así como los datos adicionales que se solicitan sobre dicho evento o actividad. Concretamente, debe cumplimentar los siguientes apartados:

                  - Nombre y apellidos
                  - DNI (o, en su defecto, pasaporte o número del documento oficial de identificación personal expedido por su país de origen).
                  - Vinculación con el proyecto: Debe indicar si la persona que realiza el viaje es el investigador principal responsable del proyecto, o bien miembro del equipo de investigación o del equipo de trabajo, investigador invitado, …
                  - Categoría profesional: Debe expresar la categoría profesional o vínculo laboral con la US o entidad para la que trabaje.
                  - Entidad para la que trabaja (sólo personal externo a la US): Si trabaja para otra entidad, indique el nombre de dicha entidad.
                  - Actividad o evento al que asiste: Indique brevemente la actividad, evitando descripciones genéricas, tales como “desarrollo de investigación”. La actividad o evento debe coincidir con el que se menciona en la comisión de servicios (o autorización de desplazamiento) y en la memoria justificativa que más adelante se señala.
                  - Itinerario: Debe señalar el recorrido completo que llevará a cabo. Por ejemplo: Sevilla-Madrid-Berlín-Sevilla.
                  - Lugar de destino donde se desarrollará el evento o actividad que motiva el viaje.
                  - Fecha de inicio y fin del evento o actividad a la que asiste (no confundir con la fecha del viaje, aunque evidentemente debe estar comprendida dentro de la misma).

                  Llegados a este punto, el siguiente paso es cumplimentar los datos necesarios para calcular las dietas.

                  \#\# 3.3.2 Datos necesarios para el cálculo de dietas de alojamiento y manutención

                  Se entiende por dietas los gastos de alojamiento y manutención.

                  Dependiendo del itinerario del viaje, debe cumplimentar los datos que se solicitan en el apartado \textbf{DIETAS EN TERRITORIO NACIONAL} y/o \textbf{DIETAS EN TERRITORIO EXTRANJERO}.

                  Una vez dentro del apartado correspondiente, los datos a cumplimentar son los siguientes, siguiendo el orden de la liquidación:

                  - Fecha de inicio y fin del viaje. Hora de salida y hora de llegada.

                  La fecha de inicio y fin del viaje debe coincidir o, al menos, encontrarse dentro del período recogido en la comisión de servicios o autorización de desplazamiento. Si el viaje comenzara en una fecha anterior a la indicada en la comisión/autorización aprobada, o finalizara en una fecha posterior, los gastos de aquellos días anteriores y/o posteriores no podrán ser abonados, por lo que será motivo de reparo por la Subdirección de Gestión Económica de la Investigación, que le instará a la subsanación correspondiente de la comisión de servicios y/o autorización de desplazamiento al objeto de que dicho gasto pueda ser admitido.

                  VERSIÓN 2 OCTUBRE 2024    PÁGINA 9 | 21
                \end{scriptsize}
        \end{itemize}


  \item \textbf{P49.} ¿Qué pasos de trámite indica la norma que el estudiante debe seguir minuciosamente para lograr que le tramiten la Transferencia de unos créditos previos realizados en una carrera extinta de otra sede española y su convalidación en el archivo? \\
        \begin{itemize}
          \item \textbf{Tipo de consulta:} simple
          \item \textbf{Respuesta del Agente:}

                ⚠️ Requisitos Previos:
                Antes de iniciar el trámite, es importante tener en cuenta que el estudiante debe estar admitido y matriculado en los correspondientes estudios en la Universidad de Sevilla, salvo en los casos de cambios de estudios oficiales de Grado. Además, es fundamental revisar la normativa específica de la universidad y los plazos establecidos en el calendario académico para la presentación de solicitudes.

                A continuación, se detallan los pasos a seguir para la transferencia de créditos:

                1. \textbf{Solicitud de Transferencia de Créditos}: El estudiante debe presentar una solicitud de transferencia de créditos dirigida al Decano o Director del Centro responsable del título para el que se solicita la transferencia. Esta solicitud debe incluir la documentación justificativa de los créditos obtenidos y su contenido académico, indicando los módulos, materias o asignaturas que considere superados.

                2. \textbf{Documentación Acreditativa}: El estudiante debe aportar la documentación acreditativa necesaria, que incluye:
                - \textbf{Certificación Académica Oficial} por traslado de expediente, en el caso de estudios cursados en otra universidad.
                - Cualquier otra documentación que se considere relevante para el proceso de transferencia de créditos.

                3. \textbf{Presentación de la Solicitud}: La solicitud y la documentación acreditativa deben presentarse en la \textbf{Secretaría del Centro} responsable del título para el que se solicita la transferencia, dentro de los plazos indicados en el calendario académico de cada curso.

                4. \textbf{Evaluación de la Solicitud}: La universidad evaluará la solicitud y la documentación presentada. Si es necesario, se podrán solicitar adicionales o aclaraciones.

                5. \textbf{Resolución}: La universidad emitirá una resolución sobre la solicitud de transferencia de créditos. Si la resolución es favorable, se indicarán los créditos que se transfieren y cómo se reflejarán en el expediente académico del estudiante.

                6. \textbf{Regularización de la Matrícula}: Una vez notificada la resolución, el estudiante dispondrá de un plazo de 15 días hábiles para regularizar su matrícula, si procede. Esto puede incluir la ampliación de la matrícula o la anulación de asignaturas no reconocidas.

                7. \textbf{Abono de Precios Públicos}: Si corresponde, el estudiante deberá abonar los precios públicos por los créditos transferidos dentro del plazo establecido.

                8. \textbf{Incorporación al Expediente}: Una vez completados los trámites anteriores, los créditos transferidos se incorporarán al expediente académico del estudiante, reflejándose en los documentos académicos oficiales acreditativos de los estudios seguidos por el mismo.

                Es importante seguir estos pasos de manera minuciosa y dentro de los plazos establecidos para asegurar el correcto trámite de la transferencia de créditos. Además, se recomienda consultar con la \textbf{Secretaría Virtual} o el personal de la universidad para cualquier duda o aclaración adicional.
          \item \textbf{Contexto utilizado:}
                \begin{scriptsize}
                  FUENTE: Normativa Reconocimiento y Transferencia de creditos.pdf | PÁGINA: 10
                  \# Textos consolidados

                  15.7. De los acuerdos de estudios que se establezcan se enviará copia a los Servicios Centrales del Rectorado que corresponda.
                  15.8. Con carácter general lo dispuesto en estas normas será de aplicación a la movilidad para dobles titulaciones sin perjuicio de las previsiones contenidas en los convenios respectivos.
                  15.9. Resultarán igualmente de aplicación las normas que eventualmente se aprueben por los órganos nacionales o internacionales competentes para cada programa específico de movilidad.

                  \#\# CAPÍTULO V
                  \#\#\# Transferencia de créditos
                  \textbf{Artículo 16. Definición.}
                  La transferencia de créditos implica que, en los documentos académicos oficiales acreditativos de las enseñanzas seguidas por cada estudiante, se incluirán la totalidad de los créditos obtenidos en enseñanzas universitarias oficiales cursadas con anterioridad, en ésta u otra universidad, que no hayan conducido a la finalización de sus estudios con la consiguiente obtención de un título oficial.

                  \textbf{Artículo 17. Aplicación.}
                  Los créditos correspondientes a materias o asignaturas previamente superadas por el estudiante, en enseñanzas universitarias no concluidas y que no puedan ser objeto de reconocimiento, serán transferidos a su expediente en los estudios a los que ha accedido con la calificación de origen y se reflejarán en los documentos académicos oficiales acreditativos de los estudios seguidos por el mismo, así como en el Suplemento Europeo al Título.

                  \#\# CAPÍTULO VI
                  \#\#\# Tramitación
                  \textbf{Artículo 18. Solicitudes de reconocimiento de créditos.}
                  18.1. Los expedientes de reconocimiento de créditos se tramitarán a solicitud del interesado, quién deberá aportar la documentación justificativa de los créditos obtenidos y su contenido académico, indicando los módulos, materias o asignaturas que considere superados.
                  18.2. Será requisito imprescindible estar admitido y matriculado en los correspondientes estudios, salvo en los casos de cambios de estudios oficiales de Grado, según lo dispuesto en la Resolución Rectoral por la que se regula la admisión a los títulos de Grado de la Universidad de Sevilla de los estudiantes que han iniciado anteriormente otros estudios universitarios.
                  18.3. Las solicitudes de reconocimiento de créditos tendrán su origen en materias o asignaturas realmente cursadas y superadas, en ningún caso se referirán a materias o asignaturas previamente reconocidas, convalidadas o adaptadas.
                  18.4. Las solicitudes se presentarán en la Secretaría del Centro responsable del título para el que se solicita el reconocimiento dentro de los plazos indicados en el calendario académico de cada curso.

                  \textbf{Artículo 19. Documentación acreditativa.}
                  19.1. En el caso de estudios universitarios cursados, estudios superiores no universitarios u otros estudios no oficiales, se aportará la siguiente documentación:

                  Este documento es de carácter informativo y no tiene valor jurídico.
                  10

                  ---

                  FUENTE: Normativa Reconocimiento y Transferencia de creditos.pdf | PÁGINA: 13
                  \# Textos consolidados

                  que en caso desestimatorio debe ser motivada académicamente, deberá dictarse y notificarse en un plazo máximo de tres meses desde la presentación de la solicitud.

                  20.6. El vencimiento del plazo sin haberse notificado resolución expresa legitimará al interesado para entenderla desestimada por silencio administrativo.

                  20.7. Contra las resoluciones del Decano o Director del Centro se podrá interponer recurso de alzada ante el Rector de la Universidad de Sevilla.

                  20.8. A efecto de la tramitación del procedimiento se declaran inhábiles los periodos no lectivos previstos en el calendario académico de cada curso.

                  \#\# Artículo 21. Solicitudes de transferencia de créditos.

                  Los expedientes de transferencia de créditos se tramitarán a petición del interesado. A estos efectos, los estudiantes que se incorporen a un nuevo estudio, mediante escrito dirigido al Decano o Director del Centro y en los plazos que se establezcan en el calendario académico de cada curso, indicarán si han cursado anteriormente otros estudios universitarios oficiales sin haberlos finalizado, aportando, en caso de no tratarse de estudios de la Universidad de Sevilla, la documentación justificativa que proceda de entre la contemplada en el artículo 19.1.

                  \#\# Artículo 22. Efectos del reconocimiento y transferencia de créditos.

                  22.1. En el proceso de reconocimiento quedarán reflejados de forma explícita aquellos módulos, materias o asignaturas que no deberán ser cursadas por el estudiante. Se entenderá en este caso que dichos módulos, materias o asignaturas ya han sido superadas, no serán susceptibles de nueva evaluación y se reflejarán en el expediente del estudiante como, módulos, materias o asignaturas reconocidas, indicándose el origen del reconocimiento.

                  22.2. En todo caso, el reconocimiento de créditos se referirá, al menos, a unidades de matrícula completas, es decir, no se podrá realizar el reconocimiento parcial de una asignatura.

                  22.3. Cuando la resolución del procedimiento dé lugar al reconocimiento de créditos optativos, el número de créditos reconocidos se minorará del número de créditos optativos exigido por el correspondiente plan de estudios y se reflejará en el expediente del estudiante como créditos optativos reconocidos, indicándose el origen del reconocimiento. En todo caso, el número de créditos optativos reconocidos no podrá superar el número de créditos exigido por el plan de estudios en cuestión.

                  22.4. En los casos procedentes, tras el proceso de reconocimiento de créditos, se permitirá a los interesados la ampliación de su matrícula en los términos recogidos en las Normas de Matrícula de cada curso académico.

                  22.5. La calificación de las asignaturas o, en su caso, de los créditos superados como consecuencia de un proceso de reconocimiento será equivalente a la calificación de las materias o asignaturas que han dado origen a éste. En caso necesario, se realizará la media ponderada cuando varias materias o asignaturas conlleven el reconocimiento de una sola en la titulación de destino.

                  22.6. Las calificaciones se reflejarán en el expediente académico en los términos recogidos en el Real Decreto 1125/2003, de 5 de septiembre, por el que se establece el sistema europeo de créditos y el sistema de calificaciones en las titulaciones universitarias de carácter oficial y validez en todo el territorio nacional.

                  22.7. No obstante lo anterior, cuando en el expediente académico de origen sólo se haga referencia a las calificaciones cualitativas, se transformarán en calificaciones numéricas, teniendo en cuenta la siguiente tabla de equivalencias:

                  Este documento es de carácter informativo y no tiene valor jurídico.

                  ---

                  FUENTE: R.R. NORMAS DE MATRICULA GRADO Y MASTER UNIVERSITARIO 2025-2026.pdf | PÁGINA: 20
                  \# UNIVERSIDAD D SEVII

                  C)EST

                  No será obligatorio matricular las asignaturas objeto de una solicitud de reconocimiento en el periodo ordinario de matrícula, salvo las de formación básica, si se realiza matrícula en otras asignaturas de un curso superior, cuya matrícula habrá de formalizarse junto con las demás asignaturas, en los términos previstos en el artículo 27 del Reglamento General de Actividades Docentes.

                  1. Una vez dictada y notificada la resolución al estudiante, se establecerá un plazo de 15 días hábiles para la regularización de su matrícula, si procede. En este plazo el estudiante podrá ampliar su matrícula, teniendo en cuenta los límites máximos previstos en el artículo 4 de estas normas.

                  Si la resolución fuese negativa, el estudiante podrá, en su caso, optar por formalizar matrícula de las asignaturas no reconocidas, en el mismo plazo indicado anteriormente, teniendo en cuenta los límites previstos en el Reglamento General de Actividades Docentes y en las presentes normas.

                  2. Si como consecuencia del reconocimiento de créditos el estudiante no alcanzara o superara el mínimo/máximo legal de matrícula, dispondrá de un plazo de 15 días desde la recepción de la notificación para regularizar su matrícula, con el apercibimiento de que, si no lo hiciere, se procedería a la anulación de su matrícula.

                  3. El reconocimiento estará condicionado al abono de los precios públicos correspondientes, en su caso, que deberá realizarse en el plazo de 15 días hábiles siguientes a la notificación. El impago en el plazo indicado se entenderá como una renuncia a los reconocimientos de créditos concedidos y estos quedarán sin efecto, previa notificación.

                  4. El estudiante podrá desistir de su solicitud de reconocimiento de créditos siempre que no se haya dictado resolución expresa por el Decano/Director.

                  5. Los estudiantes podrán solicitar la renuncia de los reconocimientos aceptados en el mismo plazo indicado de 15 días y antes de su abono e incorporación al expediente, cuando consideren que cursar de nuevo la asignatura sea beneficioso para el mejor conocimiento de la materia o si durante la tramitación del reconocimiento de créditos se ha realizado el examen y superado la asignatura por el estudiante.

                  \#\# Artículo 24. Transferencia de créditos

                  La transferencia de créditos implica que, en los documentos académicos oficiales acreditativos de las enseñanzas seguidas por cada estudiante, se incluirán la totalidad de los créditos obtenidos en enseñanzas oficiales cursadas con anterioridad, en esta u otra universidad, que no hayan conducido a la finalización de sus estudios con la consiguiente obtención de un título oficial y que no puedan ser objeto de reconocimiento.

                  \#\# Artículo 25. Solicitudes de transferencia de créditos

                  Los expedientes de transferencia de créditos se tramitarán a petición del estudiante. La solicitud de transferencia de créditos se podrá presentar durante todo el curso académico.

                  Para llevar a cabo la transferencia de créditos en titulaciones de la Universidad de Sevilla, los estudiantes deberán cerrar el expediente de la titulación abandonada y, solo en el caso de que hayan cursado los estudios en otra universidad, presentarán en la secretaría del centro el resguardo del traslado de expediente para que el centro de destino pueda incluir en el expediente académico del estudiante los créditos obtenidos en la titulación de origen.

                  La acreditación documental de los créditos cuya transferencia se solicita deberá efectuarse mediante Certificación Académica Oficial por traslado de expediente.

                  ----

                  \textbf{Código Seguro De Verificación}: eqOjhZcmL0p6VJlZiTH/Yw==
                  \textbf{Fecha}: 27/06/2025
                  \textbf{Firmado Por}: MIGUEL ANGEL CASTRO ARROYO
                  \textbf{Url De Verificación}: Verificación
                  \textbf{Página}: 16/50

                  ---

                  FUENTE: Normativa Reconocimiento y Transferencia de creditos.pdf | PÁGINA: 4
                  \# Textos consolidados

                  superado íntegramente un determinado módulo, el reconocimiento se llevará a cabo por materias o asignaturas en función de las competencias y conocimientos asociados a las mismas.

                  7) En el caso de títulos oficiales de Grado que habiliten para el ejercicio de profesiones reguladas en España, para los que el Gobierno haya establecido las condiciones a las que han de adecuarse los planes de estudios, se reconocerán los créditos de los módulos definidos en la correspondiente norma reguladora que hayan sido superados por el estudiante. En caso de no haberse superado íntegramente un determinado módulo, el reconocimiento se llevará a cabo por materias o asignaturas en función de las competencias y conocimientos asociados a las mismas.

                  \#\# Artículo 5. A partir de títulos de Máster Universitario.

                  5.1. Las solicitudes de reconocimiento de créditos basadas en estudios superados en títulos oficiales de Máster Universitario (tanto los regulados por el RD 56/2005, como por el RD 1393/2007) o periodo de formación específico del Doctorado se resolverán teniendo en cuenta la adecuación entre las competencias y conocimientos adquiridos en las materias superadas y los previstos en el plan de estudios del título de destino, indicándose las materias o asignaturas que se considerarán superadas por el interesado y que, por lo tanto, no estarán obligados a cursar.

                  5.2. En ningún caso podrán ser objeto de reconocimiento los Trabajos Fin de Grado.

                  \#\# Artículo 6. A partir de títulos de la anterior ordenación universitaria.

                  6.1. Las solicitudes de reconocimiento de créditos basadas en estudios superados en títulos universitarios de la anterior ordenación universitaria, Arquitecto Técnico, Diplomado, Ingeniero Técnico, Maestro, Licenciado, Arquitecto, Ingeniero o periodo de docencia del doctorado, se resolverán teniendo en cuenta la adecuación entre las competencias y conocimientos adquiridos en las materias superadas y los previstos en el plan de estudios del título de destino, indicándose las materias o asignaturas que se considerarán superadas por el interesado y que, por lo tanto, no estarán obligados a cursar.

                  6.2. Para la resolución de estas solicitudes se tendrán en cuenta los siguientes criterios:

                  1) Cuando las competencias y conocimientos no estén explicitados o no puedan deducirse del plan de estudios de origen del estudiante se tomarán como referencia el número de créditos y/o los contenidos de las materias o asignaturas cursadas.

                  2) En el caso de títulos en proceso de extinción por la implantación de los nuevos títulos de Grado, la adaptación de los estudiantes a estos últimos se basará en el reconocimiento de créditos previsto en la tabla de adaptación incluida en la correspondiente memoria de verificación del título de Grado en cuestión.

                  3) En los procesos de adaptación de estudiantes a los nuevos planes de los títulos de Grado deberá garantizarse que la situación académica de aquellos no resulte perjudicada. A tal efecto, las materias, asignaturas o créditos superados que no tengan equivalencia en las correspondientes al plan de estudios de Grado se incorporarán en el expediente del estudiante como créditos genéricos de carácter optativo. Si, aún así resultaran excedentes, los créditos restantes se podrán incorporar al expediente como créditos transferidos, a petición del interesado y siempre que se trate de materias o asignaturas completas.

                  4) El número de créditos reconocidos con cargo a la optatividad no podrá superar el número de créditos optativos exigido por el plan de estudios del título de destino.

                  5) En ningún caso podrán ser objeto de reconocimiento los Trabajos Fin de Grado.

                  Este documento es de carácter informativo y no tiene valor jurídico.

                  ---

                  FUENTE: Normativa Reconocimiento y Transferencia de creditos.pdf | PÁGINA: 2
                  \# Textos consolidados

                  \#\# Artículo 19. Documentación acreditativa

                  \#\# Artículo 20. Órganos responsables

                  \#\# Artículo 21. Solicitudes de transferencia de créditos

                  \#\# Artículo 22. Efectos del reconocimiento y transferencia de créditos

                  \# DISPOSICIONES ADICIONALES

                  \#\# Disposición Adicional Primera. Títulos conjuntos y dobles titulaciones

                  \#\# Disposición Adicional Segunda. Reconocimiento parcial de estudios extranjeros

                  \#\# Disposición adicional Tercera. Aplicabilidad a los Centros Adscritos

                  \#\# Disposición Adicional Cuarta. Cita en género femenino de los preceptos de estas normas

                  \# DISPOSICIÓN DEROGATORIA

                  \# DISPOSICIONES FINALES

                  \#\# Disposición final Primera. Título competencial

                  \#\# Disposición final Segunda. Habilitación para el desarrollo normativo

                  \#\# Disposición final Tercera. Entrada en vigor

                  \#\# INTRODUCCIÓN

                  El Real Decreto 1393/2007, de 29 de octubre, por el que se establece la ordenación de las enseñanzas universitarias oficiales indica que, con objeto de hacer efectiva la movilidad de estudiantes, tanto dentro del territorio nacional como fuera de él, las universidades elaborarán y harán pública su normativa sobre el sistema de reconocimiento y transferencia de créditos, con sujeción a los criterios generales establecidos en el mismo.

                  La Universidad de Sevilla, a fin de dar cumplimiento al mencionado precepto, aprobó mediante Acuerdo 5.1/C.G. 30-09-2008 las Normas Básicas sobre Reconocimiento y Transferencia de Créditos de aplicación a los estudios universitarios oficiales de Grado y Máster. Posteriormente, el Real Decreto 861/2010, de 2 de julio, por el que se modifica el anterior, modifica sustancialmente el apartado correspondiente al régimen de reconocimiento y transferencia de créditos introduciendo nuevas posibilidades de reconocimiento académico, especialmente a partir de la experiencia laboral y profesional y a partir de estudios cursados en títulos propios.

                  Por todo ello, el Consejo de Gobierno de la Universidad de Sevilla acuerda modificar las Normas Básicas aprobadas por el Acuerdo 5.1/C.G. 30-09-2008, que quedarán establecidas según las siguientes normas reguladoras:

                  \# CAPÍTULO I

                  \#\# Disposiciones generales

                  \#\#\# Artículo 1. Objeto.

                  Las presentes normas tienen por objeto establecer los criterios generales y el procedimiento para el reconocimiento y la transferencia de créditos en las enseñanzas universitarias de Grado y Máster previstas en el Real Decreto 1393/2007, de 29 de octubre, por el que se establece la ordenación de las enseñanzas universitarias oficiales, modificado por el Real Decreto 861/2010, de 2 de julio.

                  Este documento es de carácter informativo y no tiene valor jurídico.
                \end{scriptsize}
        \end{itemize}


  \item \textbf{P50.} En el portal de Currículum Nacional, ¿qué pasos en orden o acciones exactas debe ejecutar cualquier usuario para exportar e integrar adecuadamente y en última instancia su CVN en la solicitud principal creada para Bareprof? \\
        \begin{itemize}
          \item \textbf{Tipo de consulta:} simple
          \item \textbf{Respuesta del Agente:}

                ⚠️ Requisitos Previos:
                Antes de comenzar, asegúrate de tener acceso al portal de Currículum Nacional y de haber creado previamente tu CVN en el editor CVN-FECYT.

                Para exportar e integrar adecuadamente tu CVN en la solicitud principal creada para Bareprof, sigue estos pasos:

                1. \textbf{Accede al editor CVN-FECYT}: Ingresa a la dirección [http://cvn.fecyt.es/editor](http://cvn.fecyt.es/editor) y accede con tu correo electrónico y contraseña. Si no estás registrado, crea una cuenta.
                2. \textbf{Editar o importar tu CVN}: Una vez dentro, puedes editar tu CVN existente o importar uno nuevo si lo necesitas. Asegúrate de introducir solo los méritos y campos requeridos por Bareprof.
                3. \textbf{Descargar y publicar tu CVN}: En el menú principal del editor CVN-FECYT, selecciona la opción \textbf{Descargar y publicar mi CVN}. Esto te llevará a una pantalla donde podrás generar y descargar tu CVN en formato PDF.
                4. \textbf{Selecciona la opción COMPLETO}: En la pantalla de descarga, selecciona la opción \textbf{COMPLETO} y pulsa el botón \textbf{Generar}. Esto generará un CVN completo que incluya todos tus méritos.
                5. \textbf{Descarga el CVN}: Una vez generado, podrás descargar tu CVN en formato PDF. Asegúrate de guardar el archivo en un lugar accesible.
                6. \textbf{Accede a Bareprof}: Ingresa a la plataforma Bareprof y busca la sección donde debes subir tu CVN. Generalmente, esto se encuentra en la sección de \textbf{Requisitos de admisión} o \textbf{Justificación de méritos}.
                7. \textbf{Sube tu CVN a Bareprof}: Selecciona el archivo CVN que descargaste en el paso 5 y subelo a la plataforma Bareprof. Asegúrate de que el archivo no supere el tamaño máximo permitido (generalmente 2 MB).
                8. \textbf{Verifica la información}: Una vez subido, Bareprof procesará la información de tu CVN y mostrará los méritos leídos por la aplicación. Verifica que todos tus méritos estén correctamente registrados y que no haya errores.

                Recuerda que es importante seguir las instrucciones del baremo y del documento de Preguntas Frecuentes para asegurarte de que tu CVN esté correctamente cumplimentado y que tus méritos sean evaluados correctamente por las comisiones de contratación.
          \item \textbf{Contexto utilizado:}
                \begin{scriptsize}
                  FUENTE: 27830\_guia-ayuda-candidato-r04-3-.pdf | PÁGINA: 16
                  \# Guía de ayuda al candidato

                  Presentar un número mayor de méritos no mejorará la puntuación, siendo recomendable una adecuada selección de los méritos a presentar. La forma de introducir en el CVN los distintos méritos se describe en el apartado 6 (Cumplimentación del CVN para su uso en Bareprof) de este documento. Se recomienda visionar los vídeos de ayuda disponibles en la plataforma Bareprof.

                  Una vez haya introducido los méritos en el CVN, debe proceder a la descarga de este para su posterior incorporación a Bareprof. Para ello, seleccione la opción Descargar y publicar mi CVN en el menú principal del editor CVN-FECYT (véase la Figura 1). Se accede entonces a una pantalla donde se muestra un listado de los CVN generados con anterioridad y se puede generar el nuevo. Seleccione para ello la opción COMPLETO y pulse el botón Generar.

                  \#\# Descargar y hacer público mi CVN

                  Pulsando sobre el botón Generar puede descargar una copia de su CVN.

                  \vspace{0.4cm}
                  \begin{table}[h!]\centering\renewcommand{\arraystretch}{1.3}
                    \begin{tabular}{|p{0.45\textwidth}|p{0.45\textwidth}|}\hline
                      COMPLETO                   & Generar                                                                 \\\hline
                      Histórico de CVN Generados & Listado con los diez últimos CVNs generados desde su cuenta de usuario. \\\hline
                    \end{tabular}\end{table}
                  \vspace{0.4cm}

                  \vspace{0.4cm}
                  \begin{table}[h!]\centering\renewcommand{\arraystretch}{1.3}
                    \begin{tabular}{|p{0.45\textwidth}|p{0.45\textwidth}|p{0.45\textwidth}|p{0.45\textwidth}|p{0.45\textwidth}|}\hline
                      Nombre de fichero         & Tipo CVN & N de páginas & Fecha de generación & Hacer público \\\hline
                      cvn 20220928211538503.pdf & COMPLETO & 13           & 28/09/2022-21:15:56 & Publicar      \\\hline
                    \end{tabular}\end{table}
                  \vspace{0.4cm}

                  Para descargar el CVN generado (o cualquier otro de los que se muestran en el listado), se debe clicar sobre el nombre que aparece en la primera columna. Almacene el fichero en disco para posteriormente incorporarlo a Bareprof.

                  \#\# 5 Aspectos de interés sobre la baremación

                  \#\#\# 5.1 Méritos a considerar para su incorporación a Bareprof

                  El baremo de contratación de la Universidad de Sevilla solo contempla un conjunto limitado de todos los posibles méritos que conforman la vida académica y profesional de un candidato. No puede pretenderse, por tanto, incorporar todos los méritos para su uso en Bareprof, pues es posible que algunos no estén recogidos en el baremo. El candidato debe analizar sus méritos y seleccionar solo aquellos que correspondan a epígrafes del baremo.

                  A la hora de preparar el CVN, no hay ningún inconveniente en incluir en él todos los méritos de que se disponga, ya que Bareprof ignorará todos aquellos registros correspondientes a méritos no contemplados en el baremo. Muchos de los epígrafes del baremo establecen un número máximo de aportaciones (por ejemplo, se valorarán como máximo 5 libros). Si se supera ese número, Bareprof ordenará los ítems aportados y seleccionará los primeros, hasta completar el número máximo. El criterio de ordenación se indica para cada tipo de mérito en los distintos subapartados del apartado 6 de esta guía. Así, por ejemplo, los libros se ordenan descendentemente según la fecha de publicación, de tal modo que se seleccionarán los 5 de publicación más reciente.

                  Revisión 4 (marzo, 2025)

                  ---

                  FUENTE: 27830\_guia-ayuda-candidato-r04-3-.pdf | PÁGINA: 2
                  1. \textbf{INTRODUCCIÓN} .........................................................................................................................................................3
                  2. \textbf{DOCUMENTOS DE REFERENCIA} .................................................................................................................................3
                  3. \textbf{UTILIZACIÓN DE LA PLATAFORMA DE BAREMACIÓN DE LA UNIVERSIDAD DE SEVILLA} ..............................................4
                  3.1. \textbf{ACCESO A BAREPROF} ....................................................................................................................................................... 4
                  3.2. \textbf{OBTENCIÓN DE AYUDA} .................................................................................................................................................... 6
                  3.3. \textbf{INSCRIPCIÓN EN EL CONCURSO} .......................................................................................................................................... 8
                  3.4. \textbf{REQUISITOS DE ADMISIÓN} ................................................................................................................................................ 9
                  3.5. \textbf{JUSTIFICACIÓN DE MÉRITOS} ............................................................................................................................................ 11
                  3.6. \textbf{DATOS ADICIONALES} ..................................................................................................................................................... 12
                  3.7. \textbf{FINALIZACIÓN DEL PROCESO DE INSCRIPCIÓN} ...................................................................................................................... 13
                  4. \textbf{GENERACIÓN DEL CVN} .............................................................................................................................................14
                  5. \textbf{ASPECTOS DE INTERÉS SOBRE LA BAREMACIÓN} ......................................................................................................15
                  5.1. \textbf{MÉRITOS A CONSIDERAR PARA SU INCORPORACIÓN A BAREPROF} ............................................................................................ 15
                  5.2. \textbf{REUBICACIÓN DE MÉRITOS} ............................................................................................................................................. 17
                  5.3. \textbf{DUPLICIDAD DE MÉRITOS} ............................................................................................................................................... 17
                  5.4. \textbf{PUNTUACIONES ESCALADAS} ............................................................................................................................................ 18
                  6. \textbf{CUMPLIMENTACIÓN DEL CVN PARA SU USO EN BAREPROF} ....................................................................................19
                  6.1. \textbf{DATOS GENERALES} ........................................................................................................................................................ 19
                  6.2. \textbf{RESUMEN DEL CURRÍCULUM VITAE} ................................................................................................................................... 20
                  6.3. \textbf{RENDIMIENTO ACADÉMICO} ............................................................................................................................................. 20
                  6.3.1. \textbf{Expediente académico} ..................................................................................................................................................... 20
                  6.3.1.1. \textbf{Titulación de grado o equivalente} ..................................................................................................................... 20
                  6.3.1.2. \textbf{Titulación de máster universitario oficial} .......................................................................................................... 22

                  ---

                  FUENTE: 27830\_guia-ayuda-candidato-r04-3-.pdf | PÁGINA: 4
                  \# Guía de ayuda al candidato

                  \#\# 1. Introducción

                  La Universidad de Sevilla cuenta con una plataforma web para la baremación de los candidatos participantes en los concursos públicos que se convoquen para la contratación de Profesorado Ayudante Doctor, Profesorado Asociado y Profesorado Sustituto. Dicha plataforma, denominada Bareprof, permite gestionar de forma integrada todo el proceso de admisión al concurso y posterior valoración de méritos. A ella tienen acceso los candidatos que participan en tales convocatorias, las comisiones de contratación encargadas de baremar las solicitudes de los candidatos y los técnicos del Vicerrectorado de Personal Docente e Investigador responsables de la gestión de las convocatorias. El candidato hará uso de la plataforma Bareprof para inscribirse en la plaza o plazas a las que desee concursar. Al acceder a la plataforma, el candidato tendrá acceso a la información de las plazas que se encuentran en periodo de inscripción, con indicación del número de días restantes para la conclusión del plazo. El proceso de inscripción en una plaza por medio de Bareprof se describe en el apartado 3 de esta guía.

                  Durante el proceso de inscripción, el candidato debe presentar, junto a otra documentación requerida por la convocatoria, aquellos méritos que desea sean baremados por la Comisión de Contratación de la plaza a la que opta. La introducción de méritos por parte del candidato en Bareprof se realiza de forma indirecta. Así, los méritos deberán ser previamente registrados en el Curriculum Vitae Normalizado del candidato, cuya creación y mantenimiento se realiza a través del editor CVN de la Fundación Española para la Ciencia y la Tecnología (FECYT), dependiente del Ministerio de Ciencia, Innovación y Universidades. En el apartado 4 de esta guía se proporciona una breve descripción del proceso de creación del CVN.

                  Cada mérito contemplado en el baremo de la Universidad de Sevilla debe ser introducido en el CVN siguiendo las instrucciones dadas en el apartado 5 de esta guía. Este es un proceso que requiere especial atención por parte del candidato y en el que es muy importante que se sigan de forma precisa las indicaciones dadas para cada mérito, ya que, de lo contrario, el mérito puede llegar a no ser baremado por la Comisión al no ser reconocido por Bareprof.

                  \#\# 2. Documentos de referencia

                  Se relacionan a continuación los principales documentos relacionados con el procedimiento de baremación para la contratación de Profesorado Ayudante Doctor, Profesorado Asociado y Profesorado Sustituto, que se encuentran disponibles en el Boletín Oficial de la Universidad de Sevilla (http://www.us.es/bous). Deberán además tenerse en cuenta todas aquellas instrucciones aclaratorias del baremo que pudieran dictarse.

                  - Normativa sobre Procedimiento de Contratación de Profesorado Ayudante Doctor, Profesorado Asociado y Profesorado Interino, junto con el Baremo establecido para dichas figuras contractuales (BOUS 12/2024, de 11 de octubre).
                    [https://www.us.es/bous-numeros/numero-12-11-de-octubre-de-2024](https://www.us.es/bous-numeros/numero-12-11-de-octubre-de-2024)

                  - Instrucción de 17-10-19 sobre los requisitos para el reconocimiento de las estancias de investigación dentro del mérito preferente en los concursos de profesorado ayudante doctor, asociado y sustituto interino (BOUS 14/2019, de 21 de noviembre).
                    [https://www.us.es/bous-numeros/numero-14-21-de-noviembre-de-2019](https://www.us.es/bous-numeros/numero-14-21-de-noviembre-de-2019)

                  - Resolución rectoral, de 24 marzo de 2021, por la que se emite una interpretación sobre la determinación del umbral mínimo aplicable en los procedimientos selectivos para plazas de.

                  ---

                  FUENTE: 27830\_guia-ayuda-candidato-r04-3-.pdf | PÁGINA: 10
                  \# Guía de ayuda al candidato

                  \#\# 3.4 Requisitos de admisión

                  En este paso del proceso de inscripción, el candidato deberá subir a la plataforma Bareprof los documentos (en formato pdf) que acreditan el cumplimiento de los requisitos de admisión establecidos en la convocatoria del concurso público.

                  \vspace{0.4cm}
                  \begin{table}[h!]\centering\renewcommand{\arraystretch}{1.3}
                    \begin{tabular}{|p{0.45\textwidth}|p{0.45\textwidth}|p{0.45\textwidth}|p{0.45\textwidth}|p{0.45\textwidth}|p{0.45\textwidth}|}\hline
                      D                        & B                                & a                                        & T                              &                        &          \\\hline
                      Solicitud de inscripción & REQUISITOS DE ADMISIÓN           & Justificación de méritos                 & Datos adicionales              &                        &          \\\hline
                      CVN (pdf max: 2 MB)      & Título académico (pdf max: 2 MB) & Solicitud de inscripción (pdf max: 2 MB) & Pago de tasas (pdf max: 1 MB)¹ &                        &          \\\hline
                      Seleccionar archivo...   & Examinar                         & Seleccionar archivo...                   & Examinar                       & Seleccionar archivo... & Examinar \\\hline
                    \end{tabular}\end{table}
                  \vspace{0.4cm}

                  - CVN: Una vez seleccionado el archivo que contiene el CVN, generado por medio del editor CVN-FECYT, Bareprof procesará la información que contiene y mostrará finalmente el Visor de CVN, en el que podrá ver los méritos leídos por la aplicación y los errores detectados en el proceso.

                  Visor de CVN

                  A continuación se muestran los datos que usted ha introducido en los campos leídos por Bareprof, solo aquellos que estén enmarcados en color rojo se pondrán a disposición de las comisiones de contratación. Solo los datos leídos por Bareprof serán evaluados por las comisiones de contratación. Recuerde también que los méritos incluidos en apartados del baremo a los que no corresponden por su tipología, no serán evaluados por las comisiones de contratación, procediendo a su no validación. Se le recuerda que es necesario que siga detenidamente las instrucciones del baremo y el documento de Preguntas Frecuentes en la elaboración de su CVN.

                  - Méritos con campos obligatorios sin cumplimentar
                  - Méritos con campos de texto libre no puntuables
                  - Méritos con información no indexada

                  - Rendimiento académico
                  - Actividad docente
                  - Actividad investigadora
                  - Transferencia del conocimiento
                  - Actividad profesional externa
                  - Mérito preferente
                  - Resumen CV
                  - Resumen lectura CVN
                  - Expediente académico
                  - Titulación universitaria (grado o equivalente)

                  \#\#\# Mérito

                  - Nombre del título: Ingeniero Industrial Especialidad Electricidad
                  - Titulación universitaria: Titulado Superior
                  - Nota media del expediente:
                  - Fecha de titulación: 17/06/2019

                  \#\#\# Titulación de máster universitario oficial

                  - Mérito
                  - Tipo de formación: Máster
                  - Titulación de posgrado: Máster Universitario en Ingeniería de la Energía
                  - Fecha de titulación: 22/06/2020
                  - Calificación obtenida: 300 horas

                  Los méritos que aparecen recuadrados en rojo contienen algún campo obligatorio que no ha sido cumplimentado, marcándose también en ese color los campos en los que se han detectado las incidencias. Es muy importante que proceda a su corrección y vuelva a subir el CVN a Bareprof, ya que estos méritos no serán puestos a disposición de la comisión para su baremación. En color azul aparecen recuadrados los méritos que contienen opciones o literales no admitidos en alguno de sus campos, por lo que podrían no ser valorados por la comisión de contratación.

                  Revisión 4 (marzo, 2025) 9

                  ---

                  FUENTE: 27830\_guia-ayuda-candidato-r04-3-.pdf | PÁGINA: 15
                  \# Guía de ayuda al candidato

                  \#\# 4 Generación del CVN

                  El sistema de baremación de la Universidad de Sevilla hace uso del Curriculum Vitae Normalizado (CVN) de la Fundación Española para la Ciencia y Tecnología (FECYT) como herramienta para la introducción de méritos por parte de los candidatos. Para la generación y mantenimiento del CVN, FECYT proporciona un editor accesible vía web en la dirección [http://cvn.fecyt.es/editor](http://cvn.fecyt.es/editor). Si es la primera vez que lo utiliza, deberá registrarse, requiriéndose como datos obligatorios para ello su nombre y primer apellido, así como una dirección de correo electrónico.

                  \#\#\# Acceso al editor

                  - Correo electrónico
                  - Contraseña
                  - ¿Ha olvidado la contraseña?
                  - No cerrar sesión

                  \textbf{Accede al editor}
                  Si no está registrado, pulse aquí

                  Una vez haya introducido su correo electrónico y contraseña, se accede al siguiente menú de opciones:

                  - \textbf{Accesos directos}
                  - Editar CVN
                  - Utilizar curriculo existente (importar)
                  - Descargar y publicar mi CVN
                  - Cambiar contraseña
                  - Borrar mis datos CVN
                  - Donde introduzco este mérito?

                  \textbf{Figura 1. Menú principal del editor CVN-FECYT}

                  Puede introducir sus méritos usando la opción Editar CVN o bien importar un CVN del que disponga con anterioridad por medio de la opción Utilizar curriculum existente (importar); este puede haber sido generado en una plataforma distinta como, por ejemplo, el Sistema de Información Científica de Andalucía (SICA). Si solo va a utilizar el CVN de FECYT para los concursos de la Universidad de Sevilla, conviene que solo introduzca los méritos y campos requeridos por Bareprof. Además, debe tener en cuenta que ciertos bloques del sistema de baremación tienen limitado el número de ítems a.

                  Revisión 4 (marzo, 2025)
                \end{scriptsize}
        \end{itemize}


\end{itemize}

\end{document}